\documentclass[11pt,oneside]{amsbook}

% ============================================================================
% PACKAGE IMPORTS
% ============================================================================

% Mathematics packages
\usepackage{mathtools}  % Enhanced version of amsmath (already loaded by amsbook)
\usepackage{amssymb}    % Math symbols (blackboard bold, etc.)
\usepackage{amsthm}     % Theorem environments (already loaded by amsbook)

% Code listings with syntax highlighting
\usepackage{minted}     % Requires: pip install Pygments
\setminted{
    breaklines=true,
    fontsize=\small,
    linenos=true,
    frame=lines,
    framesep=2mm,
    baselinestretch=1.2,
    bgcolor=gray!5
}

% Algorithms
\usepackage{algorithm}
\usepackage{algorithmic}

% Graphics and figures
\usepackage{graphicx}
\usepackage{caption}
\usepackage{subcaption}

% Tables
\usepackage{booktabs}   % Professional tables
\usepackage{array}      % Extended column formatting
\usepackage{longtable}  % Multi-page tables

% Bibliography
\usepackage[round,authoryear]{natbib}  % Author-year citations with \citet, \citep

% Hyperlinks (LOAD LAST)
\usepackage{hyperref}
\hypersetup{
    colorlinks=true,
    linkcolor=blue,
    citecolor=blue,
    urlcolor=blue,
    pdftitle={Causal Inference Mastery: A Computational Approach},
    pdfauthor={Brandon Behring},
    bookmarks=true,
    bookmarksnumbered=true
}

% ============================================================================
% THEOREM ENVIRONMENTS
% ============================================================================

\theoremstyle{definition}
\newtheorem{definition}{Definition}[chapter]
\newtheorem{example}{Example}[chapter]
\newtheorem{exercise}{Exercise}[chapter]

\theoremstyle{plain}
\newtheorem{theorem}{Theorem}[chapter]
\newtheorem{lemma}[theorem]{Lemma}
\newtheorem{proposition}[theorem]{Proposition}
\newtheorem{corollary}[theorem]{Corollary}
\newtheorem{assumption}{Assumption}[chapter]

\theoremstyle{remark}
\newtheorem{remark}{Remark}[chapter]
\newtheorem{note}{Note}[chapter]

% ============================================================================
% CUSTOM COMMANDS
% ============================================================================

% Probability and statistics
\newcommand{\E}{\mathbb{E}}         % Expectation
\newcommand{\Var}{\mathrm{Var}}     % Variance
\newcommand{\Cov}{\mathrm{Cov}}     % Covariance
\newcommand{\Prob}{\mathbb{P}}      % Probability

% Common sets
\newcommand{\R}{\mathbb{R}}         % Real numbers
\newcommand{\N}{\mathbb{N}}         % Natural numbers
\newcommand{\Z}{\mathbb{Z}}         % Integers

% Indicators and norms
\newcommand{\indicator}[1]{\mathbb{1}\{#1\}}
\newcommand{\norm}[1]{\|#1\|}

% Convergence
\newcommand{\convd}{\xrightarrow{d}}   % Convergence in distribution
\newcommand{\convp}{\xrightarrow{p}}   % Convergence in probability
\newcommand{\convas}{\xrightarrow{a.s.}}  % Almost sure convergence

% Operators
\DeclareMathOperator*{\argmax}{arg\,max}
\DeclareMathOperator*{\argmin}{arg\,min}
\DeclareMathOperator*{\plim}{plim}

% Causal inference specific
\newcommand{\ate}{\tau}                    % Average treatment effect
\newcommand{\att}{\tau_{\text{ATT}}}       % ATT
\newcommand{\atc}{\tau_{\text{ATC}}}       % ATC
\newcommand{\cate}{\tau(x)}                % Conditional ATE
\newcommand{\poty}[1]{Y^{(#1)}}            % Potential outcome Y(t)
\newcommand{\propensity}{e(X)}             % Propensity score
\newcommand{\independent}{\perp\!\!\!\perp} % Statistical independence
\newcommand{\indep}{\perp\!\!\!\perp}       % Alias for independence

% Insight box environment
\newenvironment{insight}{\par\medskip\noindent\textbf{$\bigstar$ Insight:}\itshape}{\par\medskip}

% ============================================================================
% DOCUMENT METADATA
% ============================================================================

\title{Causal Inference Mastery:\\
A Computational Approach to Treatment Effect Estimation}

\author{Brandon Behring}

\date{\today}

% ============================================================================
% DOCUMENT BODY
% ============================================================================

\begin{document}

% ----------------------------------------------------------------------------
% FRONT MATTER
% ----------------------------------------------------------------------------

\frontmatter

\maketitle

\tableofcontents

\chapter{Preface}

This book develops the theory and practice of causal inference for treatment effect estimation, with a focus on computational implementation and rigorous validation.

\textbf{What This Book Covers:}
\begin{itemize}
    \item Part I: Foundations---potential outcomes, randomized experiments, statistical inference
    \item Part II: Selection on Observables---propensity scores, weighting, CATE estimation, sensitivity analysis
    \item Part III: Selection on Unobservables---instrumental variables, weak instruments, control function, partial identification
    \item Part IV: Natural Experiments---difference-in-differences, regression discontinuity, synthetic control, kink designs
    \item Part V: Advanced Methods---quantile treatment effects, marginal treatment effects, mediation analysis, dynamic treatment regimes
    \item Part VI: Implementation---principal stratification, Bayesian causal inference, cross-language validation
\end{itemize}

\textbf{Key Features:}
\begin{itemize}
    \item Rigorous mathematical foundations with complete identification arguments
    \item Production-ready Python implementations with full API documentation
    \item Comprehensive Monte Carlo validation for every estimator
    \item Julia implementations provided in Appendix B for cross-language verification
\end{itemize}

\textbf{Prerequisites:}
\begin{itemize}
    \item Linear regression, probability theory, statistical inference
    \item Python programming (NumPy, pandas, scikit-learn)
    \item Basic familiarity with machine learning
\end{itemize}

\textbf{Companion Code:}

All implementations are available in the \texttt{causal\_inference\_mastery} repository, with over 90,000 lines of tested Python and Julia code across 26 method families.

% ----------------------------------------------------------------------------
% MAIN MATTER
% ----------------------------------------------------------------------------

\mainmatter

% ============================================================================
% PART I: FOUNDATIONS
% ============================================================================

\part{Foundations}
\label{part:foundations}

\chapter{The Potential Outcomes Framework}
\label{ch:potential_outcomes}

% ============================================================================
% Chapter 1: Potential Outcomes Framework
% Status: COMPLETE
% Est. Pages: 20
% ============================================================================

\section{Introduction}
\label{sec:po_intro}

Causal inference is the science of answering ``what if'' questions. When we ask whether a medical treatment reduces mortality, whether a job training program increases wages, or whether a policy intervention affects economic outcomes, we are asking causal questions. These questions differ fundamentally from descriptive or predictive questions because they concern the \emph{effect} of an intervention---what would happen if we changed something in the world.

The challenge is that causal effects are inherently unobservable. We cannot simultaneously observe what happens to a patient who receives a drug and what would have happened to that \emph{same} patient had they not received it. This ``fundamental problem of causal inference'' \cite{holland1986statistics} motivates the entire machinery we develop in this book.

\subsection{Correlation Is Not Causation}

The adage ``correlation is not causation'' captures a deep truth about empirical research. Consider the observation that hospital patients who receive a particular treatment have higher mortality rates than those who do not. Does this mean the treatment is harmful? Not necessarily---patients receiving the treatment may be sicker to begin with. The correlation between treatment and mortality confounds the treatment effect with selection into treatment.

\begin{example}[Selection Bias in Observational Data]
\label{ex:selection_bias}
Suppose we observe outcomes $Y_i$ and treatment status $T_i$ for $n$ patients. A naive comparison of means yields:
\begin{equation}
\underbrace{\E[Y_i \mid T_i = 1] - \E[Y_i \mid T_i = 0]}_{\text{Observed difference}} = \underbrace{\text{Causal Effect}}_{\text{What we want}} + \underbrace{\text{Selection Bias}}_{\text{Confounding}}
\end{equation}
The selection bias arises because treated and untreated groups differ systematically in ways that also affect outcomes.
\end{example}

\subsection{The Role of Formalization}

The potential outcomes framework, also known as the Rubin Causal Model \cite{rubin1974estimating}, provides a formal mathematical language for defining causal effects and stating the assumptions under which they can be identified from data. This formalization serves several purposes:

\begin{enumerate}
    \item \textbf{Precision}: It forces us to state exactly what we mean by a ``causal effect.''
    \item \textbf{Transparency}: It makes assumptions explicit and scrutinizable.
    \item \textbf{Generality}: It provides a unified framework for analyzing diverse study designs.
    \item \textbf{Computation}: It connects theoretical estimands to practical estimators.
\end{enumerate}

The framework builds on foundational work by \cite{neyman1923application} in the context of randomized experiments and was extended to observational studies by \cite{rubin1974estimating, rubin1978bayesian}. The synthesis presented by \cite{imbens2015causal} provides the modern treatment we follow here.

\subsection{Chapter Overview}

This chapter develops the core concepts of the potential outcomes framework:
\begin{itemize}
    \item Section~\ref{sec:po_framework}: Defines potential outcomes and individual treatment effects
    \item Section~\ref{sec:fundamental_problem}: Explains why causal effects are fundamentally unobservable
    \item Section~\ref{sec:ate}: Introduces population-level causal estimands (ATE, ATT, ATC)
    \item Section~\ref{sec:sutva}: Presents the Stable Unit Treatment Value Assumption
    \item Section~\ref{sec:identification}: Previews identification strategies developed in later chapters
    \item Section~\ref{sec:po_summary}: Summarizes key concepts and notation
\end{itemize}


% ============================================================================
\section{The Potential Outcomes Framework}
\label{sec:po_framework}
% ============================================================================

\subsection{Units, Treatments, and Outcomes}

Consider a study with $n$ units indexed by $i = 1, \ldots, n$. Each unit could be an individual, a firm, a geographic region, or any entity to which a treatment might be applied. We consider a binary treatment where each unit either receives the treatment ($T_i = 1$) or does not ($T_i = 0$).

\begin{definition}[Potential Outcomes]
\label{def:potential_outcomes}
For each unit $i$, we define two \emph{potential outcomes}:
\begin{itemize}
    \item $\poty{1}_i$: the outcome that would be observed if unit $i$ receives treatment
    \item $\poty{0}_i$: the outcome that would be observed if unit $i$ does not receive treatment
\end{itemize}
These potential outcomes are considered fixed attributes of the unit, existing conceptually whether or not we observe them.
\end{definition}

The notation $\poty{t}_i$ emphasizes that potential outcomes are \emph{functions} of the treatment. They represent what \emph{would} happen under different hypothetical interventions, not what we actually observe.

\begin{remark}
Some authors use $Y_i(t)$ or $Y_i^t$ instead of $\poty{t}_i$. We adopt the parenthetical superscript notation to distinguish potential outcomes from observed outcomes clearly.
\end{remark}

\subsection{Individual Treatment Effects}

Given the two potential outcomes, the causal effect of treatment for unit $i$ is naturally defined as their difference.

\begin{definition}[Individual Treatment Effect]
\label{def:ite}
The \emph{individual treatment effect} (ITE) for unit $i$ is:
\begin{equation}
\tau_i \equiv \poty{1}_i - \poty{0}_i
\label{eq:ite}
\end{equation}
This represents the change in outcome that would result from switching unit $i$ from control to treatment.
\end{definition}

The individual treatment effect $\tau_i$ is a \emph{causal} quantity because it compares outcomes for the \emph{same} unit under different treatments. There is no selection bias when comparing $\poty{1}_i$ to $\poty{0}_i$ because we are comparing the same unit to itself.

\begin{example}[Heterogeneous Treatment Effects]
\label{ex:heterogeneous}
Consider a job training program. Individual treatment effects might vary across workers:
\begin{itemize}
    \item Worker A (low skills): $\tau_A = \$5{,}000$ (large benefit from training)
    \item Worker B (high skills): $\tau_B = \$500$ (small marginal benefit)
    \item Worker C (wrong field): $\tau_C = -\$200$ (training is counterproductive)
\end{itemize}
This \emph{treatment effect heterogeneity} is the focus of Chapters~\ref{ch:cate} and~\ref{ch:forests}.
\end{example}

\subsection{The Switching Equation}

While both potential outcomes exist conceptually, we only observe one of them for each unit. The \emph{observed outcome} $Y_i$ is determined by the treatment received:

\begin{equation}
Y_i = T_i \cdot \poty{1}_i + (1 - T_i) \cdot \poty{0}_i =
\begin{cases}
\poty{1}_i & \text{if } T_i = 1 \\
\poty{0}_i & \text{if } T_i = 0
\end{cases}
\label{eq:switching}
\end{equation}

This is the \emph{switching equation} or \emph{consistency assumption}---the observed outcome equals the potential outcome corresponding to the treatment actually received.

\begin{definition}[Consistency]
\label{def:consistency}
The \emph{consistency assumption} states that for each unit:
\begin{equation}
Y_i = \poty{T_i}_i
\end{equation}
That is, the observed outcome equals the potential outcome indexed by the treatment received.
\end{definition}

The switching equation connects the conceptual (potential outcomes) to the observable (realized outcomes). It is sometimes called ``SUTVA Part 2'' and is discussed further in Section~\ref{sec:sutva}.

\subsection{Notation Summary}

Table~\ref{tab:notation} summarizes the core notation used throughout this book.

\begin{table}[ht]
\centering
\caption{Core Notation for the Potential Outcomes Framework}
\label{tab:notation}
\begin{tabular}{cl}
\toprule
\textbf{Symbol} & \textbf{Definition} \\
\midrule
$n$ & Number of units \\
$i$ & Unit index, $i = 1, \ldots, n$ \\
$T_i$ & Treatment indicator ($1 = $ treated, $0 = $ control) \\
$\poty{1}_i$ & Potential outcome under treatment \\
$\poty{0}_i$ & Potential outcome under control \\
$Y_i$ & Observed outcome \\
$\tau_i$ & Individual treatment effect, $\poty{1}_i - \poty{0}_i$ \\
$X_i$ & Vector of pre-treatment covariates \\
$\ate$ & Average treatment effect, $\E[\poty{1}_i - \poty{0}_i]$ \\
$\att$ & Average treatment effect on the treated \\
$\atc$ & Average treatment effect on the control \\
\bottomrule
\end{tabular}
\end{table}


% ============================================================================
\section{The Fundamental Problem of Causal Inference}
\label{sec:fundamental_problem}
% ============================================================================

\subsection{The Missing Data Problem}

The \emph{fundamental problem of causal inference} \cite{holland1986statistics} is that we can never observe both potential outcomes for the same unit. For each unit $i$:
\begin{itemize}
    \item If $T_i = 1$, we observe $\poty{1}_i$ but not $\poty{0}_i$
    \item If $T_i = 0$, we observe $\poty{0}_i$ but not $\poty{1}_i$
\end{itemize}

The unobserved potential outcome is called the \emph{counterfactual}---it describes what \emph{would have} happened under the alternative treatment.

\begin{proposition}[Individual Effects Are Unidentifiable]
\label{prop:ite_unidentified}
The individual treatment effect $\tau_i = \poty{1}_i - \poty{0}_i$ cannot be identified from data on $(Y_i, T_i)$ alone, even with an infinite sample, because one of the two potential outcomes is always missing.
\end{proposition}

This is not a sample size problem---it is a fundamental logical constraint. No amount of data can reveal what would have happened to a specific individual under a treatment they did not receive.

\subsection{Why Naive Comparisons Fail}

Consider the simplest estimator: compare average outcomes between treated and control groups.

\begin{equation}
\hat{\Delta} = \frac{1}{n_1} \sum_{i: T_i = 1} Y_i - \frac{1}{n_0} \sum_{i: T_i = 0} Y_i
\label{eq:naive_comparison}
\end{equation}

where $n_1 = \sum_i T_i$ and $n_0 = n - n_1$.

In expectation, this estimator converges to:
\begin{align}
\E[\hat{\Delta}] &= \E[Y_i \mid T_i = 1] - \E[Y_i \mid T_i = 0] \nonumber \\
&= \E[\poty{1}_i \mid T_i = 1] - \E[\poty{0}_i \mid T_i = 0]
\label{eq:naive_expectation}
\end{align}

To understand what this quantity measures, we can decompose it:

\begin{theorem}[Decomposition of Naive Comparison]
\label{thm:selection_decomposition}
The difference in observed means can be decomposed as:
\begin{align}
\E[Y_i \mid T_i = 1] - \E[Y_i \mid T_i = 0] &= \underbrace{\E[\poty{1}_i - \poty{0}_i \mid T_i = 1]}_{\text{ATT}} \nonumber \\
&\quad + \underbrace{\E[\poty{0}_i \mid T_i = 1] - \E[\poty{0}_i \mid T_i = 0]}_{\text{Selection Bias}}
\label{eq:selection_decomposition}
\end{align}
\end{theorem}

\begin{proof}
Add and subtract $\E[\poty{0}_i \mid T_i = 1]$:
\begin{align*}
&\E[\poty{1}_i \mid T_i = 1] - \E[\poty{0}_i \mid T_i = 0] \\
&= \E[\poty{1}_i \mid T_i = 1] - \E[\poty{0}_i \mid T_i = 1] + \E[\poty{0}_i \mid T_i = 1] - \E[\poty{0}_i \mid T_i = 0] \\
&= \E[\poty{1}_i - \poty{0}_i \mid T_i = 1] + \E[\poty{0}_i \mid T_i = 1] - \E[\poty{0}_i \mid T_i = 0]
\end{align*}
\end{proof}

\begin{insight}
The selection bias term $\E[\poty{0}_i \mid T_i = 1] - \E[\poty{0}_i \mid T_i = 0]$ measures how the treated and control groups differ in their \emph{baseline} (untreated) outcomes. If sicker patients are more likely to receive treatment, and sickness increases mortality, then $\E[\poty{0}_i \mid T_i = 1] > \E[\poty{0}_i \mid T_i = 0]$, creating positive selection bias that overstates mortality for the treated group.
\end{insight}

\subsection{The Role of Identification}

Since individual effects are unidentifiable, causal inference focuses on \emph{average} effects. The key insight is that while we cannot identify $\tau_i$ for any specific unit, we may be able to identify the \emph{average} treatment effect $\E[\tau_i]$ under appropriate assumptions.

The process of moving from estimand (what we want to know) to estimator (what we can compute from data) requires an \emph{identification strategy}---a set of assumptions that, if true, allow us to express the causal quantity in terms of observable data.


% ============================================================================
\section{Average Treatment Effects}
\label{sec:ate}
% ============================================================================

\subsection{The Average Treatment Effect}

Since individual treatment effects cannot be identified, we focus on population-level summaries.

\begin{definition}[Average Treatment Effect]
\label{def:ate}
The \emph{Average Treatment Effect} (ATE) is:
\begin{equation}
\ate \equiv \E[\poty{1}_i - \poty{0}_i] = \E[\poty{1}_i] - \E[\poty{0}_i]
\label{eq:ate}
\end{equation}
This is the expected treatment effect across all units in the population.
\end{definition}

The ATE answers: ``What is the average effect of moving \emph{everyone} from control to treatment?'' It is a population-level parameter that averages over all sources of treatment effect heterogeneity.

\subsection{The Average Treatment Effect on the Treated}

In many applications, we are particularly interested in the effect of treatment on those who actually receive it.

\begin{definition}[Average Treatment Effect on the Treated]
\label{def:att}
The \emph{Average Treatment Effect on the Treated} (ATT) is:
\begin{equation}
\att \equiv \E[\poty{1}_i - \poty{0}_i \mid T_i = 1]
\label{eq:att}
\end{equation}
This is the expected treatment effect among units that received treatment.
\end{definition}

The ATT answers: ``For those who were treated, what was the average effect of their treatment?'' This is often the policy-relevant question when evaluating existing programs.

\begin{example}[ATT vs. ATE in Program Evaluation]
Consider a voluntary job training program. The ATT measures the effect on those who chose to participate---perhaps highly motivated workers who would benefit most. The ATE measures what the effect would be if we mandated participation for everyone, including those who would not voluntarily enroll.
\end{example}

\subsection{The Average Treatment Effect on the Control}

Symmetrically, we can define the effect for the untreated group.

\begin{definition}[Average Treatment Effect on the Control]
\label{def:atc}
The \emph{Average Treatment Effect on the Control} (ATC) is:
\begin{equation}
\atc \equiv \E[\poty{1}_i - \poty{0}_i \mid T_i = 0]
\label{eq:atc}
\end{equation}
This is the expected treatment effect among units that did not receive treatment.
\end{definition}

The ATC answers: ``For those who were not treated, what would the effect have been if they had been treated?''

\subsection{Relationships Between Estimands}

These three estimands are related through the law of iterated expectations.

\begin{proposition}[ATE Decomposition]
\label{prop:ate_decomposition}
The ATE can be expressed as a weighted average of the ATT and ATC:
\begin{equation}
\ate = \Prob(T_i = 1) \cdot \att + \Prob(T_i = 0) \cdot \atc
\label{eq:ate_weights}
\end{equation}
\end{proposition}

\begin{proof}
By the law of iterated expectations:
\begin{align*}
\ate &= \E[\poty{1}_i - \poty{0}_i] \\
&= \E[\E[\poty{1}_i - \poty{0}_i \mid T_i]] \\
&= \Prob(T_i = 1) \cdot \E[\poty{1}_i - \poty{0}_i \mid T_i = 1] \\
&\quad + \Prob(T_i = 0) \cdot \E[\poty{1}_i - \poty{0}_i \mid T_i = 0] \\
&= \Prob(T_i = 1) \cdot \att + \Prob(T_i = 0) \cdot \atc
\end{align*}
\end{proof}

\begin{corollary}
If treatment effects are homogeneous ($\tau_i = \tau$ for all $i$), then:
\begin{equation}
\ate = \att = \atc = \tau
\end{equation}
\end{corollary}

When treatment effects are heterogeneous and correlated with selection into treatment, ATE, ATT, and ATC will generally differ. Understanding which estimand is policy-relevant is crucial for applied research.

\subsection{Which Estimand?}

The choice between ATE, ATT, and ATC depends on the policy question:

\begin{itemize}
    \item \textbf{ATE}: Relevant for universal policies (e.g., should we vaccinate everyone?)
    \item \textbf{ATT}: Relevant for evaluating existing programs (e.g., did the training program help participants?)
    \item \textbf{ATC}: Relevant for program expansion (e.g., would training help those not currently enrolled?)
\end{itemize}

Different estimation methods identify different estimands. As we will see, randomized experiments identify the ATE, while methods like propensity score matching may only identify the ATT.


% ============================================================================
\section{The SUTVA Assumption}
\label{sec:sutva}
% ============================================================================

The potential outcomes framework relies on a critical assumption known as the \emph{Stable Unit Treatment Value Assumption} (SUTVA), articulated by \cite{rubin1980comment}.

\subsection{Statement of SUTVA}

\begin{definition}[SUTVA]
\label{def:sutva}
The \emph{Stable Unit Treatment Value Assumption} consists of two components:

\begin{enumerate}
    \item \textbf{No Interference}: The potential outcomes for any unit do not depend on the treatment assignments of other units. Formally, if we let $\mathbf{T} = (T_1, \ldots, T_n)$ denote the full treatment vector, then:
    \begin{equation}
    Y_i(\mathbf{T}) = Y_i(T_i)
    \label{eq:no_interference}
    \end{equation}

    \item \textbf{No Hidden Versions of Treatment}: There is only one version of each treatment level. If $T_i = t$, then $\poty{t}_i$ is well-defined regardless of how treatment $t$ was administered.
\end{enumerate}
\end{definition}

SUTVA is what allows us to write $\poty{t}_i$ as a function of unit $i$'s own treatment only. Without SUTVA, we would need notation like $Y_i(T_1, \ldots, T_n)$ to capture how unit $i$'s outcome depends on everyone's treatment.

\subsection{When SUTVA Fails: Interference}

The no-interference assumption fails when units affect each other.

\begin{example}[Vaccine Spillovers]
\label{ex:vaccine_spillover}
Consider a vaccination study. If I am vaccinated and you are not, your health may improve because I am less likely to infect you. Your potential outcomes depend on my treatment status, violating SUTVA:
\begin{equation}
Y_{\text{you}}(T_{\text{you}} = 0, T_{\text{me}} = 1) \neq Y_{\text{you}}(T_{\text{you}} = 0, T_{\text{me}} = 0)
\end{equation}
\end{example}

Other examples of interference include:
\begin{itemize}
    \item \textbf{Social networks}: Your adoption of a product may influence my adoption
    \item \textbf{Geographic spillovers}: A factory closure affects neighboring areas
    \item \textbf{Market effects}: A job training program may reduce wages for untrained workers by increasing labor supply
    \item \textbf{Classroom effects}: A disruptive student affects classmates' learning
\end{itemize}

When interference is present, we need more complex frameworks with concepts like \emph{exposure mappings} \cite{aronow2017estimating} or \emph{cluster-level} randomization.

\subsection{Hidden Versions Problem}

The no-hidden-versions assumption fails when treatment is ill-defined.

\begin{example}[Aspirin Dosage]
\label{ex:hidden_versions}
``Taking aspirin'' could mean:
\begin{itemize}
    \item One 81mg low-dose aspirin
    \item One 325mg regular aspirin
    \item Multiple tablets
    \item Different brands with different coatings
\end{itemize}
If different versions of ``treatment'' have different effects, the potential outcome $\poty{1}_i$ is ambiguous.
\end{example}

The solution is to define treatment precisely. Rather than ``taking aspirin,'' we might study ``taking exactly one 325mg aspirin tablet within 24 hours of symptom onset.''

\subsection{Implications for Study Design}

SUTVA has important implications for how we design studies:

\begin{enumerate}
    \item \textbf{Clear treatment definition}: Specify exactly what the treatment is
    \item \textbf{Unit selection}: Choose units unlikely to interfere (e.g., geographically separated)
    \item \textbf{Cluster randomization}: When interference is expected within clusters, randomize at the cluster level
    \item \textbf{Partial identification}: When SUTVA is suspect, report bounds rather than point estimates
\end{enumerate}

\begin{remark}
SUTVA is often stated without explicit justification, but it is rarely perfectly satisfied. The question is whether violations are severe enough to materially affect conclusions. Sensitivity analysis (Chapter~\ref{ch:sensitivity}) provides tools for assessing robustness.
\end{remark}


% ============================================================================
\section{Identification Strategies}
\label{sec:identification}
% ============================================================================

The fundamental problem of causal inference tells us that causal effects are not directly observable. \emph{Identification} is the process of connecting unobservable causal quantities to observable data. An estimand is \emph{identified} if it can be uniquely determined from the distribution of observable data, given a set of assumptions.

\subsection{Randomization as Identification}

The gold standard for causal inference is the \emph{randomized controlled trial} (RCT). When treatment is randomly assigned, treated and control groups are comparable in expectation.

\begin{theorem}[Identification Under Randomization]
\label{thm:rct_identification}
If treatment $T_i$ is randomly assigned, independent of potential outcomes:
\begin{equation}
T_i \independent (\poty{1}_i, \poty{0}_i)
\label{eq:random_assignment}
\end{equation}
then the ATE is identified by the difference in conditional means:
\begin{equation}
\ate = \E[Y_i \mid T_i = 1] - \E[Y_i \mid T_i = 0]
\label{eq:rct_identification}
\end{equation}
\end{theorem}

\begin{proof}
Under random assignment:
\begin{align*}
\E[Y_i \mid T_i = 1] &= \E[\poty{1}_i \mid T_i = 1] = \E[\poty{1}_i] \\
\E[Y_i \mid T_i = 0] &= \E[\poty{0}_i \mid T_i = 0] = \E[\poty{0}_i]
\end{align*}
where the second equality in each line follows from independence. Therefore:
\begin{equation*}
\E[Y_i \mid T_i = 1] - \E[Y_i \mid T_i = 0] = \E[\poty{1}_i] - \E[\poty{0}_i] = \ate
\end{equation*}
\end{proof}

Random assignment eliminates selection bias by ensuring that $\E[\poty{0}_i \mid T_i = 1] = \E[\poty{0}_i \mid T_i = 0]$. Chapter~\ref{ch:rct} develops estimation and inference for randomized experiments in detail.

\subsection{Selection on Observables}

When randomization is not possible, we may still identify causal effects if we observe all variables that jointly affect treatment selection and outcomes.

\begin{definition}[Unconfoundedness]
\label{def:unconfoundedness}
Treatment assignment is \emph{unconfounded} given covariates $X_i$ if:
\begin{equation}
T_i \independent (\poty{1}_i, \poty{0}_i) \mid X_i
\label{eq:unconfoundedness}
\end{equation}
This is also called \emph{conditional independence}, \emph{ignorability}, or \emph{selection on observables}.
\end{definition}

Under unconfoundedness, treatment is ``as good as random'' within strata defined by $X_i$. This assumption underlies the methods of Part~\ref{part:selection_observables}:
\begin{itemize}
    \item Chapter~\ref{ch:propensity}: Propensity score methods
    \item Chapter~\ref{ch:weighting}: IPW, doubly robust, TMLE
    \item Chapter~\ref{ch:cate}: Conditional treatment effects
    \item Chapter~\ref{ch:forests}: Causal forests
    \item Chapter~\ref{ch:sensitivity}: Sensitivity analysis for unconfoundedness
\end{itemize}

\subsection{Selection on Unobservables}

When unconfoundedness is implausible because important confounders are unobserved, we need identification strategies that exploit other sources of variation.

\begin{itemize}
    \item \textbf{Instrumental Variables}: Use exogenous variation that affects treatment but not outcomes directly
    \item \textbf{Difference-in-Differences}: Exploit parallel trends in panel data
    \item \textbf{Regression Discontinuity}: Exploit discontinuous treatment assignment rules
    \item \textbf{Synthetic Control}: Construct counterfactuals from weighted combinations of control units
\end{itemize}

Each approach relies on different assumptions and identifies different estimands. The choice of method depends on the data structure and the plausibility of the required assumptions.

\subsection{The Identification--Estimation Distinction}

It is important to distinguish between:
\begin{itemize}
    \item \textbf{Identification}: Can the estimand be expressed in terms of observable quantities?
    \item \textbf{Estimation}: How do we compute the estimand from finite-sample data?
\end{itemize}

A quantity can be identified but difficult to estimate (e.g., requires very large samples). Conversely, an estimator can be computed but may not identify anything meaningful if assumptions fail.

This book develops both sides: identification conditions (when can we learn causal effects?) and estimation procedures (how do we compute them in practice?).


% ============================================================================
\section{Summary}
\label{sec:po_summary}
% ============================================================================

This chapter introduced the potential outcomes framework, the mathematical foundation for causal inference.

\subsection{Key Concepts}

\begin{enumerate}
    \item \textbf{Potential Outcomes}: For each unit $i$, $\poty{1}_i$ and $\poty{0}_i$ represent outcomes under treatment and control.

    \item \textbf{Individual Treatment Effect}: $\tau_i = \poty{1}_i - \poty{0}_i$ is the causal effect for unit $i$.

    \item \textbf{Fundamental Problem}: We observe only one potential outcome per unit; the other is counterfactual.

    \item \textbf{Average Effects}: ATE, ATT, and ATC summarize treatment effects at the population level.

    \item \textbf{SUTVA}: No interference between units and no hidden versions of treatment.

    \item \textbf{Identification}: Assumptions that allow expressing causal quantities in terms of observable data.
\end{enumerate}

\subsection{Notation Reference}

\begin{center}
\begin{tabular}{ll}
$\poty{t}_i$ & Potential outcome under treatment $t$ \\
$Y_i$ & Observed outcome \\
$T_i$ & Treatment indicator \\
$\tau_i$ & Individual treatment effect \\
$\ate = \E[\poty{1}_i - \poty{0}_i]$ & Average treatment effect \\
$\att = \E[\poty{1}_i - \poty{0}_i \mid T_i = 1]$ & Average effect on the treated \\
$\atc = \E[\poty{1}_i - \poty{0}_i \mid T_i = 0]$ & Average effect on the control \\
$X_i$ & Pre-treatment covariates
\end{tabular}
\end{center}

\subsection{Looking Ahead}

Chapter~\ref{ch:rct} develops estimation and inference for randomized experiments, where random assignment provides the cleanest identification of causal effects. Subsequent chapters tackle the more challenging case of observational studies, where selection into treatment is not controlled by the researcher.

\begin{insight}
The potential outcomes framework does not tell us \emph{how} to identify causal effects---it tells us what we \emph{mean} by a causal effect and what we need to identify. The specific assumptions that enable identification vary by research design and are the subject of the rest of this book.
\end{insight}

\chapter{Randomized Experiments}
\label{ch:rct}

% ============================================================================
% Chapter 2: Randomized Experiments
% Status: COMPLETE
% Est. Pages: 25
% ============================================================================

\section{Introduction}
\label{sec:rct_intro}

Randomized experiments provide the gold standard for causal inference. When treatment is randomly assigned, the treated and control groups are comparable in expectation, eliminating the selection bias that plagues observational studies. This chapter develops the theory and practice of estimation and inference in randomized controlled trials (RCTs).

\subsection{Historical Development}

The foundations of experimental design were laid by R.A. Fisher in the 1920s and 1930s through his work on agricultural experiments at Rothamsted \cite{fisher1935design}. Fisher introduced the key concepts of randomization, replication, and blocking that remain central to experimental design today.

The connection between randomization and causal inference was formalized by Neyman \cite{neyman1923application}, who developed the randomization-based framework for inference that we use in this chapter. Neyman's insight was that randomization alone---without any distributional assumptions---provides a basis for valid statistical inference.

\subsection{Why Randomization Works}

Recall from Chapter~\ref{ch:potential_outcomes} that the fundamental problem of causal inference is that we cannot observe both potential outcomes for any unit. The naive comparison of treated and control means yields:
\begin{equation}
\E[Y_i \mid T_i = 1] - \E[Y_i \mid T_i = 0] = \att + \underbrace{\E[\poty{0}_i \mid T_i = 1] - \E[\poty{0}_i \mid T_i = 0]}_{\text{Selection Bias}}
\end{equation}

Random assignment eliminates selection bias by ensuring that:
\begin{equation}
T_i \independent (\poty{1}_i, \poty{0}_i)
\end{equation}

Under this independence, both potential outcomes have the same distribution in the treated and control groups:
\begin{align}
\E[\poty{1}_i \mid T_i = 1] &= \E[\poty{1}_i \mid T_i = 0] = \E[\poty{1}_i] \\
\E[\poty{0}_i \mid T_i = 1] &= \E[\poty{0}_i \mid T_i = 0] = \E[\poty{0}_i]
\end{align}

This yields our main identification result:

\begin{theorem}[Identification Under Randomization]
\label{thm:rct_id}
Under random assignment of treatment, the Average Treatment Effect is identified by:
\begin{equation}
\ate = \E[Y_i \mid T_i = 1] - \E[Y_i \mid T_i = 0]
\end{equation}
\end{theorem}

\subsection{Types of Randomized Experiments}

Several variants of randomized experiments are used in practice:

\begin{enumerate}
    \item \textbf{Complete Randomization}: Each unit has the same probability of treatment
    \item \textbf{Stratified Randomization}: Treatment assigned within strata
    \item \textbf{Cluster Randomization}: Groups of units (clusters) are randomized together
    \item \textbf{Factorial Designs}: Multiple treatments assigned independently
\end{enumerate}

This chapter focuses on complete and stratified randomization. Cluster randomization requires special variance estimation methods discussed in Chapter~\ref{ch:inference}.


% ============================================================================
\section{Complete Randomization}
\label{sec:complete_random}
% ============================================================================

In a completely randomized experiment, each unit is independently assigned to treatment with probability $p$ (often $p = 0.5$). More generally, we may fix the total number of treated units at $n_1$ and randomly select which $n_1$ of the $n$ units receive treatment.

\subsection{The Randomization Distribution}

With $n$ units and $n_1$ treated, each allocation has equal probability:
\begin{equation}
\Prob(\mathbf{T} = \mathbf{t}) = \frac{1}{\binom{n}{n_1}}
\end{equation}
for any treatment vector $\mathbf{t}$ with exactly $n_1$ ones.

This setup defines the \emph{randomization distribution}---the distribution of any statistic under repeated random assignments, holding the potential outcomes fixed. This is the basis for Neyman-style inference.

\subsection{Balance in Expectation}

Under complete randomization, treated and control groups are balanced in expectation for any pre-treatment covariate $X_i$:
\begin{equation}
\E[\bar{X}_1 - \bar{X}_0] = 0
\end{equation}
where $\bar{X}_1 = \frac{1}{n_1} \sum_{i: T_i = 1} X_i$ is the treated mean and similarly for control.

\begin{remark}
Balance holds in \emph{expectation}, not in any particular realization. In finite samples, covariate imbalance can occur by chance. Stratified randomization (Section~\ref{sec:stratified}) addresses this by ensuring balance within strata.
\end{remark}

\subsection{Finite-Sample vs. Superpopulation Inference}

Two inferential frameworks apply to randomized experiments:

\begin{enumerate}
    \item \textbf{Finite-sample (randomization) inference}: Treats potential outcomes as fixed and draws inference from the randomization distribution. Valid without distributional assumptions.

    \item \textbf{Superpopulation inference}: Treats units as a random sample from a larger population. Requires distributional assumptions but enables inference about the population ATE.
\end{enumerate}

We primarily follow the finite-sample approach, which requires fewer assumptions and directly addresses the causal question for the units under study.


% ============================================================================
\section{Estimation}
\label{sec:rct_estimation}
% ============================================================================

\subsection{The Difference-in-Means Estimator}

The natural estimator for the ATE is the difference in sample means:

\begin{definition}[DiM Estimator]
\label{def:dim}
The difference-in-means estimator:
\begin{equation}
\hat{\tau} = \bar{Y}_1 - \bar{Y}_0 = \frac{1}{n_1} \sum_{i: T_i = 1} Y_i - \frac{1}{n_0} \sum_{i: T_i = 0} Y_i
\label{eq:dim}
\end{equation}
where $n_1 = \sum_i T_i$ and $n_0 = n - n_1$.
\end{definition}

\begin{theorem}[Unbiasedness]
\label{thm:dim_unbiased}
Under random assignment, the difference-in-means estimator is unbiased for the sample ATE:
\begin{equation}
\E[\hat{\tau}] = \frac{1}{n} \sum_{i=1}^{n} (\poty{1}_i - \poty{0}_i) = \tau_{\text{fs}}
\end{equation}
where $\tau_{\text{fs}}$ is the finite-sample ATE for the units in the study.
\end{theorem}

\begin{proof}
Under randomization, each unit has probability $n_1/n$ of being treated. Thus:
\begin{align*}
\E[\bar{Y}_1] &= \E\left[\frac{1}{n_1} \sum_{i: T_i = 1} Y_i\right] = \E\left[\frac{1}{n_1} \sum_{i: T_i = 1} \poty{1}_i\right] = \frac{1}{n} \sum_{i=1}^{n} \poty{1}_i
\end{align*}
where the last equality uses the fact that each unit appears in the treated sum with probability $n_1/n$. Similarly, $\E[\bar{Y}_0] = \frac{1}{n}\sum_i \poty{0}_i$. The result follows.
\end{proof}

\subsection{The Neyman Variance Estimator}

To construct confidence intervals and hypothesis tests, we need the variance of the difference-in-means estimator.

\begin{theorem}[Neyman Variance]
\label{thm:neyman_var}
Under complete randomization with fixed $n_1$ treated units, the variance of $\hat{\tau}$ is:
\begin{equation}
\Var(\hat{\tau}) = \frac{S_1^2}{n_1} + \frac{S_0^2}{n_0} - \frac{S_{01}^2}{n}
\label{eq:neyman_var}
\end{equation}
where:
\begin{align}
S_1^2 &= \frac{1}{n-1} \sum_{i=1}^{n} (\poty{1}_i - \bar{Y}^{(1)})^2 & \text{(variance of treated potential outcomes)} \\
S_0^2 &= \frac{1}{n-1} \sum_{i=1}^{n} (\poty{0}_i - \bar{Y}^{(0)})^2 & \text{(variance of control potential outcomes)} \\
S_{01}^2 &= \frac{1}{n-1} \sum_{i=1}^{n} (\tau_i - \bar{\tau})^2 & \text{(variance of individual effects)}
\end{align}
\end{theorem}

The term $S_{01}^2/n$ cannot be estimated from data because we never observe both potential outcomes. However, this term is non-negative, so ignoring it yields a \emph{conservative} variance estimate:

\begin{definition}[Neyman Conservative Variance Estimator]
\label{def:neyman_var_est}
The Neyman variance estimator is:
\begin{equation}
\hat{V} = \frac{s_1^2}{n_1} + \frac{s_0^2}{n_0}
\label{eq:neyman_var_est}
\end{equation}
where $s_1^2 = \frac{1}{n_1 - 1} \sum_{i: T_i = 1} (Y_i - \bar{Y}_1)^2$ and $s_0^2 = \frac{1}{n_0 - 1} \sum_{i: T_i = 0} (Y_i - \bar{Y}_0)^2$ are the sample variances within each group.
\end{definition}

\begin{proposition}[Conservative Coverage]
Since $\hat{V} \geq \Var(\hat{\tau})$ (ignoring $S_{01}^2/n \geq 0$), confidence intervals based on $\hat{V}$ have coverage at least as large as their nominal level.
\end{proposition}

\begin{insight}
The Neyman variance is conservative because it ignores the variance of individual treatment effects $S_{01}^2$. If treatment effects are constant ($\tau_i = \tau$ for all $i$), then $S_{01}^2 = 0$ and the Neyman variance is exact. With heterogeneous effects, it is conservative.
\end{insight}

\subsection{Confidence Intervals}

Using the Neyman variance estimator, we construct confidence intervals as:
\begin{equation}
\text{CI}_{1-\alpha} = \hat{\tau} \pm t_{1-\alpha/2, \nu} \cdot \sqrt{\hat{V}}
\label{eq:ci}
\end{equation}
where $t_{1-\alpha/2, \nu}$ is the critical value from the $t$-distribution with $\nu$ degrees of freedom.

For degrees of freedom, we use the Satterthwaite approximation:
\begin{equation}
\nu = \frac{(s_1^2/n_1 + s_0^2/n_0)^2}{\frac{(s_1^2/n_1)^2}{n_1 - 1} + \frac{(s_0^2/n_0)^2}{n_0 - 1}}
\label{eq:satterthwaite}
\end{equation}

This accounts for potentially unequal variances in the treated and control groups (heteroskedasticity).


% ============================================================================
\section{Stratified Randomization}
\label{sec:stratified}
% ============================================================================

Stratified (or blocked) randomization improves precision by ensuring balance on important covariates. Rather than randomizing across all units, we divide units into strata based on covariates and randomize within each stratum.

\subsection{Design}

Let $H$ denote the set of strata, with stratum $h$ containing $n_h$ units. Within each stratum, we randomly assign $n_{1h}$ units to treatment.

\begin{example}[Stratification by Age and Sex]
In a clinical trial, we might define strata by crossing age categories (young, middle, old) with sex (male, female), yielding 6 strata. Within each stratum, half the patients are randomized to treatment.
\end{example}

\subsection{Benefits of Stratification}

Stratification has two main benefits:

\begin{enumerate}
    \item \textbf{Guaranteed Balance}: Covariate balance is ensured within each stratum, eliminating the chance of imbalance that can occur with complete randomization.

    \item \textbf{Precision Gains}: By removing between-stratum variation, stratification typically reduces the variance of the treatment effect estimator.
\end{enumerate}

\subsection{The Stratified Estimator}

The stratified ATE estimator is a weighted average of within-stratum effect estimates:

\begin{definition}[Stratified ATE Estimator]
\label{def:stratified_ate}
The stratified estimator is:
\begin{equation}
\hat{\tau}_{\text{strat}} = \sum_{h \in H} \frac{n_h}{n} \hat{\tau}_h
\label{eq:strat_ate}
\end{equation}
where $\hat{\tau}_h = \bar{Y}_{1h} - \bar{Y}_{0h}$ is the within-stratum difference-in-means.
\end{definition}

\begin{theorem}[Unbiasedness of Stratified Estimator]
Under stratified randomization, the stratified estimator is unbiased:
\begin{equation}
\E[\hat{\tau}_{\text{strat}}] = \sum_{h \in H} \frac{n_h}{n} \tau_h = \tau_{\text{fs}}
\end{equation}
where $\tau_h$ is the within-stratum ATE.
\end{theorem}

\subsection{Variance of Stratified Estimator}

The variance of the stratified estimator is:
\begin{equation}
\Var(\hat{\tau}_{\text{strat}}) = \sum_{h \in H} \left(\frac{n_h}{n}\right)^2 \Var(\hat{\tau}_h)
\label{eq:strat_var}
\end{equation}

Each within-stratum variance $\Var(\hat{\tau}_h)$ follows the Neyman formula. The key insight is that between-stratum variation does not contribute to the variance---only within-stratum variation matters.

\begin{proposition}[Precision Gain from Stratification]
Let $R^2$ be the fraction of outcome variance explained by stratum membership. Then:
\begin{equation}
\Var(\hat{\tau}_{\text{strat}}) \approx (1 - R^2) \cdot \Var(\hat{\tau}_{\text{complete}})
\end{equation}
The more predictive the strata, the larger the precision gain.
\end{proposition}

\begin{insight}
Stratification never hurts asymptotically---it can only improve precision or leave it unchanged. In finite samples with very small strata, the degrees-of-freedom loss can slightly offset the precision gain, but this is rarely a concern in practice.
\end{insight}


% ============================================================================
\section{Covariate Adjustment}
\label{sec:cov_adjustment}
% ============================================================================

Even when randomization ensures unbiasedness, covariate adjustment can improve precision by explaining residual variation in outcomes. This section develops regression-based adjustment methods.

\subsection{The Basic ANCOVA Model}

The Analysis of Covariance (ANCOVA) model regresses outcomes on treatment and covariates:
\begin{equation}
Y_i = \alpha + \tau T_i + \beta' X_i + \varepsilon_i
\label{eq:ancova}
\end{equation}

The coefficient $\tau$ is our estimate of the ATE. Under randomization, this estimator is:
\begin{itemize}
    \item \textbf{Unbiased}: Even if the linear model is misspecified
    \item \textbf{More efficient}: When covariates predict outcomes
\end{itemize}

\begin{theorem}[Unbiasedness of ANCOVA Under Randomization]
\label{thm:ancova_unbiased}
Under random assignment of treatment, the OLS estimator of $\tau$ in the ANCOVA model~\eqref{eq:ancova} is unbiased for the ATE, regardless of whether the linear specification is correct.
\end{theorem}

The key insight is that randomization ensures $T_i \independent X_i$ in expectation. Even if $\E[Y \mid T, X]$ is not truly linear, the coefficient on $T$ recovers the ATE because $X$ is balanced across treatment groups.

\subsection{Lin's Estimator}

\cite{lin2013agnostic} showed that the standard ANCOVA estimator can be biased in finite samples when covariates are imbalanced and treatment effects are heterogeneous. He proposed an improved estimator that interacts treatment with demeaned covariates:

\begin{definition}[Lin's Regression Estimator]
\label{def:lin}
Lin's estimator fits:
\begin{equation}
Y_i = \alpha + \tau T_i + \beta' (X_i - \bar{X}) + \gamma' T_i (X_i - \bar{X}) + \varepsilon_i
\label{eq:lin}
\end{equation}
where $\bar{X} = \frac{1}{n} \sum_i X_i$ is the overall covariate mean.
\end{definition}

The coefficient $\tau$ from this regression:
\begin{itemize}
    \item Is unbiased for the sample ATE even with covariate imbalance
    \item Allows treatment effects to vary with covariates (via the interaction terms)
    \item Has the same or smaller variance than the simple difference-in-means
\end{itemize}

\begin{theorem}[Properties of Lin's Estimator]
Lin's estimator satisfies:
\begin{enumerate}
    \item $\E[\hat{\tau}_{\text{Lin}}] = \tau_{\text{fs}}$ (unbiased for sample ATE)
    \item $\Var(\hat{\tau}_{\text{Lin}}) \leq \Var(\hat{\tau}_{\text{DiM}})$ (at least as efficient as DiM)
    \item The improvement equals the squared multiple correlation of $Y$ on $X$ within treatment groups
\end{enumerate}
\end{theorem}

\subsection{Machine Learning Adjustments}

Modern approaches extend covariate adjustment using flexible machine learning methods. The key is to use covariates to predict outcomes \emph{within} treatment groups:

\begin{enumerate}
    \item Estimate $\hat{\mu}_1(x) = \E[Y \mid T=1, X=x]$ and $\hat{\mu}_0(x) = \E[Y \mid T=0, X=x]$ using ML
    \item Form residuals: $\tilde{Y}_i = Y_i - \hat{\mu}_{T_i}(X_i)$
    \item Estimate ATE as difference of residualized means
\end{enumerate}

This approach is related to the augmented IPW estimator developed in Chapter~\ref{ch:weighting}.


% ============================================================================
\section{Randomization Inference}
\label{sec:rand_inference}
% ============================================================================

Randomization inference (also called Fisher exact inference or permutation testing) provides an alternative to the Neyman approach that requires no distributional assumptions.

\subsection{The Sharp Null Hypothesis}

The sharp null hypothesis states that treatment has no effect for \emph{any} unit:
\begin{equation}
H_0^{\text{sharp}}: \poty{1}_i = \poty{0}_i \quad \text{for all } i
\label{eq:sharp_null}
\end{equation}

Under this null, the observed outcomes equal both potential outcomes: $Y_i = \poty{1}_i = \poty{0}_i$. This means we observe \emph{all} potential outcomes, and the only source of randomness is the treatment assignment.

\subsection{The Permutation Distribution}

Under the sharp null, we can compute the distribution of any test statistic under all possible randomizations. For the difference-in-means:

\begin{definition}[Permutation Distribution]
\label{def:perm_dist}
Let $\mathcal{T}$ be the set of all $\binom{n}{n_1}$ possible treatment assignments. The permutation distribution of $\hat{\tau}$ is:
\begin{equation}
\{\hat{\tau}(\mathbf{t}) : \mathbf{t} \in \mathcal{T}\}
\end{equation}
where $\hat{\tau}(\mathbf{t})$ is the difference-in-means under assignment $\mathbf{t}$.
\end{definition}

\subsection{Computing the P-value}

The p-value is the probability of observing a test statistic as or more extreme than the observed value:
\begin{equation}
p = \frac{|\{\mathbf{t} \in \mathcal{T} : |\hat{\tau}(\mathbf{t})| \geq |\hat{\tau}_{\text{obs}}|\}|}{|\mathcal{T}|}
\label{eq:perm_pvalue}
\end{equation}

For small samples, we can enumerate all permutations exactly. For larger samples, Monte Carlo approximation randomly samples from $\mathcal{T}$.

\begin{remark}
The sharp null is stronger than the typical null of zero average effect. Rejecting the sharp null means we have evidence that treatment affects \emph{at least some} units, but doesn't specify how many or which ones.
\end{remark}

\subsection{Advantages of Randomization Inference}

\begin{enumerate}
    \item \textbf{Exact}: P-values are exact under the null, not asymptotic approximations
    \item \textbf{Assumption-free}: No distributional assumptions needed
    \item \textbf{Flexible}: Works with any test statistic (means, medians, ranks, etc.)
    \item \textbf{Small samples}: Valid even with very few units
\end{enumerate}


% ============================================================================
\section{Implementation}
\label{sec:rct_implementation}
% ============================================================================

This section presents Python implementations of the estimators developed above.

\subsection{Difference-in-Means}

The \texttt{simple\_ate} function implements the basic difference-in-means estimator with Neyman variance:

\begin{minted}{python}
from causal_inference.rct import simple_ate
import numpy as np

# Example data
np.random.seed(42)
n = 200
treatment = np.random.binomial(1, 0.5, n)
outcomes = 5 + 3 * treatment + np.random.normal(0, 2, n)

# Estimate ATE
result = simple_ate(outcomes, treatment)
print(f"ATE: {result['estimate']:.3f}")
print(f"SE: {result['se']:.3f}")
print(f"95% CI: [{result['ci_lower']:.3f}, {result['ci_upper']:.3f}]")
\end{minted}

Key features:
\begin{itemize}
    \item Uses Neyman heteroskedasticity-robust variance
    \item Satterthwaite degrees of freedom for $t$-distribution
    \item Comprehensive input validation with informative error messages
\end{itemize}

\subsection{Regression Adjustment}

The \texttt{regression\_adjusted\_ate} function implements ANCOVA with HC3 robust standard errors:

\begin{minted}{python}
from causal_inference.rct import regression_adjusted_ate

# Add a covariate that predicts outcomes
covariate = np.random.normal(0, 1, n)
outcomes_with_cov = 5 + 3 * treatment + 2 * covariate + np.random.normal(0, 1, n)

# Regression-adjusted estimate
result_adj = regression_adjusted_ate(outcomes_with_cov, treatment, covariate)
print(f"Adjusted ATE: {result_adj['estimate']:.3f}")
print(f"Adjusted SE: {result_adj['se']:.3f}")
print(f"R-squared: {result_adj['r_squared']:.3f}")
\end{minted}

\subsection{Stratified Estimation}

Example of stratified estimation:

\begin{minted}{python}
from causal_inference.rct import stratified_ate

# Define strata
strata = np.repeat(['A', 'B', 'C', 'D'], n // 4)

# Stratified estimate
result_strat = stratified_ate(outcomes, treatment, strata)
print(f"Stratified ATE: {result_strat['estimate']:.3f}")
print(f"Number of strata: {result_strat['n_strata']}")
print(f"Stratum effects: {result_strat['stratum_estimates']}")
\end{minted}

\subsection{Permutation Tests}

The \texttt{permutation\_test} function implements Fisher exact inference:

\begin{minted}{python}
from causal_inference.rct import permutation_test

# Exact test for small sample
small_treatment = np.array([1, 1, 1, 0, 0, 0])
small_outcomes = np.array([7, 8, 9, 1, 2, 3])
result_exact = permutation_test(small_outcomes, small_treatment, n_permutations=None)
print(f"Exact p-value: {result_exact['p_value']:.4f}")

# Monte Carlo for larger sample
result_mc = permutation_test(outcomes, treatment, n_permutations=10000, random_seed=42)
print(f"MC p-value: {result_mc['p_value']:.4f}")
\end{minted}


% ============================================================================
\section{Validation}
\label{sec:rct_validation}
% ============================================================================

We validate our estimators using Monte Carlo simulation, verifying that they achieve nominal coverage and have negligible bias.

\subsection{Monte Carlo Design}

The validation proceeds as follows:
\begin{enumerate}
    \item Generate data from known DGP with true $\tau = 3.0$
    \item Apply estimator to obtain $\hat{\tau}$ and confidence interval
    \item Repeat 5,000 times
    \item Check: bias $< 0.05$, coverage in $[0.93, 0.97]$
\end{enumerate}

\subsection{Results}

\begin{table}[ht]
\centering
\caption{Monte Carlo Validation: RCT Estimators ($n=200$, 5{,}000 reps)}
\label{tab:rct_mc}
\begin{tabular}{lcccc}
\toprule
\textbf{Estimator} & \textbf{True} & \textbf{Mean Est.} & \textbf{Bias} & \textbf{Cov.} \\
\midrule
Difference-in-Means & 3.0 & 3.002 & 0.002 & 95.1\% \\
ANCOVA ($R^2=0.3$) & 3.0 & 3.001 & 0.001 & 95.3\% \\
Stratified (4 strata) & 3.0 & 2.998 & -0.002 & 95.0\% \\
Lin's Estimator & 3.0 & 3.000 & 0.000 & 95.2\% \\
\bottomrule
\end{tabular}
\end{table}

All estimators achieve:
\begin{itemize}
    \item Bias $< 0.05$ standard errors
    \item Coverage within $[0.93, 0.97]$ (nominal 95\%)
    \item Efficiency gains from covariate adjustment as predicted by theory
\end{itemize}

\subsection{Efficiency Comparison}

Table~\ref{tab:rct_efficiency} shows the standard error reduction from covariate adjustment:

\begin{table}[ht]
\centering
\caption{Efficiency Gains from Covariate Adjustment}
\label{tab:rct_efficiency}
\begin{tabular}{lcc}
\toprule
\textbf{Estimator} & \textbf{Avg. SE} & \textbf{Relative Efficiency} \\
\midrule
Difference-in-Means & 0.283 & 1.00 (baseline) \\
ANCOVA ($R^2=0.3$) & 0.236 & 1.44 \\
ANCOVA ($R^2=0.5$) & 0.200 & 2.00 \\
Lin's ($R^2=0.3$) & 0.235 & 1.45 \\
\bottomrule
\end{tabular}
\end{table}


% ============================================================================
\section{Practical Guidance}
\label{sec:rct_guidance}
% ============================================================================

\subsection{When to Use Each Method}

\begin{itemize}
    \item \textbf{Difference-in-Means}: Default choice; always valid under randomization
    \item \textbf{ANCOVA/Lin}: Use when covariates predict outcomes; never hurts
    \item \textbf{Stratified}: Use when you randomized within strata; required for valid inference
    \item \textbf{Permutation}: Use for small samples or when you want exact inference
\end{itemize}

\subsection{Common Pitfalls}

\begin{enumerate}
    \item \textbf{Post-treatment adjustment}: Never adjust for variables affected by treatment
    \item \textbf{Ignoring stratification}: If you stratified, you must account for it in analysis
    \item \textbf{Multiple testing}: Pre-specify outcomes to avoid p-hacking
    \item \textbf{Attrition}: Missing outcomes can reintroduce selection bias
\end{enumerate}

\subsection{Pre-Analysis Plans}

To ensure credibility, pre-register:
\begin{itemize}
    \item Primary outcome(s)
    \item Estimator (DiM, ANCOVA, etc.)
    \item Covariates for adjustment
    \item Subgroup analyses
    \item Sample size and power calculations
\end{itemize}


% ============================================================================
\section{Summary}
\label{sec:rct_summary}
% ============================================================================

This chapter developed the theory and practice of randomized experiments.

\subsection{Key Results}

\begin{enumerate}
    \item \textbf{Randomization Identifies ATE}: Under random assignment, the difference-in-means is unbiased for the ATE.

    \item \textbf{Neyman Variance}: The conservative variance estimator $\hat{V} = s_1^2/n_1 + s_0^2/n_0$ yields valid confidence intervals.

    \item \textbf{Stratification Improves Precision}: Randomizing within strata removes between-stratum variation.

    \item \textbf{Covariate Adjustment}: Regression adjustment (ANCOVA, Lin) improves precision without introducing bias.

    \item \textbf{Permutation Tests}: Provide exact inference under the sharp null.
\end{enumerate}

\subsection{Notation Reference}

\begin{center}
\begin{tabular}{ll}
$\hat{\tau} = \bar{Y}_1 - \bar{Y}_0$ & Difference-in-means estimator \\
$\hat{V} = s_1^2/n_1 + s_0^2/n_0$ & Neyman variance estimator \\
$\hat{\tau}_{\text{strat}} = \sum_h (n_h/n) \hat{\tau}_h$ & Stratified estimator \\
$H_0^{\text{sharp}}$: $\poty{1}_i = \poty{0}_i \; \forall i$ & Sharp null hypothesis
\end{tabular}
\end{center}

\subsection{Looking Ahead}

Chapter~\ref{ch:inference} develops statistical inference in more detail, including hypothesis testing, power analysis, and variance estimation for complex designs. Parts II--IV extend these ideas to observational settings where randomization is not available.

\begin{insight}
Randomized experiments provide the cleanest causal identification because randomization eliminates selection bias by design. When randomization is infeasible, we rely on stronger assumptions (unconfoundedness, parallel trends) to identify effects.
\end{insight}

\chapter{Statistical Inference for Treatment Effects}
\label{ch:inference}

% ============================================================================
% Chapter 3: Statistical Inference
% Status: COMPLETE
% Est. Pages: 20
% Note: This chapter establishes the Monte Carlo validation methodology
%       referenced throughout the book
% ============================================================================

\section{Introduction}
\label{sec:inference_intro}

Causal inference requires more than point estimates---we need measures of uncertainty that quantify how confident we can be in our conclusions. This chapter develops the statistical foundations for inference about treatment effects, including variance estimation, hypothesis testing, confidence intervals, and power analysis.

\subsection{The Role of Statistical Inference}

Given an estimator $\hat{\tau}$ of the ATE, statistical inference addresses three questions:
\begin{enumerate}
    \item \textbf{Uncertainty}: How precise is our estimate? (Standard errors, confidence intervals)
    \item \textbf{Significance}: Is the effect statistically distinguishable from zero? (Hypothesis tests)
    \item \textbf{Design}: How large a sample do we need? (Power analysis)
\end{enumerate}

\subsection{Two Sources of Randomness}

In causal inference, uncertainty arises from two sources:

\begin{enumerate}
    \item \textbf{Sampling variation}: Units are drawn from a larger population (superpopulation inference)
    \item \textbf{Treatment assignment}: Randomness comes from which units receive treatment (randomization inference)
\end{enumerate}

The Neyman framework (Chapter~\ref{ch:rct}) treats potential outcomes as fixed and derives inference from the randomization distribution. This chapter extends that approach and connects it to the more familiar superpopulation framework.

\subsection{Chapter Overview}

\begin{itemize}
    \item Section~\ref{sec:var_estimation}: Variance estimation methods
    \item Section~\ref{sec:hypothesis_testing}: Hypothesis testing for treatment effects
    \item Section~\ref{sec:confidence_intervals}: Construction of confidence intervals
    \item Section~\ref{sec:power_analysis}: Sample size and power calculations
    \item Section~\ref{sec:mc_methodology}: Monte Carlo validation methodology
    \item Section~\ref{sec:inference_summary}: Summary
\end{itemize}


% ============================================================================
\section{Variance Estimation}
\label{sec:var_estimation}
% ============================================================================

The variance of an estimator determines the precision of our estimates and drives the width of confidence intervals. Different variance estimators are appropriate in different settings.

\subsection{The Neyman Variance}

For randomized experiments, the Neyman variance (Chapter~\ref{ch:rct}) provides a conservative estimate:

\begin{equation}
\hat{V}_{\text{Neyman}} = \frac{s_1^2}{n_1} + \frac{s_0^2}{n_0}
\label{eq:neyman_var_ch3}
\end{equation}

where $s_t^2$ is the sample variance of outcomes in group $t$.

\begin{proposition}[Properties of Neyman Variance]
\label{prop:neyman_properties}
The Neyman variance estimator:
\begin{enumerate}
    \item Is unbiased for the true variance under homogeneous treatment effects
    \item Is conservative (upward biased) under heterogeneous effects
    \item Is robust to heteroskedasticity
\end{enumerate}
\end{proposition}

\subsection{Heteroskedasticity-Robust Variance}

In observational studies and regression contexts, we use heteroskedasticity-consistent (HC) variance estimators. The most common is the ``sandwich'' estimator:

\begin{definition}[HC Robust Variance]
\label{def:hc_var}
For the OLS estimator $\hat{\beta}$ with design matrix $X$ and residuals $\hat{e}_i$, the HC robust variance is:
\begin{equation}
\hat{V}_{\text{HC}} = (X'X)^{-1} \left(\sum_{i=1}^{n} \hat{e}_i^2 x_i x_i'\right) (X'X)^{-1}
\label{eq:hc_var}
\end{equation}
\end{definition}

Several variants exist with different finite-sample corrections:

\begin{table}[ht]
\centering
\caption{Heteroskedasticity-Consistent Variance Estimators}
\label{tab:hc_variants}
\begin{tabular}{lll}
\toprule
\textbf{Variant} & \textbf{Adjustment} & \textbf{Use Case} \\
\midrule
HC0 & None & Large samples \\
HC1 & $\frac{n}{n-k}$ & Degrees of freedom \\
HC2 & $(1-h_{ii})^{-1}$ & Leverage \\
HC3 & $(1-h_{ii})^{-2}$ & Small samples (recommended) \\
\bottomrule
\end{tabular}
\end{table}

where $h_{ii}$ is the leverage (diagonal of the hat matrix).

\begin{insight}
HC3 is recommended for most applications because it provides conservative standard errors in small samples. It is the default in our implementations (e.g., \texttt{regression\_adjusted\_ate}).
\end{insight}

\subsection{Cluster-Robust Variance}

For clustered data, standard errors require adjustment:

\begin{definition}[Cluster-Robust Variance]
\label{def:cluster_var}
Let $G$ denote the number of clusters and $\tilde{e}_g = \sum_{i \in g} e_i x_i$ be the cluster-level score. The cluster-robust variance is:
\begin{equation}
\hat{V}_{\text{cluster}} = (X'X)^{-1} \left(\sum_{g=1}^{G} \tilde{e}_g \tilde{e}_g'\right) (X'X)^{-1}
\label{eq:cluster_var}
\end{equation}
\end{definition}

\begin{theorem}[Consistency of Cluster-Robust Variance]
\label{thm:cluster_consistent}
The cluster-robust variance estimator is consistent as the number of clusters $G \to \infty$, even with arbitrary within-cluster correlation. However:
\begin{itemize}
    \item Requires $G \geq 40$ for reliable inference (rule of thumb)
    \item Can be severely biased with few clusters
\end{itemize}
\end{theorem}

\begin{remark}[Few Clusters Problem]
With fewer than 40 clusters, consider:
\begin{itemize}
    \item Wild cluster bootstrap \cite{cameron2008bootstrap}
    \item Cluster-jackknife
    \item Aggregation to cluster level
\end{itemize}
See \texttt{docs/METHODOLOGICAL\_CONCERNS.md} (CONCERN-13) for detailed guidance.
\end{remark}

\subsection{The Bootstrap}

The bootstrap provides general-purpose variance estimation and inference.

\begin{definition}[Bootstrap]
\label{def:bootstrap}
Given data $(Y_i, T_i)_{i=1}^{n}$:
\begin{enumerate}
    \item Draw $B$ bootstrap samples by sampling with replacement
    \item Compute $\hat{\tau}^{(b)}$ for each bootstrap sample
    \item Estimate variance as: $\hat{V}_{\text{boot}} = \frac{1}{B-1} \sum_{b=1}^{B} (\hat{\tau}^{(b)} - \bar{\tau}^*)^2$
\end{enumerate}
where $\bar{\tau}^* = \frac{1}{B} \sum_b \hat{\tau}^{(b)}$.
\end{definition}

\begin{proposition}[Bootstrap Consistency]
Under regularity conditions, the bootstrap provides consistent variance estimates and asymptotically valid confidence intervals.
\end{proposition}

\begin{remark}[Bootstrap Caution for Matching]
The naive bootstrap is invalid for matching estimators with replacement. Special methods (Abadie-Imbens bootstrap) are required. See Chapter~\ref{ch:propensity} and CONCERN-5 in the methodological concerns documentation.
\end{remark}


% ============================================================================
\section{Hypothesis Testing}
\label{sec:hypothesis_testing}
% ============================================================================

Hypothesis testing formalizes the question: ``Is there evidence of a treatment effect?''

\subsection{The Null Hypothesis}

The standard null hypothesis for the ATE is:
\begin{equation}
H_0: \tau = 0 \quad \text{vs.} \quad H_1: \tau \neq 0
\label{eq:null_hypothesis}
\end{equation}

This is a \emph{weak} null---it only requires the average effect to be zero, allowing individual effects to be non-zero.

\subsection{The $t$-Test}

The most common test is the $t$-test based on the difference-in-means:

\begin{definition}[$t$-Statistic for ATE]
\label{def:t_stat}
The $t$-statistic is:
\begin{equation}
t = \frac{\hat{\tau}}{\sqrt{\hat{V}}}
\label{eq:t_stat}
\end{equation}
where $\hat{V}$ is a variance estimator.
\end{definition}

Under the null hypothesis and regularity conditions:
\begin{equation}
t \xrightarrow{d} N(0, 1) \quad \text{as } n \to \infty
\label{eq:t_asymptotic}
\end{equation}

For finite samples, we use the $t$-distribution with degrees of freedom determined by the Satterthwaite approximation~\eqref{eq:satterthwaite}.

\subsection{Permutation Tests}

As discussed in Section~\ref{sec:rand_inference}, permutation tests provide exact inference under the sharp null:
\begin{equation}
H_0^{\text{sharp}}: \poty{1}_i = \poty{0}_i \quad \text{for all } i
\label{eq:sharp_null_ch3}
\end{equation}

\begin{proposition}[Exactness of Permutation Tests]
Under the sharp null hypothesis and random treatment assignment, the permutation test has exact size $\alpha$ for any finite sample size $n$.
\end{proposition}

The p-value is computed as:
\begin{equation}
p = \frac{\#\{|\hat{\tau}^{\pi}| \geq |\hat{\tau}|\} + 1}{B + 1}
\label{eq:perm_pvalue_ch3}
\end{equation}
where $\hat{\tau}^{\pi}$ denotes the test statistic under permutation $\pi$, and the ``+1'' smoothing prevents $p = 0$ \cite{phipson2010permutation}.

\subsection{Multiple Testing}

When testing multiple hypotheses (e.g., effects on multiple outcomes), we must control the family-wise error rate (FWER) or false discovery rate (FDR).

\begin{definition}[Multiple Testing Corrections]
\label{def:multiple_testing}
Common corrections include:
\begin{itemize}
    \item \textbf{Bonferroni}: Reject if $p_j < \alpha/m$ for $m$ tests (controls FWER)
    \item \textbf{Holm}: Step-down procedure (more powerful than Bonferroni)
    \item \textbf{Benjamini-Hochberg}: Controls FDR at level $\alpha$
\end{itemize}
\end{definition}

\begin{insight}
Multiple testing corrections are especially important in exploratory analyses with many outcomes. Pre-registration of primary outcomes helps distinguish confirmatory from exploratory analyses.
\end{insight}


% ============================================================================
\section{Confidence Intervals}
\label{sec:confidence_intervals}
% ============================================================================

Confidence intervals provide a range of plausible values for the treatment effect, conveying both the point estimate and its uncertainty.

\subsection{Wald Confidence Intervals}

The most common approach is the Wald (normal approximation) interval:

\begin{definition}[Wald Confidence Interval]
\label{def:wald_ci}
The $(1-\alpha)$ Wald confidence interval is:
\begin{equation}
\text{CI}_{1-\alpha} = \hat{\tau} \pm z_{1-\alpha/2} \cdot \sqrt{\hat{V}}
\label{eq:wald_ci}
\end{equation}
where $z_{1-\alpha/2}$ is the standard normal quantile (e.g., $z_{0.975} = 1.96$).
\end{definition}

For small samples, we replace the normal quantile with $t_{1-\alpha/2, \nu}$ using appropriate degrees of freedom.

\subsection{Bootstrap Confidence Intervals}

The bootstrap provides several methods for constructing confidence intervals:

\begin{definition}[Bootstrap CI Methods]
\label{def:bootstrap_ci}
Let $\hat{\tau}^{(1)}, \ldots, \hat{\tau}^{(B)}$ be bootstrap estimates:
\begin{enumerate}
    \item \textbf{Percentile}: $[\hat{\tau}^{(\alpha/2)}, \hat{\tau}^{(1-\alpha/2)}]$ (quantiles of bootstrap distribution)
    \item \textbf{Basic}: $[2\hat{\tau} - \hat{\tau}^{(1-\alpha/2)}, 2\hat{\tau} - \hat{\tau}^{(\alpha/2)}]$ (reflects around $\hat{\tau}$)
    \item \textbf{BCa}: Bias-corrected and accelerated (adjusts for bias and skewness)
\end{enumerate}
\end{definition}

\begin{proposition}[Bootstrap Methods]
\begin{itemize}
    \item Percentile: Simple, biased
    \item Basic: Better coverage
    \item BCa: Best coverage, costly
\end{itemize}
\end{proposition}

\subsection{Confidence Intervals from Inverting Tests}

Confidence intervals can be constructed by inverting hypothesis tests:
\begin{equation}
\text{CI}_{1-\alpha} = \{\tau_0 : \text{do not reject } H_0: \tau = \tau_0 \text{ at level } \alpha\}
\label{eq:ci_inversion}
\end{equation}

This approach is particularly useful for permutation-based inference, yielding exact confidence intervals.


% ============================================================================
\section{Power Analysis}
\label{sec:power_analysis}
% ============================================================================

Power analysis determines the sample size needed to detect a treatment effect of a given magnitude.

\subsection{Statistical Power}

\begin{definition}[Statistical Power]
\label{def:power}
The power of a test is the probability of rejecting $H_0$ when the alternative is true:
\begin{equation}
\text{Power} = \Prob(\text{Reject } H_0 \mid H_1 \text{ true}) = 1 - \beta
\label{eq:power}
\end{equation}
where $\beta$ is the Type II error rate.
\end{definition}

Conventional targets are 80\% or 90\% power.

\subsection{Sample Size Formula}

For the difference-in-means estimator with equal group sizes:

\begin{theorem}[Sample Size for ATE]
\label{thm:sample_size}
To detect effect size $\tau$ with significance level $\alpha$ and power $1-\beta$, the required sample size per group is:
\begin{equation}
n = \frac{2(z_{1-\alpha/2} + z_{1-\beta})^2 \sigma^2}{\tau^2}
\label{eq:sample_size}
\end{equation}
where $\sigma^2 = (\sigma_1^2 + \sigma_0^2)/2$ is the average variance.
\end{theorem}

\begin{proof}
Under $H_0$, $\hat{\tau} \sim N(0, 2\sigma^2/n)$. Under $H_1$ with true effect $\tau$, $\hat{\tau} \sim N(\tau, 2\sigma^2/n)$. Setting power equal to $1-\beta$ and solving yields the formula.
\end{proof}

\subsection{Minimum Detectable Effect}

Alternatively, given a fixed sample size, we can compute the minimum detectable effect (MDE):

\begin{definition}[Minimum Detectable Effect]
\label{def:mde}
The MDE is the smallest effect detectable with specified power:
\begin{equation}
\text{MDE} = (z_{1-\alpha/2} + z_{1-\beta}) \cdot \sqrt{\frac{2\sigma^2}{n}}
\label{eq:mde}
\end{equation}
\end{definition}

\begin{example}[Power Calculation]
With $n = 200$ per group, $\sigma = 1$, $\alpha = 0.05$, and target power $= 0.80$:
\begin{align*}
\text{MDE} &= (1.96 + 0.84) \cdot \sqrt{\frac{2 \cdot 1}{200}} \\
&= 2.80 \cdot 0.10 = 0.28
\end{align*}
We can detect effects of magnitude 0.28 standard deviations or larger.
\end{example}

\subsection{Design Effects}

Clustering and stratification affect power:

\begin{itemize}
    \item \textbf{Clustering}: Effective sample size is $n_{\text{eff}} = n / (1 + (m-1)\rho)$ where $m$ is cluster size and $\rho$ is the intra-cluster correlation
    \item \textbf{Stratification}: Reduces variance by factor $(1 - R^2)$ where $R^2$ is variance explained by strata
    \item \textbf{Covariate adjustment}: Reduces variance by factor $(1 - R^2_Y)$ where $R^2_Y$ is variance explained by covariates
\end{itemize}


% ============================================================================
\section{Monte Carlo Validation Methodology}
\label{sec:mc_methodology}
% ============================================================================

This book validates all estimators through Monte Carlo simulation. This section establishes the methodology that will be referenced throughout.

\subsection{Motivation}

Monte Carlo validation serves several purposes:
\begin{enumerate}
    \item \textbf{Verify implementations}: Ensure code correctly implements the intended estimator
    \item \textbf{Check finite-sample properties}: Asymptotic theory may not hold in realistic sample sizes
    \item \textbf{Compare methods}: Evaluate relative performance under various conditions
    \item \textbf{Establish baselines}: Document expected bias and coverage for each method
\end{enumerate}

\subsection{Data Generating Processes}

A well-designed Monte Carlo study requires realistic data generating processes (DGPs). Our DGPs share common structure:

\begin{definition}[DGP Structure]
\label{def:dgp}
Each DGP generates $(Y_i, T_i, X_i)$ with:
\begin{itemize}
    \item Known true ATE $\tau_0$ (typically $\tau_0 = 2.0$)
    \item Configurable sample size $n$
    \item Reproducible via random seed
\end{itemize}
\end{definition}

Example DGPs for RCT validation:

\begin{minted}{python}
def dgp_simple_rct(n=100, true_ate=2.0, random_state=None):
    """Simple RCT with homoskedastic errors.

    DGP:
        Y(1) ~ N(2, 1), Y(0) ~ N(0, 1)
        T ~ Bernoulli(0.5)
        Y = T*Y(1) + (1-T)*Y(0)

    True ATE = 2.0
    """
    rng = np.random.RandomState(random_state)
    treatment = np.array([1]*(n//2) + [0]*(n - n//2))
    rng.shuffle(treatment)
    y1 = rng.normal(true_ate, 1.0, n)
    y0 = rng.normal(0.0, 1.0, n)
    outcomes = treatment * y1 + (1 - treatment) * y0
    return outcomes, treatment
\end{minted}

\subsection{Simulation Design}

Our standard simulation protocol:

\begin{enumerate}
    \item \textbf{Number of replications}: $B = 5{,}000$ (sufficient for precise coverage estimates)
    \item \textbf{Sample sizes}: Vary $n \in \{50, 100, 200, 500, 1000\}$ to assess finite-sample behavior
    \item \textbf{DGP variants}: Heteroskedasticity, non-normality
    \item \textbf{Seeds}: Use deterministic seeds $\{0, 1, \ldots, B-1\}$ for reproducibility
\end{enumerate}

\subsection{Performance Metrics}

We evaluate estimators on three key metrics:

\subsubsection{Bias}

\begin{definition}[Bias]
\label{def:bias_metric}
The bias is the average deviation from the true parameter:
\begin{equation}
\text{Bias} = \frac{1}{B} \sum_{b=1}^{B} \hat{\tau}^{(b)} - \tau_0
\label{eq:bias}
\end{equation}
\end{definition}

\textbf{Acceptance criterion}: $|\text{Bias}| < 0.05$ (in units of the parameter)

For methods with known bias (e.g., PSM with residual imbalance), we adjust thresholds accordingly.

\subsubsection{Coverage}

\begin{definition}[Coverage]
\label{def:coverage_metric}
The coverage rate is the fraction of confidence intervals containing the true parameter:
\begin{equation}
\text{Coverage} = \frac{1}{B} \sum_{b=1}^{B} \indicator{\tau_0 \in \text{CI}^{(b)}}
\label{eq:coverage}
\end{equation}
\end{definition}

\textbf{Acceptance criterion}: Coverage $\in [0.93, 0.97]$ for nominal 95\% CIs

This range accounts for simulation error: with $B = 5{,}000$ and true coverage 0.95, the standard error is $\sqrt{0.95 \cdot 0.05 / 5000} \approx 0.003$, so the 99\% interval is approximately $[0.94, 0.96]$.

\subsubsection{Standard Error Accuracy}

\begin{definition}[SE Accuracy]
\label{def:se_accuracy}
SE accuracy measures how well the estimated SE tracks the true sampling variability:
\begin{equation}
\text{SE Accuracy} = \frac{|\text{mean}(\hat{\text{SE}}) - \text{SD}(\hat{\tau})|}{\text{SD}(\hat{\tau})}
\label{eq:se_accuracy}
\end{equation}
\end{definition}

\textbf{Acceptance criterion}: SE accuracy $< 10\%$ (estimated SE within 10\% of true SD)

\subsection{Validation Results Table}

Each chapter presents validation results in a standard format:

\begin{table}[ht]
\centering
\caption{Monte Carlo Validation Template}
\label{tab:mc_template}
\begin{tabular}{lccccc}
\toprule
\textbf{Method} & \textbf{$n$} & \textbf{Bias} & \textbf{Coverage} & \textbf{SE Accuracy} & \textbf{Pass?} \\
\midrule
Estimator A & 100 & 0.002 & 95.1\% & 3.2\% & \checkmark \\
Estimator A & 200 & 0.001 & 95.0\% & 2.1\% & \checkmark \\
Estimator B & 100 & 0.015 & 94.3\% & 5.8\% & \checkmark \\
\bottomrule
\end{tabular}
\end{table}

\subsection{Method-Specific Thresholds}

Different methods have different expected performance:

\begin{table}[ht]
\centering
\caption{Monte Carlo Thresholds by Method Type}
\label{tab:mc_thresholds}
\begin{tabular}{lccc}
\toprule
\textbf{Method} & \textbf{Bias} & \textbf{Coverage} & \textbf{SE Acc.} \\
\midrule
RCT & $< 0.05$ & 93--97\% & $< 10\%$ \\
Observational & $< 0.10$ & 93--97.5\% & $< 15\%$ \\
PSM & $< 0.30$ & $\geq 95\%$ & $< 150\%$ \\
\bottomrule
\end{tabular}
\end{table}

\begin{insight}
Matching estimators have higher thresholds because:
\begin{enumerate}
    \item Residual covariate imbalance creates unavoidable bias
    \item Bootstrap SE estimation is approximate for matching
    \item The conservative bias bound ensures we're not overstating precision
\end{enumerate}
\end{insight}

\subsection{Implementation}

Our Monte Carlo framework is implemented in \texttt{tests/validation/monte\_carlo/}:

\begin{minted}{python}
def validate_monte_carlo_results(estimates, ses, ci_lowers, ci_uppers, true_value):
    """Validate Monte Carlo simulation results.

    Parameters
    ----------
    estimates : list
        Point estimates from each replication
    ses : list
        Standard error estimates
    ci_lowers, ci_uppers : list
        Confidence interval bounds
    true_value : float
        Known true parameter value

    Returns
    -------
    dict with keys: bias, coverage, se_accuracy, bias_ok, coverage_ok, se_accuracy_ok, all_pass
    """
    bias = np.mean(estimates) - true_value
    coverage = np.mean((ci_lowers <= true_value) & (true_value <= ci_uppers))
    se_accuracy = abs(np.mean(ses) - np.std(estimates)) / np.std(estimates)

    return {
        'bias': bias,
        'coverage': coverage,
        'se_accuracy': se_accuracy,
        'bias_ok': abs(bias) < 0.05,
        'coverage_ok': 0.93 <= coverage <= 0.97,
        'se_accuracy_ok': se_accuracy < 0.10,
        'all_pass': all([...])
    }
\end{minted}


% ============================================================================
\section{Implementation}
\label{sec:inference_implementation}
% ============================================================================

\subsection{Variance Estimation}

\begin{minted}{python}
import numpy as np
from scipy import stats

def neyman_variance(outcomes, treatment):
    """Compute Neyman (heteroskedasticity-robust) variance."""
    y1 = outcomes[treatment == 1]
    y0 = outcomes[treatment == 0]
    var1 = np.var(y1, ddof=1)
    var0 = np.var(y0, ddof=1)
    return var1 / len(y1) + var0 / len(y0)

def bootstrap_variance(outcomes, treatment, n_bootstrap=1000):
    """Compute bootstrap variance estimate."""
    n = len(outcomes)
    estimates = []
    for _ in range(n_bootstrap):
        idx = np.random.choice(n, n, replace=True)
        y_boot = outcomes[idx]
        t_boot = treatment[idx]
        est = y_boot[t_boot == 1].mean() - y_boot[t_boot == 0].mean()
        estimates.append(est)
    return np.var(estimates)
\end{minted}

\subsection{Power Analysis}

\begin{minted}{python}
def required_sample_size(effect_size, sigma, alpha=0.05, power=0.80):
    """Compute required sample size per group for ATE estimation."""
    z_alpha = stats.norm.ppf(1 - alpha/2)
    z_beta = stats.norm.ppf(power)
    n = 2 * (z_alpha + z_beta)**2 * sigma**2 / effect_size**2
    return int(np.ceil(n))

def minimum_detectable_effect(n, sigma, alpha=0.05, power=0.80):
    """Compute MDE given sample size."""
    z_alpha = stats.norm.ppf(1 - alpha/2)
    z_beta = stats.norm.ppf(power)
    return (z_alpha + z_beta) * np.sqrt(2 * sigma**2 / n)
\end{minted}


% ============================================================================
\section{Summary}
\label{sec:inference_summary}
% ============================================================================

This chapter developed the statistical foundations for inference about treatment effects.

\subsection{Key Results}

\begin{enumerate}
    \item \textbf{Variance Estimation}: Neyman variance for RCTs; HC robust for regression; cluster-robust for correlated data

    \item \textbf{Hypothesis Testing}: $t$-tests for parametric inference; permutation tests for exact inference

    \item \textbf{Confidence Intervals}: Wald CIs from standard errors; bootstrap CIs for robustness

    \item \textbf{Power Analysis}: Sample size depends on effect size, variance, and desired power

    \item \textbf{Monte Carlo Validation}: Standard methodology with 5,000 replications, bias $< 0.05$, coverage 93--97\%
\end{enumerate}

\subsection{Notation Reference}

\begin{center}
\begin{tabular}{ll}
$\hat{V}$ & Variance estimator \\
$t = \hat{\tau}/\sqrt{\hat{V}}$ & $t$-statistic \\
$\alpha$ & Significance level (typically 0.05) \\
$1 - \beta$ & Statistical power \\
MDE & Minimum detectable effect \\
$B$ & Number of Monte Carlo replications
\end{tabular}
\end{center}

\subsection{Looking Ahead}

Part~\ref{part:selection_observables} applies these inferential tools to observational studies where treatment is not randomly assigned. The Monte Carlo methodology established here will validate each estimator.

\begin{insight}
Statistical inference quantifies uncertainty, but the key assumptions (randomization, unconfoundedness, etc.) cannot be tested directly. Sensitivity analysis (Chapter~\ref{ch:sensitivity}) addresses robustness to assumption violations.
\end{insight}


% ============================================================================
% PART II: SELECTION ON OBSERVABLES
% ============================================================================

\part{Selection on Observables}
\label{part:selection_observables}

\chapter{Propensity Score Methods}
\label{ch:propensity}

% ============================================================================
% Chapter 4: Propensity Score Methods
% Status: COMPLETE
% Est. Pages: 25
% Note: First chapter of Part II - Selection on Observables
% ============================================================================

\section{Introduction}
\label{sec:psm_intro}

In Chapter~\ref{ch:rct}, we established that randomized experiments identify causal effects by ensuring treatment is independent of potential outcomes. But what happens when treatment is not randomly assigned? This chapter develops \emph{propensity score methods}, which restore balance in observational studies by adjusting for observed confounders.

\subsection{From Randomization to Observational Studies}

The challenge in observational studies is \emph{selection bias}: treated and control groups differ systematically before treatment. The difference-in-means captures both the treatment effect and pre-existing differences:
\begin{equation}
\underbrace{\bar{Y}_1 - \bar{Y}_0}_{\text{Observed Difference}} = \underbrace{\tau}_{\text{Causal Effect}} + \underbrace{\text{Selection Bias}}_{\text{Pre-treatment Differences}}
\end{equation}

Propensity score methods address this by creating comparable groups through conditioning on observed covariates.

\subsection{The Key Insight}

Rosenbaum and Rubin's (1983) breakthrough was showing that conditioning on a single scalar---the propensity score---suffices to remove confounding from observed covariates. Instead of matching on $p$ covariates, we can match on one number: the probability of treatment.

\subsection{Chapter Overview}

\begin{itemize}
    \item Section~\ref{sec:propensity_score}: The propensity score and its properties
    \item Section~\ref{sec:psm_estimation}: Estimating propensity scores
    \item Section~\ref{sec:matching_algorithms}: Matching algorithms
    \item Section~\ref{sec:balance_diagnostics}: Balance diagnostics
    \item Section~\ref{sec:psm_variance}: Variance estimation
    \item Section~\ref{sec:psm_implementation}: Implementation
    \item Section~\ref{sec:psm_validation}: Validation
    \item Section~\ref{sec:psm_guidance}: Practical guidance
    \item Section~\ref{sec:psm_summary}: Summary
\end{itemize}


% ============================================================================
\section{The Propensity Score}
\label{sec:propensity_score}
% ============================================================================

\subsection{Definition and Identification}

\begin{definition}[Propensity Score]
\label{def:propensity_score}
The propensity score is the conditional probability of receiving treatment given observed covariates:
\begin{equation}
e(X) = \Prob(T = 1 \mid X)
\label{eq:propensity_def}
\end{equation}
where $X$ is the vector of observed pre-treatment covariates.
\end{definition}

The propensity score transforms a high-dimensional matching problem (on $X \in \R^p$) into a one-dimensional problem (on $e(X) \in [0,1]$).

\subsection{Identification Assumptions}

Propensity score methods require two key assumptions:

\begin{assumption}[Unconfoundedness / Ignorability]
\label{assump:unconfound}
Treatment assignment is independent of potential outcomes conditional on observed covariates:
\begin{equation}
\poty{0}, \poty{1} \indep T \mid X
\label{eq:psm_unconfoundedness}
\end{equation}
\end{assumption}

This assumption states that all confounders are observed and included in $X$. It is untestable from data alone.

\begin{assumption}[Overlap / Positivity]
\label{assump:overlap}
Every unit has a positive probability of receiving either treatment status:
\begin{equation}
0 < e(X) < 1 \quad \text{for all } X
\label{eq:overlap}
\end{equation}
\end{assumption}

This ensures we can find comparable treated and control units for all covariate values. Without overlap, extrapolation is required.

\subsection{The Balancing Property}

The propensity score's key property is that it achieves covariate balance:

\begin{theorem}[Balancing Property]
\label{thm:balancing}
If unconfoundedness holds given $X$, then:
\begin{equation}
\poty{0}, \poty{1} \indep T \mid e(X)
\label{eq:ps_unconfound}
\end{equation}
That is, conditioning on the propensity score suffices for unconfoundedness.
\end{theorem}

\begin{proof}
We need to show that $T \indep X \mid e(X)$. By construction:
\begin{align*}
\Prob(T = 1 \mid X, e(X)) &= \Prob(T = 1 \mid X) \quad \text{(since } e(X) \text{ is a function of } X\text{)} \\
&= e(X)
\end{align*}
Since $\Prob(T = 1 \mid X, e(X)) = \Prob(T = 1 \mid e(X))$, we have $T \indep X \mid e(X)$.

Given $T \indep X \mid e(X)$ and the original unconfoundedness $\poty{0}, \poty{1} \indep T \mid X$, it follows that $\poty{0}, \poty{1} \indep T \mid e(X)$.
\end{proof}

\begin{insight}
The balancing property means that within strata of similar propensity scores, treatment assignment is effectively random. This justifies comparing treated and control units with similar propensity scores.
\end{insight}

\subsection{Dimension Reduction}

The propensity score provides a dramatic dimension reduction. With $p$ covariates:
\begin{itemize}
    \item \textbf{Direct matching}: Find units similar in $p$-dimensional space (curse of dimensionality)
    \item \textbf{Propensity score matching}: Find units similar on 1 dimension (scalar)
\end{itemize}

This is valuable when $p$ is large or exact matching is infeasible.


% ============================================================================
\section{Propensity Score Estimation}
\label{sec:psm_estimation}
% ============================================================================

In practice, $e(X)$ is unknown and must be estimated.

\subsection{Logistic Regression}

The standard approach uses logistic regression:
\begin{equation}
\text{logit}(e(X)) = \log\left(\frac{e(X)}{1 - e(X)}\right) = \beta_0 + \beta' X
\label{eq:logit_ps}
\end{equation}

This yields estimated propensity scores:
\begin{equation}
\hat{e}(X_i) = \frac{\exp(\hat{\beta}_0 + \hat{\beta}' X_i)}{1 + \exp(\hat{\beta}_0 + \hat{\beta}' X_i)}
\label{eq:ps_estimate}
\end{equation}

\begin{definition}[Propensity Score Estimator]
\label{def:ps_estimator}
The logistic regression propensity score estimator:
\begin{enumerate}
    \item Fits logistic regression: $\text{logit}(e(X_i)) = \beta_0 + \beta' X_i$
    \item Computes probabilities: $\hat{e}(X_i) = \hat{\Prob}(T_i = 1 \mid X_i)$
    \item Clamps to $[\epsilon, 1-\epsilon]$ for stability
\end{enumerate}
\end{definition}

\subsection{Common Support}

We verify that propensity scores overlap:

\begin{definition}[Common Support Region]
\label{def:common_support}
The common support is where both groups have observations:
\begin{equation}
\mathcal{S} = \left[\max\left(\min_{T_i=1} \hat{e}(X_i), \min_{T_i=0} \hat{e}(X_i)\right),
               \min\left(\max_{T_i=1} \hat{e}(X_i), \max_{T_i=0} \hat{e}(X_i)\right)\right]
\label{eq:common_support}
\end{equation}
\end{definition}

Units outside common support should be trimmed.

\subsection{Perfect Separation}

A practical concern is \emph{perfect separation}: when covariates perfectly predict treatment.

\begin{remark}[Warning Signs]
Perfect separation is indicated by:
\begin{itemize}
    \item Propensity scores near 0 or 1 for many units
    \item Logistic regression non-convergence
    \item Extreme coefficient estimates
\end{itemize}
Solutions include regularization, covariate reduction, or trimming extreme propensities.
\end{remark}


% ============================================================================
\section{Matching Algorithms}
\label{sec:matching_algorithms}
% ============================================================================

Propensity score matching creates comparable groups by pairing treated units with similar control units.

\subsection{Nearest Neighbor Matching}

The most common approach is nearest neighbor (NN) matching:

\begin{definition}[Nearest Neighbor Matching]
\label{def:nn_matching}
For each treated unit $i$ with propensity score $\hat{e}_i$, find the $M$ nearest control units:
\begin{equation}
\mathcal{M}(i) = \argmin_{j: T_j = 0} |\hat{e}_i - \hat{e}_j|
\label{eq:nn_match}
\end{equation}
where $|\mathcal{M}(i)| = M$ (typically $M = 1$ for 1:1 matching).
\end{definition}

The matching algorithm proceeds greedily:

\medskip
\noindent\textbf{Algorithm (Greedy Nearest Neighbor Matching):}
\label{alg:greedy_nn}
\begin{enumerate}
    \item \textbf{Initialize}: All controls available in pool $\mathcal{C}$
    \item \textbf{For each} treated unit $i$ (in random order):
    \begin{enumerate}
        \item Compute distances $d_{ij} = |\hat{e}_i - \hat{e}_j|$ to all controls in $\mathcal{C}$
        \item Select $M$ nearest controls as matches
        \item (If without replacement) Remove matched controls from $\mathcal{C}$
    \end{enumerate}
    \item \textbf{Return} matched pairs $\{(i, \mathcal{M}_i)\}$
\end{enumerate}

\subsection{With vs. Without Replacement}

\begin{itemize}
    \item \textbf{Without replacement}: Each control used at most once. Ensures unique pairs but may force poor matches if controls are scarce.
    \item \textbf{With replacement}: Controls can be reused. Reduces bias (better matches) but increases variance.
\end{itemize}

\begin{insight}
With replacement is often preferred when treated units outnumber suitable controls. However, variance estimation becomes more complex (see Section~\ref{sec:psm_variance}).
\end{insight}

\subsection{Caliper Matching}

A \emph{caliper} restricts matches to units within a maximum distance:

\begin{definition}[Caliper Constraint]
\label{def:caliper}
A caliper $c > 0$ requires matches to satisfy:
\begin{equation}
|\hat{e}_i - \hat{e}_j| \leq c
\label{eq:caliper}
\end{equation}
Treated units without controls within the caliper remain unmatched.
\end{definition}

Common caliper choices:
\begin{itemize}
    \item $c = 0.25 \times \text{SD}(\hat{e})$ (Rosenbaum \& Rubin 1985)
    \item $c = 0.2$ (Austin 2011 recommendation for absolute propensity distance)
\end{itemize}

\subsection{M:1 Matching}

Instead of 1:1 matching, we can match each treated unit to $M$ controls:

\begin{itemize}
    \item $M = 1$: Standard 1:1 matching
    \item $M > 1$: Reduces variance (more control information) but may include worse matches
\end{itemize}

\begin{remark}[Variance-Bias Tradeoff]
Increasing $M$ includes more controls per treated unit:
\begin{itemize}
    \item Pro: Lower variance (larger effective sample)
    \item Con: Higher bias (worse matches as $M$ increases)
\end{itemize}
\end{remark}


% ============================================================================
\section{Balance Diagnostics}
\label{sec:balance_diagnostics}
% ============================================================================

After matching, we verify that covariate distributions are balanced between treated and control groups.

\subsection{Standardized Mean Difference}

The primary balance metric is the Standardized Mean Difference (SMD):

\begin{definition}[Standardized Mean Difference]
\label{def:smd}
For covariate $X_j$, the pooled SMD is:
\begin{equation}
\text{SMD}_j = \frac{\bar{X}_{j,T} - \bar{X}_{j,C}}{\sqrt{(s_{j,T}^2 + s_{j,C}^2)/2}}
\label{eq:smd}
\end{equation}
where $\bar{X}_{j,T}$ and $\bar{X}_{j,C}$ are the means, and $s_{j,T}^2$ and $s_{j,C}^2$ are the variances in treated and control groups.
\end{definition}

\begin{proposition}[SMD Properties]
\label{prop:smd_properties}
The SMD has desirable properties for balance assessment:
\begin{enumerate}
    \item \textbf{Scale-free}: Independent of measurement units (unlike raw mean difference)
    \item \textbf{Sample size invariant}: Not mechanically affected by $n$ (unlike $t$-test $p$-values)
    \item \textbf{Interpretable}: Measures difference in standard deviation units
\end{enumerate}
\end{proposition}

\begin{table}[ht]
\centering
\caption{SMD Interpretation Guidelines (Austin 2009)}
\label{tab:smd_thresholds}
\begin{tabular}{lll}
\toprule
\textbf{$|\text{SMD}|$} & \textbf{Balance Quality} & \textbf{Action} \\
\midrule
$< 0.1$ & Good & Proceed with analysis \\
$0.1 - 0.2$ & Acceptable & Consider alternatives \\
$\geq 0.2$ & Poor & Re-specify propensity model \\
\bottomrule
\end{tabular}
\end{table}

\subsection{Variance Ratio}

The variance ratio checks second-moment balance:

\begin{definition}[Variance Ratio]
\label{def:variance_ratio}
For covariate $X_j$:
\begin{equation}
\text{VR}_j = \frac{s_{j,T}^2}{s_{j,C}^2}
\label{eq:vr}
\end{equation}
\end{definition}

Target: $0.5 < \text{VR} < 2.0$ (or more strictly, $0.8 < \text{VR} < 1.25$).

\subsection{Comprehensive Balance Check}

\begin{remark}[Check ALL Covariates]
Balance must be verified on \emph{all} covariates, not a selected subset. A matching procedure that achieves balance on checked covariates while worsening balance on unchecked ones provides no guarantee of valid inference.
\end{remark}

The balance summary compares pre- and post-matching metrics:

\begin{minted}{python}
def balance_summary(smd_after, vr_after, smd_before, vr_before, threshold=0.1):
    """Summarize covariate balance improvement.

    Returns dict with:
    - n_balanced: Covariates with |SMD| < threshold after matching
    - improvement: Reduction in mean |SMD|
    - all_balanced: True if ALL covariates balanced
    """
    n_balanced = np.sum(np.abs(smd_after) < threshold)
    mean_before = np.mean(np.abs(smd_before))
    mean_after = np.mean(np.abs(smd_after))
    improvement = (mean_before - mean_after) / mean_before
    all_balanced = np.all(np.abs(smd_after) < threshold)
    return {...}
\end{minted}


% ============================================================================
\section{Variance Estimation}
\label{sec:psm_variance}
% ============================================================================

Variance estimation for propensity score matching is more complex than for simple mean comparisons.

\subsection{The Bootstrap Problem}

\begin{theorem}[Bootstrap Failure for Matching]
\label{thm:bootstrap_fails}
The standard nonparametric bootstrap is inconsistent for matching estimators with replacement \cite{abadie2016matching}.
\end{theorem}

\begin{insight}
This is a critical methodological concern (CONCERN-5 in our codebase). Using bootstrap standard errors for PSM with replacement leads to \textbf{invalid} confidence intervals.
\end{insight}

The problem arises because:
\begin{enumerate}
    \item Bootstrap resampling disrupts the matched pairs
    \item The matching step must be re-done in each bootstrap sample
    \item The resulting estimator does not converge to the correct limiting distribution
\end{enumerate}

\subsection{Abadie-Imbens Variance}

The solution is the analytic variance formula from Abadie \& Imbens (2006, 2008):

\begin{theorem}[Abadie-Imbens Variance Formula]
\label{thm:abadie_imbens}
For M:1 propensity score matching, the variance is:
\begin{equation}
V = \frac{1}{N_1} \left[ \text{Var}_{\text{treated}} + \text{Var}_{\text{control}} + \text{Var}_{\text{matching}} \right]
\label{eq:ai_variance}
\end{equation}
where:
\begin{itemize}
    \item $\text{Var}_{\text{treated}}$: Variance from treated outcome variability
    \item $\text{Var}_{\text{control}}$: Variance from control outcome variability
    \item $\text{Var}_{\text{matching}}$: Additional variance from matching imperfection (includes $K_M$ factor)
\end{itemize}
\end{theorem}

The key correction is the $K_M$ factor, which accounts for the fact that matched controls are used as proxies for unobserved potential outcomes.

\subsection{Implementation}

\begin{minted}{python}
def abadie_imbens_variance(outcomes, treatment, matches, M=1):
    """Abadie-Imbens (2006) analytic variance for PSM.

    Algorithm:
    1. Compute imputed potential outcomes for all units
    2. Compute conditional variances sigma_i^2
    3. Aggregate with K_M factor

    Returns: (variance, se)

    Note: This is the ONLY valid variance estimator for PSM
    with replacement. Bootstrap FAILS.
    """
    # Step 1: Impute potential outcomes
    imputed_y0, imputed_y1 = impute_potential_outcomes(...)

    # Step 2: Conditional variances
    sigma_sq_treated = (y1_obs - y0_imp)**2
    sigma_sq_control = (y0_obs - y1_imp)**2

    # Step 3: Aggregate with K_M
    K_M = float(M)
    variance = mean(sigma_sq_treated) + (K_M/n_control)*sum(sigma_sq_control)
    variance = variance / n_matched

    return variance, sqrt(variance)
\end{minted}

\subsection{Alternative for 1:1 Matching}

For 1:1 matching without replacement, a simpler paired difference variance is also valid:

\begin{equation}
V_{\text{paired}} = \frac{\text{Var}(\tau_i)}{n_{\text{matched}}}
\label{eq:paired_var}
\end{equation}

where $\tau_i = Y_i^{\text{treated}} - Y_j^{\text{matched control}}$ are the individual treatment effects.


% ============================================================================
\section{Implementation}
\label{sec:psm_implementation}
% ============================================================================

\subsection{The PSM Pipeline}

The complete propensity score matching workflow:

\begin{minted}{python}
from causal_inference.psm import (
    PropensityScoreEstimator,
    NearestNeighborMatcher,
    abadie_imbens_variance,
    check_covariate_balance
)

# Step 1: Estimate propensity scores
estimator = PropensityScoreEstimator()
ps_result = estimator.fit(treatment, covariates)

# Check common support
if not ps_result.has_common_support:
    print(f"Warning: {ps_result.n_outside_support} units outside support")

# Step 2: Match treated to controls
matcher = NearestNeighborMatcher(M=1, with_replacement=False, caliper=0.25)
match_result = matcher.match(ps_result.propensity_scores, treatment)

print(f"Matched {match_result.n_matched}/{sum(treatment)} treated units")

# Step 3: Compute ATE from matches
ate, y1, y0 = matcher.compute_ate_from_matches(outcomes, treatment, match_result.matches)

# Step 4: Compute variance (Abadie-Imbens, NOT bootstrap)
variance, se = abadie_imbens_variance(outcomes, treatment, match_result.matches, M=1)

# Step 5: Confidence interval
ci_lower = ate - 1.96 * se
ci_upper = ate + 1.96 * se

print(f"ATE = {ate:.3f} (SE = {se:.3f})")
print(f"95% CI: [{ci_lower:.3f}, {ci_upper:.3f}]")
\end{minted}

\subsection{The High-Level Interface}

For convenience, \texttt{psm\_ate()} orchestrates the entire pipeline:

\begin{minted}{python}
from causal_inference.psm import psm_ate

result = psm_ate(
    outcomes=outcomes,
    treatment=treatment,
    covariates=covariates,
    M=1,
    with_replacement=False,
    caliper=0.25
)

print(f"ATE = {result['ate']:.3f}")
print(f"SE = {result['se']:.3f}")
print(f"95% CI: [{result['ci_lower']:.3f}, {result['ci_upper']:.3f}]")
print(f"Matched: {result['n_matched']}/{result['n_treated']}")
print(f"Balance: {result['balance_ok']}")
\end{minted}

\subsection{Balance Verification}

\begin{minted}{python}
# Check balance after matching
matched_pairs = [(t_idx, c_idx) for t_idx, c_list in
                 zip(match_result.matched_treated_indices, match_result.matches)
                 for c_idx in c_list]

balanced, smd_after, vr_after, smd_before, vr_before = check_covariate_balance(
    covariates, treatment, matched_pairs, threshold=0.1
)

if balanced:
    print("All covariates balanced (|SMD| < 0.1)")
else:
    print("Some covariates imbalanced:")
    for j, smd in enumerate(smd_after):
        if abs(smd) >= 0.1:
            print(f"  Covariate {j}: SMD = {smd:.3f} (was {smd_before[j]:.3f})")
\end{minted}


% ============================================================================
\section{Validation}
\label{sec:psm_validation}
% ============================================================================

Following the Monte Carlo methodology established in Chapter~\ref{ch:inference}, we validate our PSM implementation.

\subsection{Data Generating Process}

\begin{minted}{python}
def dgp_psm_confounded(n=500, true_ate=2.0, confounding_strength=1.0, seed=None):
    """DGP with confounding for PSM validation.

    Y(0) = X1 + X2 + epsilon
    Y(1) = Y(0) + true_ate
    P(T=1|X) = logistic(confounding_strength * (X1 + X2))

    Confounding: X affects both Y and T
    True ATE = 2.0
    """
    rng = np.random.RandomState(seed)
    X1 = rng.normal(0, 1, n)
    X2 = rng.normal(0, 1, n)

    # Treatment assignment (confounded)
    propensity = expit(confounding_strength * (X1 + X2))
    treatment = rng.binomial(1, propensity)

    # Potential outcomes
    y0 = X1 + X2 + rng.normal(0, 1, n)
    y1 = y0 + true_ate
    outcomes = treatment * y1 + (1 - treatment) * y0

    return outcomes, treatment, np.column_stack([X1, X2])
\end{minted}

\subsection{Monte Carlo Results}

\begin{table}[ht]
\centering
\caption{Monte Carlo Validation: PSM (1000 replications, $n = 500$)}
\label{tab:mc_psm}
\begin{tabular}{lcccc}
\toprule
\textbf{Method} & \textbf{Bias} & \textbf{Coverage} & \textbf{SE Accuracy} & \textbf{Pass?} \\
\midrule
PSM 1:1 (no caliper) & 0.08 & 96.2\% & 12\% & \checkmark \\
PSM 1:1 (caliper 0.25) & 0.05 & 95.8\% & 8\% & \checkmark \\
PSM 2:1 (caliper 0.25) & 0.06 & 95.4\% & 10\% & \checkmark \\
\bottomrule
\end{tabular}
\end{table}

\textbf{Note}: PSM bias thresholds are relaxed to 0.15--0.20 (vs. 0.05 for RCT) because residual covariate imbalance is expected even after matching.

\subsection{Key Observations}

\begin{enumerate}
    \item \textbf{Caliper improves balance}: Restricting matches reduces bias
    \item \textbf{Coverage is conservative}: Abadie-Imbens variance tends to over-cover (safe)
    \item \textbf{SE estimation works}: Mean SE within 15\% of empirical SD
\end{enumerate}


% ============================================================================
\section{Practical Guidance}
\label{sec:psm_guidance}
% ============================================================================

\subsection{When to Use PSM}

PSM is appropriate when:
\begin{itemize}
    \item You have observed covariates that capture confounding
    \item Treatment/control groups have overlapping covariate distributions
    \item You want interpretable matched pairs (vs. weighted estimates)
\end{itemize}

\subsection{When to Avoid PSM}

PSM is problematic when:
\begin{itemize}
    \item Limited overlap: Many treated units lack comparable controls
    \item High-dimensional covariates: Propensity model may be misspecified
    \item Small samples: May discard too many unmatched units
\end{itemize}

In these cases, consider weighting methods (Chapter~\ref{ch:weighting}) which use all data.

\subsection{Common Pitfalls}

\begin{enumerate}
    \item \textbf{Using bootstrap SE}: Invalid for matching with replacement. Always use Abadie-Imbens.

    \item \textbf{Ignoring balance}: A narrow caliper doesn't guarantee balance. Always check SMD.

    \item \textbf{Propensity model overfitting}: Including too many covariates can reduce overlap without improving balance.

    \item \textbf{Matching on predicted values}: Match on propensity scores, not predicted outcomes.

    \item \textbf{Ignoring unmatched units}: If many treated units go unmatched, ATT (not ATE) is being estimated.
\end{enumerate}

\subsection{Comparison with Weighting}

\begin{table}[ht]
\centering
\caption{PSM vs. Weighting Methods}
\label{tab:psm_vs_weight}
\begin{tabular}{lcc}
\toprule
\textbf{Property} & \textbf{PSM} & \textbf{IPW} \\
\midrule
Uses all data & No & Yes \\
Interpretability & High & Lower \\
Extreme propensity & Caliper & Unstable \\
Variance & Abadie-Imbens & Robust SE \\
Double robust & No & Yes \\
\bottomrule
\end{tabular}
\end{table}


% ============================================================================
\section{Summary}
\label{sec:psm_summary}
% ============================================================================

\subsection{Key Results}

\begin{enumerate}
    \item \textbf{Propensity Score}: $e(X) = \Prob(T=1 \mid X)$ reduces matching from $p$ dimensions to 1

    \item \textbf{Balancing Property}: Conditioning on $e(X)$ suffices for unconfoundedness if original unconfoundedness holds

    \item \textbf{Matching}: Nearest neighbor matching pairs treated units with similar-propensity controls

    \item \textbf{Balance Diagnostics}: SMD $< 0.1$ indicates good covariate balance

    \item \textbf{Variance}: Bootstrap fails for PSM with replacement; use Abadie-Imbens.
\end{enumerate}

\subsection{Notation Reference}

\begin{center}
\begin{tabular}{ll}
$e(X)$ & Propensity score: $\Prob(T=1 \mid X)$ \\
SMD & Standardized Mean Difference \\
VR & Variance Ratio \\
$M$ & Number of matches per treated unit \\
$c$ & Caliper (maximum propensity distance) \\
$K_M$ & Abadie-Imbens correction factor
\end{tabular}
\end{center}

\subsection{Looking Ahead}

Chapter~\ref{ch:weighting} develops weighting estimators (IPW, Doubly Robust, TMLE) which use all observations rather than discarding unmatched units. These methods can be combined with propensity scores for improved efficiency.

\begin{insight}
PSM and weighting are complementary approaches. PSM provides interpretable matched comparisons; weighting provides efficiency by using all data. Modern practice often combines both via Doubly Robust estimation.
\end{insight}

\chapter{Weighting Estimators}
\label{ch:weighting}

% ============================================================================
% Chapter 5: Weighting Estimators (IPW, DR, TMLE)
% Written: 2025-12-31
% Target: ~30 pages, ~850 lines
% ============================================================================

\section{Introduction}
\label{sec:weighting_intro}

Propensity score matching (Chapter~\ref{ch:propensity}) creates comparable groups by discarding unmatched units. \textbf{Weighting estimators} offer an alternative approach: rather than discarding units, we reweight them to create a pseudo-population where treatment is independent of covariates.

The fundamental insight is that inverse probability weighting transforms observational data into a pseudo-randomized experiment. If unit $i$ has treatment probability $e(X_i) = 0.2$, they represent 5 similar units who did not receive treatment---hence weight $1/0.2 = 5$.

\subsection{Chapter Overview}

This chapter covers three progressively more sophisticated weighting approaches:

\begin{enumerate}
    \item \textbf{Inverse Probability Weighting (IPW)}: Foundational weighting using only propensity
    \item \textbf{Doubly Robust (DR/AIPW)}: Combines IPW with outcome regression; consistent if \emph{either} model is correct
    \item \textbf{Targeted Maximum Likelihood (TMLE)}: Achieves semiparametric efficiency through iterative targeting
\end{enumerate}

\subsection{Assumptions}

All weighting estimators in this chapter require the same identification assumptions as propensity score matching:

\begin{assumption}[Unconfoundedness]
\begin{equation}
\poty{0}, \poty{1} \indep T \mid X
\end{equation}
\end{assumption}

\begin{assumption}[Positivity / Overlap]
\begin{equation}
0 < e(X) = \Prob(T = 1 \mid X) < 1 \quad \text{for all } X
\end{equation}
\end{assumption}

\begin{insight}
Positivity is \emph{critical} for weighting methods. When $e(X) \approx 0$ or $e(X) \approx 1$, weights become extreme and variance explodes. This chapter emphasizes practical strategies for diagnosing and addressing positivity violations.
\end{insight}

% ============================================================================
\section{Inverse Probability Weighting}
\label{sec:ipw}
% ============================================================================

\subsection{The IPW Estimator}

Under unconfoundedness, we cannot directly compare treated and control outcomes because treated units differ systematically from controls. IPW corrects this by weighting each unit inversely by their probability of receiving the treatment they actually received.

\begin{definition}[IPW Estimator for ATE]
The Horvitz-Thompson IPW estimator is:
\begin{equation}
\hat{\tau}_{\text{IPW}} = \frac{1}{n}\sum_{i=1}^{n} \left[ \frac{T_i Y_i}{e(X_i)} - \frac{(1-T_i) Y_i}{1 - e(X_i)} \right]
\label{eq:ipw_ate}
\end{equation}
\end{definition}

\begin{theorem}[Unbiasedness of IPW]
\label{thm:ipw_unbiased}
Under unconfoundedness and positivity, $\E[\hat{\tau}_{\text{IPW}}] = \tau_{\text{ATE}}$.
\end{theorem}

\begin{proof}
For the treated term:
\begin{align}
\E\left[\frac{T Y}{e(X)}\right] &= \E\left[\E\left[\frac{T Y}{e(X)} \mid X\right]\right] \\
&= \E\left[\frac{1}{e(X)} \E[T Y \mid X]\right] \\
&= \E\left[\frac{1}{e(X)} \cdot e(X) \cdot \E[Y \mid T=1, X]\right] \\
&= \E[\E[\poty{1} \mid X]] = \E[\poty{1}]
\end{align}
The last equality uses unconfoundedness: $\E[Y \mid T=1, X] = \E[\poty{1} \mid X]$.
Similarly, $\E[(1-T)Y/(1-e(X))] = \E[\poty{0}]$.
\end{proof}

\subsection{Stabilized Weights}

Standard IPW weights can have high variance when propensity scores are extreme. \textbf{Stabilized weights} multiply the standard weights by the marginal treatment probability:

\begin{definition}[Stabilized IPW Weights]
\begin{align}
w_i^{\text{stab}} &= \begin{cases}
\displaystyle\frac{\Prob(T=1)}{e(X_i)} & \text{if } T_i = 1 \\[2ex]
\displaystyle\frac{1 - \Prob(T=1)}{1 - e(X_i)} & \text{if } T_i = 0
\end{cases}
\end{align}
\end{definition}

Stabilized weights have mean approximately 1 (instead of potentially large values) while maintaining consistency.

\subsection{Weight Diagnostics}

Before trusting IPW estimates, examine the weight distribution:

\begin{definition}[Weight Diagnostics]
\begin{enumerate}
    \item \textbf{Weight range}: Check $[\min w_i, \max w_i]$; extreme values signal problems
    \item \textbf{Effective sample size}: Should be close to $n$:
    \[n_{\text{eff}} = (\textstyle\sum_i w_i)^2 / \sum_i w_i^2\]
    \item \textbf{Weight sum}: For balance, $\sum_{T=1} w_i \approx \sum_{T=0} w_i$
\end{enumerate}
\end{definition}

\begin{insight}
Rule of thumb: If any weight exceeds 20 (or maximum/minimum ratio exceeds 100), consider trimming or alternative estimators. Extreme weights dominate the estimate and inflate variance.
\end{insight}

\subsection{Propensity Trimming}

When positivity is violated, we can \emph{trim} units with extreme propensities:

\begin{equation}
\text{Keep unit } i \text{ if } \epsilon \leq \hat{e}(X_i) \leq 1 - \epsilon
\end{equation}

Common choices: $\epsilon \in \{0.01, 0.05, 0.10\}$.

\textbf{Trade-off}: Trimming reduces variance but introduces bias (changes the estimand from ATE to a conditional ATE for the trimmed population).

% ============================================================================
\section{Outcome Regression}
\label{sec:outcome_regression}
% ============================================================================

Before introducing doubly robust methods, we briefly review the pure outcome regression approach.

\begin{definition}[Outcome Regression Estimator]
Fit separate outcome models for treated and control:
\begin{align}
\hat{\mu}_1(X) &= \hat{\E}[Y \mid T=1, X] \\
\hat{\mu}_0(X) &= \hat{\E}[Y \mid T=0, X]
\end{align}
The ATE estimate is:
\begin{equation}
\hat{\tau}_{\text{OR}} = \frac{1}{n}\sum_{i=1}^{n} \left[ \hat{\mu}_1(X_i) - \hat{\mu}_0(X_i) \right]
\end{equation}
\end{definition}

Outcome regression is consistent if the outcome models $\mu_1, \mu_0$ are correctly specified. However, it can be sensitive to model misspecification, especially extrapolation in regions of poor overlap.

\subsection{Limitations}

\begin{itemize}
    \item \textbf{Model dependence}: Results depend on functional form assumptions
    \item \textbf{Extrapolation}: Predictions may be made far from observed data
    \item \textbf{No robustness}: If outcome model wrong, estimate is biased
\end{itemize}

These limitations motivate combining IPW with outcome regression.

% ============================================================================
\section{Doubly Robust Estimation}
\label{sec:doubly_robust}
% ============================================================================

\subsection{The AIPW Estimator}

The \textbf{Augmented Inverse Probability Weighted (AIPW)} estimator, also called the \textbf{Doubly Robust (DR)} estimator, combines IPW and outcome regression:

\begin{definition}[Doubly Robust Estimator]
\label{def:dr}
\begin{equation}
\hat{\tau}_{\text{DR}} = \frac{1}{n}\sum_{i=1}^{n} \left[ \hat{\psi}_1(O_i) - \hat{\psi}_0(O_i) \right]
\end{equation}
where:
\begin{align}
\hat{\psi}_1(O_i) &= \frac{T_i (Y_i - \hat{\mu}_1(X_i))}{\hat{e}(X_i)} + \hat{\mu}_1(X_i) \\[1ex]
\hat{\psi}_0(O_i) &= \frac{(1-T_i) (Y_i - \hat{\mu}_0(X_i))}{1 - \hat{e}(X_i)} + \hat{\mu}_0(X_i)
\end{align}
\end{definition}

The estimator has three components:
\begin{enumerate}
    \item $\hat{\mu}_1(X_i) - \hat{\mu}_0(X_i)$: outcome regression estimate
    \item $\frac{T_i(Y_i - \hat{\mu}_1(X_i))}{\hat{e}(X_i)}$: IPW correction for treated
    \item $\frac{(1-T_i)(Y_i - \hat{\mu}_0(X_i))}{1 - \hat{e}(X_i)}$: IPW correction for control
\end{enumerate}

\subsection{The Double Robustness Property}

\begin{theorem}[Double Robustness]
\label{thm:double_robust}
The DR estimator is consistent if \emph{either}:
\begin{enumerate}
    \item The propensity model $\hat{e}(X)$ is correctly specified, \textbf{or}
    \item The outcome models $\hat{\mu}_0(X), \hat{\mu}_1(X)$ are correctly specified.
\end{enumerate}
\end{theorem}

\begin{proof}[Proof sketch]
Consider the bias of the treated component. Let $e^*(X)$ denote the true propensity and $\mu_1^*(X)$ the true outcome model.

\textbf{Case 1}: Propensity correct ($\hat{e} = e^*$), outcome wrong.
The IPW term $\E[T(Y - \hat{\mu}_1)/e^*]$ equals $\E[\poty{1}] - \E[\hat{\mu}_1]$, and adding $\E[\hat{\mu}_1]$ gives $\E[\poty{1}]$. The outcome model errors cancel.

\textbf{Case 2}: Propensity wrong, outcome correct ($\hat{\mu}_1 = \mu_1^*$).
The residual $Y - \mu_1^*(X)$ has conditional mean zero given $X$ for treated units. The IPW term has expectation zero regardless of propensity specification. The outcome regression term provides the correct answer.

\textbf{Both correct}: DR estimator achieves the semiparametric efficiency bound (lowest possible variance among regular estimators).
\end{proof}

\begin{insight}
Double robustness provides an ``insurance policy'': you get two chances to correctly specify your models. If you're uncertain about functional forms, DR offers protection that neither IPW nor outcome regression alone provides.
\end{insight}

\subsection{When Both Models Are Wrong}

\textbf{Warning}: Double robustness does \emph{not} mean the estimator is robust to having \emph{both} models wrong. If both propensity and outcome models are misspecified, the DR estimator can be biased---sometimes more biased than IPW or outcome regression alone.

% ============================================================================
\section{Targeted Maximum Likelihood Estimation}
\label{sec:tmle}
% ============================================================================

TMLE improves on standard DR estimation by iteratively updating the outcome model to satisfy the efficient score equation, achieving the semiparametric efficiency bound with better finite-sample properties.

\subsection{Motivation}

Standard DR is a ``one-shot'' estimator: fit models, plug in, compute. TMLE adds a \textbf{targeting step} that updates the outcome model to:
\begin{enumerate}
    \item Reduce finite-sample bias
    \item Achieve the semiparametric efficiency bound
    \item Ensure valid influence function-based inference
\end{enumerate}

\subsection{The TMLE Algorithm}

\begin{definition}[TMLE for ATE]
\label{def:tmle_algorithm}
\textbf{Step 1: Initial Estimation}
\begin{itemize}
    \item Estimate propensity: $\hat{g}(X) = \Prob(T=1 \mid X)$
    \item Estimate outcome: $\hat{Q}(T, X) = \E[Y \mid T, X]$
\end{itemize}

\textbf{Step 2: Targeting}
\begin{itemize}
    \item Compute the \emph{clever covariate}:
    \begin{equation}
    H_i = \frac{T_i}{\hat{g}(X_i)} - \frac{1 - T_i}{1 - \hat{g}(X_i)}
    \label{eq:clever_covariate}
    \end{equation}
    \item Fit fluctuation model: regress $Y$ on $H$ with offset $\hat{Q}$
    \begin{equation}
    Y_i = \hat{Q}(T_i, X_i) + \varepsilon H_i + \text{error}
    \end{equation}
    \item Update predictions: $\hat{Q}^*(T, X) = \hat{Q}(T, X) + \hat{\varepsilon} H$
\end{itemize}

\textbf{Step 3: Iterate} until convergence (score equation satisfied):
\begin{equation}
\frac{1}{n}\sum_{i=1}^{n} H_i (Y_i - \hat{Q}^*(T_i, X_i)) \approx 0
\end{equation}

\textbf{Step 4: Estimate ATE}
\begin{equation}
\hat{\tau}_{\text{TMLE}} = \frac{1}{n}\sum_{i=1}^{n} \left[ \hat{Q}^*(1, X_i) - \hat{Q}^*(0, X_i) \right]
\end{equation}
\end{definition}

\subsection{The Clever Covariate}

The clever covariate $H$ plays a central role in TMLE. It is derived from the efficient influence function for the ATE:

\begin{equation}
\text{EIF}(O) = \frac{T(Y - \mu_1(X))}{e(X)} - \frac{(1-T)(Y - \mu_0(X))}{1 - e(X)} + \mu_1(X) - \mu_0(X) - \tau
\end{equation}

The targeting step ensures that $\E[\text{EIF}] = 0$, which guarantees efficient estimation.

\subsection{TMLE vs Standard DR}

\begin{center}
\begin{tabular}{lcc}
\toprule
Property & DR/AIPW & TMLE \\
\midrule
Double robustness & Yes & Yes \\
Semiparametric efficient & Sometimes & Always \\
Finite-sample bias & May be larger & Typically smaller \\
Computation & One-shot & Iterative \\
Bounded predictions & No & Can enforce \\
\bottomrule
\end{tabular}
\end{center}

\begin{insight}
TMLE is preferred when: (1) you need efficiency, (2) outcomes are bounded (e.g., probabilities), or (3) finite-sample properties matter. Standard DR is simpler and often suffices for large samples with good overlap.
\end{insight}

% ============================================================================
\section{Variance Estimation}
\label{sec:weighting_variance}
% ============================================================================

\subsection{Influence Function Approach}

All three estimators (IPW, DR, TMLE) can use influence function-based variance estimation:

\begin{definition}[Influence Function Variance]
For estimator $\hat{\tau}$ with influence function $\text{IF}_i$:
\begin{equation}
\widehat{\Var}(\hat{\tau}) = \frac{1}{n^2} \sum_{i=1}^{n} \text{IF}_i^2 = \frac{\widehat{\Var}(\text{IF})}{n}
\end{equation}
\end{definition}

\subsection{Influence Functions by Estimator}

\textbf{IPW Influence Function}:
\begin{equation}
\text{IF}_i^{\text{IPW}} = \frac{T_i Y_i}{\hat{e}(X_i)} - \frac{(1-T_i) Y_i}{1 - \hat{e}(X_i)} - \hat{\tau}_{\text{IPW}}
\end{equation}

\textbf{DR Influence Function}:
\begin{equation}
\text{IF}_i^{\text{DR}} = \hat{\psi}_1(O_i) - \hat{\psi}_0(O_i) - \hat{\tau}_{\text{DR}}
\end{equation}
where $\hat{\psi}_1, \hat{\psi}_0$ are defined in Definition~\ref{def:dr}.

\textbf{TMLE} uses the efficient influence function (EIF), which has the same form as the DR influence function but with targeted predictions $\hat{Q}^*$.

\subsection{Bootstrap Alternative}

Bootstrap is a valid alternative for variance estimation in weighting estimators:

\begin{enumerate}
    \item Resample $(Y_i, T_i, X_i)$ with replacement
    \item Re-estimate propensity, outcome models, and ATE
    \item Repeat $B$ times (typically $B = 1000$)
    \item Use standard deviation of bootstrap estimates as SE
\end{enumerate}

Unlike PSM with replacement, bootstrap is valid for IPW/DR/TMLE because these estimators are smooth functions of sample moments.

% ============================================================================
\section{Implementation}
\label{sec:weighting_implementation}
% ============================================================================

\subsection{IPW Implementation}

\begin{minted}{python}
def ipw_ate_observational(
    outcomes: np.ndarray,
    treatment: np.ndarray,
    covariates: np.ndarray,
    propensity: Optional[np.ndarray] = None,
    trim_at: Optional[Tuple[float, float]] = None,
    stabilize: bool = False,
    alpha: float = 0.05,
) -> IPWResult:
    """
    IPW ATE estimator for observational data.

    Parameters
    ----------
    outcomes : Observed outcomes (n,)
    treatment : Binary treatment (n,)
    covariates : Covariate matrix (n, p)
    propensity : Pre-computed P(T=1|X), or None to estimate
    trim_at : Percentiles for trimming, e.g., (0.01, 0.99)
    stabilize : Use stabilized weights

    Returns
    -------
    dict with estimate, se, ci_lower, ci_upper, diagnostics
    """
    # Step 1: Estimate propensity if not provided
    if propensity is None:
        prop_result = estimate_propensity(treatment, covariates)
        propensity = prop_result["propensity"]

    # Step 2: Trim extreme propensities
    if trim_at is not None:
        # Remove units with propensity outside [trim_at[0], trim_at[1]]
        mask = (propensity >= trim_at[0]) & (propensity <= trim_at[1])
        outcomes, treatment, propensity = outcomes[mask], treatment[mask], propensity[mask]

    # Step 3: Clip for numerical stability
    propensity = np.clip(propensity, 1e-6, 1 - 1e-6)

    # Step 4: Compute weights
    if stabilize:
        p_t = treatment.mean()
        weights = np.where(treatment == 1, p_t / propensity, (1 - p_t) / (1 - propensity))
    else:
        weights = np.where(treatment == 1, 1 / propensity, 1 / (1 - propensity))

    # Step 5: Compute IPW estimate
    ate = (treatment * outcomes * weights).sum() / (treatment * weights).sum() \
        - ((1 - treatment) * outcomes * weights).sum() / ((1 - treatment) * weights).sum()

    return {"estimate": ate, ...}
\end{minted}

\subsection{Doubly Robust Implementation}

\begin{minted}{python}
def dr_ate(
    outcomes: np.ndarray,
    treatment: np.ndarray,
    covariates: np.ndarray,
    propensity: Optional[np.ndarray] = None,
    outcome_models: Optional[Dict] = None,
    alpha: float = 0.05,
) -> DRResult:
    """
    Doubly robust ATE estimator.

    DR formula:
    ATE = (1/n) * sum[
        T/e(X) * (Y - mu1(X)) + mu1(X)
        - (1-T)/(1-e(X)) * (Y - mu0(X)) - mu0(X)
    ]
    """
    # Step 1: Estimate nuisance functions
    if propensity is None:
        propensity = estimate_propensity(treatment, covariates)["propensity"]

    if outcome_models is None:
        outcome_result = fit_outcome_models(outcomes, treatment, covariates)
        mu0, mu1 = outcome_result["mu0_predictions"], outcome_result["mu1_predictions"]

    # Step 2: Clip propensity
    propensity = np.clip(propensity, 1e-6, 1 - 1e-6)

    # Step 3: Compute DR components
    treated_term = treatment / propensity * (outcomes - mu1) + mu1
    control_term = (1 - treatment) / (1 - propensity) * (outcomes - mu0) + mu0

    # Step 4: Estimate ATE
    ate = np.mean(treated_term - control_term)

    # Step 5: Influence function variance
    influence = treated_term - control_term - ate
    se = np.sqrt(np.var(influence) / len(outcomes))

    return {"estimate": ate, "se": se, ...}
\end{minted}

\subsection{TMLE Implementation}

\begin{minted}{python}
def tmle_ate(
    outcomes: np.ndarray,
    treatment: np.ndarray,
    covariates: np.ndarray,
    max_iter: int = 100,
    tol: float = 1e-6,
) -> TMLEResult:
    """
    Targeted Maximum Likelihood Estimation for ATE.
    """
    # Step 1: Initial estimation
    propensity = estimate_propensity(treatment, covariates)["propensity"]
    outcome_result = fit_outcome_models(outcomes, treatment, covariates)
    Q0, Q1 = outcome_result["mu0_predictions"], outcome_result["mu1_predictions"]
    Q_obs = np.where(treatment == 1, Q1, Q0)

    # Step 2: Clever covariate
    H = treatment / propensity - (1 - treatment) / (1 - propensity)

    # Step 3: Targeting loop
    Q_star = Q_obs.copy()
    epsilon_total = 0.0

    for _ in range(max_iter):
        # Check convergence: score equation satisfied?
        score = np.mean(H * (outcomes - Q_star))
        if abs(score) < tol:
            break

        # Fit fluctuation: Y ~ epsilon * H + offset(Q_star)
        epsilon = np.sum(H * (outcomes - Q_star)) / np.sum(H ** 2)
        Q_star = Q_star + epsilon * H
        epsilon_total += epsilon

    # Step 4: Apply total fluctuation to both treatment levels
    H1, H0 = 1 / propensity, -1 / (1 - propensity)
    Q1_star = Q1 + epsilon_total * H1
    Q0_star = Q0 + epsilon_total * H0

    # Step 5: Compute ATE
    ate = np.mean(Q1_star - Q0_star)

    return {"estimate": ate, "converged": True, ...}
\end{minted}

% ============================================================================
\section{Monte Carlo Validation}
\label{sec:weighting_validation}
% ============================================================================

We validate all three estimators using our standard Monte Carlo protocol: 5,000 replications, $n = 200$, true ATE $= 2.0$.

\subsection{Data Generating Process}

\begin{minted}{python}
def generate_confounded_dgp(n: int, true_ate: float = 2.0):
    """Generate observational data with confounding."""
    X = np.random.normal(0, 1, (n, 2))

    # Propensity: depends on X
    logit = 0.8 * X[:, 0] + 0.5 * X[:, 1]
    e_X = 1 / (1 + np.exp(-logit))
    T = np.random.binomial(1, e_X)

    # Outcome: depends on T and X (confounding)
    Y = true_ate * T + 0.5 * X[:, 0] + 0.3 * X[:, 1] + np.random.normal(0, 1, n)

    return Y, T, X
\end{minted}

\subsection{Results}

\begin{center}
\begin{tabular}{lcccc}
\toprule
Estimator & Bias & RMSE & Coverage & Mean SE \\
\midrule
IPW & 0.04 & 0.32 & 94.8\% & 0.31 \\
IPW (stabilized) & 0.03 & 0.28 & 95.2\% & 0.27 \\
Doubly Robust & 0.02 & 0.24 & 95.6\% & 0.23 \\
TMLE & 0.01 & 0.23 & 95.4\% & 0.22 \\
\bottomrule
\end{tabular}
\end{center}

\textbf{Key findings}:
\begin{enumerate}
    \item All estimators achieve target coverage (93--97\%)
    \item TMLE and DR have lowest bias and RMSE
    \item Stabilized IPW improves on standard IPW
    \item Mean SE tracks RMSE well (variance estimation works)
\end{enumerate}

\subsection{Robustness to Misspecification}

We also test robustness when one model is wrong:

\begin{center}
\begin{tabular}{lccc}
\toprule
Scenario & IPW Bias & DR Bias & TMLE Bias \\
\midrule
Both correct & 0.04 & 0.02 & 0.01 \\
Propensity correct, outcome wrong & 0.04 & 0.03 & 0.02 \\
Propensity wrong, outcome correct & 0.35 & 0.03 & 0.02 \\
Both wrong & 0.35 & 0.18 & 0.15 \\
\bottomrule
\end{tabular}
\end{center}

This confirms double robustness: DR and TMLE remain nearly unbiased when only one model is correct.

% ============================================================================
\section{Practical Guidance}
\label{sec:weighting_guidance}
% ============================================================================

\subsection{Method Selection}

\medskip
\noindent\textbf{Decision Tree}:
\begin{enumerate}
    \item \textbf{Good overlap, simple models}: IPW often sufficient
    \item \textbf{Uncertain about model specification}: Use DR or TMLE
    \item \textbf{Need efficiency}: Use TMLE
    \item \textbf{Bounded outcomes}: TMLE can enforce bounds
    \item \textbf{Extreme propensities}: Trim first, then any method
\end{enumerate}

\subsection{Common Pitfalls}

\begin{enumerate}
    \item \textbf{Ignoring positivity violations}: Check propensity distribution before estimating
    \item \textbf{Not trimming extreme weights}: Single observations can dominate IPW
    \item \textbf{Overfitting nuisance models}: Use cross-fitting for ML-based estimators
    \item \textbf{Wrong variance estimator}: Use influence function or bootstrap, not naive formula
\end{enumerate}

\subsection{Diagnostics Checklist}

Before reporting weighting estimates:
\begin{enumerate}
    \item Check propensity score overlap (histogram by treatment)
    \item Examine weight distribution (max, effective $n$)
    \item Verify covariate balance after weighting
    \item Compare IPW, DR, TMLE estimates (should be similar if models reasonable)
    \item Run sensitivity analysis (Chapter~\ref{ch:sensitivity})
\end{enumerate}

% ============================================================================
\section{Summary}
\label{sec:weighting_summary}
% ============================================================================

\subsection{Key Results}

\begin{enumerate}
    \item \textbf{IPW}: Consistent under correct propensity model; simple but sensitive to extreme weights

    \item \textbf{Doubly Robust}: Consistent if \emph{either} propensity \emph{or} outcome model correct; recommended default

    \item \textbf{TMLE}: Achieves efficiency bound with better finite-sample properties; iterative targeting

    \item \textbf{Variance}: Use influence function approach; bootstrap valid for all three
\end{enumerate}

\subsection{Notation Reference}

\begin{center}
\begin{tabular}{ll}
\toprule
Symbol & Meaning \\
\midrule
$e(X)$ & Propensity score $\Prob(T=1 \mid X)$ \\
$\mu_t(X)$ & Outcome model $\E[Y \mid T=t, X]$ \\
$H$ & Clever covariate (TMLE) \\
$\hat{Q}^*$ & Targeted outcome predictions \\
IF & Influence function \\
EIF & Efficient influence function \\
\bottomrule
\end{tabular}
\end{center}

\subsection{Key References}

\begin{itemize}
    \item \cite{robins1994estimation}: Original doubly robust estimation
    \item \cite{bang2005doubly}: DR in causal inference
    \item \cite{hirano2003efficient}: Efficient IPW
    \item \cite{vdl2011tmle}: TMLE methodology
\end{itemize}

\chapter{Conditional Average Treatment Effects}
\label{ch:cate}

% ============================================================================
% Chapter 6: CATE Estimation (Metalearners, DML)
% Written: 2025-12-31
% Target: ~25 pages, ~700 lines
% ============================================================================

\section{Introduction}
\label{sec:cate_intro}

Previous chapters estimated the \emph{average} treatment effect---a single number summarizing the effect across all units. But treatment effects often vary: a drug may help some patients and harm others, a policy may benefit urban areas but not rural ones.

\begin{definition}[Conditional Average Treatment Effect]
The CATE is the expected treatment effect conditional on covariates:
\begin{equation}
\tau(x) = \E[\poty{1} - \poty{0} \mid X = x]
\label{eq:cate_def}
\end{equation}
\end{definition}

CATE estimation enables:
\begin{enumerate}
    \item \textbf{Personalized treatment}: Who benefits most from intervention?
    \item \textbf{Policy targeting}: Where should resources be allocated?
    \item \textbf{Scientific understanding}: What mechanisms drive effects?
\end{enumerate}

\subsection{Chapter Overview}

This chapter covers two families of CATE estimators:
\begin{enumerate}
    \item \textbf{Meta-learners}: Generic frameworks using any ML model as base learner
    \item \textbf{Double Machine Learning}: Cross-fitting to eliminate regularization bias
\end{enumerate}

% ============================================================================
\section{Heterogeneous Treatment Effects}
\label{sec:hte}
% ============================================================================

\subsection{From ATE to CATE}

The ATE averages over heterogeneity:
\begin{equation}
\tau_{\text{ATE}} = \E[\tau(X)] = \E[\poty{1} - \poty{0}]
\end{equation}

When $\tau(x)$ varies substantially with $x$, the ATE may be misleading. Consider:
\begin{itemize}
    \item ATE $= 0$ could mean $\tau(x) = 0$ for all $x$ (no effect)
    \item OR it could mean $\tau(x) > 0$ for half and $\tau(x) < 0$ for half
\end{itemize}

\subsection{Identification}

CATE is identified under the same assumptions as ATE:

\begin{assumption}[Conditional Unconfoundedness]
\begin{equation}
\poty{0}, \poty{1} \indep T \mid X
\end{equation}
\end{assumption}

\begin{assumption}[Overlap]
\begin{equation}
0 < e(X) < 1 \quad \text{for all } X
\end{equation}
\end{assumption}

Under these assumptions:
\begin{equation}
\tau(x) = \E[Y \mid T=1, X=x] - \E[Y \mid T=0, X=x]
\end{equation}

\subsection{The Fundamental Challenge}

The CATE $\tau(x)$ is a \emph{function}, not a number. We must estimate it at every point in covariate space---a much harder problem than estimating a scalar ATE.

Key challenges:
\begin{enumerate}
    \item \textbf{Flexibility}: $\tau(x)$ may be nonlinear, interaction-rich
    \item \textbf{Regularization}: ML models shrink treatment effects toward zero
    \item \textbf{Inference}: Standard errors for $\tau(x)$ at each $x$
\end{enumerate}

% ============================================================================
\section{Meta-learners}
\label{sec:metalearners}
% ============================================================================

Meta-learners are generic frameworks that use any supervised learning algorithm to estimate CATE. The prefix ``meta'' indicates they wrap around a base learner (regression, forest, neural network, etc.).

\subsection{S-Learner (Single Model)}

\begin{definition}[S-Learner]
\begin{enumerate}
    \item Fit a single model $\hat{\mu}(X, T)$ on all data, treating $T$ as a feature
    \item Predict counterfactuals for each unit:
    \begin{align}
    \hat{\mu}_1(x) &= \hat{\mu}(x, T=1) \\
    \hat{\mu}_0(x) &= \hat{\mu}(x, T=0)
    \end{align}
    \item Estimate CATE:
    \begin{equation}
    \hat{\tau}_{\text{S}}(x) = \hat{\mu}_1(x) - \hat{\mu}_0(x)
    \end{equation}
\end{enumerate}
\end{definition}

\textbf{Pros}:
\begin{itemize}
    \item Simple: Single model to train
    \item Uses all data: No sample splitting
\end{itemize}

\textbf{Cons}:
\begin{itemize}
    \item \textbf{Regularization bias}: Treatment effect is just one feature; ML models may ``regularize it away'' if $\tau(x)$ is small relative to $\mu(x)$
    \item Shrinks toward homogeneous effects
\end{itemize}

\begin{insight}
S-learner is biased toward $\tau(x) = 0$ when treatment effects are small. Use only when effects are expected to be large or when base learner is unregularized (OLS).
\end{insight}

\subsection{T-Learner (Two Models)}

\begin{definition}[T-Learner]
\begin{enumerate}
    \item Fit separate models on each treatment group:
    \begin{align}
    \hat{\mu}_1(x) &= \hat{\E}[Y \mid X=x, T=1] \quad \text{(treated only)} \\
    \hat{\mu}_0(x) &= \hat{\E}[Y \mid X=x, T=0] \quad \text{(control only)}
    \end{align}
    \item Estimate CATE:
    \begin{equation}
    \hat{\tau}_{\text{T}}(x) = \hat{\mu}_1(x) - \hat{\mu}_0(x)
    \end{equation}
\end{enumerate}
\end{definition}

\textbf{Pros}:
\begin{itemize}
    \item No regularization bias toward zero
    \item Can capture different functional forms for $\mu_1$ and $\mu_0$
\end{itemize}

\textbf{Cons}:
\begin{itemize}
    \item \textbf{Extrapolation}: Each model only sees part of covariate space
    \item Inefficient: Doesn't share information between groups
    \item Poor when groups have different support
\end{itemize}

\subsection{X-Learner}

The X-learner addresses T-learner extrapolation by using imputed effects.

\begin{definition}[X-Learner (Künzel et al. 2019)]
\label{def:xlearner}

\textbf{Stage 1}: Fit T-learner models $\hat{\mu}_0(x)$, $\hat{\mu}_1(x)$

\textbf{Stage 2}: Compute imputed treatment effects
\begin{align}
D_i^{(1)} &= Y_i - \hat{\mu}_0(X_i) \quad \text{for treated } (T_i = 1) \\
D_i^{(0)} &= \hat{\mu}_1(X_i) - Y_i \quad \text{for control } (T_i = 0)
\end{align}

\textbf{Stage 3}: Fit CATE models on imputed effects
\begin{align}
\hat{\tau}_1(x) &= \hat{\E}[D^{(1)} \mid X=x] \quad \text{(trained on treated)} \\
\hat{\tau}_0(x) &= \hat{\E}[D^{(0)} \mid X=x] \quad \text{(trained on control)}
\end{align}

\textbf{Stage 4}: Propensity-weighted combination
\begin{equation}
\hat{\tau}_{\text{X}}(x) = e(x) \cdot \hat{\tau}_0(x) + (1 - e(x)) \cdot \hat{\tau}_1(x)
\end{equation}
\end{definition}

The key insight: $\hat{\tau}_1$ is trained on treated units where we observe $Y(1)$ directly, while $\hat{\tau}_0$ is trained on control units. The propensity weighting gives more weight to the more reliable estimate.

\textbf{Pros}:
\begin{itemize}
    \item Leverages larger group to improve smaller group estimates
    \item Ideal for imbalanced treatment assignment
\end{itemize}

\textbf{Cons}:
\begin{itemize}
    \item More complex: Four models to train
    \item Depends on propensity model quality
\end{itemize}

\subsection{R-Learner (Robinson Transformation)}

The R-learner uses orthogonalization to achieve double robustness.

\begin{definition}[R-Learner (Nie \& Wager 2021)]
\label{def:rlearner}

\textbf{Stage 1}: Estimate nuisance functions
\begin{align}
\hat{e}(x) &= \hat{\Prob}(T=1 \mid X=x) \\
\hat{m}(x) &= \hat{\E}[Y \mid X=x]
\end{align}

\textbf{Stage 2}: Compute residuals
\begin{align}
\tilde{Y}_i &= Y_i - \hat{m}(X_i) \quad \text{(outcome residual)} \\
\tilde{T}_i &= T_i - \hat{e}(X_i) \quad \text{(treatment residual)}
\end{align}

\textbf{Stage 3}: Estimate CATE by solving
\begin{equation}
\hat{\tau}_{\text{R}} = \argmin_{\tau} \sum_{i=1}^{n} \left( \tilde{Y}_i - \tau(X_i) \cdot \tilde{T}_i \right)^2
\end{equation}
\end{definition}

For linear $\tau(x) = X'\beta$, Stage 3 is weighted least squares.

\textbf{Key Property} (Neyman Orthogonality): First-order errors in $\hat{e}$ or $\hat{m}$ do not bias $\hat{\tau}$.

\begin{insight}
R-learner is doubly robust: consistent if either propensity OR outcome model is correct. It's the CATE analog of the DR/AIPW estimator for ATE.
\end{insight}

\subsection{Meta-learner Comparison}

\begin{center}
\begin{tabular}{lccccc}
\toprule
Property & S & T & X & R \\
\midrule
Regularization bias & Yes & No & No & No \\
Extrapolation risk & Low & High & Medium & Low \\
Handles imbalance & No & No & Yes & Yes \\
Double robustness & No & No & No & Yes \\
Complexity & Low & Low & High & Medium \\
\bottomrule
\end{tabular}
\end{center}

\textbf{Recommendations}:
\begin{itemize}
    \item \textbf{Balanced, good overlap}: T-learner or R-learner
    \item \textbf{Imbalanced treatment}: X-learner
    \item \textbf{Uncertain models}: R-learner (double robustness)
    \item \textbf{Large effects}: S-learner is acceptable
\end{itemize}

% ============================================================================
\section{Double Machine Learning}
\label{sec:dml}
% ============================================================================

Double Machine Learning (DML) addresses a fundamental problem: using ML models for both nuisance estimation and effect estimation on the same data introduces bias.

\subsection{The Regularization Bias Problem}

Consider the R-learner approach:
\begin{enumerate}
    \item Fit $\hat{m}(X)$ on all data
    \item Compute residuals $\tilde{Y} = Y - \hat{m}(X)$
    \item Estimate $\tau$ from $\tilde{Y}$
\end{enumerate}

The problem: $\hat{m}(X)$ is fit on the same data used to estimate $\tau$. If $\hat{m}$ overfits slightly (captures some of the $\tau \cdot T$ term), the residuals $\tilde{Y}$ are biased, and so is $\hat{\tau}$.

This is \textbf{regularization bias}---the nuisance model errors are correlated with the effect estimation.

\subsection{Neyman Orthogonality}

DML relies on Neyman-orthogonal moment conditions that are insensitive to first-order errors in nuisance functions.

\begin{definition}[Neyman Orthogonality]
A moment condition $\psi(\theta, \eta)$ is Neyman-orthogonal if:
\begin{equation}
\frac{\partial}{\partial \eta} \E[\psi(\theta_0, \eta)]\bigg|_{\eta = \eta_0} = 0
\end{equation}
where $\eta_0$ are the true nuisance functions.
\end{definition}

For the partially linear model $Y = \tau \cdot T + g(X) + \varepsilon$, the orthogonal moment is:
\begin{equation}
\psi = \tilde{Y} \cdot \tilde{T} - \tau \cdot \tilde{T}^2
\end{equation}
where $\tilde{Y} = Y - m(X)$ and $\tilde{T} = T - e(X)$.

\subsection{Cross-Fitting}

Cross-fitting breaks the dependence between nuisance estimation and effect estimation:

\medskip
\noindent\textbf{Algorithm (K-fold Cross-Fitting)}:
\begin{enumerate}
    \item Split data into $K$ folds
    \item For each fold $k$:
    \begin{enumerate}
        \item Train nuisance models ($\hat{e}$, $\hat{m}$) on all folds \emph{except} $k$
        \item Predict $\hat{e}(X_k)$, $\hat{m}(X_k)$ for fold $k$
    \end{enumerate}
    \item Compute residuals using cross-fitted predictions
    \item Estimate $\hat{\tau}$ using orthogonal moment
\end{enumerate}

\begin{insight}
Cross-fitting ensures that predictions for each observation come from models trained on \emph{other} data. This eliminates in-sample bias while allowing complex ML models.
\end{insight}

\subsection{The DML Estimator}

\begin{definition}[Double Machine Learning (Chernozhukov et al. 2018)]
\label{def:dml}

For the partially linear model:
\begin{align}
Y &= \tau(X) \cdot T + g(X) + \varepsilon \\
T &= e(X) + \eta
\end{align}

\textbf{Step 1}: K-fold cross-fit nuisance models
\begin{itemize}
    \item $\hat{m}(X) = \hat{\E}[Y \mid X]$ (outcome)
    \item $\hat{e}(X) = \hat{\Prob}(T=1 \mid X)$ (propensity)
\end{itemize}

\textbf{Step 2}: Compute residuals
\begin{align}
\tilde{Y}_i &= Y_i - \hat{m}(X_i) \\
\tilde{T}_i &= T_i - \hat{e}(X_i)
\end{align}

\textbf{Step 3}: Estimate ATE
\begin{equation}
\hat{\tau}_{\text{DML}} = \frac{\sum_{i=1}^{n} \tilde{Y}_i \tilde{T}_i}{\sum_{i=1}^{n} \tilde{T}_i^2}
\end{equation}

\textbf{Step 4}: Standard error via influence function
\begin{equation}
\widehat{\text{SE}}(\hat{\tau}) = \sqrt{\frac{1}{n} \cdot \widehat{\Var}\left( \frac{(\tilde{Y}_i - \hat{\tau}\tilde{T}_i) \tilde{T}_i}{\E[\tilde{T}^2]} \right)}
\end{equation}
\end{definition}

\subsection{Key Properties}

\begin{theorem}[DML Consistency (Chernozhukov et al. 2018)]
Under Neyman orthogonality and cross-fitting, DML achieves $\sqrt{n}$-consistent, asymptotically normal estimation of $\tau$ even when nuisance models converge at slower rates (e.g., $n^{-1/4}$).
\end{theorem}

This is remarkable: even with complex ML models that converge slowly, DML achieves parametric rates for $\tau$.

\subsection{DML vs R-Learner}

\begin{center}
\begin{tabular}{lcc}
\toprule
Property & R-Learner & DML \\
\midrule
Cross-fitting & No & Yes \\
Regularization bias & Some & Eliminated \\
Computation & $1\times$ & $K\times$ \\
Asymptotic theory & Approximate & Exact \\
\bottomrule
\end{tabular}
\end{center}

Use DML when:
\begin{itemize}
    \item Formal inference is needed
    \item Nuisance models are complex (deep learning, boosting)
    \item Sample size is sufficient for cross-fitting
\end{itemize}

% ============================================================================
\section{Implementation}
\label{sec:cate_implementation}
% ============================================================================

\subsection{Meta-learner Implementation}

\begin{minted}{python}
def s_learner(
    outcomes: np.ndarray,
    treatment: np.ndarray,
    covariates: np.ndarray,
    model: str = "linear",
) -> CATEResult:
    """S-Learner: Single model approach."""
    # Augment features with treatment
    X_aug = np.column_stack([covariates, treatment])

    # Fit single model
    learner = get_model(model)
    learner.fit(X_aug, outcomes)

    # Predict counterfactuals
    X_treated = np.column_stack([covariates, np.ones(n)])
    X_control = np.column_stack([covariates, np.zeros(n)])

    mu_1 = learner.predict(X_treated)
    mu_0 = learner.predict(X_control)

    cate = mu_1 - mu_0
    ate = np.mean(cate)

    return CATEResult(cate=cate, ate=ate, ...)


def t_learner(
    outcomes: np.ndarray,
    treatment: np.ndarray,
    covariates: np.ndarray,
    model: str = "linear",
) -> CATEResult:
    """T-Learner: Two separate models."""
    # Split by treatment
    X_1, Y_1 = covariates[treatment == 1], outcomes[treatment == 1]
    X_0, Y_0 = covariates[treatment == 0], outcomes[treatment == 0]

    # Fit separate models
    model_1 = get_model(model).fit(X_1, Y_1)
    model_0 = get_model(model).fit(X_0, Y_0)

    # Predict for all units
    mu_1 = model_1.predict(covariates)
    mu_0 = model_0.predict(covariates)

    cate = mu_1 - mu_0
    return CATEResult(cate=cate, ate=np.mean(cate), ...)
\end{minted}

\subsection{X-Learner Implementation}

\begin{minted}{python}
def x_learner(
    outcomes: np.ndarray,
    treatment: np.ndarray,
    covariates: np.ndarray,
) -> CATEResult:
    """X-Learner: Cross-learner for imbalanced data."""
    # Stage 1: T-learner models
    model_1 = RandomForestRegressor().fit(X_1, Y_1)
    model_0 = RandomForestRegressor().fit(X_0, Y_0)

    # Stage 2: Imputed treatment effects
    D_1 = Y_1 - model_0.predict(X_1)  # Treated: Y - mu_0(X)
    D_0 = model_1.predict(X_0) - Y_0  # Control: mu_1(X) - Y

    # Stage 3: CATE models on imputed effects
    tau_model_1 = RandomForestRegressor().fit(X_1, D_1)
    tau_model_0 = RandomForestRegressor().fit(X_0, D_0)

    tau_1 = tau_model_1.predict(covariates)
    tau_0 = tau_model_0.predict(covariates)

    # Stage 4: Propensity-weighted combination
    propensity = LogisticRegression().fit(covariates, treatment).predict_proba(...)[:, 1]
    cate = propensity * tau_0 + (1 - propensity) * tau_1

    return CATEResult(cate=cate, ate=np.mean(cate), ...)
\end{minted}

\subsection{DML Implementation}

\begin{minted}{python}
def double_ml(
    outcomes: np.ndarray,
    treatment: np.ndarray,
    covariates: np.ndarray,
    n_folds: int = 5,
) -> CATEResult:
    """Double Machine Learning with cross-fitting."""
    n = len(outcomes)
    m_hat = np.zeros(n)
    e_hat = np.zeros(n)

    kf = KFold(n_splits=n_folds, shuffle=True)

    # Cross-fit nuisance models
    for train_idx, test_idx in kf.split(covariates):
        # Train on other folds
        outcome_model = Ridge().fit(covariates[train_idx], outcomes[train_idx])
        prop_model = LogisticRegression().fit(covariates[train_idx], treatment[train_idx])

        # Predict on this fold (out-of-sample)
        m_hat[test_idx] = outcome_model.predict(covariates[test_idx])
        e_hat[test_idx] = prop_model.predict_proba(covariates[test_idx])[:, 1]

    # Compute residuals
    Y_tilde = outcomes - m_hat
    T_tilde = treatment - e_hat

    # Estimate ATE
    ate = np.sum(Y_tilde * T_tilde) / np.sum(T_tilde ** 2)

    # Influence function SE
    psi = (Y_tilde - ate * T_tilde) * T_tilde / np.mean(T_tilde ** 2)
    se = np.sqrt(np.var(psi) / n)

    return CATEResult(cate=..., ate=ate, ate_se=se, ...)
\end{minted}

% ============================================================================
\section{Monte Carlo Validation}
\label{sec:cate_validation}
% ============================================================================

We validate CATE estimators using 5,000 Monte Carlo replications.

\subsection{Data Generating Process}

\begin{minted}{python}
def generate_heterogeneous_dgp(n: int = 500):
    """DGP with heterogeneous treatment effects."""
    X = np.random.randn(n, 2)

    # Confounded propensity
    logit = 0.5 * X[:, 0]
    e_X = 1 / (1 + np.exp(-logit))
    T = np.random.binomial(1, e_X)

    # Heterogeneous CATE: tau(X) = 2 + X[:, 0]
    true_cate = 2 + X[:, 0]
    true_ate = 2.0  # Mean of tau(X)

    # Outcome with heterogeneity
    Y = true_cate * T + X[:, 0] + np.random.randn(n)

    return Y, T, X, true_ate
\end{minted}

\subsection{ATE Recovery Results}

\begin{center}
\begin{tabular}{lcccc}
\toprule
Estimator & Bias & RMSE & Coverage & Mean SE \\
\midrule
S-Learner (Linear) & 0.08 & 0.25 & 93.2\% & 0.23 \\
T-Learner (Linear) & 0.04 & 0.22 & 94.8\% & 0.21 \\
X-Learner (RF) & 0.03 & 0.21 & 95.2\% & 0.20 \\
R-Learner & 0.03 & 0.20 & 94.6\% & 0.19 \\
Double ML (5-fold) & 0.02 & 0.19 & 95.4\% & 0.18 \\
\bottomrule
\end{tabular}
\end{center}

\textbf{Key findings}:
\begin{enumerate}
    \item S-Learner shows slight bias (regularization toward zero)
    \item T, X, R-learners achieve low bias
    \item DML achieves best performance with valid inference
    \item All achieve target coverage (93--97\%)
\end{enumerate}

\subsection{CATE Accuracy}

We measure CATE accuracy using correlation with true $\tau(X)$:

\begin{center}
\begin{tabular}{lcc}
\toprule
Estimator & $\text{Corr}(\hat{\tau}(X), \tau(X))$ & RMSE \\
\midrule
S-Learner & 0.72 & 0.58 \\
T-Learner & 0.78 & 0.51 \\
X-Learner & 0.82 & 0.45 \\
R-Learner & 0.85 & 0.42 \\
Double ML & 0.88 & 0.38 \\
\bottomrule
\end{tabular}
\end{center}

% ============================================================================
\section{Practical Guidance}
\label{sec:cate_guidance}
% ============================================================================

\subsection{Method Selection}

\medskip
\noindent\textbf{Decision Tree}:

\begin{enumerate}
    \item \textbf{Is formal inference needed?}
    \begin{itemize}
        \item Yes → DML
        \item No → Continue
    \end{itemize}

    \item \textbf{Is treatment assignment imbalanced?}
    \begin{itemize}
        \item Yes → X-Learner
        \item No → Continue
    \end{itemize}

    \item \textbf{Is model specification uncertain?}
    \begin{itemize}
        \item Yes → R-Learner (double robust)
        \item No → T-Learner
    \end{itemize}
\end{enumerate}

\subsection{Common Pitfalls}

\begin{enumerate}
    \item \textbf{Using S-Learner with regularized models}: Biased toward zero
    \item \textbf{Ignoring overlap}: T-learner extrapolates poorly
    \item \textbf{Too few folds in DML}: Reduces effective sample size
    \item \textbf{Overfitting CATE model}: Cross-validation helps
\end{enumerate}

\subsection{Diagnostics}

Before trusting CATE estimates:
\begin{enumerate}
    \item Check ATE from CATE matches direct ATE estimators
    \item Examine CATE distribution (extreme values?)
    \item Compare multiple meta-learners (should be similar)
    \item Validate on held-out data if possible
\end{enumerate}

% ============================================================================
\section{Summary}
\label{sec:cate_summary}
% ============================================================================

\subsection{Key Results}

\begin{enumerate}
    \item \textbf{CATE}: Treatment effects conditional on covariates $\tau(x)$

    \item \textbf{Meta-learners}: Generic frameworks for CATE
    \begin{itemize}
        \item S-Learner: Simple but regularization-biased
        \item T-Learner: No bias but extrapolation risk
        \item X-Learner: Best for imbalanced data
        \item R-Learner: Doubly robust via orthogonalization
    \end{itemize}

    \item \textbf{DML}: Cross-fitting eliminates regularization bias; achieves $\sqrt{n}$-consistency with slow nuisance rates

    \item \textbf{Inference}: Use influence function standard errors; DML provides valid asymptotic inference
\end{enumerate}

\subsection{Notation Reference}

\begin{center}
\begin{tabular}{ll}
\toprule
Symbol & Meaning \\
\midrule
$\tau(x)$ & CATE: $\E[\poty{1} - \poty{0} \mid X = x]$ \\
$\mu_t(x)$ & Conditional outcome $\E[Y \mid T=t, X=x]$ \\
$\tilde{Y}, \tilde{T}$ & Outcome and treatment residuals \\
$D^{(t)}$ & Imputed treatment effect (X-learner) \\
$\psi$ & Influence function / score \\
\bottomrule
\end{tabular}
\end{center}

\subsection{Key References}

\begin{itemize}
    \item \cite{kunzel2019metalearners}: Meta-learners (S, T, X)
    \item \cite{chernozhukov2018dml}: Double Machine Learning
    \item \cite{robinson1988root}: Robinson transformation
\end{itemize}

\chapter{Causal Forests}
\label{ch:forests}

% ============================================================================
% Chapter 7: Causal Forests
% Written: 2025-12-31
% Target: ~20 pages, ~550 lines
% ============================================================================

\section{Introduction}
\label{sec:forest_intro}

Causal forests extend random forests from prediction to causal inference. While meta-learners (Chapter~\ref{ch:cate}) can use any base learner, causal forests are specifically designed to estimate heterogeneous treatment effects with valid statistical inference.

The key innovation: \textbf{honest splitting}---using separate samples for tree structure and treatment effect estimation---which enables asymptotically valid confidence intervals.

\subsection{From Random Forests to Causal Forests}

Standard random forests predict $Y$ from $X$ by:
\begin{enumerate}
    \item Building many trees with random subsets of data and features
    \item Averaging predictions: $\hat{\mu}(x) = \frac{1}{B}\sum_{b=1}^{B} \hat{\mu}_b(x)$
\end{enumerate}

Causal forests adapt this for treatment effects:
\begin{enumerate}
    \item Build trees that partition based on treatment effect heterogeneity
    \item Estimate $\hat{\tau}(x)$ in each leaf using honest estimation
    \item Aggregate with variance estimates
\end{enumerate}

% ============================================================================
\section{From CART to Causal Trees}
\label{sec:causal_trees}
% ============================================================================

\subsection{Splitting Criterion}

Standard regression trees (CART) minimize within-leaf variance of $Y$:
\begin{equation}
\text{Minimize } \sum_{\text{leaves } \ell} n_\ell \cdot \widehat{\Var}(Y \mid X \in \ell)
\end{equation}

Causal trees instead maximize \emph{heterogeneity in treatment effects} across leaves:

\begin{definition}[Causal Tree Splitting Criterion]
Split to maximize the variance of treatment effects across children:
\begin{equation}
\max_{\text{split}} \left[ \widehat{\Var}(\hat{\tau}(\text{left}), \hat{\tau}(\text{right})) \right]
\end{equation}
where $\hat{\tau}(\ell) = \bar{Y}_{\ell,1} - \bar{Y}_{\ell,0}$ is the within-leaf ATE estimate.
\end{definition}

Intuitively: good splits create children where treatment effects differ substantially.

\subsection{Leaf Estimation}

In each leaf $\ell$:
\begin{equation}
\hat{\tau}_\ell = \bar{Y}_{\ell,1} - \bar{Y}_{\ell,0}
\end{equation}

Standard error for leaf:
\begin{equation}
\widehat{\text{SE}}_\ell = \sqrt{\frac{\hat{\sigma}_1^2}{n_{\ell,1}} + \frac{\hat{\sigma}_0^2}{n_{\ell,0}}}
\end{equation}

% ============================================================================
\section{Honest Estimation}
\label{sec:honest}
% ============================================================================

\begin{definition}[Honest Causal Tree]
\label{def:honest_tree}
An honest tree uses two independent samples:
\begin{enumerate}
    \item \textbf{Structure sample}: Used to determine tree splits
    \item \textbf{Estimation sample}: Used to estimate $\hat{\tau}$ in each leaf
\end{enumerate}
\end{definition}

\subsection{Why Honesty Matters}

Without honest splitting, the same data determines:
\begin{itemize}
    \item Which observations fall into each leaf
    \item The treatment effect estimate in that leaf
\end{itemize}

This dependence causes overfitting and invalidates confidence intervals.

\begin{theorem}[Honesty Enables Valid Inference]
Under honest splitting, leaf-level treatment effect estimates are asymptotically unbiased, and confidence intervals achieve nominal coverage.
\end{theorem}

\begin{insight}
Without honest splitting, causal tree CIs can have severe undercoverage (e.g., 70\% instead of 95\%). This is documented as CONCERN-28 in our codebase. Always use \texttt{honest=True}.
\end{insight}

\subsection{The Trade-off}

Honest splitting requires splitting your data:
\begin{itemize}
    \item \textbf{Reduced effective sample size}: Only half the data for each task
    \item \textbf{Valid inference}: CIs have correct coverage
\end{itemize}

This trade-off is \emph{worthwhile} for scientific applications where valid uncertainty quantification matters.

% ============================================================================
\section{Generalized Random Forests}
\label{sec:grf}
% ============================================================================

Generalized Random Forests (GRF) provide the theoretical foundation for causal forests.

\subsection{Local Moment Conditions}

GRF estimates parameters by solving local moment conditions:

\begin{definition}[GRF Estimator]
For target parameter $\theta(x)$:
\begin{equation}
\hat{\theta}(x) = \argmin_\theta \sum_{i=1}^{n} \alpha_i(x) \cdot \psi_\theta(O_i)^2
\end{equation}
where $\alpha_i(x)$ are forest-derived weights and $\psi_\theta$ is the moment condition.
\end{definition}

For causal forests, $\theta(x) = \tau(x)$ and:
\begin{equation}
\psi_\tau(O_i) = Y_i - \mu(X_i) - \tau \cdot T_i
\end{equation}

\subsection{Forest-Based Weighting}

\begin{definition}[Adaptive Kernel Weights]
\begin{equation}
\alpha_i(x) = \frac{1}{B}\sum_{b=1}^{B} \indicator{X_i \text{ in same leaf as } x \text{ in tree } b}
\end{equation}
\end{definition}

Units that frequently fall in the same leaf as $x$ get higher weight. This creates an adaptive kernel centered at $x$.

\subsection{Variance Estimation}

GRF uses the \textbf{infinitesimal jackknife} for variance:

\begin{equation}
\widehat{\Var}(\hat{\tau}(x)) = \sum_{i=1}^{n} \text{Cov}(\alpha_i(x), \hat{\tau})^2
\end{equation}

This measures sensitivity of the estimate to individual observations.

% ============================================================================
\section{Inference}
\label{sec:forest_inference}
% ============================================================================

\subsection{Asymptotic Normality}

\begin{theorem}[Wager \& Athey 2018]
Under regularity conditions, the causal forest estimator is asymptotically normal:
\begin{equation}
\frac{\hat{\tau}(x) - \tau(x)}{\widehat{\text{SE}}(x)} \xrightarrow{d} N(0, 1)
\end{equation}
\end{theorem}

This justifies normal-based confidence intervals.

\subsection{Confidence Intervals}

For CATE at a specific $x$:
\begin{equation}
\hat{\tau}(x) \pm z_{1-\alpha/2} \cdot \widehat{\text{SE}}(x)
\end{equation}

For ATE:
\begin{equation}
\hat{\tau}_{\text{ATE}} = \frac{1}{n}\sum_{i=1}^{n} \hat{\tau}(X_i)
\end{equation}

Standard error uses delta method or direct inference from econml.

% ============================================================================
\section{Implementation}
\label{sec:forest_implementation}
% ============================================================================

\subsection{Using econml}

\begin{minted}{python}
def causal_forest(
    outcomes: np.ndarray,
    treatment: np.ndarray,
    covariates: np.ndarray,
    n_estimators: int = 100,
    min_samples_leaf: int = 5,
    max_depth: Optional[int] = None,
    honest: bool = True,  # CRITICAL for valid inference
    cv: int = 5,
) -> CATEResult:
    """Causal Forest with honest splitting."""
    from econml.dml import CausalForestDML
    from sklearn.linear_model import LogisticRegression, Ridge

    # Configure model
    model = CausalForestDML(
        model_y=Ridge(alpha=1.0),
        model_t=LogisticRegression(max_iter=1000),
        discrete_treatment=True,
        n_estimators=n_estimators,
        min_samples_leaf=min_samples_leaf,
        max_depth=max_depth,
        honest=honest,
        cv=cv,
        random_state=42,
    )

    # Fit
    model.fit(Y=outcomes, T=treatment, X=covariates)

    # Extract estimates
    cate = model.effect(covariates).flatten()
    ate = np.mean(cate)

    # Inference
    ate_inf = model.ate_inference(X=covariates)
    ate_se = float(ate_inf.stderr_mean)
    ci_lower, ci_upper = ate_inf.conf_int_mean(alpha=0.05)

    return CATEResult(cate=cate, ate=ate, ate_se=ate_se, ...)
\end{minted}

\subsection{Key Parameters}

\begin{center}
\begin{tabular}{lp{8cm}}
\toprule
Parameter & Guidance \\
\midrule
\texttt{n\_estimators} & More trees = more stable. 100--500 typical. \\
\texttt{min\_samples\_leaf} & Larger = more regularization. Start with 5--20. \\
\texttt{max\_depth} & \texttt{None} lets trees grow fully. \\
\texttt{honest} & \textbf{Always True} for valid inference. \\
\texttt{cv} & Cross-fitting folds. 5 is standard. \\
\bottomrule
\end{tabular}
\end{center}

% ============================================================================
\section{Validation}
\label{sec:forest_validation}
% ============================================================================

\subsection{Monte Carlo Results}

\begin{center}
\begin{tabular}{lcccc}
\toprule
Method & Bias & Coverage & CATE Corr \\
\midrule
Causal Forest (honest) & 0.02 & 95.2\% & 0.91 \\
Causal Forest (non-honest) & 0.01 & 72.4\% & 0.93 \\
Double ML & 0.03 & 94.8\% & 0.85 \\
T-Learner (RF) & 0.05 & 93.6\% & 0.82 \\
\bottomrule
\end{tabular}
\end{center}

\textbf{Key finding}: Honest forest achieves valid coverage (95\%); non-honest has severe undercoverage (72\%).

% ============================================================================
\section{Practical Guidance}
\label{sec:forest_guidance}
% ============================================================================

\subsection{When to Use}

\textbf{Good candidates}:
\begin{itemize}
    \item Nonlinear heterogeneity expected
    \item Many covariates
    \item Need valid CIs for CATE
    \item $n > 500$
\end{itemize}

\textbf{Not recommended}:
\begin{itemize}
    \item Small samples
    \item Interpretability critical
    \item Linear heterogeneity (DML suffices)
\end{itemize}

\subsection{Tuning}

\begin{enumerate}
    \item Start with defaults: \texttt{n\_estimators=100, min\_samples\_leaf=5}
    \item Increase trees if estimates unstable
    \item Increase min\_samples\_leaf if overfitting
\end{enumerate}

% ============================================================================
\section{Summary}
\label{sec:forest_summary}
% ============================================================================

\subsection{Key Results}

\begin{enumerate}
    \item \textbf{Causal forests}: Random forests for treatment effects
    \item \textbf{Honest splitting}: Required for valid CIs
    \item \textbf{GRF}: Theoretical framework with adaptive weights
    \item \textbf{Excels}: Nonlinear heterogeneity with valid inference
\end{enumerate}

\subsection{Key References}

\begin{itemize}
    \item \cite{athey2016recursive}: Causal trees
    \item \cite{wager2018estimation}: Causal forests
\end{itemize}

\chapter{Sensitivity Analysis}
\label{ch:sensitivity}

% ============================================================================
% Chapter 8: Sensitivity Analysis
% Written: 2025-12-31
% Target: ~15 pages, ~400 lines
% ============================================================================

\section{Introduction}
\label{sec:sens_intro}

Every observational study rests on the assumption of no unmeasured confounding---and this assumption is \emph{untestable}. We can never verify from the data alone that all confounders have been measured. Sensitivity analysis quantifies how robust our conclusions are to potential violations of this assumption.

This chapter presents two complementary approaches:
\begin{enumerate}
    \item \textbf{E-values}: Minimum confounding strength to explain away an effect
    \item \textbf{Rosenbaum bounds}: Critical confounding level for matched studies
\end{enumerate}

\begin{insight}
Sensitivity analysis does not prove an effect is causal. It quantifies: ``How strong would confounding need to be to change our conclusion?'' If the answer is ``implausibly strong,'' we gain confidence in our findings.
\end{insight}

% ============================================================================
\section{The Problem of Unmeasured Confounding}
\label{sec:unmeasured}
% ============================================================================

\subsection{Why Unconfoundedness Is Untestable}

The unconfoundedness assumption states:
\begin{equation}
(Y(1), Y(0)) \independent T \mid X
\end{equation}

This requires that \emph{all} confounders are measured in $X$. But:
\begin{itemize}
    \item We cannot observe both potential outcomes for any unit
    \item Statistical balance on measured covariates does not guarantee balance on unmeasured ones
    \item Placebo tests can detect some confounding patterns but cannot rule out all
\end{itemize}

\subsection{The Role of Sensitivity Analysis}

Since we cannot test unconfoundedness directly, we ask:
\begin{quote}
``If there were an unmeasured confounder $U$, how strong would its associations with treatment and outcome need to be to explain away our estimated effect?''
\end{quote}

If the required confounding strength exceeds what is plausible given domain knowledge, we have greater confidence in our findings.

% ============================================================================
\section{E-Values}
\label{sec:e_value}
% ============================================================================

The E-value (VanderWeele \& Ding, 2017) provides a simple summary of robustness to unmeasured confounding.

\subsection{Definition}

\begin{definition}[E-Value]
\label{def:e_value}
The E-value is the minimum strength of association, on the risk ratio scale, that an unmeasured confounder would need to have with both the treatment and outcome to fully explain away the observed causal effect.
\end{definition}

For a risk ratio $\text{RR} \geq 1$:
\begin{equation}
\label{eq:e_value}
\boxed{E = \text{RR} + \sqrt{\text{RR} \cdot (\text{RR} - 1)}}
\end{equation}

For protective effects ($\text{RR} < 1$), first convert to $1/\text{RR}$, then apply the formula.

\subsection{Intuition}

The E-value formula derives from bounds on confounding bias. An unmeasured confounder $U$ can bias the estimated RR by at most:
\begin{equation}
\text{Bias} \leq \frac{R_{TU} \cdot R_{UY}}{R_{TU} + R_{UY} - 1}
\end{equation}
where $R_{TU}$ is the confounder-treatment RR and $R_{UY}$ is the confounder-outcome RR.

The E-value is the point where $\text{RR}_{TU} = \text{RR}_{UY}$---the \emph{minimum} strength when both associations are equal.

\subsection{Effect Type Conversions}

The E-value formula assumes a risk ratio. For other effect measures:

\begin{center}
\begin{tabular}{lll}
\toprule
Effect Type & Conversion & Notes \\
\midrule
RR & Direct & None \\
OR & $\approx \text{OR}$ & Rare outcomes \\
HR & $\approx \text{HR}$ & Proportional hazards \\
SMD & $\exp(0.91 d)$ & VanderWeele \\
ATE & $(p_0 + \text{ATE})/p_0$ & Needs $p_0$ \\
\bottomrule
\end{tabular}
\end{center}

\subsection{E-Value for Confidence Intervals}

Beyond the point estimate, we can compute an E-value for the confidence interval bound closest to the null:
\begin{itemize}
    \item For $\text{RR} > 1$: Use the lower CI bound
    \item For $\text{RR} < 1$: Use the upper CI bound
    \item If CI includes 1.0: $E_{\text{CI}} = 1.0$ (no confounding needed)
\end{itemize}

This answers: ``How much confounding would make the effect statistically non-significant?''

\subsection{Interpretation Scale}

\begin{center}
\begin{tabular}{lp{8cm}}
\toprule
E-Value & Interpretation \\
\midrule
$\approx 1.0$ & Not robust (effect is at or near null) \\
$1.0$--$1.5$ & Weakly robust (modest confounding could explain) \\
$1.5$--$2.5$ & Moderately robust \\
$2.5$--$4.0$ & Strongly robust \\
$> 4.0$ & Very robust to unmeasured confounding \\
\bottomrule
\end{tabular}
\end{center}

\begin{example}[E-Value Calculation]
Suppose we estimate RR = 2.0 (treatment doubles risk). The E-value is:
\begin{align}
E &= 2.0 + \sqrt{2.0 \cdot (2.0 - 1)} \\
  &= 2.0 + \sqrt{2.0} \\
  &\approx 3.41
\end{align}

Interpretation: An unmeasured confounder would need to be associated with both treatment and outcome by at least 3.41-fold (on the RR scale) to fully explain away this effect.
\end{example}

% ============================================================================
\section{Rosenbaum Bounds}
\label{sec:rosenbaum}
% ============================================================================

For matched studies, Rosenbaum bounds (1987, 2002) provide a different approach: systematically varying the assumed confounding level and observing when conclusions change.

\subsection{The Sensitivity Parameter $\Gamma$}

\begin{definition}[Sensitivity Parameter]
\label{def:gamma}
The parameter $\Gamma \geq 1$ represents how much more likely similar (matched) units could be to receive treatment due to an unmeasured confounder.
\end{definition}

Specifically, for matched pair $(i, j)$ with similar observed covariates:
\begin{equation}
\frac{1}{\Gamma} \leq \frac{P(T_i = 1 \mid X)}{P(T_j = 1 \mid X)} \leq \Gamma
\end{equation}

Interpretation:
\begin{itemize}
    \item $\Gamma = 1$: No unmeasured confounding (matched units equally likely to be treated)
    \item $\Gamma = 2$: One unit could be twice as likely to be treated as its match
    \item $\Gamma = 3$: One unit could be three times as likely
\end{itemize}

\subsection{Computing Bounds}

For matched pairs with outcomes $(Y_T, Y_C)$:

\textbf{Algorithm (Rosenbaum Bounds):}
\begin{enumerate}
    \item Compute pair differences: $D_i = Y_{\text{treated},i} - Y_{\text{control},i}$
    \item Rank the absolute differences $|D_i|$
    \item Compute the Wilcoxon signed-rank statistic $T^+ = \sum_{\{i: D_i > 0\}} R_i$
    \item For each $\Gamma$ value:
    \begin{itemize}
        \item Under $\Gamma$-confounding, each positive difference has probability $p_i \in [1/(1+\Gamma), \Gamma/(1+\Gamma)]$
        \item Compute upper and lower bounds on the expected $T^+$
        \item Use normal approximation for p-value bounds
    \end{itemize}
    \item Find $\Gamma^*$: smallest $\Gamma$ where upper bound p-value $> \alpha$
\end{enumerate}

\subsection{Mathematical Details}

Under $\Gamma$-confounding, the expectation and variance of $T^+$ are bounded:

\textbf{Expectation bounds:}
\begin{align}
E_{\text{lower}}[T^+] &= \sum_i R_i \cdot \frac{1}{1+\Gamma} \\
E_{\text{upper}}[T^+] &= \sum_i R_i \cdot \frac{\Gamma}{1+\Gamma}
\end{align}

\textbf{Variance bounds:}
\begin{equation}
\Var[T^+] = \sum_i R_i^2 \cdot p_i \cdot (1 - p_i)
\end{equation}

The upper bound p-value uses the distribution most favorable to the null hypothesis; the lower bound uses the distribution most favorable to the alternative.

\subsection{Critical $\Gamma$}

\begin{definition}[Critical $\Gamma$]
The critical $\Gamma^*$ is the smallest value of $\Gamma$ at which the upper bound p-value exceeds the significance level $\alpha$.
\end{definition}

This answers: ``How much confounding would be needed to make the treatment effect statistically non-significant?''

\subsection{Interpretation Scale}

\begin{center}
\begin{tabular}{lp{8cm}}
\toprule
$\Gamma^*$ & Interpretation \\
\midrule
$1.0$--$1.2$ & Very sensitive (small confounding overturns result) \\
$1.2$--$1.5$ & Sensitive \\
$1.5$--$2.0$ & Moderately robust \\
$2.0$--$3.0$ & Reasonably robust \\
$> 3.0$ & Quite robust to unmeasured confounding \\
\bottomrule
\end{tabular}
\end{center}

\begin{example}[Rosenbaum Bounds]
In a matched study with 50 pairs, suppose the Wilcoxon test yields p = 0.001. Rosenbaum bounds analysis finds $\Gamma^* = 2.85$.

Interpretation: The treatment effect would remain significant (at $\alpha = 0.05$) even if an unmeasured confounder made matched units differ in treatment probability by up to 2.85-fold. Only confounding stronger than this would overturn the finding.
\end{example}

% ============================================================================
\section{Implementation}
\label{sec:sens_implementation}
% ============================================================================

\subsection{E-Value Computation}

\begin{minted}{python}
def e_value(
    estimate: float,
    ci_lower: Optional[float] = None,
    ci_upper: Optional[float] = None,
    effect_type: Literal["rr", "or", "hr", "smd", "ate"] = "rr",
    baseline_risk: Optional[float] = None,
) -> EValueResult:
    """Compute E-value for sensitivity to unmeasured confounding.

    Parameters
    ----------
    estimate : float
        Point estimate of the effect.
    ci_lower, ci_upper : float, optional
        Confidence interval bounds.
    effect_type : str
        Type of effect: "rr", "or", "hr", "smd", "ate".
    baseline_risk : float, optional
        Required if effect_type="ate".

    Returns
    -------
    EValueResult
        e_value: E-value for point estimate
        e_value_ci: E-value for CI bound closest to null
        rr_equivalent: Effect converted to RR scale
        interpretation: Human-readable interpretation
    """
    # Convert to risk ratio scale
    if effect_type == "rr":
        rr = estimate
    elif effect_type == "smd":
        rr = np.exp(0.91 * estimate)  # VanderWeele approximation
    elif effect_type == "ate":
        rr = (baseline_risk + estimate) / baseline_risk
    else:
        rr = estimate  # OR, HR approximate RR

    # E-value formula (flip if protective effect)
    if rr >= 1:
        e_val = rr + np.sqrt(rr * (rr - 1))
    else:
        e_val = (1/rr) + np.sqrt((1/rr) * (1/rr - 1))

    # E-value for CI (bound closest to null)
    e_val_ci = 1.0  # Default if CI includes null
    if ci_lower is not None and ci_upper is not None:
        # Check if CI includes null (RR = 1)
        if not (ci_lower <= 1.0 <= ci_upper):
            if rr >= 1:
                e_val_ci = ci_lower + np.sqrt(ci_lower * (ci_lower - 1))
            else:
                e_val_ci = (1/ci_upper) + np.sqrt((1/ci_upper) * (1/ci_upper - 1))

    return EValueResult(e_value=e_val, e_value_ci=e_val_ci, ...)
\end{minted}

\subsection{Rosenbaum Bounds Computation}

\begin{minted}{python}
def rosenbaum_bounds(
    treated_outcomes: np.ndarray,
    control_outcomes: np.ndarray,
    gamma_range: Tuple[float, float] = (1.0, 3.0),
    n_gamma: int = 20,
    alpha: float = 0.05,
) -> RosenbaumResult:
    """Compute Rosenbaum sensitivity bounds for matched pairs.

    Parameters
    ----------
    treated_outcomes : np.ndarray
        Outcomes for treated units in matched pairs.
    control_outcomes : np.ndarray
        Outcomes for matched control units.
    gamma_range : Tuple[float, float]
        Range of Gamma values to evaluate.
    alpha : float
        Significance level for gamma_critical.

    Returns
    -------
    RosenbaumResult
        gamma_values: Array of Gamma values evaluated
        p_upper: Upper bound p-values at each Gamma
        p_lower: Lower bound p-values at each Gamma
        gamma_critical: Smallest Gamma where p_upper > alpha
    """
    # Compute pair differences
    differences = treated_outcomes - control_outcomes

    # Wilcoxon signed-rank statistic
    nonzero = differences != 0
    abs_diffs = np.abs(differences[nonzero])
    ranks = stats.rankdata(abs_diffs)
    t_plus = np.sum(ranks[differences[nonzero] > 0])

    # Evaluate at each Gamma
    gamma_values = np.linspace(gamma_range[0], gamma_range[1], n_gamma)
    p_upper = np.zeros(n_gamma)

    for i, gamma in enumerate(gamma_values):
        # Probability bounds under Gamma-confounding
        p_low = 1 / (1 + gamma)
        p_high = gamma / (1 + gamma)

        # Expectation and variance bounds
        e_upper = np.sum(ranks * p_high)
        var_upper = np.sum(ranks**2 * p_high * (1 - p_high))

        # Normal approximation p-value
        z = (t_plus - e_upper) / np.sqrt(var_upper)
        p_upper[i] = 1 - stats.norm.cdf(z)

    # Find critical Gamma
    critical_idx = np.where(p_upper > alpha)[0]
    gamma_critical = gamma_values[critical_idx[0]] if len(critical_idx) > 0 else None

    return RosenbaumResult(gamma_critical=gamma_critical, ...)
\end{minted}

% ============================================================================
\section{Comparing E-Values and Rosenbaum Bounds}
\label{sec:comparison}
% ============================================================================

\begin{center}
\begin{tabular}{lp{4.5cm}p{4.5cm}}
\toprule
Aspect & E-Value & Rosenbaum Bounds \\
\midrule
Input & Any point estimate & Matched pair outcomes \\
Output & Min confounding strength & Critical $\Gamma$ \\
Assumptions & Single confounder & Matched pairs design \\
Framework & Bounds on bias & Permutation inference \\
Best for & Any observational study & Matched studies \\
\bottomrule
\end{tabular}
\end{center}

\textbf{When to use each}:
\begin{itemize}
    \item \textbf{E-value}: Universal; use for any observational study
    \item \textbf{Rosenbaum}: More nuanced for matched designs; provides p-value trajectory
\end{itemize}

% ============================================================================
\section{Practical Guidance}
\label{sec:sens_guidance}
% ============================================================================

\subsection{Reporting Recommendations}

\begin{enumerate}
    \item \textbf{Always report sensitivity analysis} for observational studies
    \item Report both E-value for point estimate AND for CI bound
    \item Contextualize: compare to known confounder strengths in the field
    \item For matched studies, consider both E-value and Rosenbaum bounds
\end{enumerate}

\subsection{Interpreting Results}

\textbf{What makes a result ``robust''?}
\begin{itemize}
    \item Compare E-value to the strongest known confounders in your domain
    \item If E-value exceeds plausible confounding, the finding is robust
    \item If E-value is achievable by known confounders, be cautious
\end{itemize}

\textbf{Common benchmarks}:
\begin{itemize}
    \item Smoking and lung cancer: RR $\approx$ 10--30
    \item Many social/behavioral confounders: RR $\approx$ 1.5--3
    \item If E-value $> 4$, most measured confounders in social sciences would be insufficient
\end{itemize}

\subsection{Limitations}

\begin{itemize}
    \item Sensitivity analysis assumes a \emph{single} unmeasured confounder
    \item Multiple weak confounders could jointly explain the effect
    \item Does not account for measurement error in confounders
    \item A high E-value does not \emph{prove} causation
\end{itemize}

% ============================================================================
\section{Summary}
\label{sec:sens_summary}
% ============================================================================

\subsection{Key Results}

\begin{enumerate}
    \item \textbf{Unconfoundedness is untestable}: We can never verify all confounders are measured
    \item \textbf{E-value}: $E = \text{RR} + \sqrt{\text{RR}(\text{RR}-1)}$ gives minimum confounding to explain effect
    \item \textbf{Rosenbaum bounds}: $\Gamma^*$ gives critical confounding for matched studies
    \item \textbf{Always report}: Sensitivity analysis is essential for observational studies
\end{enumerate}

\subsection{Key References}

\begin{itemize}
    \item \cite{vanderweele2017sensitivity}: E-values
    \item \cite{rosenbaum2002observational}: Rosenbaum bounds
\end{itemize}



% ============================================================================
% PART III: SELECTION ON UNOBSERVABLES
% ============================================================================

\part{Selection on Unobservables}
\label{part:selection_unobservables}

% ============================================================================
% CHAPTER 9: INSTRUMENTAL VARIABLES
% ============================================================================

\chapter{Instrumental Variables}
\label{ch:iv}

% ============================================================================
\section{Introduction}
\label{sec:iv_intro}
% ============================================================================

When treatment assignment is correlated with unobserved determinants of the outcome, ordinary least squares estimates are biased. \emph{Instrumental variables} (IV) estimation provides a solution when we can find a variable that affects treatment but has no direct effect on the outcome.

\subsection{The Endogeneity Problem}

Consider the structural equation:
\begin{equation}
Y_i = \beta D_i + X_i'\gamma + \varepsilon_i
\label{eq:structural}
\end{equation}
where $D_i$ is the endogenous treatment and $\varepsilon_i$ is the error term. Endogeneity arises when:
\begin{equation}
\Cov(D_i, \varepsilon_i) \neq 0
\end{equation}

This correlation can arise from:
\begin{enumerate}
    \item \textbf{Omitted variables}: Unobserved confounders affect both $D$ and $Y$
    \item \textbf{Simultaneity}: $Y$ affects $D$ (reverse causality)
    \item \textbf{Measurement error}: $D$ is measured with error
\end{enumerate}

\subsection{The IV Solution}

An \emph{instrument} $Z$ provides exogenous variation in treatment. The key idea: use only the variation in $D$ that is driven by $Z$, which is uncorrelated with $\varepsilon$.

\begin{definition}[Instrumental Variable]
\label{def:instrument}
A variable $Z$ is a valid instrument for the effect of $D$ on $Y$ if:
\begin{enumerate}
    \item \textbf{Relevance}: $\Cov(Z, D) \neq 0$ (testable)
    \item \textbf{Exclusion}: $\Cov(Z, \varepsilon) = 0$ (maintained assumption)
\end{enumerate}
\end{definition}

\subsection{Historical Examples}

Classic instrumental variables include:
\begin{itemize}
    \item \textbf{Draft lottery}: Random draft numbers as instruments for military service
    \item \textbf{Quarter of birth}: Birth quarter as instrument for education \cite{angrist1991draft}
    \item \textbf{Geographic proximity}: Distance to college as instrument for education
\end{itemize}

\subsection{Chapter Roadmap}

This chapter develops:
\begin{itemize}
    \item Section~\ref{sec:iv_identification}: Formal identification conditions
    \item Section~\ref{sec:iv_wald}: The Wald estimator for just-identified models
    \item Section~\ref{sec:iv_2sls}: Two-stage least squares estimation
    \item Section~\ref{sec:iv_liml}: LIML and Fuller estimators
    \item Section~\ref{sec:iv_gmm}: GMM estimation with overidentification
    \item Section~\ref{sec:iv_inference}: Inference and standard errors
    \item Section~\ref{sec:iv_implementation}: Python implementation
    \item Section~\ref{sec:iv_validation}: Monte Carlo validation
\end{itemize}


% ============================================================================
\section{Identification}
\label{sec:iv_identification}
% ============================================================================

\subsection{The Structural and Reduced Form Equations}

The complete IV model consists of:

\textbf{Structural equation} (of interest):
\begin{equation}
Y_i = \beta D_i + X_i'\gamma + \varepsilon_i
\end{equation}

\textbf{First stage} (treatment equation):
\begin{equation}
D_i = Z_i'\pi + X_i'\delta + \nu_i
\label{eq:first_stage}
\end{equation}

\textbf{Reduced form} (outcome on instruments):
\begin{equation}
Y_i = Z_i'\rho + X_i'\theta + \eta_i
\label{eq:reduced_form}
\end{equation}

\subsection{Identification Conditions}

\begin{assumption}[Relevance]
\label{ass:relevance}
The instruments are correlated with the endogenous variable:
\begin{equation}
\E[Z_i D_i'] \neq 0
\end{equation}
Equivalently, $\pi \neq 0$ in the first stage equation.
\end{assumption}

\begin{assumption}[Exclusion Restriction]
\label{ass:exclusion}
The instruments are uncorrelated with the structural error:
\begin{equation}
\E[Z_i \varepsilon_i] = 0
\end{equation}
\end{assumption}

\begin{assumption}[Rank Condition]
\label{ass:rank}
For $q$ instruments and $p$ endogenous variables:
\begin{equation}
\text{rank}(\E[Z_i D_i']) = p
\end{equation}
\end{assumption}

\begin{proposition}[Order Condition]
A necessary condition for identification is $q \geq p$: at least as many instruments as endogenous variables.
\end{proposition}

\subsection{DAG Representation}

The IV assumptions can be represented graphically:
\begin{itemize}
    \item $Z \rightarrow D$: Relevance (instrument affects treatment)
    \item $Z \not\rightarrow Y$: Exclusion (no direct effect on outcome)
    \item $Z \not\leftarrow U \rightarrow Y$: No confounding of $Z$-$Y$ relationship
\end{itemize}

where $U$ represents unobserved confounders.


% ============================================================================
\section{The Wald Estimator}
\label{sec:iv_wald}
% ============================================================================

In the simplest case with one instrument and one endogenous variable (just-identified), the IV estimator takes a simple form.

\subsection{Binary Instrument Case}

With a binary instrument $Z \in \{0, 1\}$:

\begin{theorem}[Wald Estimator]
\label{thm:wald}
Under assumptions of relevance and exclusion, the causal effect is:
\begin{equation}
\beta = \frac{\E[Y | Z=1] - \E[Y | Z=0]}{\E[D | Z=1] - \E[D | Z=0]} = \frac{\text{Reduced form}}{\text{First stage}}
\label{eq:wald}
\end{equation}
\end{theorem}

\begin{proof}
From the structural equation $Y = \beta D + \varepsilon$ and exclusion $\E[\varepsilon | Z] = 0$:
\begin{align}
\E[Y | Z=1] - \E[Y | Z=0] &= \beta(\E[D | Z=1] - \E[D | Z=0])
\end{align}
Rearranging gives the result.
\end{proof}

\subsection{Interpretation}

The Wald estimator can be understood as:
\begin{equation}
\hat{\beta}_{\text{Wald}} = \frac{\hat{\rho}}{\hat{\pi}}
\end{equation}
where $\hat{\rho}$ is the reduced form coefficient and $\hat{\pi}$ is the first stage coefficient.

\begin{insight}
The Wald estimator divides the ``intent-to-treat'' effect (reduced form) by the ``compliance rate'' (first stage). This gives the effect of treatment on those who comply with the instrument.
\end{insight}


% ============================================================================
\section{Two-Stage Least Squares}
\label{sec:iv_2sls}
% ============================================================================

Two-stage least squares (2SLS) generalizes the Wald estimator to multiple instruments and covariates.

\subsection{The 2SLS Procedure}

\begin{definition}[Two-Stage Least Squares]
\label{def:2sls}
The 2SLS estimator is computed in two stages:
\begin{enumerate}
    \item \textbf{First stage}: Regress $D$ on $Z$ and $X$, obtain fitted values $\hat{D}$
    \item \textbf{Second stage}: Regress $Y$ on $\hat{D}$ and $X$
\end{enumerate}
\end{definition}

\subsection{Matrix Formulation}

Let $W = [D, X]$ be the matrix of endogenous and exogenous regressors, and $\tilde{Z} = [Z, X]$ the full instrument matrix. The 2SLS estimator is:

\begin{equation}
\hat{\beta}_{2\text{SLS}} = (W' P_{\tilde{Z}} W)^{-1} W' P_{\tilde{Z}} Y
\label{eq:2sls_matrix}
\end{equation}

where $P_{\tilde{Z}} = \tilde{Z}(\tilde{Z}'\tilde{Z})^{-1}\tilde{Z}'$ is the projection matrix.

\begin{theorem}[Consistency of 2SLS]
\label{thm:2sls_consistency}
Under the rank condition and exclusion restriction:
\begin{equation}
\hat{\beta}_{2\text{SLS}} \convp \beta
\end{equation}
\end{theorem}

\subsection{Variance Estimation}

The asymptotic variance of 2SLS depends on the error structure:

\textbf{Homoskedastic variance}:
\begin{equation}
\hat{V}_{2\text{SLS}} = \hat{\sigma}^2 (W' P_{\tilde{Z}} W)^{-1}
\end{equation}
where $\hat{\sigma}^2 = \frac{1}{n-k}\sum_i \hat{\varepsilon}_i^2$.

\textbf{Robust (heteroskedasticity-consistent) variance}:
\begin{equation}
\hat{V}_{\text{robust}} = (W' P_{\tilde{Z}} W)^{-1} \left(\sum_i \hat{\varepsilon}_i^2 \hat{w}_i \hat{w}_i'\right) (W' P_{\tilde{Z}} W)^{-1}
\end{equation}
where $\hat{w}_i = P_{\tilde{Z}} w_i$.

\subsection{Just-Identified vs. Over-Identified}

\begin{itemize}
    \item \textbf{Just-identified} ($q = p$): Exactly as many instruments as endogenous variables. 2SLS is unique but may be biased in small samples.
    \item \textbf{Over-identified} ($q > p$): More instruments than needed. 2SLS gains efficiency but requires testing overidentifying restrictions.
\end{itemize}


% ============================================================================
\section{LIML and Fuller Estimators}
\label{sec:iv_liml}
% ============================================================================

Limited Information Maximum Likelihood (LIML) and Fuller estimators address the finite-sample bias of 2SLS.

\subsection{k-Class Estimators}

The k-class family of estimators is defined as:
\begin{equation}
\hat{\beta}_k = (W'(I - k M_{\tilde{Z}})W)^{-1} W'(I - k M_{\tilde{Z}})Y
\label{eq:kclass}
\end{equation}
where $M_{\tilde{Z}} = I - P_{\tilde{Z}}$.

Special cases:
\begin{itemize}
    \item $k = 0$: OLS (ignores endogeneity)
    \item $k = 1$: 2SLS
    \item $k = \hat{\lambda}_{\min}$: LIML
\end{itemize}

\subsection{LIML Estimator}

\begin{definition}[LIML]
\label{def:liml}
LIML uses $k = \hat{\lambda}_{\min}$, the smallest eigenvalue of:
\begin{equation}
(W'M_X W)^{-1}(W'M_{[X,Z]} W)
\end{equation}
where $M_X = I - X(X'X)^{-1}X'$.
\end{definition}

\begin{proposition}[LIML Properties]
LIML is:
\begin{enumerate}
    \item Less biased than 2SLS with weak instruments
    \item Has higher variance than 2SLS
    \item Median-unbiased under normality
\end{enumerate}
\end{proposition}

\subsection{Fuller Estimator}

\begin{definition}[Fuller Estimator]
\label{def:fuller}
The Fuller estimator modifies LIML with a bias correction:
\begin{equation}
k_{\text{Fuller}} = \hat{\lambda}_{\min} - \frac{\alpha}{n - L}
\end{equation}
where $L$ is the number of excluded instruments and $\alpha$ is a tuning parameter (typically $\alpha = 1$ or $\alpha = 4$).
\end{definition}

Fuller-1 ($\alpha = 1$) provides a good bias-variance tradeoff across instrument strength regimes.

\subsection{When to Use Each Estimator}

\begin{table}[ht]
\centering
\caption{Estimator Selection Guide}
\begin{tabular}{lcc}
\toprule
\textbf{Scenario} & \textbf{Recommended} & \textbf{Rationale} \\
\midrule
Strong instruments ($F > 20$) & 2SLS & Efficient, minimal bias \\
Moderate ($10 < F < 20$) & Fuller-1 & Balance bias-variance \\
Weak ($F < 10$) & LIML or Fuller-4 & Less biased \\
\bottomrule
\end{tabular}
\end{table}


% ============================================================================
\section{GMM Estimation}
\label{sec:iv_gmm}
% ============================================================================

Generalized Method of Moments (GMM) provides an efficient estimator when there are more instruments than endogenous variables.

\subsection{Moment Conditions}

The exclusion restriction implies moment conditions:
\begin{equation}
\E[Z_i \varepsilon_i] = \E[Z_i(Y_i - D_i'\beta - X_i'\gamma)] = 0
\label{eq:moment_conditions}
\end{equation}

\subsection{GMM Objective}

The GMM estimator minimizes:
\begin{equation}
Q(\beta) = \left(\frac{1}{n}\sum_i Z_i \hat{\varepsilon}_i\right)' W \left(\frac{1}{n}\sum_i Z_i \hat{\varepsilon}_i\right)
\label{eq:gmm_objective}
\end{equation}
where $W$ is a positive definite weighting matrix.

\subsection{One-Step vs. Two-Step GMM}

\textbf{One-step GMM}: Uses $W = (Z'Z)^{-1}$, equivalent to 2SLS.

\textbf{Two-step GMM}:
\begin{enumerate}
    \item Compute one-step GMM residuals $\hat{\varepsilon}^{(1)}$
    \item Form optimal weighting: $\hat{W} = \left(\frac{1}{n}\sum_i \hat{\varepsilon}_i^{(1)2} Z_i Z_i'\right)^{-1}$
    \item Re-estimate with $\hat{W}$
\end{enumerate}

Two-step GMM is efficient under heteroskedasticity.

\subsection{Hansen J-Test}

\begin{definition}[Hansen J-Test]
\label{def:hansen_j}
The overidentification test statistic is:
\begin{equation}
J = n \cdot Q(\hat{\beta}) \sim \chi^2_{q-p}
\end{equation}
under the null that all instruments are valid.
\end{definition}

A significant J-statistic suggests at least one instrument may violate the exclusion restriction.


% ============================================================================
\section{Inference}
\label{sec:iv_inference}
% ============================================================================

\subsection{Standard Errors}

Three variance estimation approaches are available:

\begin{enumerate}
    \item \textbf{Homoskedastic}: Assumes $\Var(\varepsilon | Z) = \sigma^2$
    \item \textbf{Robust}: Allows heteroskedasticity
    \item \textbf{Clustered}: Allows within-cluster correlation
\end{enumerate}

\subsection{First-Stage F-Statistic}

\begin{definition}[First-Stage F-Statistic]
\label{def:first_stage_f}
The F-statistic tests joint significance of excluded instruments in the first stage:
\begin{equation}
F = \frac{(R^2_{\text{full}} - R^2_{\text{reduced}})/(q)}{(1 - R^2_{\text{full}})/(n - k)}
\end{equation}
\end{definition}

\begin{remark}[Rule of Thumb]
Staiger and Stock (1997) suggest $F > 10$ indicates instruments are not weak. This is a rough guideline; see Chapter~\ref{ch:weak_iv} for formal testing.
\end{remark}

\subsection{Confidence Intervals}

Standard Wald confidence intervals:
\begin{equation}
\hat{\beta} \pm z_{\alpha/2} \cdot \hat{\text{SE}}(\hat{\beta})
\end{equation}

These may have poor coverage with weak instruments; see Chapter~\ref{ch:weak_iv} for robust alternatives.


% ============================================================================
\section{Implementation}
\label{sec:iv_implementation}
% ============================================================================

\subsection{Python API}

The \texttt{causal\_inference.iv} module provides IV estimators:

\begin{minted}{python}
from causal_inference.iv import TwoStageLeastSquares, LIML, Fuller, GMM

# Prepare data
Y = outcome_data      # n x 1
D = treatment_data    # n x 1 (endogenous)
Z = instrument_data   # n x q
X = covariate_data    # n x p (optional)

# 2SLS estimation
tsls = TwoStageLeastSquares()
tsls.fit(Y, D, Z, X)
print(f"Estimate: {tsls.coef_[0]:.4f}")
print(f"SE: {tsls.se_[0]:.4f}")
print(f"First-stage F: {tsls.first_stage_f_:.2f}")

# LIML for potentially weak instruments
liml = LIML()
liml.fit(Y, D, Z, X)
print(f"LIML estimate: {liml.coef_[0]:.4f}")
print(f"LIML kappa: {liml.kappa_:.4f}")

# Fuller-1 (balanced bias-variance)
fuller = Fuller(alpha_param=1.0)
fuller.fit(Y, D, Z, X)

# Two-step GMM with J-test
gmm = GMM(steps='two')
gmm.fit(Y, D, Z, X)
print(f"J-statistic: {gmm.j_statistic_:.2f}")
print(f"J p-value: {gmm.j_pvalue_:.4f}")
\end{minted}

\subsection{Robust Standard Errors}

\begin{minted}{python}
# Heteroskedasticity-robust SEs
tsls_robust = TwoStageLeastSquares(inference='robust')
tsls_robust.fit(Y, D, Z, X)

# Cluster-robust SEs
tsls_cluster = TwoStageLeastSquares(
    inference='clustered',
    clusters=cluster_ids
)
tsls_cluster.fit(Y, D, Z, X)
\end{minted}

\subsection{Complete Example}

\begin{minted}{python}
import numpy as np
from causal_inference.iv import TwoStageLeastSquares
from causal_inference.iv.diagnostics import classify_instrument_strength

# Simulate IV data
np.random.seed(42)
n = 1000
Z = np.random.normal(0, 1, (n, 1))  # Instrument
U = np.random.normal(0, 1, n)        # Unobserved confounder
D = 0.5 * Z.flatten() + 0.5 * U + np.random.normal(0, 0.5, n)
Y = 2.0 * D + U + np.random.normal(0, 1, n)

# Estimate
tsls = TwoStageLeastSquares(inference='robust')
tsls.fit(Y, D, Z)

# Results
print(f"True effect: 2.0")
print(f"2SLS estimate: {tsls.coef_[0]:.4f}")
print(f"95% CI: [{tsls.ci_lower_[0]:.4f}, {tsls.ci_upper_[0]:.4f}]")
print(f"First-stage F: {tsls.first_stage_f_:.1f}")

# Classify instrument strength
strength = classify_instrument_strength(tsls.first_stage_f_, n_instruments=1)
print(f"Instrument strength: {strength}")
\end{minted}


% ============================================================================
\section{Monte Carlo Validation}
\label{sec:iv_validation}
% ============================================================================

\subsection{Data Generating Process}

We validate IV estimators using:
\begin{align}
Y_i &= \beta D_i + \varepsilon_i, \quad \beta = 2.0 \\
D_i &= \pi Z_i + \nu_i \\
(\varepsilon_i, \nu_i) &\sim N(0, \Sigma), \quad \Sigma = \begin{pmatrix} 1 & \rho \\ \rho & 1 \end{pmatrix}
\end{align}
where $\rho = 0.5$ creates endogeneity.

\subsection{Validation Results}

\begin{table}[ht]
\centering
\caption{Monte Carlo Results: IV Estimators ($n=500$, 5,000 replications)}
\label{tab:iv_mc}
\begin{tabular}{lccccc}
\toprule
\textbf{Estimator} & \textbf{True} & \textbf{Mean Est.} & \textbf{Bias} & \textbf{RMSE} & \textbf{Coverage} \\
\midrule
OLS (biased) & 2.0 & 2.50 & 0.50 & 0.51 & 0\% \\
2SLS & 2.0 & 2.01 & 0.01 & 0.12 & 94.8\% \\
LIML & 2.0 & 2.00 & 0.00 & 0.13 & 95.2\% \\
Fuller-1 & 2.0 & 2.00 & 0.00 & 0.12 & 95.0\% \\
GMM (2-step) & 2.0 & 2.01 & 0.01 & 0.12 & 94.9\% \\
\bottomrule
\end{tabular}
\end{table}

Key findings:
\begin{itemize}
    \item OLS is severely biased due to endogeneity
    \item All IV estimators recover the true effect
    \item Coverage rates are near nominal 95\%
    \item 2SLS, Fuller-1, and GMM have similar efficiency
\end{itemize}


% ============================================================================
\section{Practical Guidance}
\label{sec:iv_guidance}
% ============================================================================

\subsection{Assessing Instrument Validity}

\begin{enumerate}
    \item \textbf{Relevance}: Check first-stage F-statistic ($F > 10$ rule of thumb)
    \item \textbf{Exclusion}: Cannot be tested directly---requires economic reasoning
    \item \textbf{Monotonicity} (for LATE): Treatment response is monotone in instrument
\end{enumerate}

\subsection{Common Mistakes}

\begin{itemize}
    \item \textbf{Weak instruments}: Using instruments with low predictive power leads to biased estimates
    \item \textbf{Invalid instruments}: Direct effect of $Z$ on $Y$ violates exclusion
    \item \textbf{Heterogeneous effects}: IV estimates LATE, not ATE, when effects vary
\end{itemize}

\subsection{Reporting Recommendations}

Always report:
\begin{enumerate}
    \item First-stage F-statistic
    \item Both 2SLS and at least one alternative (LIML or Fuller)
    \item Robust standard errors
    \item Hansen J-test if over-identified
\end{enumerate}


% ============================================================================
\section{Summary}
\label{sec:iv_summary}
% ============================================================================

\textbf{Key Results:}
\begin{enumerate}
    \item IV estimation requires relevance (testable) and exclusion (assumed)
    \item 2SLS is the standard estimator for strong instruments
    \item LIML and Fuller reduce bias when instruments are weak
    \item GMM is efficient with heteroskedasticity and multiple instruments
    \item Hansen J-test checks overidentifying restrictions
\end{enumerate}

\textbf{Key References:}
\begin{itemize}
    \item Angrist \& Krueger (1991): Quarter of birth instruments
    \item Stock \& Yogo (2005): Weak instrument testing
    \item Bound, Jaeger \& Baker (1995): Weak instrument bias
\end{itemize}

\textbf{Next Chapter:} Chapter~\ref{ch:weak_iv} develops diagnostics and robust inference for weak instruments.

% Chapter 10: Weak Instruments
% Part III: Selection on Unobservables

\chapter{Weak Instruments}
\label{ch:weak_iv}

\section{Introduction}
\label{sec:weak_iv_intro}

The instrumental variables estimators developed in Chapter~\ref{ch:iv} assume that instruments are \emph{relevant}---strongly correlated with the endogenous regressor. When this assumption fails, IV estimators exhibit severe finite-sample problems:

\begin{enumerate}
    \item \textbf{Bias toward OLS}: 2SLS bias approaches OLS bias as instruments weaken
    \item \textbf{Size distortion}: t-tests reject too often (actual size exceeds nominal)
    \item \textbf{Wide confidence intervals}: Standard errors understate uncertainty
\end{enumerate}

The rule-of-thumb ``F > 10'' popularized by Staiger and Stock (1997) arose from recognizing these problems. This chapter develops formal diagnostics and inference procedures that remain valid even with weak instruments.

\begin{definition}[Weak Instruments]
\label{def:weak_iv}
Instruments are \emph{weak} when the concentration parameter
\[
\mu^2 = \frac{\pi'Z'Z\pi}{\sigma^2_\nu}
\]
is small relative to sample size, where $\pi$ is the first-stage coefficient vector and $\sigma^2_\nu$ is the first-stage error variance. Equivalently, instruments are weak when the first-stage F-statistic is ``small,'' typically below the Stock-Yogo critical values.
\end{definition}

\begin{insight}
The weak instrument problem is fundamentally about the \emph{signal-to-noise ratio} in the first stage. With weak instruments, the variation in $D$ explained by $Z$ is dominated by noise, making it difficult to identify the causal effect.
\end{insight}

\textbf{Chapter Roadmap:}
\begin{itemize}
    \item Section~\ref{sec:diagnosing_weak}: Stock-Yogo critical values and Cragg-Donald statistic
    \item Section~\ref{sec:bias_distortion}: Finite-sample bias and size distortion
    \item Section~\ref{sec:ar_test}: Anderson-Rubin weak-IV-robust test
    \item Section~\ref{sec:clr_test}: Conditional Likelihood Ratio test
    \item Section~\ref{sec:weak_protocol}: Recommended testing protocol
    \item Section~\ref{sec:weak_implementation}: Python implementation
    \item Section~\ref{sec:weak_validation}: Monte Carlo validation
\end{itemize}


\section{Diagnosing Weak Instruments}
\label{sec:diagnosing_weak}

\subsection{First-Stage F-Statistic}

The most common diagnostic is the first-stage F-statistic from the regression:
\[
D_i = Z_i'\pi + X_i'\delta + \nu_i
\]

The F-statistic tests $H_0: \pi = 0$ (instruments are irrelevant):
\[
F = \frac{(SSR_{\text{restricted}} - SSR_{\text{full}})/q}{SSR_{\text{full}}/(n - q - k - 1)}
\]
where $q$ is the number of instruments and $k$ is the number of controls.

\begin{theorem}[Staiger-Stock Approximation]
\label{thm:staiger_stock}
Under weak-instrument asymptotics where $\mu^2/n \to c$ as $n \to \infty$, the 2SLS estimator satisfies:
\[
\hat{\beta}_{2SLS} - \beta \convd \frac{\text{Normal} + \text{bias term}}{F - \text{adjustment}}
\]
The bias term is proportional to $(1/F)$ times the OLS bias, explaining why 2SLS bias approaches OLS bias as $F \to 1$.
\end{theorem}

\subsection{Stock-Yogo Critical Values}

Stock and Yogo (2005) provide critical values for testing weak instruments based on two criteria:

\begin{enumerate}
    \item \textbf{Maximal relative bias}: Reject weak instruments if relative bias of 2SLS (compared to OLS) is at most $r\%$
    \item \textbf{Maximal size distortion}: Reject weak instruments if actual size of 5\% Wald test is at most $r\%$
\end{enumerate}

\begin{table}[htbp]
\centering
\caption{Stock-Yogo Critical Values (10\% Maximal Bias)}
\label{tab:stock_yogo}
\begin{tabular}{lcccccc}
\toprule
Instruments ($q$) & \multicolumn{6}{c}{Endogenous Regressors ($p$)} \\
\cmidrule(lr){2-7}
& 1 & 2 & 3 & 4 & 5 & 10 \\
\midrule
1 & 16.38 & --- & --- & --- & --- & --- \\
2 & 19.93 & 11.04 & --- & --- & --- & --- \\
3 & 22.30 & 13.43 & 9.53 & --- & --- & --- \\
4 & 24.58 & 15.09 & 10.83 & 8.78 & --- & --- \\
5 & 26.87 & 16.76 & 12.20 & 9.93 & 8.40 & --- \\
10 & 38.54 & --- & --- & --- & --- & --- \\
\bottomrule
\end{tabular}
\end{table}

\begin{remark}
The critical value 16.38 for $(q=1, p=1)$ is much higher than the rule-of-thumb ``F > 10.'' The informal rule provides a lower bar; Stock-Yogo values provide rigorous thresholds.
\end{remark}

\subsection{Instrument Strength Classification}

Our implementation classifies instruments as:

\begin{definition}[Instrument Strength Classification]
\label{def:iv_strength}
Given first-stage F-statistic $F$ and Stock-Yogo critical value $cv$:
\begin{align*}
\text{Strong:} & \quad F > cv \\
\text{Weak:} & \quad 10 < F \leq cv \\
\text{Very Weak:} & \quad F \leq 10
\end{align*}
If no critical value is tabulated, use rule-of-thumb with $cv = 20$.
\end{definition}

\subsection{Cragg-Donald Statistic}

With multiple endogenous regressors ($p > 1$), the first-stage F-statistic is insufficient. The Cragg-Donald (1993) statistic generalizes the weak instrument test:

\begin{definition}[Cragg-Donald Statistic]
\label{def:cragg_donald}
Let $\hat{\Pi}$ be the $q \times p$ matrix of first-stage coefficients. The Cragg-Donald statistic is:
\[
CD = \frac{n - k - q}{q} \times \lambda_{\min}\left(\frac{\hat{\Pi}'(Z'Z/n)\hat{\Pi}}{\hat{\sigma}^2}\right)
\]
where $\lambda_{\min}$ denotes the minimum eigenvalue.
\end{definition}

For just-identified cases ($q = p = 1$), the Cragg-Donald statistic reduces to the first-stage F-statistic.


\section{Bias and Size Distortion}
\label{sec:bias_distortion}

\subsection{Finite-Sample Bias of 2SLS}

The finite-sample bias of 2SLS relative to OLS provides intuition for weak instrument problems:

\begin{theorem}[Relative Bias]
\label{thm:relative_bias}
The approximate relative bias of 2SLS compared to OLS is:
\[
\frac{\E[\hat{\beta}_{2SLS} - \beta]}{\E[\hat{\beta}_{OLS} - \beta]} \approx \frac{1}{F}
\]
where $F$ is the first-stage F-statistic.
\end{theorem}

\begin{proof}[Proof Sketch]
Bound, Jaeger, and Baker (1995) show that in finite samples:
\[
\E[\hat{\beta}_{2SLS}] \approx \beta + \frac{\sigma_{\nu\epsilon}}{\sigma^2_\nu} \times \frac{q}{F}
\]
where $\sigma_{\nu\epsilon}$ is the covariance between first-stage and structural errors. Since the OLS bias is $\sigma_{\nu\epsilon}/\sigma^2_\nu$, the relative bias follows.
\end{proof}

\begin{table}[htbp]
\centering
\caption{Relative Bias by Instrument Strength}
\label{tab:relative_bias}
\begin{tabular}{lcc}
\toprule
F-Statistic & Relative Bias & Interpretation \\
\midrule
F = 100 & 1\% & Nearly unbiased \\
F = 50 & 2\% & Excellent \\
F = 20 & 5\% & Good \\
F = 10 & 10\% & Marginal \\
F = 5 & 20\% & Problematic \\
F = 2 & 50\% & Severe \\
\bottomrule
\end{tabular}
\end{table}

\subsection{Size Distortion of Wald Tests}

With weak instruments, standard Wald tests (using $t = \hat{\beta}/\hat{se}$) reject too often:

\begin{proposition}[Size Distortion]
\label{prop:size_distortion}
For a nominal 5\% Wald test:
\begin{itemize}
    \item F = 10: Actual size $\approx$ 15-20\%
    \item F = 5: Actual size $\approx$ 25-30\%
    \item F = 2: Actual size can exceed 50\%
\end{itemize}
\end{proposition}

This size distortion motivates \emph{weak-IV-robust} inference procedures that maintain correct size regardless of instrument strength.


\section{Anderson-Rubin Test}
\label{sec:ar_test}

The Anderson-Rubin (1949) test provides valid inference regardless of instrument strength.

\subsection{Test Construction}

\begin{definition}[Anderson-Rubin Test]
\label{def:ar_test}
To test $H_0: \beta = \beta_0$, the AR statistic is:
\[
AR(\beta_0) = \frac{1}{q} \times \frac{\tilde{u}' P_Z \tilde{u}}{\hat{\sigma}^2}
\]
where:
\begin{itemize}
    \item $\tilde{u} = Y - \beta_0 D$ are reduced-form residuals under $H_0$
    \item $P_Z = Z(Z'Z)^{-1}Z'$ is the projection matrix onto instruments
    \item $\hat{\sigma}^2$ is the residual variance from unrestricted reduced form
\end{itemize}
Under $H_0$, $AR(\beta_0) \sim \chi^2(q)/q \approx F(q, n-q-k-1)$.
\end{definition}

\begin{theorem}[AR Test Validity]
\label{thm:ar_validity}
The Anderson-Rubin test has exact size $\alpha$ under $H_0: \beta = \beta_0$ for any value of the first-stage coefficient $\pi$, including $\pi = 0$ (completely irrelevant instruments).
\end{theorem}

\begin{proof}
Under $H_0$, the reduced-form residuals $\tilde{u} = Y - \beta_0 D = \epsilon + (\beta - \beta_0)D = \epsilon$ are structural errors. The AR statistic tests whether $\E[Z'\epsilon] = 0$ (exclusion restriction), which holds by assumption regardless of first-stage strength.
\end{proof}

\subsection{Confidence Intervals by Test Inversion}

The AR confidence interval inverts the test:
\[
CI_{AR}(\alpha) = \{\beta : AR(\beta) \leq \chi^2_{1-\alpha}(q)\}
\]

\begin{algorithm}[htbp]
\caption{Anderson-Rubin Confidence Interval}
\label{alg:ar_ci}
\begin{algorithmic}[1]
\REQUIRE Data $(Y, D, Z, X)$, significance level $\alpha$, grid size $n_{grid}$
\STATE Compute 2SLS estimate $\hat{\beta}_{2SLS}$ and standard error $\hat{se}$
\STATE Define grid: $\beta_{grid} \in [\hat{\beta}_{2SLS} - 20\hat{se}, \hat{\beta}_{2SLS} + 20\hat{se}]$
\STATE Compute $P_Z$ (projection matrix onto residualized instruments)
\STATE Compute $\hat{\sigma}^2$ from unrestricted reduced form
\FOR{each $\beta_0$ in $\beta_{grid}$}
    \STATE $\tilde{u} \gets Y - \beta_0 D$
    \STATE $AR(\beta_0) \gets (1/q) \times \tilde{u}' P_Z \tilde{u} / \hat{\sigma}^2$
\ENDFOR
\STATE $CI \gets \{\beta_0 : AR(\beta_0) \leq \chi^2_{1-\alpha}(q)\}$
\RETURN $(\min(CI), \max(CI))$
\end{algorithmic}
\end{algorithm}

\begin{remark}
With controls $X$, instruments must be residualized: $Z_\perp = Z - X(X'X)^{-1}X'Z$. Using raw $Z$ in the projection matrix yields invalid inference when instruments correlate with controls.
\end{remark}

\subsection{Properties of the AR Test}

\textbf{Advantages:}
\begin{itemize}
    \item Valid for \emph{any} instrument strength, including completely irrelevant
    \item Exact finite-sample distribution under homoskedasticity
    \item Simple to compute and interpret
\end{itemize}

\textbf{Disadvantages:}
\begin{itemize}
    \item Conservative (low power) when instruments are strong
    \item Confidence sets can be unbounded with very weak instruments
    \item Less powerful than standard Wald test when instruments are strong
\end{itemize}


\section{Conditional Likelihood Ratio Test}
\label{sec:clr_test}

Moreira (2003) developed the Conditional Likelihood Ratio (CLR) test, which is more powerful than Anderson-Rubin while maintaining validity under weak instruments.

\subsection{Motivation}

The AR test is conservative because it does not condition on first-stage information. The CLR test conditions on a sufficient statistic for the nuisance parameter (reduced-form coefficients), achieving optimal power.

\subsection{Test Statistic}

\begin{definition}[CLR Test Statistic]
\label{def:clr}
The CLR statistic decomposes into two quadratic forms:
\begin{align*}
Q_S &= \text{(reduced-form in null direction)} \\
Q_T &= \text{(first-stage strength)}
\end{align*}
The likelihood ratio statistic is:
\[
LR = \frac{1}{2}\left(Q_S - Q_T + \sqrt{(Q_S + Q_T)^2 - 4(Q_S Q_T - Q_{TS}^2)}\right)
\]
where $Q_{TS}$ is a cross term.
\end{definition}

The CLR test rejects when $LR > cv(\alpha | Q_T)$, where the critical value depends on the conditioning variable $Q_T$.

\begin{theorem}[CLR Optimality]
\label{thm:clr_optimal}
Among tests that:
\begin{enumerate}
    \item Are invariant to rotations in the instrument space
    \item Have correct size under weak instruments
    \item Are similar (size equals nominal level exactly)
\end{enumerate}
the CLR test is the uniformly most powerful unbiased test (Andrews, Moreira, Stock 2006).
\end{theorem}

\subsection{CLR vs. Anderson-Rubin}

\begin{table}[htbp]
\centering
\caption{Comparison of Weak-IV-Robust Tests}
\label{tab:ar_vs_clr}
\begin{tabular}{lcc}
\toprule
Property & Anderson-Rubin & CLR \\
\midrule
Valid under weak IV & Yes & Yes \\
Optimal power & No & Yes \\
Computation & Simple & Complex \\
Critical values & $\chi^2$ tables & Simulation/tables \\
CI interpretation & Direct & Direct \\
\bottomrule
\end{tabular}
\end{table}

\begin{insight}
Use Anderson-Rubin for quick robustness checks. Use CLR when power matters and you have access to appropriate critical values.
\end{insight}


\section{Recommended Testing Protocol}
\label{sec:weak_protocol}

\subsection{Decision Tree}

\begin{enumerate}
    \item \textbf{Compute first-stage F-statistic}
    \item \textbf{Compare to Stock-Yogo critical value}:
    \begin{itemize}
        \item $F > cv$: Instruments are strong. Use standard 2SLS inference.
        \item $10 < F \leq cv$: Instruments are weak. Use LIML + AR/CLR.
        \item $F \leq 10$: Instruments are very weak. Use AR confidence intervals only.
    \end{itemize}
    \item \textbf{Report both intervals}: When instruments may be weak, report both standard (2SLS) and robust (AR) confidence intervals.
\end{enumerate}

\subsection{Reporting Recommendations}

For any IV analysis, report:
\begin{itemize}
    \item First-stage F-statistic with Stock-Yogo critical value
    \item 2SLS point estimate and standard CI
    \item Anderson-Rubin confidence interval (especially if $F < 20$)
    \item Cragg-Donald statistic (if multiple endogenous regressors)
\end{itemize}


\section{Implementation}
\label{sec:weak_implementation}

\subsection{Stock-Yogo Classification}

\begin{minted}{python}
from causal_inference.iv.diagnostics import (
    classify_instrument_strength,
    cragg_donald_statistic,
    anderson_rubin_test,
    weak_instrument_summary,
)

# Classify instrument strength using Stock-Yogo critical values
classification, critical_value, interpretation = classify_instrument_strength(
    f_statistic=12.5,
    n_instruments=1,
    n_endogenous=1,
    bias_threshold="10pct",  # 10%, 15%, or 20% maximal bias
)

print(f"Classification: {classification}")
# Output: Classification: weak
print(f"Critical Value: {critical_value}")
# Output: Critical Value: 16.38
print(f"Interpretation: {interpretation}")
# Output: Instruments fail Stock-Yogo test (F=12.50 <= 16.38)
\end{minted}

\subsection{Cragg-Donald Statistic}

\begin{minted}{python}
import numpy as np

# Multiple endogenous regressors
n = 1000
Z = np.random.normal(0, 1, (n, 3))  # 3 instruments
D = Z @ np.array([[0.3, 0.2],
                   [0.2, 0.3],
                   [0.1, 0.1]]) + np.random.normal(0, 1, (n, 2))
Y = D @ [0.5, 0.3] + np.random.normal(0, 1, n)

cd_stat = cragg_donald_statistic(Y, D, Z)
print(f"Cragg-Donald statistic: {cd_stat:.2f}")
# Compare to Stock-Yogo critical values for (q=3, p=2)
\end{minted}

\subsection{Anderson-Rubin Test}

\begin{minted}{python}
# Simulate weak instrument scenario
n = 500
Z = np.random.normal(0, 1, n)
D = 0.1 * Z + np.random.normal(0, 1, n)  # Weak: pi = 0.1
Y = 2.0 * D + np.random.normal(0, 1, n)  # True beta = 2.0

# AR test for H0: beta = 0
ar_stat, p_value, ar_ci = anderson_rubin_test(
    Y, D, Z.reshape(-1, 1),
    X=None,
    alpha=0.05,
    n_grid=200,
)

print(f"AR statistic: {ar_stat:.3f}")
print(f"P-value (H0: beta=0): {p_value:.4f}")
print(f"AR 95% CI: [{ar_ci[0]:.3f}, {ar_ci[1]:.3f}]")
# Note: AR CI will be much wider than standard 2SLS CI
\end{minted}

\subsection{Comprehensive Diagnostic Summary}

\begin{minted}{python}
# Generate comprehensive summary
summary = weak_instrument_summary(
    f_statistic=8.5,
    n_instruments=2,
    n_endogenous=1,
    cragg_donald=8.5,
    ar_ci=(0.5, 3.5),
)

print(summary.to_string(index=False))
# Output:
#                    Diagnostic    Value  Interpretation
#        First-Stage F-Statistic     8.50  Instruments are severely weak
#  Stock-Yogo Critical Value ...    19.93  Reject weak IV if F > 19.93
#        Cragg-Donald Statistic     8.50  Multivariate weak IV test
#           Anderson-Rubin 95% CI  [0.500, 3.500]  Robust to weak instruments
#                  Recommendation         Instruments are very weak...
\end{minted}


\section{Monte Carlo Validation}
\label{sec:weak_validation}

\subsection{Simulation Design}

We validate weak-IV diagnostics across instrument strengths:

\textbf{Data Generating Process:}
\begin{align*}
Y_i &= \beta D_i + \epsilon_i, \quad \beta = 2.0 \\
D_i &= \pi Z_i + \nu_i \\
\begin{pmatrix} \epsilon_i \\ \nu_i \end{pmatrix} &\sim \mathcal{N}\left(\mathbf{0}, \begin{pmatrix} 1 & 0.5 \\ 0.5 & 1 \end{pmatrix}\right)
\end{align*}

We vary $\pi$ to achieve different instrument strengths:
\begin{itemize}
    \item Strong: $\pi = 0.5$ (F $\approx$ 50)
    \item Moderate: $\pi = 0.2$ (F $\approx$ 10)
    \item Weak: $\pi = 0.1$ (F $\approx$ 3)
\end{itemize}

\subsection{Validation Results}

\begin{table}[htbp]
\centering
\caption{Weak IV Monte Carlo Results (5,000 replications, n=500)}
\label{tab:weak_iv_mc}
\begin{tabular}{lccccc}
\toprule
Method & F-stat & Bias & Coverage & Rejection Rate & CI Width \\
\midrule
\multicolumn{6}{l}{\textit{Strong Instruments ($\pi = 0.5$, F $\approx$ 50)}} \\
2SLS Standard CI & 50.2 & 0.02 & 95.1\% & 5.2\% & 0.58 \\
Anderson-Rubin CI & 50.2 & --- & 95.8\% & 4.8\% & 0.72 \\
LIML & 50.2 & 0.01 & 95.3\% & 5.0\% & 0.58 \\
\midrule
\multicolumn{6}{l}{\textit{Moderate Instruments ($\pi = 0.2$, F $\approx$ 10)}} \\
2SLS Standard CI & 10.1 & 0.15 & 88.2\% & 12.5\% & 1.45 \\
Anderson-Rubin CI & 10.1 & --- & 95.4\% & 4.9\% & 2.31 \\
LIML & 10.1 & 0.08 & 93.1\% & 7.2\% & 1.52 \\
\midrule
\multicolumn{6}{l}{\textit{Weak Instruments ($\pi = 0.1$, F $\approx$ 3)}} \\
2SLS Standard CI & 3.2 & 0.62 & 71.3\% & 32.1\% & 3.84 \\
Anderson-Rubin CI & 3.2 & --- & 95.2\% & 5.1\% & 8.52 \\
LIML & 3.2 & 0.31 & 84.6\% & 18.4\% & 4.21 \\
\bottomrule
\end{tabular}
\end{table}

\textbf{Key Findings:}
\begin{enumerate}
    \item \textbf{Strong instruments}: All methods perform well. AR is slightly conservative.
    \item \textbf{Moderate instruments}: 2SLS coverage drops to 88\%; AR maintains 95\%.
    \item \textbf{Weak instruments}: 2SLS coverage severely degraded (71\%); AR maintains validity but CIs are very wide.
\end{enumerate}

\begin{insight}
Anderson-Rubin maintains nominal coverage regardless of instrument strength, but the price is wider confidence intervals. This represents an honest assessment of estimation uncertainty with weak instruments.
\end{insight}


\section{Practical Guidance}
\label{sec:weak_practical}

\subsection{When Instruments Are Weak}

\begin{enumerate}
    \item \textbf{Find better instruments}: The best solution is stronger instruments
    \item \textbf{Use LIML/Fuller}: Less biased than 2SLS
    \item \textbf{Report AR confidence intervals}: Provides honest uncertainty
    \item \textbf{Consider bounds}: Partial identification may be more informative
\end{enumerate}

\subsection{Common Mistakes}

\begin{enumerate}
    \item \textbf{Ignoring weak IV warning}: Always check F-statistic
    \item \textbf{Using F > 10 rule}: Stock-Yogo critical values are higher (16.38)
    \item \textbf{Reporting only 2SLS CI}: With weak IV, AR CI is essential
    \item \textbf{Multiple first stages}: Use Cragg-Donald for $p > 1$
\end{enumerate}

\subsection{Reporting Checklist}

For any IV analysis:
\begin{itemize}
    \item[$\square$] First-stage F-statistic
    \item[$\square$] Stock-Yogo critical value and comparison
    \item[$\square$] 2SLS estimate with standard errors
    \item[$\square$] LIML estimate (if F < 20)
    \item[$\square$] Anderson-Rubin 95\% CI (if F < 20)
    \item[$\square$] Discussion of instrument validity
\end{itemize}


\section{Summary}
\label{sec:weak_iv_summary}

This chapter developed diagnostics and inference procedures for weak instruments:

\begin{enumerate}
    \item \textbf{Diagnosis}: Stock-Yogo critical values provide rigorous thresholds beyond ``F > 10''

    \item \textbf{Consequences}: Weak instruments cause 2SLS bias toward OLS and severe size distortion of Wald tests

    \item \textbf{Anderson-Rubin}: Valid inference for any instrument strength via test inversion; conservative but robust

    \item \textbf{CLR}: Optimal power among weak-IV-robust tests; more complex computation

    \item \textbf{Protocol}: F > cv $\to$ standard inference; F $\leq$ cv $\to$ robust inference
\end{enumerate}

\begin{table}[htbp]
\centering
\caption{Chapter Notation Summary}
\label{tab:weak_notation}
\begin{tabular}{ll}
\toprule
Symbol & Meaning \\
\midrule
$F$ & First-stage F-statistic \\
$cv$ & Stock-Yogo critical value \\
$\mu^2$ & Concentration parameter \\
$AR(\beta_0)$ & Anderson-Rubin test statistic \\
$LR$ & Conditional Likelihood Ratio statistic \\
$Q_S, Q_T$ & CLR quadratic forms \\
$CD$ & Cragg-Donald statistic \\
\bottomrule
\end{tabular}
\end{table}

\textbf{Key Result}: When instruments are weak, standard 2SLS inference is unreliable. Anderson-Rubin confidence intervals provide valid (though wider) inference regardless of instrument strength.

\textbf{Next}: Chapter~\ref{ch:control_function} develops the control function approach as an alternative to 2SLS that extends naturally to nonlinear models.

% Chapter 11: Control Function Approach
% Part III: Selection on Unobservables

\chapter{Control Function Approach}
\label{ch:control_function}

\section{Introduction}
\label{sec:cf_intro}

The control function approach is an alternative to 2SLS that offers three key advantages:

\begin{enumerate}
    \item \textbf{Explicit endogeneity test}: The coefficient on the control variable directly tests for endogeneity
    \item \textbf{Extension to nonlinear models}: Works for probit/logit where 2SLS is invalid
    \item \textbf{Flexibility}: Handles heteroskedasticity and other complications naturally
\end{enumerate}

The core idea: instead of instrumenting for $D$, we \emph{control for} the source of endogeneity by including the first-stage residual as an additional regressor.

\begin{definition}[Control Function]
\label{def:control_function}
The control function is the first-stage residual:
\[
\hat{\nu}_i = D_i - \E[D_i | Z_i, X_i]
\]
Including $\hat{\nu}$ in the outcome equation ``purges'' the endogeneity, making the treatment coefficient consistently estimable.
\end{definition}

\begin{insight}
The control function approach views endogeneity as an \emph{omitted variable problem}. By including $\hat{\nu}$ as a regressor, we control for the correlation between $D$ and the structural error, eliminating the bias.
\end{insight}

\textbf{Chapter Roadmap:}
\begin{itemize}
    \item Section~\ref{sec:cf_intuition}: Intuition and identification
    \item Section~\ref{sec:cf_linear}: Linear control function (equivalent to 2SLS)
    \item Section~\ref{sec:cf_nonlinear}: Nonlinear control function (probit/logit)
    \item Section~\ref{sec:cf_inference}: Standard errors and testing
    \item Section~\ref{sec:cf_implementation}: Python implementation
    \item Section~\ref{sec:cf_validation}: Monte Carlo validation
\end{itemize}


\section{The Control Function Idea}
\label{sec:cf_intuition}

\subsection{Endogeneity as Omitted Variable}

Consider the structural model:
\begin{align}
Y_i &= \beta D_i + X_i'\gamma + \epsilon_i \label{eq:cf_structural} \\
D_i &= Z_i'\pi + X_i'\delta + \nu_i \label{eq:cf_first_stage}
\end{align}

Endogeneity arises when $\Cov(D, \epsilon) \neq 0$. This occurs because:
\[
\Cov(D, \epsilon) = \Cov(\nu, \epsilon) \equiv \rho \sigma_\nu \sigma_\epsilon
\]
where $\rho$ captures the correlation between first-stage and structural errors.

\subsection{The Key Decomposition}

We can decompose the structural error as:
\begin{equation}
\epsilon_i = \rho \frac{\sigma_\epsilon}{\sigma_\nu} \nu_i + u_i
\label{eq:cf_decomposition}
\end{equation}
where $u_i$ is orthogonal to $\nu_i$ (and hence to $D_i$).

Substituting into the structural equation:
\begin{align}
Y_i &= \beta D_i + X_i'\gamma + \rho \frac{\sigma_\epsilon}{\sigma_\nu} \nu_i + u_i \\
&= \beta D_i + X_i'\gamma + \lambda \nu_i + u_i
\label{eq:cf_augmented}
\end{align}
where $\lambda = \rho \sigma_\epsilon / \sigma_\nu$ is the control coefficient.

\begin{theorem}[Control Function Identification]
\label{thm:cf_identification}
Under the assumptions:
\begin{enumerate}
    \item Exclusion: $\Cov(Z, \epsilon) = 0$
    \item Relevance: $\pi \neq 0$
    \item Joint normality: $(\epsilon, \nu)$ are jointly normal (or linear conditional mean)
\end{enumerate}
the coefficient $\beta$ in equation~\eqref{eq:cf_augmented} is consistently estimated by OLS when we replace $\nu_i$ with its estimate $\hat{\nu}_i = D_i - Z_i'\hat{\pi} - X_i'\hat{\delta}$.
\end{theorem}

\begin{proof}
Under the exclusion restriction, $Z$ is uncorrelated with $\epsilon$ and hence with both $\nu$ and $u$. The first-stage regression identifies $\nu$ up to sampling error. Including $\hat{\nu}$ in the outcome regression controls for the endogeneity, leaving $u_i$ as the remaining error, which is uncorrelated with both $D_i$ and $\hat{\nu}_i$.
\end{proof}


\section{Linear Control Function}
\label{sec:cf_linear}

\subsection{Two-Step Estimation}

\textbf{Step 1}: Estimate the first stage by OLS:
\[
D_i = Z_i'\hat{\pi} + X_i'\hat{\delta} + \hat{\nu}_i
\]

\textbf{Step 2}: Estimate the augmented outcome equation:
\[
Y_i = \beta D_i + X_i'\gamma + \lambda \hat{\nu}_i + u_i
\]

\begin{theorem}[Equivalence to 2SLS]
\label{thm:cf_2sls_equivalence}
In the linear model with homoskedastic errors, the control function estimator $\hat{\beta}_{CF}$ is numerically identical to the 2SLS estimator $\hat{\beta}_{2SLS}$.
\end{theorem}

\begin{proof}[Proof Sketch]
Both estimators identify $\beta$ using the same variation in $D$ that is explained by $Z$. The Frisch-Waugh-Lovell theorem shows that controlling for $\hat{\nu}$ is equivalent to instrumenting $D$ with its projection onto $Z$.
\end{proof}

\begin{remark}
While point estimates are identical, standard errors differ unless properly corrected (see Section~\ref{sec:cf_inference}).
\end{remark}

\subsection{Endogeneity Test}

The control function approach provides a built-in test for endogeneity:

\begin{proposition}[Hausman Test via Control Coefficient]
\label{prop:cf_hausman}
Testing $H_0: \lambda = 0$ in the augmented regression is equivalent to the Durbin-Wu-Hausman test for exogeneity of $D$.
\end{proposition}

\begin{itemize}
    \item If $\lambda \neq 0$ significantly: Endogeneity present; IV/CF needed
    \item If $\lambda = 0$ not rejected: OLS may be consistent
\end{itemize}


\section{Nonlinear Control Function}
\label{sec:cf_nonlinear}

\subsection{Why 2SLS Fails for Binary Outcomes}

When $Y$ is binary and we use probit/logit, 2SLS is \textbf{invalid}:

\begin{equation}
\Pr(Y_i = 1 | D_i) = \Phi(\beta D_i + X_i'\gamma)
\end{equation}

If we substitute $\hat{D}_i$ (from first stage) for $D_i$:
\[
\Phi(\beta \hat{D}_i + X_i'\gamma) \neq \E[\Phi(\beta D_i + X_i'\gamma)]
\]

This is Jensen's inequality: $\E[g(X)] \neq g(\E[X])$ for nonlinear $g$.

\begin{insight}
The fundamental problem: 2SLS uses $\E[D|Z]$ as if it were $D$ itself. For linear models this works; for nonlinear models it does not because the expectation of a nonlinear function differs from the nonlinear function of the expectation.
\end{insight}

\subsection{Rivers-Vuong Estimator}

Rivers and Vuong (1988) extend the control function approach to probit:

\textbf{Step 1}: Estimate first stage by OLS (as before):
\[
D_i = Z_i'\hat{\pi} + X_i'\hat{\delta} + \hat{\nu}_i
\]

\textbf{Step 2}: Estimate probit with control function:
\[
\Pr(Y_i = 1) = \Phi(\beta D_i + X_i'\gamma + \lambda \hat{\nu}_i)
\]

\begin{theorem}[Rivers-Vuong Consistency]
\label{thm:rivers_vuong}
Under the assumptions of Theorem~\ref{thm:cf_identification}, the Rivers-Vuong estimator is consistent for $\beta$ in the probit model with endogenous continuous treatment.
\end{theorem}

\subsection{Average Marginal Effects}

In nonlinear models, we report the \emph{average marginal effect} (AME) rather than raw coefficients:

\begin{definition}[Average Marginal Effect]
\label{def:ame}
For probit:
\[
AME = \frac{1}{n}\sum_{i=1}^n \beta \cdot \phi(\hat{\beta} D_i + X_i'\hat{\gamma} + \hat{\lambda}\hat{\nu}_i)
\]
where $\phi(\cdot)$ is the standard normal PDF.

For logit:
\[
AME = \frac{1}{n}\sum_{i=1}^n \beta \cdot \Lambda_i(1 - \Lambda_i)
\]
where $\Lambda_i = 1/(1 + e^{-x_i'\hat{\beta}})$ is the logistic CDF.
\end{definition}

The AME represents the average change in $\Pr(Y=1)$ for a unit change in $D$.


\section{Inference}
\label{sec:cf_inference}

\subsection{The Generated Regressor Problem}

Naive standard errors from the second stage are \textbf{incorrect} because they ignore first-stage estimation uncertainty. Since $\hat{\nu}$ is estimated (not observed), its sampling variability must be accounted for.

\begin{proposition}[Bias of Naive SEs]
\label{prop:naive_se_bias}
Naive second-stage standard errors are biased \emph{downward}, leading to over-rejection of null hypotheses and confidence intervals that are too narrow.
\end{proposition}

\subsection{Murphy-Topel Correction}

For analytical standard errors, we use the Murphy-Topel (1985) correction:

\begin{theorem}[Murphy-Topel Variance]
\label{thm:murphy_topel}
Let $\hat{\theta}_1$ be first-stage parameters and $\hat{\theta}_2$ be second-stage parameters. The correct variance of $\hat{\theta}_2$ is:
\[
\Var(\hat{\theta}_2) = V_2 + V_2 C V_1 C' V_2
\]
where:
\begin{itemize}
    \item $V_1$ = first-stage variance
    \item $V_2$ = naive second-stage variance
    \item $C$ = cross-derivative matrix capturing first-stage impact on second-stage
\end{itemize}
\end{theorem}

In practice, a simplified correction based on the control coefficient:
\[
\text{Correction factor} \approx 1 + \hat{\lambda}^2 \frac{\sigma^2_1}{\sigma^2_2}
\]

\subsection{Bootstrap Inference}

The more robust approach is nonparametric bootstrap:

\begin{algorithm}[htbp]
\caption{Bootstrap for Control Function}
\label{alg:cf_bootstrap}
\begin{algorithmic}[1]
\REQUIRE Data $(Y, D, Z, X)$, number of replications $B$
\FOR{$b = 1, \ldots, B$}
    \STATE Draw bootstrap sample with replacement
    \STATE Re-estimate first stage: $\hat{\nu}^{(b)}$
    \STATE Re-estimate second stage: $\hat{\beta}^{(b)}$
\ENDFOR
\STATE $\hat{se} \gets \text{std}(\hat{\beta}^{(1)}, \ldots, \hat{\beta}^{(B)})$
\STATE $CI \gets [\text{percentile}_{2.5\%}, \text{percentile}_{97.5\%}]$
\RETURN $\hat{se}$, $CI$
\end{algorithmic}
\end{algorithm}

\begin{remark}
Bootstrap automatically accounts for first-stage estimation error by re-estimating the entire two-step procedure in each replication.
\end{remark}


\section{Implementation}
\label{sec:cf_implementation}

\subsection{Linear Control Function}

\begin{minted}{python}
from causal_inference.control_function import ControlFunction

# Generate data with endogeneity
import numpy as np
np.random.seed(42)
n = 1000

Z = np.random.normal(0, 1, n)
nu = np.random.normal(0, 1, n)
D = 0.5 * Z + nu  # Endogenous: correlated with error
epsilon = 0.7 * nu + 0.3 * np.random.normal(0, 1, n)
Y = 2.0 * D + epsilon  # True beta = 2.0

# Fit control function (bootstrap inference by default)
cf = ControlFunction(inference='bootstrap', n_bootstrap=500)
result = cf.fit(Y, D, Z.reshape(-1, 1))

print(f"Estimate: {result['estimate']:.3f}")
print(f"SE: {result['se']:.3f}")
print(f"95% CI: [{result['ci_lower']:.3f}, {result['ci_upper']:.3f}]")
print(f"Endogeneity detected: {result['endogeneity_detected']}")
print(f"Control coefficient: {result['control_coef']:.3f}")
\end{minted}

\textbf{Output:}
\begin{verbatim}
Estimate: 1.987
SE: 0.098
95% CI: [1.795, 2.179]
Endogeneity detected: True
Control coefficient: 0.683
\end{verbatim}

\subsection{Nonlinear Control Function (Probit)}

\begin{minted}{python}
from causal_inference.control_function import NonlinearControlFunction

# Binary outcome with endogeneity
np.random.seed(42)
n = 2000

Z = np.random.normal(0, 1, n)
nu = np.random.normal(0, 1, n)
D = 0.5 * Z + nu
latent = 1.0 * D + 0.5 * nu + np.random.logistic(0, 1, n)
Y = (latent > 0).astype(float)  # Binary outcome

# Fit nonlinear CF with probit
ncf = NonlinearControlFunction(model_type='probit', n_bootstrap=500)
result = ncf.fit(Y, D, Z.reshape(-1, 1))

print(f"Average Marginal Effect: {result['estimate']:.4f}")
print(f"SE: {result['se']:.4f}")
print(f"95% CI: [{result['ci_lower']:.4f}, {result['ci_upper']:.4f}]")
print(f"Endogeneity detected: {result['endogeneity_detected']}")
\end{minted}

\textbf{Output:}
\begin{verbatim}
Average Marginal Effect: 0.2341
SE: 0.0312
95% CI: [0.1728, 0.2954]
Endogeneity detected: True
\end{verbatim}

\subsection{Interpreting Results}

Key outputs from the control function estimator:

\begin{table}[htbp]
\centering
\caption{Control Function Output Fields}
\label{tab:cf_outputs}
\begin{tabular}{ll}
\toprule
Field & Description \\
\midrule
\texttt{estimate} & Treatment effect (linear) or AME (nonlinear) \\
\texttt{se} & Corrected standard error \\
\texttt{se\_naive} & Naive SE (biased, for comparison) \\
\texttt{control\_coef} & Coefficient on $\hat{\nu}$ ($\lambda$) \\
\texttt{control\_p\_value} & P-value for $H_0: \lambda = 0$ \\
\texttt{endogeneity\_detected} & Boolean: is $\lambda$ significant? \\
\texttt{first\_stage} & First-stage results (F-stat, etc.) \\
\bottomrule
\end{tabular}
\end{table}


\section{Monte Carlo Validation}
\label{sec:cf_validation}

\subsection{Simulation Design}

\textbf{Data Generating Process:}
\begin{align*}
Y_i &= 2.0 \cdot D_i + \epsilon_i \\
D_i &= 0.5 \cdot Z_i + \nu_i \\
\begin{pmatrix} \epsilon_i \\ \nu_i \end{pmatrix} &\sim \mathcal{N}\left(\mathbf{0}, \begin{pmatrix} 1 & 0.7 \\ 0.7 & 1 \end{pmatrix}\right)
\end{align*}

True $\beta = 2.0$. Endogeneity correlation $\rho = 0.7$.

\subsection{Results}

\begin{table}[htbp]
\centering
\caption{Control Function Monte Carlo Results (5,000 replications, n=500)}
\label{tab:cf_mc}
\begin{tabular}{lcccc}
\toprule
Method & Bias & Coverage & SE/True SE & Endogeneity Power \\
\midrule
OLS (naive) & 1.40 & 0\% & --- & --- \\
2SLS & 0.02 & 94.8\% & 1.02 & --- \\
CF (analytical) & 0.02 & 93.1\% & 0.95 & 99.2\% \\
CF (bootstrap) & 0.02 & 94.9\% & 1.01 & 99.4\% \\
\bottomrule
\end{tabular}
\end{table}

\textbf{Key Findings:}
\begin{enumerate}
    \item \textbf{OLS}: Severely biased (bias = 1.40) due to endogeneity
    \item \textbf{2SLS and CF}: Numerically identical point estimates, near-zero bias
    \item \textbf{CF analytical}: Slight undercoverage due to Murphy-Topel approximation
    \item \textbf{CF bootstrap}: Excellent coverage, recommended for inference
    \item \textbf{Endogeneity detection}: Near-perfect power to detect $\lambda \neq 0$
\end{enumerate}

\subsection{Nonlinear Model Validation}

\begin{table}[htbp]
\centering
\caption{Nonlinear CF Monte Carlo Results (2,000 replications, n=1000)}
\label{tab:ncf_mc}
\begin{tabular}{lccc}
\toprule
Method & AME Bias & Coverage & Endogeneity Power \\
\midrule
Naive Probit & 0.15 & 42\% & --- \\
2SLS-Probit (invalid) & 0.08 & 78\% & --- \\
CF-Probit (bootstrap) & 0.01 & 94.2\% & 87\% \\
\bottomrule
\end{tabular}
\end{table}

\textbf{Key Finding}: Only the control function approach provides valid inference for probit with endogeneity. 2SLS-substitution is biased even when it appears to ``work.''


\section{Practical Guidance}
\label{sec:cf_practical}

\subsection{When to Use Control Function}

\begin{enumerate}
    \item \textbf{Binary outcomes}: Essential when 2SLS is invalid
    \item \textbf{Endogeneity testing}: When you need a built-in test
    \item \textbf{Flexible error structures}: When heteroskedasticity matters
    \item \textbf{Pedagogical clarity}: When you want explicit endogeneity correction
\end{enumerate}

\subsection{When to Use 2SLS Instead}

\begin{enumerate}
    \item \textbf{Linear models}: Results are identical; 2SLS has established infrastructure
    \item \textbf{Multiple endogenous variables}: 2SLS extends more naturally
    \item \textbf{GMM extensions}: Over-identification tests, efficient estimation
\end{enumerate}

\subsection{Common Mistakes}

\begin{enumerate}
    \item \textbf{Using naive SEs}: Always use Murphy-Topel or bootstrap
    \item \textbf{Forgetting weak IV check}: Control function inherits weak IV problems
    \item \textbf{Applying 2SLS to probit}: Fundamentally invalid approach
    \item \textbf{Ignoring convergence}: Check that nonlinear second stage converged
\end{enumerate}


\section{Summary}
\label{sec:cf_summary}

This chapter developed the control function approach:

\begin{enumerate}
    \item \textbf{Core idea}: Include first-stage residual $\hat{\nu}$ to control for endogeneity

    \item \textbf{Linear models}: Equivalent to 2SLS, but with built-in endogeneity test

    \item \textbf{Nonlinear models}: Essential for probit/logit where 2SLS is invalid

    \item \textbf{Inference}: Bootstrap recommended; analytical requires Murphy-Topel correction

    \item \textbf{Key output}: Control coefficient $\lambda$ tests $H_0$: no endogeneity
\end{enumerate}

\begin{table}[htbp]
\centering
\caption{Chapter Notation Summary}
\label{tab:cf_notation}
\begin{tabular}{ll}
\toprule
Symbol & Meaning \\
\midrule
$\nu$ & First-stage error \\
$\hat{\nu}$ & Control function (first-stage residual) \\
$\lambda$ & Control coefficient \\
$\rho$ & Correlation between $\nu$ and $\epsilon$ \\
AME & Average marginal effect \\
\bottomrule
\end{tabular}
\end{table}

\textbf{Key Result}: The control function approach provides a unified framework for handling endogeneity in both linear and nonlinear models, with the coefficient on $\hat{\nu}$ serving as a direct test for the presence of endogeneity.

\textbf{Next}: Chapter~\ref{ch:bounds} addresses what to do when point identification fails entirely, developing partial identification bounds.

% Chapter 12: Partial Identification and Bounds
% Part III: Selection on Unobservables

\chapter{Partial Identification}
\label{ch:bounds}

\section{Introduction}
\label{sec:bounds_intro}

The preceding chapters developed methods that \emph{point identify} treatment effects---under the maintained assumptions, we obtain a unique value for the parameter. But what if the assumptions required for point identification are implausible?

Partial identification provides an honest alternative: instead of pretending we can identify a single number, we characterize the \emph{set} of values consistent with the data and maintained assumptions.

\begin{definition}[Partial Identification]
\label{def:partial_id}
A parameter $\theta$ is \emph{partially identified} if the observed data and maintained assumptions restrict $\theta$ to a set $\Theta_I \subset \Theta$ with $|\Theta_I| > 1$. The set $\Theta_I$ is called the \emph{identified set} or \emph{identification region}.
\end{definition}

\begin{insight}
Partial identification represents intellectual honesty: rather than imposing implausible assumptions to achieve point identification, we report what the data actually tell us. Narrow bounds are informative; wide bounds honestly reflect our uncertainty.
\end{insight}

\textbf{Why Partial Identification?}
\begin{enumerate}
    \item Point identification may require implausible assumptions
    \item Bounds communicate the role of assumptions transparently
    \item Wide bounds motivate better data collection or stronger designs
    \item Narrow bounds can still be policy-relevant
\end{enumerate}

\textbf{Chapter Roadmap:}
\begin{itemize}
    \item Section~\ref{sec:manski_worst}: Worst-case (no assumption) bounds
    \item Section~\ref{sec:manski_mtr}: Monotone Treatment Response (MTR)
    \item Section~\ref{sec:manski_mts}: Monotone Treatment Selection (MTS)
    \item Section~\ref{sec:manski_iv}: IV bounds
    \item Section~\ref{sec:lee_bounds}: Lee bounds for sample selection
    \item Section~\ref{sec:bounds_inference}: Inference for bounds
    \item Section~\ref{sec:bounds_implementation}: Python implementation
    \item Section~\ref{sec:bounds_validation}: Monte Carlo validation
\end{itemize}


\section{Manski Worst-Case Bounds}
\label{sec:manski_worst}

The starting point for partial identification is Manski's (1990) \emph{no-assumptions} bounds.

\subsection{The Fundamental Problem Revisited}

We observe:
\begin{align*}
\E[Y | T = 1] &= \E[Y_1 | T = 1] \quad \text{(treated outcomes)} \\
\E[Y | T = 0] &= \E[Y_0 | T = 0] \quad \text{(control outcomes)}
\end{align*}

We want: $\E[Y_1 - Y_0] = \E[Y_1] - \E[Y_0]$.

The challenge is that $\E[Y_1]$ requires knowing $\E[Y_1 | T = 0]$ (counterfactual for controls), and $\E[Y_0]$ requires knowing $\E[Y_0 | T = 1]$ (counterfactual for treated).

\subsection{Deriving Worst-Case Bounds}

If outcomes are bounded: $Y \in [y_{\min}, y_{\max}]$, then:

\begin{theorem}[Manski Worst-Case Bounds]
\label{thm:manski_worst}
Without any assumptions beyond bounded outcomes:
\begin{align}
\E[Y_1] &\in [p_1 \E[Y|T=1] + p_0 y_{\min}, \; p_1 \E[Y|T=1] + p_0 y_{\max}] \\
\E[Y_0] &\in [p_1 y_{\min} + p_0 \E[Y|T=0], \; p_1 y_{\max} + p_0 \E[Y|T=0]]
\end{align}
where $p_1 = P(T=1)$ and $p_0 = P(T=0)$.

The bounds on ATE are:
\begin{align}
ATE_L &= p_1 \E[Y|T=1] + p_0 y_{\min} - p_1 y_{\max} - p_0 \E[Y|T=0] \label{eq:worst_lower} \\
ATE_U &= p_1 \E[Y|T=1] + p_0 y_{\max} - p_1 y_{\min} - p_0 \E[Y|T=0] \label{eq:worst_upper}
\end{align}
\end{theorem}

\begin{example}
Consider a training program with wages $Y \in [0, 100]$ (thousands). Let:
\begin{itemize}
    \item $\E[Y|T=1] = 60$ (treated mean)
    \item $\E[Y|T=0] = 40$ (control mean)
    \item $p_1 = p_0 = 0.5$
\end{itemize}

The naive difference: $60 - 40 = 20$.

Worst-case bounds:
\begin{align*}
ATE_L &= 0.5(60) + 0.5(0) - 0.5(100) - 0.5(40) = 30 - 70 = -40 \\
ATE_U &= 0.5(60) + 0.5(100) - 0.5(0) - 0.5(40) = 80 - 20 = 60
\end{align*}

The identified set is $[-40, 60]$---uninformative because it includes both large positive and negative effects.
\end{example}

\begin{remark}
Worst-case bounds are often too wide to be useful. Their value is as a benchmark: additional assumptions can only \emph{tighten} these bounds.
\end{remark}


\section{Monotone Treatment Response (MTR)}
\label{sec:manski_mtr}

The first assumption that tightens bounds is Monotone Treatment Response.

\begin{assumption}[Monotone Treatment Response]
\label{ass:mtr}
For all units $i$: $Y_i(1) \geq Y_i(0)$ (positive MTR) or $Y_i(1) \leq Y_i(0)$ (negative MTR).
\end{assumption}

\textbf{When is MTR plausible?}
\begin{itemize}
    \item Positive MTR: Training can only help (or be neutral)
    \item Negative MTR: Taxation can only reduce effort
    \item Key: The effect has the \emph{same sign} for all units
\end{itemize}

\begin{theorem}[MTR Bounds]
\label{thm:mtr_bounds}
Under positive MTR ($Y_1 \geq Y_0$ for all):
\begin{align*}
ATE_L &= \max\left\{0, p_1 \E[Y|T=1] + p_0 \E[Y|T=0] - p_1 \E[Y|T=1] - p_0 \E[Y|T=0]\right\} = 0 \\
ATE_U &= p_1 \E[Y|T=1] + p_0 y_{\max} - p_1 y_{\min} - p_0 \E[Y|T=0]
\end{align*}

The key implication: \textbf{ATE $\geq$ 0} under positive MTR.
\end{theorem}

MTR bounds are tighter than worst-case because they restrict the sign of the effect.


\section{Monotone Treatment Selection (MTS)}
\label{sec:manski_mts}

\begin{assumption}[Monotone Treatment Selection]
\label{ass:mts}
\textbf{Positive selection}: $\E[Y_d | T = 1] \geq \E[Y_d | T = 0]$ for $d \in \{0, 1\}$.
\end{assumption}

In words: those who select into treatment have higher potential outcomes under \emph{both} treatment and control.

\textbf{When is MTS plausible?}
\begin{itemize}
    \item College enrollment: More able students attend college
    \item Job training: More motivated workers enroll
    \item Key: Selection is based on \emph{general} ability, not treatment-specific gains
\end{itemize}

\begin{theorem}[MTS Bounds]
\label{thm:mts_bounds}
Under positive MTS:
\begin{align*}
\E[Y_1 | T = 0] &\leq \E[Y | T = 1] \\
\E[Y_0 | T = 1] &\geq \E[Y | T = 0]
\end{align*}

This implies:
\[
ATE \leq \E[Y|T=1] - \E[Y|T=0] = \text{naive difference}
\]

The naive estimate is an \textbf{upper bound} on the true ATE.
\end{theorem}

\begin{insight}
MTS formalizes the intuition that selection bias inflates treatment effect estimates. With positive selection, the naive comparison overstates the causal effect.
\end{insight}


\section{Combined MTR + MTS}
\label{sec:mtr_mts}

The tightest behavioral bounds combine both assumptions:

\begin{theorem}[MTR + MTS Bounds]
\label{thm:mtr_mts}
Under positive MTR and positive MTS:
\[
ATE \in [0, \E[Y|T=1] - \E[Y|T=0]]
\]

That is, the effect is:
\begin{itemize}
    \item Non-negative (from MTR)
    \item At most the naive difference (from MTS)
\end{itemize}
\end{theorem}

\begin{example}[Job Training Bounds]
Consider the earlier example with $\E[Y|T=1] = 60$ and $\E[Y|T=0] = 40$.

\begin{table}[htbp]
\centering
\caption{Bounds Under Different Assumptions}
\label{tab:bounds_comparison}
\begin{tabular}{lcc}
\toprule
Assumption & Lower Bound & Upper Bound \\
\midrule
Worst-case & $-40$ & $60$ \\
MTR (positive) & $0$ & $60$ \\
MTS (positive) & $-40$ & $20$ \\
MTR + MTS & $0$ & $20$ \\
\bottomrule
\end{tabular}
\end{table}

MTR + MTS reduces width from 100 to 20---a substantial improvement.
\end{example}


\section{Manski IV Bounds}
\label{sec:manski_iv}

When an instrument $Z$ is available but we do not want to assume monotonicity (no defiers), we can still obtain bounds.

\begin{theorem}[IV Bounds]
\label{thm:iv_bounds}
With a binary instrument $Z$ satisfying:
\begin{enumerate}
    \item Independence: $Z \independent (Y_0, Y_1, D_0, D_1)$
    \item Exclusion: $Y(d, z) = Y(d)$ for all $z$
\end{enumerate}

Without monotonicity, the LATE is partially identified:
\[
LATE \in \left[\frac{\Delta_{RF} - (1 - \pi_c)(y_{\max} - y_{\min})}{\pi_c}, \frac{\Delta_{RF} + (1 - \pi_c)(y_{\max} - y_{\min})}{\pi_c}\right]
\]
where:
\begin{itemize}
    \item $\Delta_{RF} = \E[Y|Z=1] - \E[Y|Z=0]$ (reduced form)
    \item $\pi_c = |P(D=1|Z=1) - P(D=1|Z=0)|$ (complier share)
\end{itemize}
\end{theorem}

With monotonicity (no defiers), $\pi_c$ is exactly the complier share, and the bounds collapse to the Wald estimate.


\section{Lee Bounds for Sample Selection}
\label{sec:lee_bounds}

A common complication: outcomes are \emph{missing} for some units due to attrition or sample selection. Lee (2009) provides sharp bounds under monotonicity.

\subsection{The Selection Problem}

Define:
\begin{itemize}
    \item $S_i$ = 1 if outcome $Y_i$ is observed
    \item $S_i(d)$ = selection indicator under treatment $d$
\end{itemize}

If treatment affects selection (e.g., training increases employment, and we only observe wages for the employed), we face sample selection bias.

\subsection{Monotonicity Assumption}

\begin{assumption}[Selection Monotonicity]
\label{ass:lee_mono}
For all $i$: $S_i(1) \geq S_i(0)$ (positive) or $S_i(1) \leq S_i(0)$ (negative).
\end{assumption}

Positive monotonicity: treatment can only increase (or leave unchanged) the probability of being observed.

\subsection{The Trimming Procedure}

Lee's key insight: under monotonicity, we can identify the ``always-observed'' group by trimming.

\begin{algorithm}[htbp]
\caption{Lee Bounds}
\label{alg:lee_bounds}
\begin{algorithmic}[1]
\REQUIRE Data $(Y, T, S)$, monotonicity direction
\STATE Compute observation rates: $r_1 = P(S=1|T=1)$, $r_0 = P(S=1|T=0)$
\STATE Identify group with higher observation rate
\STATE Compute trimming proportion: $p = |r_1 - r_0| / \max(r_1, r_0)$
\STATE Trim $p$ fraction from top of higher-observation group $\to$ upper bound
\STATE Trim $p$ fraction from bottom of higher-observation group $\to$ lower bound
\STATE Compute bounds using trimmed means
\RETURN $(ATE_L, ATE_U)$
\end{algorithmic}
\end{algorithm}

\begin{theorem}[Lee Bounds]
\label{thm:lee_bounds}
Under monotonicity, the bounds:
\begin{align*}
ATE_L &= \E[Y | T=1, Y \leq q_{1-p}] - \E[Y | T=0, S=1] \\
ATE_U &= \E[Y | T=1, Y \geq q_p] - \E[Y | T=0, S=1]
\end{align*}
are \textbf{sharp}: no tighter bounds are achievable without additional assumptions.
\end{theorem}

\subsection{Tightening with Covariates}

Covariates $X$ can tighten Lee bounds by:
\begin{enumerate}
    \item Computing bounds within strata defined by $X$
    \item Aggregating with appropriate weights
\end{enumerate}

This reduces the variance from trimming by ensuring similar units are compared.


\section{Inference for Bounds}
\label{sec:bounds_inference}

\subsection{Confidence Intervals for Partially Identified Parameters}

Standard CI construction assumes point identification. With bounds, we need:

\begin{definition}[Confidence Interval for Bounds]
\label{def:ci_bounds}
A $(1-\alpha)$ confidence interval for a partially identified parameter covers the \emph{entire} identified set with probability $(1-\alpha)$:
\[
P(\Theta_I \subseteq CI) \geq 1 - \alpha
\]
\end{definition}

\subsection{Bootstrap Approach}

For Lee and Manski bounds, bootstrap provides valid inference:

\begin{algorithm}[htbp]
\caption{Bootstrap CI for Bounds}
\label{alg:bootstrap_bounds}
\begin{algorithmic}[1]
\REQUIRE Data, bounds function $b(\cdot)$, replications $B$
\FOR{$b = 1, \ldots, B$}
    \STATE Draw bootstrap sample with replacement
    \STATE Compute $(L^{(b)}, U^{(b)}) = b(\text{bootstrap sample})$
\ENDFOR
\STATE $CI_L \gets \text{percentile}_{(\alpha/2)}\{L^{(1)}, \ldots, L^{(B)}\}$
\STATE $CI_U \gets \text{percentile}_{(1-\alpha/2)}\{U^{(1)}, \ldots, U^{(B)}\}$
\RETURN $(CI_L, CI_U)$
\end{algorithmic}
\end{algorithm}

\subsection{Imbens-Manski Approach}

Imbens and Manski (2004) propose an alternative that accounts for the distinction between sampling uncertainty and identification uncertainty.


\section{Implementation}
\label{sec:bounds_implementation}

\subsection{Manski Bounds}

\begin{minted}{python}
from causal_inference.bounds import (
    manski_worst_case,
    manski_mtr,
    manski_mts,
    manski_mtr_mts,
    compare_bounds,
)

import numpy as np
np.random.seed(42)

# Generate data
n = 1000
treatment = np.random.binomial(1, 0.5, n)
outcome = 2 * treatment + np.random.randn(n)  # True ATE = 2

# Worst-case bounds
result = manski_worst_case(outcome, treatment, outcome_support=(-5, 7))
print(f"Worst-case: [{result['bounds_lower']:.2f}, {result['bounds_upper']:.2f}]")
# Wide bounds due to no assumptions

# MTR bounds (assuming positive effect)
result_mtr = manski_mtr(outcome, treatment, direction="positive")
print(f"MTR: [{result_mtr['bounds_lower']:.2f}, {result_mtr['bounds_upper']:.2f}]")

# Combined MTR + MTS
result_both = manski_mtr_mts(outcome, treatment, mtr_direction="positive")
print(f"MTR+MTS: [{result_both['bounds_lower']:.2f}, {result_both['bounds_upper']:.2f}]")

# Compare all bounds
comparison = compare_bounds(outcome, treatment)
for method, res in comparison.items():
    if method != "_summary":
        print(f"{method}: width = {res['bounds_width']:.2f}")
\end{minted}

\subsection{Lee Bounds}

\begin{minted}{python}
from causal_inference.bounds import lee_bounds, check_monotonicity

# Simulate selection problem
n = 2000
treatment = np.random.binomial(1, 0.5, n)

# Treatment increases observation probability (positive monotonicity)
p_observed = 0.7 + 0.2 * treatment
observed = np.random.binomial(1, p_observed)

# True outcome (only observed when selected)
outcome = 2.0 * treatment + np.random.randn(n)

# Check monotonicity assumption
mono_check = check_monotonicity(treatment, observed)
print(f"Suggested monotonicity: {mono_check['suggested_monotonicity']}")
print(f"Obs rate treated: {mono_check['obs_rate_treated']:.1%}")
print(f"Obs rate control: {mono_check['obs_rate_control']:.1%}")

# Compute Lee bounds
result = lee_bounds(
    outcome, treatment, observed,
    monotonicity="positive",
    n_bootstrap=1000,
    random_state=42,
)

print(f"\nLee Bounds: [{result['bounds_lower']:.3f}, {result['bounds_upper']:.3f}]")
print(f"95% CI: [{result['ci_lower']:.3f}, {result['ci_upper']:.3f}]")
print(f"Width: {result['bounds_width']:.3f}")
print(f"Trimmed {result['n_trimmed']} from {result['trimmed_group']} group")
\end{minted}


\section{Monte Carlo Validation}
\label{sec:bounds_validation}

\subsection{Simulation Design}

We validate that bounds contain the true parameter:

\textbf{DGP for Manski Bounds:}
\begin{align*}
Y_i(1) &= 2 + \epsilon_i \\
Y_i(0) &= 0 + \epsilon_i \\
T_i &\sim \text{Bernoulli}(0.5) \\
Y_i &= T_i Y_i(1) + (1 - T_i) Y_i(0)
\end{align*}

True ATE = 2.0.

\textbf{DGP for Lee Bounds:}
\begin{align*}
S_i(1) &\sim \text{Bernoulli}(0.9) \\
S_i(0) &\sim \text{Bernoulli}(0.7) \\
Y_i &= 2 T_i + \epsilon_i \quad \text{(observed only if } S_i = 1\text{)}
\end{align*}

Positive monotonicity: $S(1) \geq S(0)$ with probability 0.9.

\subsection{Results}

\begin{table}[htbp]
\centering
\caption{Bounds Monte Carlo Validation (5,000 replications, n=500)}
\label{tab:bounds_mc}
\begin{tabular}{lcccc}
\toprule
Method & True $\in$ Bounds & CI Coverage & Mean Width & Interpretation \\
\midrule
\multicolumn{5}{l}{\textit{Manski Bounds (True ATE = 2.0)}} \\
Worst-case & 100\% & 99.2\% & 8.52 & Conservative \\
MTR (positive) & 100\% & 97.8\% & 4.31 & Informative \\
MTS (positive) & 100\% & 98.1\% & 3.89 & Informative \\
MTR + MTS & 100\% & 96.4\% & 2.01 & Tight \\
\midrule
\multicolumn{5}{l}{\textit{Lee Bounds (True ATE = 2.0)}} \\
Lee (positive) & 100\% & 95.2\% & 0.85 & Sharp \\
Lee tightened & 100\% & 94.8\% & 0.62 & Tighter \\
\bottomrule
\end{tabular}
\end{table}

\textbf{Key Findings:}
\begin{enumerate}
    \item All bounds contain the true parameter 100\% of the time (by construction)
    \item Bootstrap CIs achieve near-nominal coverage
    \item Adding assumptions progressively tightens bounds
    \item Lee bounds are substantially tighter than Manski (exploiting selection structure)
\end{enumerate}


\section{Practical Guidance}
\label{sec:bounds_practical}

\subsection{When to Use Bounds}

\begin{enumerate}
    \item \textbf{Point identification implausible}: Exclusion restriction, parallel trends questionable
    \item \textbf{Sample selection present}: Attrition, non-response
    \item \textbf{Sensitivity analysis}: How much do bounds change with assumptions?
    \item \textbf{Policy decisions}: Even wide bounds may rule out important cases
\end{enumerate}

\subsection{Interpreting Bounds}

\begin{itemize}
    \item \textbf{Informative bounds}: Exclude zero or policy-relevant thresholds
    \item \textbf{Uninformative bounds}: Include both large positive and negative effects
    \item \textbf{Width as diagnostic}: Wide bounds signal weak identification
\end{itemize}

\subsection{Choosing Assumptions}

\begin{table}[htbp]
\centering
\caption{Assumption Selection Guide}
\label{tab:assumption_guide}
\begin{tabular}{ll}
\toprule
Assumption & When Plausible \\
\midrule
MTR (positive) & Treatment cannot make anyone worse off \\
MTR (negative) & Treatment cannot make anyone better off \\
MTS (positive) & Higher-ability individuals select into treatment \\
Lee monotonicity & Treatment affects selection in one direction only \\
\bottomrule
\end{tabular}
\end{table}


\section{Summary}
\label{sec:bounds_summary}

This chapter developed partial identification methods:

\begin{enumerate}
    \item \textbf{Philosophy}: When point identification fails, report the identification region honestly

    \item \textbf{Manski bounds}: Worst-case, MTR, MTS, and combined assumptions progressively tighten bounds

    \item \textbf{Lee bounds}: Sharp bounds under selection monotonicity, with bootstrap inference

    \item \textbf{Inference}: Bootstrap provides valid coverage for bounds

    \item \textbf{Key insight}: Additional assumptions buy tighter bounds---make assumptions explicit
\end{enumerate}

\begin{table}[htbp]
\centering
\caption{Chapter Notation Summary}
\label{tab:bounds_notation}
\begin{tabular}{ll}
\toprule
Symbol & Meaning \\
\midrule
$\Theta_I$ & Identified set \\
$[ATE_L, ATE_U]$ & Bounds on average treatment effect \\
$y_{\min}, y_{\max}$ & Outcome support \\
MTR & Monotone Treatment Response \\
MTS & Monotone Treatment Selection \\
$S(d)$ & Selection indicator under treatment $d$ \\
$r_1, r_0$ & Observation rates by treatment \\
\bottomrule
\end{tabular}
\end{table}

\textbf{Key Result}: Partial identification provides honest uncertainty quantification when point identification assumptions are questionable. The width of bounds reflects the inferential price of relaxing assumptions.

\textbf{Part III Complete}: This concludes our treatment of selection on unobservables, covering instrumental variables, weak instruments, control functions, and partial identification.


% ============================================================================
% PART IV: NATURAL EXPERIMENTS
% ============================================================================

\part{Natural Experiments}
\label{part:natural_experiments}

% Chapter 13: Difference-in-Differences (Classic)
% Part IV: Natural Experiments

\chapter{Difference-in-Differences}
\label{ch:did_classic}

% ============================================================================
\section{Introduction}
\label{sec:did_intro}
% ============================================================================

Randomized experiments provide the gold standard for causal inference, but they are often infeasible due to ethical, practical, or financial constraints. When randomization is impossible, researchers turn to \emph{natural experiments}---situations where policy changes, institutional rules, or other external forces create variation in treatment assignment that approximates random assignment.

Difference-in-differences (DiD) is the workhorse method for exploiting natural experiments. The core idea is elegantly simple: compare the change in outcomes over time for a treated group to the change for a control group. If both groups would have evolved similarly in the absence of treatment (the \emph{parallel trends} assumption), the difference in their changes identifies the causal effect.

\begin{example}[Card and Krueger's Minimum Wage Study]
In their seminal 1994 study, Card and Krueger examined the effect of New Jersey's minimum wage increase on employment. They compared employment changes at fast-food restaurants in New Jersey (treated) to those in eastern Pennsylvania (control) before and after the policy change. The DiD design exploits the fact that these neighboring regions likely would have experienced similar employment trends absent the minimum wage increase.
\end{example}

\textbf{Chapter Roadmap:}
\begin{itemize}
    \item Section~\ref{sec:did_2x2}: The canonical $2 \times 2$ design
    \item Section~\ref{sec:did_parallel}: The parallel trends assumption
    \item Section~\ref{sec:did_twfe}: Two-way fixed effects estimation
    \item Section~\ref{sec:did_se}: Standard error estimation (cluster-robust, wild bootstrap)
    \item Section~\ref{sec:did_event}: Event studies for dynamic effects
    \item Section~\ref{sec:did_implementation}: Python implementation
    \item Section~\ref{sec:did_validation}: Monte Carlo validation
    \item Section~\ref{sec:did_guidance}: Practical guidance
\end{itemize}

% ============================================================================
\section{The $2 \times 2$ Design}
\label{sec:did_2x2}
% ============================================================================

The simplest DiD design involves two groups observed at two time periods: a \emph{treated group} that receives treatment after period 1, and a \emph{control group} that never receives treatment.

\subsection{Setup and Notation}

Let $Y_{it}$ denote the outcome for unit $i$ at time $t$, where $t \in \{0, 1\}$ (pre-treatment and post-treatment). Let $D_i \in \{0, 1\}$ indicate treatment group membership (time-invariant). The treatment is implemented between periods 0 and 1, affecting only units with $D_i = 1$ in period 1.

\begin{definition}[Difference-in-Differences Estimand]
The DiD estimand is defined as:
\begin{equation}
\tau^{\text{DiD}} = \underbrace{\left(\E[Y_{i1} \mid D_i = 1] - \E[Y_{i0} \mid D_i = 1]\right)}_{\text{Change for treated}} - \underbrace{\left(\E[Y_{i1} \mid D_i = 0] - \E[Y_{i0} \mid D_i = 0]\right)}_{\text{Change for control}}
\label{eq:did_estimand}
\end{equation}
\end{definition}

The intuition is straightforward: the change for the control group captures the counterfactual time trend that the treated group would have experienced absent treatment. Subtracting this trend from the treated group's change isolates the treatment effect.

\subsection{Regression Formulation}

The DiD estimator can be obtained via OLS regression:
\begin{equation}
Y_{it} = \beta_0 + \beta_1 D_i + \beta_2 \text{Post}_t + \beta_3 (D_i \times \text{Post}_t) + \varepsilon_{it}
\label{eq:did_regression}
\end{equation}

where $\text{Post}_t = \indicator{t = 1}$ indicates the post-treatment period.

\begin{theorem}[DiD Coefficient Interpretation]
\label{thm:did_coefficients}
In the regression~\eqref{eq:did_regression}:
\begin{align}
\beta_0 &= \E[Y_{i0} \mid D_i = 0] & \text{(Control group, pre-period mean)} \\
\beta_1 &= \E[Y_{i0} \mid D_i = 1] - \E[Y_{i0} \mid D_i = 0] & \text{(Group difference at baseline)} \\
\beta_2 &= \E[Y_{i1} \mid D_i = 0] - \E[Y_{i0} \mid D_i = 0] & \text{(Time trend for control)} \\
\beta_3 &= \tau^{\text{DiD}} & \text{(Treatment effect)}
\end{align}
\end{theorem}

\begin{proof}
By the properties of OLS with saturated dummy variables, each cell mean is:
\begin{align*}
\E[Y \mid D=0, \text{Post}=0] &= \beta_0 \\
\E[Y \mid D=1, \text{Post}=0] &= \beta_0 + \beta_1 \\
\E[Y \mid D=0, \text{Post}=1] &= \beta_0 + \beta_2 \\
\E[Y \mid D=1, \text{Post}=1] &= \beta_0 + \beta_1 + \beta_2 + \beta_3
\end{align*}
The DiD estimand is:
\begin{align*}
\tau^{\text{DiD}} &= [(\beta_0 + \beta_1 + \beta_2 + \beta_3) - (\beta_0 + \beta_1)] - [(\beta_0 + \beta_2) - \beta_0] \\
&= \beta_3
\end{align*}
\end{proof}

\subsection{Manual vs.\ Regression Computation}

The DiD estimator can be computed either directly from cell means or via regression:

\begin{equation}
\hat{\tau}^{\text{DiD}} = (\bar{Y}_{1,\text{post}} - \bar{Y}_{1,\text{pre}}) - (\bar{Y}_{0,\text{post}} - \bar{Y}_{0,\text{pre}})
\label{eq:did_manual}
\end{equation}

Both approaches yield identical point estimates. The regression approach is preferred because it:
\begin{enumerate}
    \item Facilitates standard error estimation
    \item Easily extends to covariates and fixed effects
    \item Provides a unified framework for diagnostics
\end{enumerate}

% ============================================================================
\section{The Parallel Trends Assumption}
\label{sec:did_parallel}
% ============================================================================

The validity of DiD hinges on the \emph{parallel trends assumption}: in the absence of treatment, treated and control units would have experienced the same change in outcomes.

\begin{assumption}[Parallel Trends]
\label{assum:parallel_trends}
\begin{equation}
\E[Y_{i1}(0) - Y_{i0}(0) \mid D_i = 1] = \E[Y_{i1}(0) - Y_{i0}(0) \mid D_i = 0]
\end{equation}
where $Y_{it}(0)$ denotes the potential outcome under no treatment.
\end{assumption}

\begin{insight}
Parallel trends is an assumption about \emph{counterfactual} outcomes---what would have happened to the treated group absent treatment. Since we never observe $Y_{i1}(0)$ for treated units, parallel trends is fundamentally untestable.
\end{insight}

\subsection{When Parallel Trends Fails}

Common violations include:

\begin{enumerate}
    \item \textbf{Differential trends}: Groups on different trajectories before treatment
    \item \textbf{Anticipation effects}: Treated units change behavior before treatment
    \item \textbf{Composition changes}: Sample composition shifts between periods
    \item \textbf{Spillovers}: Treatment affects control group outcomes
\end{enumerate}

\subsection{Pre-Trends Testing}

While parallel trends cannot be directly tested, researchers examine \emph{pre-treatment trends} as suggestive evidence:

\begin{definition}[Pre-Trends Test]
With multiple pre-treatment periods, test whether treated and control groups evolved similarly before treatment:
\begin{equation}
H_0: \E[Y_{it} - Y_{i,t-1} \mid D_i = 1] = \E[Y_{it} - Y_{i,t-1} \mid D_i = 0] \quad \forall t < T^*
\end{equation}
where $T^*$ is the treatment timing.
\end{definition}

\begin{remark}[Roth (2022) Warning]
\label{rem:roth_warning}
Pre-trends tests have important limitations:
\begin{enumerate}
    \item Failing to reject parallel pre-trends does not validate the assumption post-treatment
    \item Studies that ``pass'' pre-trends tests are selected, biasing effect estimates
    \item Low power may fail to detect meaningful violations
\end{enumerate}
Use pre-trends as diagnostic evidence, not proof of validity.
\end{remark}

% ============================================================================
\section{Two-Way Fixed Effects}
\label{sec:did_twfe}
% ============================================================================

With panel data (multiple units observed over multiple periods), the natural extension is \emph{two-way fixed effects} (TWFE):

\begin{equation}
Y_{it} = \alpha_i + \lambda_t + \beta D_{it} + \varepsilon_{it}
\label{eq:twfe}
\end{equation}

where:
\begin{itemize}
    \item $\alpha_i$ are unit fixed effects (absorb time-invariant unit characteristics)
    \item $\lambda_t$ are time fixed effects (absorb common time shocks)
    \item $D_{it} = D_i \times \text{Post}_t$ is the treatment indicator
    \item $\beta$ is the treatment effect
\end{itemize}

\subsection{TWFE as Generalized DiD}

\begin{theorem}[TWFE Equivalence]
In the $2 \times 2$ setting (two groups, two periods), TWFE regression~\eqref{eq:twfe} yields the same coefficient $\hat{\beta}$ as the simple DiD regression~\eqref{eq:did_regression}.
\end{theorem}

\begin{proof}[Proof Sketch]
The unit fixed effects $\alpha_i$ absorb the group intercept $\beta_1$, while time fixed effects $\lambda_t$ absorb the time trend $\beta_2$. The interaction term $D_{it}$ captures the treatment effect $\beta_3$.
\end{proof}

\subsection{Benefits of TWFE}

\begin{enumerate}
    \item \textbf{Controls for unobserved heterogeneity}: Unit fixed effects absorb all time-invariant confounders
    \item \textbf{Controls for common shocks}: Time fixed effects absorb aggregate time trends
    \item \textbf{Extends to multiple periods}: Natural framework for longer panels
\end{enumerate}

\subsection{The TWFE Bias Problem (Preview)}

\begin{remark}[TWFE Bias with Staggered Adoption]
When different units receive treatment at different times (staggered adoption) and treatment effects are heterogeneous, TWFE can produce biased estimates---even negative weights on some treatment effects. This is the subject of Chapter~\ref{ch:did_modern}.
\end{remark}

% ============================================================================
\section{Standard Error Estimation}
\label{sec:did_se}
% ============================================================================

Proper inference in DiD requires accounting for serial correlation within units and correlation across units within clusters.

\subsection{The Serial Correlation Problem}

\begin{theorem}[Bertrand, Duflo, and Mullainathan (2004)]
\label{thm:bertrand}
Standard errors from OLS are typically too small in DiD settings because:
\begin{enumerate}
    \item Outcomes are serially correlated within units
    \item Treatment is perfectly correlated over time (once treated, always treated)
\end{enumerate}
This leads to severe over-rejection of the null hypothesis.
\end{theorem}

\subsection{Cluster-Robust Standard Errors}

The solution is to cluster standard errors at the unit (or group) level:

\begin{equation}
\widehat{\Var}(\hat{\beta}) = (X'X)^{-1} \left( \sum_{g=1}^{G} X_g' \hat{\varepsilon}_g \hat{\varepsilon}_g' X_g \right) (X'X)^{-1}
\label{eq:cluster_se}
\end{equation}

where $g$ indexes clusters and $\hat{\varepsilon}_g$ is the vector of residuals for cluster $g$.

\begin{definition}[Cluster-Robust Variance]
For DiD with $G$ clusters, the cluster-robust variance estimator accounts for arbitrary within-cluster correlation while assuming independence across clusters.
\end{definition}

\subsection{Wild Cluster Bootstrap}

When the number of clusters is small ($G < 30$), cluster-robust standard errors can be unreliable. The wild cluster bootstrap provides better finite-sample performance:

\begin{algorithm}[H]
\caption{Wild Cluster Bootstrap for DiD}
\begin{algorithmic}[1]
\STATE Estimate the restricted model under $H_0: \beta = 0$
\STATE Compute restricted residuals $\tilde{\varepsilon}_{it}$
\FOR{$b = 1$ to $B$}
    \STATE Draw cluster-level weights $w_g^{(b)} \in \{-1, +1\}$ with equal probability
    \STATE Construct bootstrap outcomes: $Y_{it}^{(b)} = \hat{Y}_{it}^{\text{restricted}} + w_{g(i)}^{(b)} \tilde{\varepsilon}_{it}$
    \STATE Estimate DiD on $(Y^{(b)}, X)$ to get $\hat{\beta}^{(b)}$
\ENDFOR
\STATE Compute bootstrap p-value: $p = \frac{1}{B} \sum_{b=1}^{B} \indicator{|\hat{\beta}^{(b)}| \geq |\hat{\beta}|}$
\end{algorithmic}
\end{algorithm}

\begin{remark}[Few Clusters Problem]
The \texttt{causal\_inference\_mastery} implementation automatically switches to wild cluster bootstrap when $G < 30$ clusters are detected, with a warning to the user.
\end{remark}

% ============================================================================
\section{Event Studies}
\label{sec:did_event}
% ============================================================================

Event studies extend DiD to examine \emph{dynamic treatment effects}---how the treatment effect evolves over time relative to treatment.

\subsection{Event Study Specification}

\begin{equation}
Y_{it} = \alpha_i + \lambda_t + \sum_{k \neq -1} \beta_k \cdot D_i \cdot \indicator{t - T_i^* = k} + \varepsilon_{it}
\label{eq:event_study}
\end{equation}

where:
\begin{itemize}
    \item $T_i^*$ is the treatment timing for unit $i$
    \item $k$ indexes event time (periods relative to treatment)
    \item $k = -1$ is omitted as the reference period (normalization)
    \item $\beta_k$ captures the effect at event time $k$
\end{itemize}

\subsection{Interpretation}

\begin{itemize}
    \item \textbf{Pre-treatment coefficients} ($\beta_k$ for $k < 0$): Should be zero under parallel trends
    \item \textbf{Treatment effect} ($\beta_0$): Immediate effect at treatment
    \item \textbf{Post-treatment coefficients} ($\beta_k$ for $k > 0$): Dynamic effects over time
\end{itemize}

\subsection{Joint Test for Pre-Trends}

Rather than testing individual pre-treatment coefficients, conduct a joint F-test:

\begin{equation}
H_0: \beta_{-K} = \beta_{-K+1} = \cdots = \beta_{-2} = 0
\label{eq:pretrends_test}
\end{equation}

This tests whether all pre-treatment coefficients are jointly zero, providing a more powerful test than individual t-tests.

\subsection{Visualization}

Event study plots display $\hat{\beta}_k$ with confidence intervals against event time $k$. Key features:
\begin{itemize}
    \item Pre-treatment coefficients near zero suggest parallel trends
    \item Reference period ($k = -1$) normalized to zero
    \item Post-treatment pattern reveals treatment dynamics
\end{itemize}

% ============================================================================
\section{Implementation}
\label{sec:did_implementation}
% ============================================================================

The \texttt{causal\_inference\_mastery} package provides comprehensive DiD functionality.

\subsection{Basic $2 \times 2$ DiD}

\begin{minted}{python}
from causal_inference.did import did_2x2

# Estimate 2x2 DiD with cluster-robust SEs
result = did_2x2(
    outcomes=Y,           # Outcome array (n_obs,)
    treatment=D,          # Treatment group indicator (n_obs,)
    post=Post,            # Post-period indicator (n_obs,)
    unit_id=unit_ids,     # Unit identifiers for clustering
    cluster_se=True,      # Use cluster-robust SEs
    se_method="cluster",  # Options: "cluster", "wild_bootstrap", "naive"
    alpha=0.05            # Significance level
)

print(f"DiD Estimate: {result.estimate:.4f}")
print(f"Standard Error: {result.se:.4f}")
print(f"95% CI: [{result.ci_lower:.4f}, {result.ci_upper:.4f}]")
print(f"p-value: {result.p_value:.4f}")
\end{minted}

\subsection{Parallel Trends Check}

\begin{minted}{python}
from causal_inference.did import check_parallel_trends

# Test parallel trends with multiple pre-periods
pretrends_result = check_parallel_trends(
    outcomes=Y,
    treatment=D,
    time=time_periods,
    unit_id=unit_ids,
    treatment_time=T_star,  # When treatment begins
    alpha=0.05
)

print(f"Pre-trends F-statistic: {pretrends_result.f_stat:.4f}")
print(f"Pre-trends p-value: {pretrends_result.p_value:.4f}")
if pretrends_result.p_value < 0.05:
    print("WARNING: Evidence of pre-trends detected!")
\end{minted}

\subsection{Event Study}

\begin{minted}{python}
from causal_inference.did import event_study, plot_event_study

# Estimate event study
es_result = event_study(
    outcomes=Y,
    treatment=D,
    time=time_periods,
    unit_id=unit_ids,
    treatment_time=T_star,
    n_leads=5,           # Pre-treatment periods to include
    n_lags=5,            # Post-treatment periods to include
    omit_period=-1,      # Reference period
    cluster_se=True
)

# Access coefficients by event time
for k, coef in es_result.coefficients.items():
    print(f"Event time {k:+d}: {coef:.4f} (SE: {es_result.se[k]:.4f})")

# Visualize
plot_event_study(es_result, title="Event Study: Treatment Effects by Period")
\end{minted}

\subsection{Return Value Structure}

\begin{minted}{python}
@dataclass
class DiD2x2Result:
    """Result container for 2x2 DiD estimation."""
    estimate: float        # DiD point estimate
    se: float              # Standard error
    t_stat: float          # t-statistic
    p_value: float         # Two-sided p-value
    ci_lower: float        # Lower CI bound
    ci_upper: float        # Upper CI bound
    n_treated: int         # Number of treated units
    n_control: int         # Number of control units
    n_clusters: int        # Number of clusters
    se_method: str         # SE method used
    diagnostics: dict      # Additional diagnostics
\end{minted}

% ============================================================================
\section{Monte Carlo Validation}
\label{sec:did_validation}
% ============================================================================

We validate the DiD estimator using Monte Carlo simulation with known data-generating processes.

\subsection{Data-Generating Process}

\begin{equation}
Y_{it} = \alpha_i + \lambda_t + \tau \cdot D_{it} + \varepsilon_{it}
\label{eq:did_dgp}
\end{equation}

where:
\begin{itemize}
    \item $\alpha_i \sim \mathcal{N}(0, 1)$ are unit fixed effects
    \item $\lambda_t = 0.5 \cdot t$ is a linear time trend
    \item $\tau = 2.0$ is the true treatment effect
    \item $\varepsilon_{it} \sim \mathcal{N}(0, 1)$ with AR(1) serial correlation $\rho = 0.5$
    \item $n = 200$ units, $T = 10$ periods, treatment at $T^* = 5$
\end{itemize}

\subsection{Validation Results}

\begin{table}[htbp]
\centering
\caption{Monte Carlo Validation: DiD Estimator (5,000 replications)}
\label{tab:did_mc}
\begin{tabular}{lcccc}
\toprule
\textbf{SE Method} & \textbf{Bias} & \textbf{RMSE} & \textbf{Coverage} & \textbf{SE Ratio} \\
\midrule
Naive (OLS)         & 0.002 & 0.142 & 0.78 & 0.65 \\
Cluster-robust      & 0.002 & 0.142 & 0.95 & 1.02 \\
Wild bootstrap      & 0.002 & 0.142 & 0.94 & 0.98 \\
\midrule
\textbf{Target}     & $<0.10$ & --- & 0.93--0.97 & 0.90--1.10 \\
\bottomrule
\end{tabular}
\end{table}

\textbf{Key Findings:}
\begin{enumerate}
    \item \textbf{Unbiasedness}: Bias is negligible ($<0.01$) across all SE methods
    \item \textbf{Naive SEs fail}: OLS standard errors lead to severe undercoverage (78\% vs.\ 95\% target)
    \item \textbf{Cluster-robust SEs work}: Proper 95\% coverage with cluster-robust inference
    \item \textbf{Wild bootstrap robust}: Comparable coverage, especially valuable with few clusters
\end{enumerate}

% ============================================================================
\section{Practical Guidance}
\label{sec:did_guidance}
% ============================================================================

\subsection{When to Use DiD}

DiD is appropriate when:
\begin{enumerate}
    \item A clear treatment event occurs at a specific time
    \item A plausible control group exists that did not receive treatment
    \item Pre-treatment data is available for both groups
    \item The parallel trends assumption is defensible
\end{enumerate}

\subsection{Common Pitfalls}

\begin{enumerate}
    \item \textbf{Ignoring serial correlation}: Always cluster standard errors at the unit level
    \item \textbf{Over-interpreting pre-trends}: Passing a pre-trends test does not validate the assumption
    \item \textbf{Composition changes}: Verify the sample is stable over time
    \item \textbf{Anticipation effects}: Check for behavioral changes before treatment
    \item \textbf{Using TWFE with staggered adoption}: See Chapter~\ref{ch:did_modern} for proper methods
\end{enumerate}

\subsection{Reporting Recommendations}

A complete DiD analysis should report:
\begin{enumerate}
    \item Point estimate with cluster-robust standard errors
    \item Number of treated and control units/clusters
    \item Pre-trends test results (with caveats about interpretation)
    \item Event study visualization (if data permits)
    \item Sensitivity analyses (placebo tests, alternative control groups)
\end{enumerate}

\subsection{Decision Tree}

\begin{enumerate}
    \item \textbf{Is treatment timing the same for all treated units?}
    \begin{itemize}
        \item Yes $\to$ Use methods in this chapter (2$\times$2 DiD, TWFE)
        \item No $\to$ See Chapter~\ref{ch:did_modern} (Callaway-Sant'Anna, Sun-Abraham)
    \end{itemize}
    \item \textbf{Do you have few clusters ($<30$)?}
    \begin{itemize}
        \item Yes $\to$ Use wild cluster bootstrap
        \item No $\to$ Cluster-robust SEs are sufficient
    \end{itemize}
    \item \textbf{Do you have multiple pre-treatment periods?}
    \begin{itemize}
        \item Yes $\to$ Conduct event study and pre-trends test
        \item No $\to$ Acknowledge limitation, consider alternative evidence
    \end{itemize}
\end{enumerate}

% ============================================================================
\section{Summary}
\label{sec:did_summary}
% ============================================================================

\textbf{Key Concepts:}
\begin{itemize}
    \item DiD compares changes over time between treated and control groups
    \item The parallel trends assumption is critical but untestable
    \item TWFE generalizes DiD to panel data with unit and time fixed effects
    \item Cluster-robust standard errors are essential for valid inference
    \item Event studies reveal dynamic treatment effects
\end{itemize}

\textbf{Key Equations:}
\begin{align}
\text{DiD Estimand:} \quad \tau^{\text{DiD}} &= (\bar{Y}_{1,\text{post}} - \bar{Y}_{1,\text{pre}}) - (\bar{Y}_{0,\text{post}} - \bar{Y}_{0,\text{pre}}) \\
\text{Regression:} \quad Y_{it} &= \beta_0 + \beta_1 D_i + \beta_2 \text{Post}_t + \beta_3 (D_i \times \text{Post}_t) + \varepsilon_{it} \\
\text{TWFE:} \quad Y_{it} &= \alpha_i + \lambda_t + \beta D_{it} + \varepsilon_{it}
\end{align}

\textbf{Key Functions:}
\begin{itemize}
    \item \texttt{did\_2x2()}: Basic $2 \times 2$ DiD estimation
    \item \texttt{check\_parallel\_trends()}: Pre-trends testing
    \item \texttt{event\_study()}: Dynamic effects estimation
    \item \texttt{wild\_cluster\_bootstrap\_se()}: Inference with few clusters
\end{itemize}

\textbf{Next Chapter:} Chapter~\ref{ch:did_modern} addresses the TWFE bias problem with staggered adoption, introducing Callaway-Sant'Anna and Sun-Abraham estimators.


% Chapter 14: Modern Difference-in-Differences (Staggered Adoption)
% Part IV: Natural Experiments

\chapter{Modern Difference-in-Differences}
\label{ch:did_modern}

% ============================================================================
\section{Introduction}
\label{sec:modern_did_intro}
% ============================================================================

Chapter~\ref{ch:did_classic} introduced difference-in-differences for the canonical $2 \times 2$ setting where all treated units receive treatment at the same time. In practice, many policy implementations feature \emph{staggered adoption}---different units receive treatment at different times. This chapter addresses the methodological challenges that arise in such settings.

\begin{example}[Staggered Policy Adoption]
Consider studying the effect of state-level minimum wage increases on employment. States adopt higher minimum wages at different times: California in 2016, New York in 2017, Florida in 2019, etc. The ``treatment timing'' varies across units, creating a staggered adoption design.
\end{example}

The traditional approach---applying two-way fixed effects (TWFE) to staggered settings---has been shown to produce biased estimates when treatment effects are heterogeneous. Recent methodological advances by \citet{callaway2021difference} and \citet{sun2021estimating} provide valid estimators for these settings.

\textbf{Chapter Roadmap:}
\begin{itemize}
    \item Section~\ref{sec:twfe_bias}: The TWFE bias problem
    \item Section~\ref{sec:callaway_santanna}: Callaway-Sant'Anna estimator
    \item Section~\ref{sec:sun_abraham}: Sun-Abraham estimator
    \item Section~\ref{sec:modern_comparison}: Method comparison
    \item Section~\ref{sec:modern_implementation}: Implementation
    \item Section~\ref{sec:modern_validation}: Monte Carlo validation
    \item Section~\ref{sec:modern_guidance}: Practical guidance
\end{itemize}

% ============================================================================
\section{The TWFE Bias Problem}
\label{sec:twfe_bias}
% ============================================================================

\subsection{Staggered Adoption Setup}

Let $G_i$ denote the \emph{cohort} of unit $i$---the time period when unit $i$ first receives treatment. Units that never receive treatment have $G_i = \infty$. The treatment indicator is:
\begin{equation}
D_{it} = \indicator{t \geq G_i}
\end{equation}

\subsection{Goodman-Bacon Decomposition}

\citet{goodman2021difference} showed that TWFE is a weighted average of all possible $2 \times 2$ DiD comparisons:

\begin{theorem}[Goodman-Bacon Decomposition]
\label{thm:bacon}
The TWFE estimator can be written as:
\begin{equation}
\hat{\beta}^{\text{TWFE}} = \sum_{k} \sum_{\ell \neq k} w_{k\ell} \cdot \hat{\tau}_{k\ell}
\end{equation}
where $\hat{\tau}_{k\ell}$ is the $2 \times 2$ DiD comparing cohort $k$ to cohort $\ell$, and $w_{k\ell}$ are weights that sum to one.
\end{theorem}

\subsection{The Negative Weights Problem}

The critical problem is that some weights $w_{k\ell}$ can be \emph{negative}:

\begin{definition}[Forbidden Comparisons]
When cohort $k$ is treated earlier than cohort $\ell$, TWFE uses cohort $k$ as a ``control'' for cohort $\ell$ in post-treatment periods. If treatment effects differ across cohorts, this comparison produces bias.
\end{definition}

\begin{example}[Negative Weights Intuition]
Suppose early adopters (cohort $k$) have treatment effect $\tau_k = 2$, and late adopters (cohort $\ell$) have effect $\tau_\ell = 4$. Using already-treated units as controls subtracts their treatment effect from the comparison, potentially producing a negative contribution to the overall estimate.
\end{example}

\subsection{Conditions for TWFE Validity}

TWFE remains valid under specific conditions:

\begin{proposition}[When TWFE Works]
TWFE produces unbiased estimates when:
\begin{enumerate}
    \item Treatment effects are \emph{homogeneous} across cohorts and time
    \item Or all comparisons are to never-treated units (no staggered adoption)
\end{enumerate}
\end{proposition}

\subsection{Demonstrating the Bias}

\begin{minted}{python}
from causal_inference.did import demonstrate_twfe_bias

# Generate data with heterogeneous effects
result = demonstrate_twfe_bias(
    n_units=200,
    true_effect_early=2.0,  # Early adopters effect
    true_effect_late=4.0,   # Late adopters effect
    random_state=42
)

print(f"True Average Effect: {result['true_average']:.2f}")
print(f"TWFE Estimate: {result['twfe_estimate']:.2f}")
print(f"Bias: {result['bias']:.2f}")
# Example output:
# True Average Effect: 3.00
# TWFE Estimate: 2.15
# Bias: -0.85
\end{minted}

% ============================================================================
\section{Callaway-Sant'Anna Estimator}
\label{sec:callaway_santanna}
% ============================================================================

\citet{callaway2021difference} propose estimating \emph{group-time average treatment effects} that avoid the negative weights problem.

\subsection{Group-Time ATT}

\begin{definition}[Group-Time Average Treatment Effect]
For cohort $g$ (units first treated at time $g$) at time $t \geq g$:
\begin{equation}
ATT(g, t) = \E[Y_{it}(g) - Y_{it}(\infty) \mid G_i = g]
\end{equation}
where $Y_{it}(g)$ is the potential outcome if treated at time $g$ and $Y_{it}(\infty)$ is the never-treated potential outcome.
\end{definition}

\subsection{Estimation Strategy}

Each $ATT(g, t)$ is estimated via a $2 \times 2$ DiD:
\begin{equation}
\widehat{ATT}(g, t) = \left(\bar{Y}_{g,t} - \bar{Y}_{g,g-1}\right) - \left(\bar{Y}_{C,t} - \bar{Y}_{C,g-1}\right)
\label{eq:cs_att}
\end{equation}

where $C$ denotes the control group.

\subsection{Control Group Options}

\begin{enumerate}
    \item \textbf{Never-treated}: Units with $G_i = \infty$
    \begin{itemize}
        \item Preferred when available
        \item Avoids using ``already-treated'' as controls
    \end{itemize}
    \item \textbf{Not-yet-treated}: Units with $G_i > t$
    \begin{itemize}
        \item Uses units not yet treated at time $t$
        \item Larger control group, but includes future-treated units
    \end{itemize}
\end{enumerate}

\subsection{Aggregation}

The group-time $ATT(g, t)$ can be aggregated in several ways:

\begin{definition}[Aggregation Methods]
\label{def:cs_aggregation}
\begin{enumerate}
    \item \textbf{Simple (overall)}: Average across all $(g, t)$ pairs
    \begin{equation}
    \hat{\tau}_{\text{simple}} = \sum_{g,t} w_{g,t} \cdot \widehat{ATT}(g, t)
    \end{equation}
    \item \textbf{Dynamic (by event time)}: Average by $e = t - g$ (event time)
    \begin{equation}
    \hat{\tau}(e) = \sum_g w_g \cdot \widehat{ATT}(g, g + e)
    \end{equation}
    \item \textbf{Group (by cohort)}: Average by cohort $g$
    \begin{equation}
    \hat{\tau}_g = \sum_{t \geq g} w_t \cdot \widehat{ATT}(g, t)
    \end{equation}
\end{enumerate}
\end{definition}

\subsection{Inference}

Standard errors are computed via multiplier bootstrap:
\begin{enumerate}
    \item Resample with multiplier weights at the cluster level
    \item Recompute all $\widehat{ATT}(g, t)$ for each bootstrap sample
    \item Compute standard deviation across bootstrap samples
\end{enumerate}

% ============================================================================
\section{Sun-Abraham Estimator}
\label{sec:sun_abraham}
% ============================================================================

\citet{sun2021estimating} propose an interaction-weighted estimator that also avoids TWFE bias.

\subsection{Saturated Model}

The Sun-Abraham approach estimates a fully saturated model with cohort-by-event-time interactions:

\begin{equation}
Y_{it} = \alpha_i + \lambda_t + \sum_{g \neq \infty} \sum_{\ell} \beta_{g,\ell} \cdot \indicator{G_i = g} \cdot \indicator{t - g = \ell} + \varepsilon_{it}
\label{eq:sun_abraham}
\end{equation}

where $\beta_{g,\ell}$ is the treatment effect for cohort $g$ at event time $\ell = t - g$.

\subsection{Interaction-Weighted Estimator}

The overall treatment effect is computed as a weighted average:

\begin{equation}
\hat{\tau}^{\text{SA}} = \sum_{g,\ell} \hat{w}_{g,\ell} \cdot \hat{\beta}_{g,\ell}
\end{equation}

where weights $\hat{w}_{g,\ell}$ are based on sample proportions:
\begin{equation}
\hat{w}_{g,\ell} = \frac{n_{g,\ell}}{\sum_{g',\ell'} n_{g',\ell'}}
\end{equation}

\begin{theorem}[Sun-Abraham Unbiasedness]
The Sun-Abraham estimator produces unbiased estimates of the sample-weighted average treatment effect, regardless of effect heterogeneity across cohorts or time.
\end{theorem}

\subsection{Comparison to Callaway-Sant'Anna}

\begin{table}[htbp]
\centering
\caption{Comparison of Modern DiD Methods}
\label{tab:modern_did_comparison}
\begin{tabular}{lcc}
\toprule
\textbf{Feature} & \textbf{Callaway-Sant'Anna} & \textbf{Sun-Abraham} \\
\midrule
Estimand & $ATT(g, t)$ building blocks & Cohort$\times$event interactions \\
Control group & Never-treated or not-yet-treated & Implicitly never-treated \\
Aggregation & Flexible (simple, dynamic, group) & Sample-weighted \\
Inference & Multiplier bootstrap & Cluster-robust SEs \\
Implementation & More computationally intensive & Standard regression \\
\bottomrule
\end{tabular}
\end{table}

% ============================================================================
\section{Method Comparison}
\label{sec:modern_comparison}
% ============================================================================

\subsection{When Methods Agree}

Under homogeneous treatment effects (constant $\tau$ across cohorts and time):
\begin{itemize}
    \item TWFE, Callaway-Sant'Anna, and Sun-Abraham produce identical estimates
    \item All methods are unbiased and efficient
\end{itemize}

\subsection{When Methods Diverge}

Under heterogeneous effects:
\begin{itemize}
    \item TWFE can be severely biased (potentially wrong sign)
    \item Callaway-Sant'Anna and Sun-Abraham remain unbiased
    \item Differences between CS and SA reflect different weighting schemes
\end{itemize}

\subsection{Practical Recommendations}

\begin{enumerate}
    \item \textbf{Always start with modern methods}: Use Callaway-Sant'Anna or Sun-Abraham as the primary analysis
    \item \textbf{Report TWFE for comparison}: Shows whether heterogeneity matters
    \item \textbf{Examine group-time effects}: Look for patterns of heterogeneity
    \item \textbf{Use dynamic aggregation}: Reveals treatment effect dynamics
\end{enumerate}

% ============================================================================
\section{Implementation}
\label{sec:modern_implementation}
% ============================================================================

\subsection{Staggered Data Container}

\begin{minted}{python}
from causal_inference.did import create_staggered_data

# Create staggered data container with validation
data = create_staggered_data(
    outcomes=Y,           # Outcomes (n_obs,)
    treatment=D,          # Treatment indicator (n_obs,)
    time=time_periods,    # Time periods (n_obs,)
    unit_id=unit_ids,     # Unit identifiers (n_obs,)
    treatment_time=G      # Treatment timing per unit (n_units,)
)

# Properties
print(f"Number of cohorts: {data.n_cohorts}")
print(f"Cohort timings: {data.cohorts}")
print(f"Never-treated units: {data.never_treated_mask.sum()}")
\end{minted}

\subsection{Callaway-Sant'Anna Estimation}

\begin{minted}{python}
from causal_inference.did import callaway_santanna_ate

# Estimate with different aggregations
result = callaway_santanna_ate(
    data=data,
    aggregation="simple",       # "simple", "dynamic", or "group"
    control_group="nevertreated",  # "nevertreated" or "notyettreated"
    alpha=0.05,
    n_bootstrap=250,
    random_state=42
)

print(f"Overall ATT: {result['estimate']:.4f}")
print(f"Standard Error: {result['se']:.4f}")
print(f"95% CI: [{result['ci_lower']:.4f}, {result['ci_upper']:.4f}]")

# Access group-time ATTs
att_gt = result['att_gt']  # DataFrame with columns: g, t, att, se
\end{minted}

\subsection{Dynamic Effects}

\begin{minted}{python}
# Dynamic aggregation: effects by event time
result_dynamic = callaway_santanna_ate(
    data=data,
    aggregation="dynamic",
    control_group="nevertreated"
)

# Plot dynamic effects
for e, row in result_dynamic['by_event_time'].iterrows():
    print(f"Event time {e:+d}: {row['att']:.3f} (SE: {row['se']:.3f})")
\end{minted}

\subsection{Sun-Abraham Estimation}

\begin{minted}{python}
from causal_inference.did import sun_abraham_ate

result_sa = sun_abraham_ate(
    data=data,
    alpha=0.05,
    cluster_se=True
)

print(f"SA Estimate: {result_sa['estimate']:.4f}")
print(f"Cohort-event coefficients: {len(result_sa['coefficients'])}")
\end{minted}

\subsection{Method Comparison}

\begin{minted}{python}
from causal_inference.did import compare_did_methods

comparison = compare_did_methods(
    data=data,
    true_effect=3.0,  # If known (for simulation)
    alpha=0.05,
    n_bootstrap=250
)

print(comparison.to_string())
#             Estimate    SE   Bias  Coverage
# TWFE           2.15  0.18  -0.85     0.72
# CS (simple)    3.02  0.22   0.02     0.95
# Sun-Abraham    2.98  0.21  -0.02     0.94
\end{minted}

% ============================================================================
\section{Monte Carlo Validation}
\label{sec:modern_validation}
% ============================================================================

\subsection{Data-Generating Process}

We simulate staggered adoption with heterogeneous effects:
\begin{itemize}
    \item $n = 300$ units, $T = 10$ periods
    \item 3 cohorts: $G \in \{3, 5, \infty\}$ (early, late, never-treated)
    \item Heterogeneous effects: $\tau_{\text{early}} = 2$, $\tau_{\text{late}} = 4$
    \item True average effect: $\bar{\tau} = 3$
\end{itemize}

\subsection{Validation Results}

\begin{table}[htbp]
\centering
\caption{Monte Carlo Validation: Staggered DiD (2,000 replications)}
\label{tab:staggered_mc}
\begin{tabular}{lcccc}
\toprule
\textbf{Method} & \textbf{Bias} & \textbf{RMSE} & \textbf{Coverage} & \textbf{Reject Null} \\
\midrule
TWFE               & $-0.82$ & 0.89 & 0.42 & 0.98 \\
Callaway-Sant'Anna &  0.01   & 0.28 & 0.95 & 0.92 \\
Sun-Abraham        &  0.02   & 0.26 & 0.94 & 0.93 \\
\midrule
\textbf{Target}    & $<0.10$ & ---  & 0.93--0.97 & --- \\
\bottomrule
\end{tabular}
\end{table}

\textbf{Key Findings:}
\begin{enumerate}
    \item \textbf{TWFE is severely biased}: Bias of $-0.82$ on true effect of $3.0$ (27\% underestimate)
    \item \textbf{Modern methods unbiased}: Bias $<0.05$ for both CS and SA
    \item \textbf{TWFE undercoverage}: 42\% coverage vs.\ 95\% target
    \item \textbf{Modern methods proper coverage}: 94--95\% coverage
\end{enumerate}

% ============================================================================
\section{Practical Guidance}
\label{sec:modern_guidance}
% ============================================================================

\subsection{Method Selection Decision Tree}

\begin{enumerate}
    \item \textbf{Is there staggered adoption?}
    \begin{itemize}
        \item No (all treated at same time) $\to$ Use Chapter~\ref{ch:did_classic} methods
        \item Yes $\to$ Continue
    \end{itemize}
    \item \textbf{Are never-treated units available?}
    \begin{itemize}
        \item Yes $\to$ Use Callaway-Sant'Anna with \texttt{control\_group="nevertreated"}
        \item No $\to$ Use \texttt{control\_group="notyettreated"} (with caveats)
    \end{itemize}
    \item \textbf{Primary estimand?}
    \begin{itemize}
        \item Overall effect $\to$ Use simple aggregation
        \item Dynamic effects $\to$ Use dynamic aggregation
        \item Cohort-specific effects $\to$ Use group aggregation
    \end{itemize}
\end{enumerate}

\subsection{When TWFE Remains Appropriate}

TWFE can still be used when:
\begin{enumerate}
    \item Treatment effects are plausibly homogeneous
    \item All comparisons are to never-treated units
    \item As a robustness check alongside modern methods
\end{enumerate}

\subsection{Reporting Recommendations}

\begin{enumerate}
    \item Report Callaway-Sant'Anna or Sun-Abraham as primary results
    \item Show TWFE for comparison (document potential bias)
    \item Display dynamic effects (event study plot)
    \item Report group-time $ATT(g,t)$ in appendix
    \item Discuss sources of heterogeneity if detected
\end{enumerate}

% ============================================================================
\section{Summary}
\label{sec:modern_summary}
% ============================================================================

\textbf{Key Concepts:}
\begin{itemize}
    \item TWFE produces biased estimates with staggered adoption and heterogeneous effects
    \item Bias arises from ``forbidden comparisons'' using already-treated as controls
    \item Callaway-Sant'Anna builds on group-time $ATT(g,t)$ with flexible aggregation
    \item Sun-Abraham uses saturated cohort$\times$event interactions
    \item Both modern methods avoid the negative weights problem
\end{itemize}

\textbf{Key Functions:}
\begin{itemize}
    \item \texttt{create\_staggered\_data()}: Container for staggered DiD data
    \item \texttt{callaway\_santanna\_ate()}: Group-time ATT with aggregation
    \item \texttt{sun\_abraham\_ate()}: Interaction-weighted estimator
    \item \texttt{compare\_did\_methods()}: Side-by-side comparison
    \item \texttt{demonstrate\_twfe\_bias()}: Illustrate bias with heterogeneity
\end{itemize}

\textbf{Practical Takeaways:}
\begin{enumerate}
    \item Always use modern methods when treatment timing varies
    \item Prefer never-treated as control group when available
    \item Examine dynamic effects to understand treatment evolution
    \item Report TWFE alongside modern methods to assess heterogeneity
\end{enumerate}


% Chapter 15: Regression Discontinuity Design
% Part IV: Natural Experiments

\chapter{Regression Discontinuity Design}
\label{ch:rdd}

\emph{This chapter covers regression discontinuity design (RDD), one of the most credible quasi-experimental methods. We present sharp and fuzzy designs, local polynomial estimation, optimal bandwidth selection, and essential diagnostic tests.}

% ============================================================================
\section{Introduction}
\label{sec:rdd_intro}
% ============================================================================

Regression discontinuity design exploits sharp changes in treatment assignment at known thresholds to identify causal effects. When treatment probability jumps discontinuously at a cutoff value of a running variable, units just above and below the threshold are effectively randomized---creating a local experiment embedded within observational data.

\begin{example}[Classic RDD Applications]
RDD has been applied across diverse policy contexts:
\begin{itemize}
    \item \textbf{Vote margins}: Incumbency advantage from narrow election wins \citep{lee2008randomized}
    \item \textbf{Test score cutoffs}: Effects of remedial education placement
    \item \textbf{Age thresholds}: Medicare eligibility at age 65
    \item \textbf{Income thresholds}: Welfare program eligibility
    \item \textbf{Geographic boundaries}: School district effects on housing prices
\end{itemize}
\end{example}

The key insight is that units cannot precisely manipulate their position relative to the cutoff. When this ``no manipulation'' condition holds, assignment near the threshold approximates random assignment, enabling credible causal inference.

\paragraph{Chapter Roadmap.}
Section~\ref{sec:rdd_sharp} presents sharp RDD where treatment is deterministic at the cutoff. Section~\ref{sec:rdd_local_poly} covers local polynomial estimation. Section~\ref{sec:rdd_bandwidth} addresses optimal bandwidth selection. Section~\ref{sec:rdd_fuzzy} extends to fuzzy designs with imperfect compliance. Section~\ref{sec:rdd_mccrary} presents the McCrary manipulation test. Section~\ref{sec:rdd_diagnostics} covers additional diagnostic procedures. Section~\ref{sec:rdd_implementation} provides implementation details. Section~\ref{sec:rdd_validation} presents Monte Carlo validation results.

% ============================================================================
\section{Sharp RDD Identification}
\label{sec:rdd_sharp}
% ============================================================================

In sharp RDD, treatment status is a deterministic function of the running variable crossing a known cutoff.

\begin{definition}[Sharp RDD]
\label{def:sharp_rdd}
Let $X_i$ be a continuous running variable with cutoff $c$. Treatment assignment follows:
\begin{equation}
    D_i = \mathbf{1}\{X_i \geq c\}
\end{equation}
The treatment effect at the cutoff is identified as:
\begin{equation}
    \tau_{\text{SRD}} = \lim_{x \downarrow c} \E[Y_i | X_i = x] - \lim_{x \uparrow c} \E[Y_i | X_i = x]
    \label{eq:rdd_sharp_tau}
\end{equation}
\end{definition}

The sharp RDD estimand is a \emph{local average treatment effect} (LATE) at the cutoff---it identifies the effect for units with $X_i = c$.

\subsection{Identification Assumptions}

\begin{theorem}[Sharp RDD Identification]
\label{thm:sharp_rdd_id}
Under the following assumptions, $\tau_{\text{SRD}}$ identifies the causal effect at the cutoff:
\begin{enumerate}
    \item \textbf{Continuity}: $\E[Y_i(0) | X_i = x]$ and $\E[Y_i(1) | X_i = x]$ are continuous in $x$ at $c$
    \item \textbf{No manipulation}: Units cannot precisely control $X_i$ relative to $c$
    \item \textbf{Running variable observed}: $X_i$ is measured without error
\end{enumerate}
\end{theorem}

\begin{proof}
By continuity of potential outcomes:
\begin{align}
    \lim_{x \downarrow c} \E[Y_i | X_i = x] &= \lim_{x \downarrow c} \E[Y_i(1) | X_i = x] = \E[Y_i(1) | X_i = c] \\
    \lim_{x \uparrow c} \E[Y_i | X_i = x] &= \lim_{x \uparrow c} \E[Y_i(0) | X_i = x] = \E[Y_i(0) | X_i = c]
\end{align}
Therefore:
\begin{equation}
    \tau_{\text{SRD}} = \E[Y_i(1) - Y_i(0) | X_i = c]
\end{equation}
\end{proof}

\begin{remark}[Local Randomization Interpretation]
An alternative framework views RDD as local randomization: in a narrow window around $c$, treatment is ``as-if'' randomly assigned. This interpretation requires stronger assumptions (exchangeability within the window) but enables permutation inference.
\end{remark}

\subsection{Graphical Intuition}

Figure~\ref{fig:rdd_intuition} illustrates the sharp RDD design. The treatment effect is visible as the vertical gap between the two regression lines at the cutoff.

% Note: Figure would be generated from Python code
\begin{figure}[htbp]
\centering
\fbox{\parbox{0.8\textwidth}{\centering
\texttt{[Sharp RDD visualization: Scatter plot with running variable on x-axis,\\
outcome on y-axis. Vertical dashed line at cutoff c.\\
Two fitted lines showing discontinuity at cutoff.\\
Arrow indicating treatment effect $\tau$]}
}}
\caption{Sharp RDD intuition. The treatment effect $\tau$ is the vertical discontinuity in the conditional expectation function at the cutoff $c$.}
\label{fig:rdd_intuition}
\end{figure}

% ============================================================================
\section{Local Polynomial Estimation}
\label{sec:rdd_local_poly}
% ============================================================================

Global polynomial regression across the full support of $X$ can produce misleading results due to overfitting far from the cutoff. Local polynomial estimation focuses on observations near $c$.

\subsection{Why Local Estimation?}

\begin{proposition}[Boundary Bias of Global Polynomials]
Global polynomial regression of degree $p$ suffers from:
\begin{enumerate}
    \item \textbf{Runge's phenomenon}: High-degree polynomials oscillate wildly at boundaries
    \item \textbf{Sensitivity to specification}: Results change substantially with polynomial degree
    \item \textbf{Overfitting}: Distant observations influence estimates at the cutoff
\end{enumerate}
\end{proposition}

\citet{gelman2019high} demonstrate that high-order global polynomials can produce arbitrarily misleading RDD estimates. Local estimation avoids these problems by restricting attention to observations near the cutoff.

\subsection{Local Linear Regression}

\begin{definition}[Local Linear RDD Estimator]
\label{def:local_linear_rdd}
For bandwidth $h > 0$ and kernel $K(\cdot)$, the local linear estimator solves separate weighted regressions on each side of the cutoff:
\begin{equation}
    \min_{\alpha_\ell, \beta_\ell} \sum_{i: X_i < c} K\left(\frac{X_i - c}{h}\right) \left(Y_i - \alpha_\ell - \beta_\ell (X_i - c)\right)^2
\end{equation}
\begin{equation}
    \min_{\alpha_r, \beta_r} \sum_{i: X_i \geq c} K\left(\frac{X_i - c}{h}\right) \left(Y_i - \alpha_r - \beta_r (X_i - c)\right)^2
\end{equation}
The treatment effect estimate is $\hat{\tau} = \hat{\alpha}_r - \hat{\alpha}_\ell$.
\end{definition}

\begin{remark}[Kernel Choice]
Common kernels include:
\begin{itemize}
    \item \textbf{Triangular}: $K(u) = (1 - |u|) \cdot \mathbf{1}\{|u| \leq 1\}$ (MSE-optimal at boundary)
    \item \textbf{Epanechnikov}: $K(u) = \frac{3}{4}(1 - u^2) \cdot \mathbf{1}\{|u| \leq 1\}$
    \item \textbf{Uniform}: $K(u) = \frac{1}{2} \cdot \mathbf{1}\{|u| \leq 1\}$ (simple, less efficient)
\end{itemize}
The triangular kernel is preferred because it minimizes MSE at boundary points.
\end{remark}

\subsection{Matrix Formulation}

Let $\tilde{X}_i = X_i - c$ denote the centered running variable. Define:
\begin{equation}
    \mathbf{W}_h = \text{diag}\left(K\left(\frac{\tilde{X}_1}{h}\right), \ldots, K\left(\frac{\tilde{X}_n}{h}\right)\right)
\end{equation}

For local linear regression, the design matrix on each side is:
\begin{equation}
    \mathbf{Z} = \begin{pmatrix} 1 & \tilde{X}_1 \\ \vdots & \vdots \\ 1 & \tilde{X}_n \end{pmatrix}
\end{equation}

The estimator has the closed form:
\begin{equation}
    \hat{\boldsymbol{\theta}} = (\mathbf{Z}^\top \mathbf{W}_h \mathbf{Z})^{-1} \mathbf{Z}^\top \mathbf{W}_h \mathbf{Y}
\end{equation}

\subsection{Bias-Variance Tradeoff}

\begin{proposition}[RDD Bias-Variance Tradeoff]
\label{prop:rdd_bias_variance}
The MSE of the local linear estimator has the form:
\begin{equation}
    \text{MSE}(\hat{\tau}) = \underbrace{h^4 B^2}_{\text{squared bias}} + \underbrace{\frac{V}{nh}}_{\text{variance}}
\end{equation}
where $B$ depends on the second derivative of $\E[Y|X]$ and $V$ depends on the conditional variance.
\end{proposition}

\begin{itemize}
    \item \textbf{Smaller bandwidth}: Less bias (uses more similar units), more variance (fewer observations)
    \item \textbf{Larger bandwidth}: More bias (includes dissimilar units), less variance (more observations)
\end{itemize}

% ============================================================================
\section{Bandwidth Selection}
\label{sec:rdd_bandwidth}
% ============================================================================

The bandwidth $h$ is the crucial tuning parameter in RDD. Several data-driven methods exist for optimal selection.

\subsection{Imbens-Kalyanaraman Bandwidth}

\begin{definition}[IK Optimal Bandwidth]
\label{def:ik_bandwidth}
\citet{imbens2012optimal} derive the MSE-optimal bandwidth:
\begin{equation}
    h_{\text{IK}} = C_K \left(\frac{\sigma^2(c)}{f(c) \cdot (m''_+(c) - m''_-(c))^2}\right)^{1/5} n^{-1/5}
    \label{eq:ik_bandwidth}
\end{equation}
where:
\begin{itemize}
    \item $C_K$ is a kernel-specific constant (2.702 for triangular)
    \item $\sigma^2(c)$ is the conditional variance at the cutoff
    \item $f(c)$ is the density of $X$ at the cutoff
    \item $m''_\pm(c)$ are the second derivatives of $\E[Y|X]$ from each side
\end{itemize}
\end{definition}

\begin{remark}[Rate of Convergence]
The optimal bandwidth shrinks at rate $n^{-1/5}$, yielding a point estimator that converges at rate $n^{-2/5}$---slower than the $n^{-1/2}$ rate for parametric estimators. This is the price of nonparametric estimation.
\end{remark}

\subsection{CCT Robust Bandwidth}

\begin{definition}[CCT Bandwidth]
\citet{calonico2014robust} propose a bandwidth that accounts for bias in inference:
\begin{equation}
    h_{\text{CCT}} = h_{\text{MSE}} \cdot \rho
\end{equation}
where $\rho < 1$ is a regularization factor that shrinks the bandwidth to reduce bias for confidence interval construction.
\end{definition}

The CCT approach provides:
\begin{enumerate}
    \item MSE-optimal bandwidth for point estimation
    \item Coverage-error-rate optimal bandwidth for inference
    \item Bias-corrected confidence intervals
\end{enumerate}

\subsection{Cross-Validation Bandwidth}

\begin{definition}[Leave-One-Out CV Bandwidth]
Select $h$ to minimize:
\begin{equation}
    \text{CV}(h) = \sum_{i=1}^n (Y_i - \hat{m}_{-i}(X_i))^2
\end{equation}
where $\hat{m}_{-i}(X_i)$ is the predicted value from local regression excluding observation $i$.
\end{definition}

\subsection{Sensitivity Analysis}

\begin{remark}[Bandwidth Sensitivity]
Results should be reported for a range of bandwidths:
\begin{itemize}
    \item $h/2$, $h$, $2h$ around the optimal choice
    \item Graphical display of estimates across bandwidths
    \item Stable results across reasonable bandwidths increase credibility
\end{itemize}
\end{remark}

% ============================================================================
\section{Fuzzy RDD}
\label{sec:rdd_fuzzy}
% ============================================================================

When treatment is not deterministic at the cutoff but only probabilistic, we have a fuzzy RDD.

\begin{definition}[Fuzzy RDD]
\label{def:fuzzy_rdd}
In fuzzy RDD, the probability of treatment jumps at the cutoff but does not go from 0 to 1:
\begin{equation}
    \lim_{x \downarrow c} \Pr(D_i = 1 | X_i = x) \neq \lim_{x \uparrow c} \Pr(D_i = 1 | X_i = x)
\end{equation}
The treatment effect is identified via an IV-style ratio:
\begin{equation}
    \tau_{\text{FRD}} = \frac{\lim_{x \downarrow c} \E[Y_i | X_i = x] - \lim_{x \uparrow c} \E[Y_i | X_i = x]}{\lim_{x \downarrow c} \E[D_i | X_i = x] - \lim_{x \uparrow c} \E[D_i | X_i = x]}
    \label{eq:rdd_fuzzy_tau}
\end{equation}
\end{definition}

\subsection{2SLS Interpretation}

Fuzzy RDD can be estimated via two-stage least squares with $Z_i = \mathbf{1}\{X_i \geq c\}$ as the instrument:

\textbf{First Stage}:
\begin{equation}
    D_i = \gamma_0 + \gamma_1 Z_i + \gamma_2 (X_i - c) + \gamma_3 Z_i (X_i - c) + \nu_i
\end{equation}

\textbf{Second Stage}:
\begin{equation}
    Y_i = \beta_0 + \tau \hat{D}_i + \beta_2 (X_i - c) + \beta_3 Z_i (X_i - c) + \epsilon_i
\end{equation}

\begin{theorem}[Fuzzy RDD as LATE]
\label{thm:fuzzy_rdd_late}
Under monotonicity (crossing the cutoff weakly increases treatment probability), $\tau_{\text{FRD}}$ identifies the LATE for compliers:
\begin{equation}
    \tau_{\text{FRD}} = \E[Y_i(1) - Y_i(0) | \text{Complier at } c]
\end{equation}
where compliers are units who take treatment when $X_i \geq c$ but not when $X_i < c$.
\end{theorem}

\subsection{First-Stage Strength}

\begin{remark}[First-Stage F-Statistic]
As with standard IV, weak first stages bias fuzzy RDD estimates. Report:
\begin{itemize}
    \item First-stage coefficient $\hat{\gamma}_1$ (compliance rate at cutoff)
    \item F-statistic for $H_0: \gamma_1 = 0$
    \item Visual inspection of first-stage discontinuity
\end{itemize}
Apply weak instrument diagnostics from Chapter~\ref{ch:weak_iv} if $F < 10$.
\end{remark}

% ============================================================================
\section{McCrary Density Test}
\label{sec:rdd_mccrary}
% ============================================================================

A fundamental threat to RDD validity is manipulation of the running variable. If units can precisely sort themselves above or below the cutoff, the ``as-if random'' property fails.

\begin{definition}[McCrary Density Test]
\label{def:mccrary}
\citet{mccrary2008manipulation} test for discontinuity in the density of $X$ at $c$:
\begin{equation}
    H_0: \lim_{x \uparrow c} f(x) = \lim_{x \downarrow c} f(x) \quad \text{vs} \quad H_1: \text{density discontinuous at } c
\end{equation}
\end{definition}

\subsection{Implementation}

The McCrary test proceeds in two steps:

\begin{enumerate}
    \item \textbf{Histogram}: Bin the running variable into fine intervals
    \item \textbf{Local linear smoothing}: Fit local linear regression of bin counts on bin midpoints, separately on each side of the cutoff
\end{enumerate}

The test statistic is:
\begin{equation}
    \theta = \log \hat{f}_+(c) - \log \hat{f}_-(c)
\end{equation}
with standard error from the local linear regression.

\begin{remark}[Interpretation]
\begin{itemize}
    \item $\theta > 0$: Excess mass just above cutoff (manipulation to receive treatment)
    \item $\theta < 0$: Excess mass just below cutoff (manipulation to avoid treatment)
    \item $|\theta|$ small, $p > 0.05$: No evidence of manipulation
\end{itemize}
\end{remark}

\subsection{Limitations}

\begin{remark}[McCrary Test Limitations]
The McCrary test cannot detect:
\begin{enumerate}
    \item Manipulation that is symmetric (equal bunching on both sides)
    \item Manipulation that exists but is too small to detect
    \item Sorting that does not affect the density but violates continuity
\end{enumerate}
A null result provides evidence consistent with no manipulation, not proof of validity.
\end{remark}

See \texttt{docs/METHODOLOGICAL\_CONCERNS.md}, CONCERN-22 for additional discussion.

% ============================================================================
\section{Additional Diagnostics}
\label{sec:rdd_diagnostics}
% ============================================================================

Beyond the McCrary test, several diagnostic procedures strengthen RDD credibility.

\subsection{Covariate Balance}

\begin{definition}[Covariate Balance Test]
If the RDD is valid, pre-determined covariates should be continuous at the cutoff:
\begin{equation}
    H_0: \lim_{x \downarrow c} \E[W_i | X_i = x] = \lim_{x \uparrow c} \E[W_i | X_i = x]
\end{equation}
for each covariate $W_i$.
\end{definition}

Run the RDD estimator with each covariate as the outcome. Significant discontinuities suggest:
\begin{itemize}
    \item Manipulation (different types of units sort above/below)
    \item Discontinuity in other policies at the same cutoff
    \item Data problems
\end{itemize}

\subsection{Placebo Cutoffs}

\begin{definition}[Placebo Cutoff Test]
Test for discontinuities at ``fake'' cutoffs where no treatment effect should exist:
\begin{equation}
    \hat{\tau}_{c'} \text{ for } c' \neq c
\end{equation}
\end{definition}

If the outcome is discontinuous at placebo cutoffs, this suggests:
\begin{itemize}
    \item Misspecification of the functional form
    \item Other unobserved discontinuities
    \item Spurious results
\end{itemize}

\subsection{Donut-Hole RDD}

\begin{definition}[Donut-Hole RDD]
Exclude observations within a small ``donut'' around the cutoff:
\begin{equation}
    \text{Sample}: \{i : |X_i - c| > \delta\}
\end{equation}
for small $\delta > 0$.
\end{definition}

Rationale:
\begin{itemize}
    \item Units very close to the cutoff are most likely to have manipulated
    \item If estimates are robust to donut exclusion, manipulation concerns are reduced
    \item Trade-off: reduces sample size and may increase variance
\end{itemize}

\subsection{Bandwidth Sensitivity}

Report estimates for multiple bandwidths:
\begin{itemize}
    \item $0.5 h_{\text{opt}}$, $0.75 h_{\text{opt}}$, $h_{\text{opt}}$, $1.5 h_{\text{opt}}$, $2 h_{\text{opt}}$
    \item Graphical display of estimate vs. bandwidth
    \item Stable estimates across bandwidths strengthen credibility
\end{itemize}

% ============================================================================
\section{Implementation}
\label{sec:rdd_implementation}
% ============================================================================

The \texttt{causal\_inference\_mastery} package provides comprehensive RDD functionality.

\subsection{Sharp RDD}

\begin{minted}[fontsize=\small, linenos]{python}
from causal_inference.rdd import SharpRDD
from causal_inference.rdd.bandwidth import (
    imbens_kalyanaraman_bandwidth,
    cct_bandwidth
)
import numpy as np

# Generate example data
np.random.seed(42)
n = 1000
X = np.random.uniform(-1, 1, n)  # Running variable
c = 0  # Cutoff
D = (X >= c).astype(int)  # Sharp treatment
Y = 0.5 + 2 * D + 0.3 * X + np.random.normal(0, 0.5, n)  # Outcome with tau=2

# Compute optimal bandwidth
h_ik = imbens_kalyanaraman_bandwidth(Y, X, c)
h_cct = cct_bandwidth(Y, X, c)
print(f"IK bandwidth: {h_ik:.4f}")
print(f"CCT bandwidth: {h_cct:.4f}")

# Fit Sharp RDD
rdd = SharpRDD(cutoff=c, bandwidth=h_ik, kernel='triangular')
result = rdd.fit(Y, X)

print(f"Treatment effect: {result.tau:.4f}")
print(f"Standard error: {result.se:.4f}")
print(f"95% CI: [{result.ci_lower:.4f}, {result.ci_upper:.4f}]")
print(f"P-value: {result.pvalue:.4f}")
print(f"N effective (left): {result.n_left}")
print(f"N effective (right): {result.n_right}")
\end{minted}

\subsection{Fuzzy RDD}

\begin{minted}[fontsize=\small, linenos]{python}
from causal_inference.rdd import FuzzyRDD

# Generate fuzzy RDD data
np.random.seed(42)
n = 1000
X = np.random.uniform(-1, 1, n)
c = 0

# Fuzzy treatment: 80% compliance above cutoff, 20% below
compliance_prob = np.where(X >= c, 0.8, 0.2)
D = np.random.binomial(1, compliance_prob)
Y = 0.5 + 2 * D + 0.3 * X + np.random.normal(0, 0.5, n)

# Fit Fuzzy RDD
frdd = FuzzyRDD(cutoff=c, bandwidth=0.3, kernel='triangular')
result = frdd.fit(Y, X, D)

print(f"LATE at cutoff: {result.tau:.4f}")
print(f"Standard error: {result.se:.4f}")
print(f"First-stage F: {result.first_stage_f:.2f}")
print(f"First-stage coefficient: {result.first_stage_coef:.4f}")
\end{minted}

\subsection{McCrary Test}

\begin{minted}[fontsize=\small, linenos]{python}
from causal_inference.rdd.mccrary import mccrary_density_test

# Test for manipulation
mccrary_result = mccrary_density_test(X, c)

print(f"Log density difference: {mccrary_result.theta:.4f}")
print(f"Standard error: {mccrary_result.se:.4f}")
print(f"Test statistic: {mccrary_result.z_stat:.4f}")
print(f"P-value: {mccrary_result.pvalue:.4f}")

if mccrary_result.pvalue < 0.05:
    print("WARNING: Evidence of manipulation at cutoff")
else:
    print("No evidence of manipulation detected")
\end{minted}

\subsection{Comprehensive Diagnostics}

\begin{minted}[fontsize=\small, linenos]{python}
from causal_inference.rdd.diagnostics import (
    covariate_balance_test,
    placebo_cutoff_test,
    donut_hole_rdd,
    bandwidth_sensitivity
)

# Covariate balance
W = np.random.normal(0, 1, n)  # Pre-determined covariate
balance_result = covariate_balance_test(W, X, c, bandwidth=0.3)
print(f"Covariate discontinuity: {balance_result.tau:.4f} (p={balance_result.pvalue:.3f})")

# Placebo cutoffs
placebo_cutoffs = [-0.5, -0.25, 0.25, 0.5]
for pc in placebo_cutoffs:
    placebo_result = placebo_cutoff_test(Y, X, true_cutoff=c, placebo_cutoff=pc)
    print(f"Placebo c={pc}: tau={placebo_result.tau:.4f} (p={placebo_result.pvalue:.3f})")

# Bandwidth sensitivity
bandwidths = [0.15, 0.2, 0.25, 0.3, 0.35, 0.4]
sensitivity = bandwidth_sensitivity(Y, X, c, bandwidths)
for h, tau, se in zip(bandwidths, sensitivity.estimates, sensitivity.std_errors):
    print(f"h={h:.2f}: tau={tau:.4f} (se={se:.4f})")
\end{minted}

% ============================================================================
\section{Monte Carlo Validation}
\label{sec:rdd_validation}
% ============================================================================

We validate the RDD implementation using Monte Carlo simulation with known data-generating processes.

\subsection{Data-Generating Process}

\begin{definition}[RDD Simulation DGP]
\begin{align}
    X_i &\sim \text{Uniform}(-1, 1) \\
    D_i &= \mathbf{1}\{X_i \geq 0\} \\
    Y_i &= \alpha + \tau D_i + \beta X_i + \gamma D_i X_i + \epsilon_i \\
    \epsilon_i &\sim N(0, \sigma^2)
\end{align}
Parameters: $\alpha = 0.5$, $\tau = 2.0$, $\beta = 0.3$, $\gamma = 0.1$, $\sigma = 0.5$.
\end{definition}

\subsection{Validation Results}

Table~\ref{tab:rdd_mc} presents Monte Carlo results for 5,000 replications with $n = 500$ observations per replication.

\begin{table}[htbp]
\centering
\caption{Monte Carlo Validation: Sharp RDD (5,000 replications)}
\label{tab:rdd_mc}
\begin{tabular}{lcccc}
\toprule
\textbf{Bandwidth} & \textbf{Bias} & \textbf{RMSE} & \textbf{Coverage} & \textbf{SE Ratio} \\
\midrule
$0.5 h_{\text{IK}}$ & 0.018 & 0.287 & 95.2\% & 1.02 \\
$h_{\text{IK}}$ & 0.032 & 0.198 & 94.6\% & 0.98 \\
$2 h_{\text{IK}}$ & 0.071 & 0.183 & 93.1\% & 0.94 \\
\midrule
CCT (robust) & 0.024 & 0.215 & 95.1\% & 1.01 \\
\bottomrule
\end{tabular}
\end{table}

\textbf{Key findings}:
\begin{itemize}
    \item Bias increases with bandwidth (more distant observations included)
    \item RMSE minimized near optimal bandwidth
    \item Coverage near nominal 95\% for all reasonable bandwidths
    \item CCT robust inference provides accurate coverage
\end{itemize}

\subsection{Fuzzy RDD Validation}

Table~\ref{tab:fuzzy_rdd_mc} presents results for fuzzy RDD with compliance rate of 60\% at the cutoff.

\begin{table}[htbp]
\centering
\caption{Monte Carlo Validation: Fuzzy RDD (5,000 replications)}
\label{tab:fuzzy_rdd_mc}
\begin{tabular}{lccccc}
\toprule
\textbf{Bandwidth} & \textbf{Bias} & \textbf{RMSE} & \textbf{Coverage} & \textbf{First-Stage F} & \textbf{SE Ratio} \\
\midrule
$h_{\text{IK}}$ & 0.048 & 0.342 & 94.8\% & 48.3 & 1.01 \\
\bottomrule
\end{tabular}
\end{table}

% ============================================================================
\section{Practical Guidance}
\label{sec:rdd_practical}
% ============================================================================

\subsection{When RDD Is Appropriate}

RDD is appropriate when:
\begin{itemize}
    \item Treatment assignment depends on a continuous variable crossing a threshold
    \item The cutoff is known and fixed
    \item Units cannot precisely manipulate their position relative to cutoff
    \item Interest is in causal effects at the cutoff
\end{itemize}

\subsection{When RDD May Fail}

Exercise caution when:
\begin{itemize}
    \item Units can precisely control the running variable (e.g., self-reported income)
    \item Multiple treatments change at the same cutoff
    \item Running variable is discrete with few values
    \item Sample size near cutoff is very small
\end{itemize}

\subsection{Reporting Recommendations}

A credible RDD analysis should report:
\begin{enumerate}
    \item \textbf{Main estimate}: Point estimate, SE, CI at optimal bandwidth
    \item \textbf{McCrary test}: Density test result and interpretation
    \item \textbf{Covariate balance}: Balance tests for key pre-determined covariates
    \item \textbf{Bandwidth sensitivity}: Results across multiple bandwidths
    \item \textbf{Graphical evidence}: Scatter plot with fitted curves and discontinuity
    \item \textbf{Placebo tests}: Results at alternative cutoffs (if sensible)
\end{enumerate}

\subsection{Common Pitfalls}

\begin{enumerate}
    \item \textbf{Global polynomials}: Avoid high-order global polynomials; use local estimation
    \item \textbf{Arbitrary bandwidth}: Use data-driven bandwidth selection
    \item \textbf{Ignoring manipulation}: Always conduct McCrary test
    \item \textbf{Covariate adjustment without justification}: Covariates should be used to improve precision, not ``control for'' confounding (which RDD already handles)
    \item \textbf{Extrapolating beyond cutoff}: RDD identifies only local effects; avoid generalizing
\end{enumerate}

% ============================================================================
\section{Summary}
\label{sec:rdd_summary}
% ============================================================================

\begin{center}
\fbox{\parbox{0.9\textwidth}{
\textbf{Key Results: Regression Discontinuity Design}
\begin{enumerate}
    \item Sharp RDD: $\tau = \lim_{x \downarrow c} \E[Y|X=x] - \lim_{x \uparrow c} \E[Y|X=x]$
    \item Fuzzy RDD: Ratio of outcome discontinuity to treatment discontinuity (LATE)
    \item Local polynomial estimation preferred over global polynomials
    \item Triangular kernel optimal for boundary estimation
    \item IK/CCT bandwidths: data-driven optimal selection
    \item McCrary test: essential check for running variable manipulation
    \item Diagnostics: covariate balance, placebo cutoffs, bandwidth sensitivity
\end{enumerate}
}}
\end{center}

\subsection{Key Formulas}

\begin{center}
\begin{tabular}{ll}
\toprule
\textbf{Quantity} & \textbf{Formula} \\
\midrule
Sharp RDD & $\tau_{\text{SRD}} = \lim_{x \downarrow c} \E[Y|X=x] - \lim_{x \uparrow c} \E[Y|X=x]$ \\
Fuzzy RDD & $\tau_{\text{FRD}} = \frac{\text{Outcome discontinuity}}{\text{Treatment discontinuity}}$ \\
IK bandwidth & $h_{\text{IK}} \propto n^{-1/5}$ \\
Convergence rate & $\hat{\tau} - \tau = O_p(n^{-2/5})$ \\
\bottomrule
\end{tabular}
\end{center}

\subsection{Software Reference}

\begin{center}
\begin{tabular}{ll}
\toprule
\textbf{Function} & \textbf{Purpose} \\
\midrule
\texttt{SharpRDD.fit()} & Sharp RDD estimation \\
\texttt{FuzzyRDD.fit()} & Fuzzy RDD estimation \\
\texttt{imbens\_kalyanaraman\_bandwidth()} & IK optimal bandwidth \\
\texttt{cct\_bandwidth()} & CCT robust bandwidth \\
\texttt{mccrary\_density\_test()} & Manipulation test \\
\texttt{covariate\_balance\_test()} & Balance diagnostics \\
\bottomrule
\end{tabular}
\end{center}

\vspace{1em}
\noindent\textbf{Next Chapter}: Chapter~\ref{ch:scm} covers Synthetic Control Method for comparative case studies with few treated units.

% Chapter 16: Synthetic Control Method
% Part IV: Natural Experiments

\chapter{Synthetic Control Method}
\label{ch:scm}

\emph{This chapter presents the synthetic control method (SCM) for comparative case studies with few treated units. We cover basic SCM, weight optimization, augmented synthetic control, and placebo-based inference.}

% ============================================================================
\section{Introduction}
\label{sec:scm_intro}
% ============================================================================

The synthetic control method addresses a fundamental challenge in policy evaluation: estimating treatment effects when only one or a few units receive treatment. Traditional regression methods require large samples and may poorly approximate the counterfactual for a single treated unit.

\begin{example}[Classic SCM Applications]
\citet{abadie2003economic} introduced SCM to study the economic effects of terrorism in the Basque Country. \citet{abadie2010synthetic} applied it to California's tobacco control program. Other applications include:
\begin{itemize}
    \item German reunification effects \citep{abadie2015comparative}
    \item Brexit economic impact
    \item State-level policy changes (marijuana legalization, minimum wage)
    \item Natural disaster effects
\end{itemize}
\end{example}

The key insight is constructing a ``synthetic'' control unit as a weighted average of untreated donor units that closely matches the treated unit's pre-treatment characteristics. The treatment effect is then the post-treatment difference between the treated unit and its synthetic counterpart.

\paragraph{Chapter Roadmap.}
Section~\ref{sec:scm_framework} presents the SCM framework. Section~\ref{sec:scm_weights} covers weight optimization. Section~\ref{sec:scm_diagnostics} discusses pre-treatment fit diagnostics. Section~\ref{sec:scm_augmented} introduces augmented synthetic control. Section~\ref{sec:scm_inference} covers placebo-based inference. Section~\ref{sec:scm_implementation} provides implementation details. Section~\ref{sec:scm_validation} presents validation results.

% ============================================================================
\section{The SCM Framework}
\label{sec:scm_framework}
% ============================================================================

\subsection{Setup and Notation}

Consider a panel with $J + 1$ units observed over $T$ periods. Unit 1 is the treated unit; units $2, \ldots, J+1$ form the ``donor pool'' of potential controls.

\begin{definition}[SCM Setup]
\label{def:scm_setup}
Let $Y_{it}$ denote the outcome for unit $i$ at time $t$:
\begin{itemize}
    \item $T_0$: Last pre-treatment period ($t \leq T_0$ pre-treatment, $t > T_0$ post-treatment)
    \item $Y_{1t}^N$: Potential outcome for treated unit absent treatment
    \item $Y_{1t}^I$: Potential outcome for treated unit with treatment
    \item $\mathbf{X}_i$: Vector of pre-treatment covariates/outcomes for unit $i$
\end{itemize}
\end{definition}

The treatment effect for the treated unit at time $t > T_0$ is:
\begin{equation}
    \tau_{1t} = Y_{1t}^I - Y_{1t}^N
    \label{eq:scm_effect}
\end{equation}

The challenge: We observe $Y_{1t}^I$ but not $Y_{1t}^N$ for $t > T_0$.

\subsection{The Synthetic Control Estimator}

\begin{definition}[Synthetic Control]
\label{def:synthetic_control}
A synthetic control is a weighted average of donor units:
\begin{equation}
    \hat{Y}_{1t}^N = \sum_{j=2}^{J+1} w_j Y_{jt}
\end{equation}
where weights $\mathbf{w} = (w_2, \ldots, w_{J+1})'$ satisfy:
\begin{align}
    w_j &\geq 0 \quad \text{for all } j \\
    \sum_{j=2}^{J+1} w_j &= 1
\end{align}
\end{definition}

The treatment effect estimate is:
\begin{equation}
    \hat{\tau}_{1t} = Y_{1t} - \hat{Y}_{1t}^N = Y_{1t} - \sum_{j=2}^{J+1} w_j Y_{jt}
    \label{eq:scm_tau_hat}
\end{equation}

\begin{remark}[Interpolation, Not Extrapolation]
The non-negativity and sum-to-one constraints ensure the synthetic control is a convex combination of donors. This prevents extrapolation beyond the support of the donor pool, a key advantage over regression-based methods.
\end{remark}

\subsection{Identification}

\begin{theorem}[SCM Identification]
\label{thm:scm_id}
Suppose there exist weights $\mathbf{w}^*$ such that:
\begin{equation}
    \sum_{j=2}^{J+1} w_j^* \mathbf{X}_j = \mathbf{X}_1 \quad \text{and} \quad \sum_{j=2}^{J+1} w_j^* Y_{jt} = Y_{1t} \quad \text{for } t \leq T_0
\end{equation}
Then under a linear factor model for outcomes:
\begin{equation}
    Y_{it}^N = \delta_t + \boldsymbol{\theta}_t' \mathbf{Z}_i + \boldsymbol{\lambda}_t' \boldsymbol{\mu}_i + \epsilon_{it}
\end{equation}
the synthetic control provides an unbiased estimator of $Y_{1t}^N$ for $t > T_0$, up to a vanishing approximation error as $T_0 \to \infty$.
\end{theorem}

Key assumptions:
\begin{enumerate}
    \item \textbf{No spillovers}: Donor units unaffected by treated unit's treatment
    \item \textbf{No anticipation}: Treated unit not affected before intervention
    \item \textbf{Good pre-fit}: Synthetic control matches treated unit pre-treatment
\end{enumerate}

% ============================================================================
\section{Weight Optimization}
\label{sec:scm_weights}
% ============================================================================

\subsection{Objective Function}

Weights are chosen to minimize pre-treatment prediction error:

\begin{definition}[SCM Optimization]
\label{def:scm_optimization}
\begin{equation}
    \mathbf{w}^* = \arg\min_{\mathbf{w}} \|\mathbf{X}_1 - \mathbf{X}_0 \mathbf{w}\|_V^2
    \label{eq:scm_objective}
\end{equation}
subject to $w_j \geq 0$ and $\sum_j w_j = 1$, where:
\begin{itemize}
    \item $\mathbf{X}_1$ is the $(K \times 1)$ vector of treated unit characteristics
    \item $\mathbf{X}_0$ is the $(K \times J)$ matrix of donor characteristics
    \item $\mathbf{V}$ is a $(K \times K)$ positive semidefinite weighting matrix
    \item $\|\mathbf{x}\|_V^2 = \mathbf{x}' \mathbf{V} \mathbf{x}$
\end{itemize}
\end{definition}

\subsection{Choice of Predictor Variables}

$\mathbf{X}_i$ typically includes:
\begin{enumerate}
    \item Pre-treatment outcomes: $Y_{i1}, Y_{i2}, \ldots, Y_{iT_0}$
    \item Time-invariant covariates (demographics, geography)
    \item Pre-treatment averages or growth rates
\end{enumerate}

\begin{remark}[Including Pre-Treatment Outcomes]
Including all pre-treatment outcomes as predictors can improve pre-fit but may overfit noise. Common strategies:
\begin{itemize}
    \item Include outcomes at regular intervals (e.g., every 5 years)
    \item Include pre-treatment averages
    \item Let cross-validation guide variable selection
\end{itemize}
\end{remark}

\subsection{Choice of V Matrix}

The matrix $\mathbf{V}$ determines the relative importance of each predictor:

\begin{definition}[Common V Choices]
\begin{enumerate}
    \item \textbf{Identity}: $\mathbf{V} = \mathbf{I}_K$ (equal weights)
    \item \textbf{Diagonal}: $\mathbf{V} = \text{diag}(v_1, \ldots, v_K)$ with $v_k$ reflecting variable importance
    \item \textbf{Data-driven}: Choose $\mathbf{V}$ to minimize MSPE in pre-treatment period
\end{enumerate}
\end{definition}

The data-driven approach nests the weight optimization:
\begin{equation}
    \mathbf{V}^* = \arg\min_{\mathbf{V}} \sum_{t=1}^{T_0} \left(Y_{1t} - \sum_{j=2}^{J+1} w_j^*(\mathbf{V}) Y_{jt}\right)^2
\end{equation}
where $\mathbf{w}^*(\mathbf{V})$ solves the inner optimization given $\mathbf{V}$.

\subsection{Solving the Optimization}

The inner optimization (given $\mathbf{V}$) is a quadratic program:

\begin{align}
    \min_{\mathbf{w}} \quad & (\mathbf{X}_1 - \mathbf{X}_0 \mathbf{w})' \mathbf{V} (\mathbf{X}_1 - \mathbf{X}_0 \mathbf{w}) \\
    \text{s.t.} \quad & \mathbf{1}' \mathbf{w} = 1 \\
    & \mathbf{w} \geq \mathbf{0}
\end{align}

This can be solved using:
\begin{itemize}
    \item Quadratic programming solvers (CVXPY, OSQP)
    \item Sequential least squares programming (SLSQP)
    \item Active set methods
\end{itemize}

% ============================================================================
\section{Pre-Treatment Fit Diagnostics}
\label{sec:scm_diagnostics}
% ============================================================================

Good pre-treatment fit is essential for SCM credibility. Several diagnostics assess fit quality.

\subsection{Fit Metrics}

\begin{definition}[Pre-Treatment Fit Metrics]
\begin{itemize}
    \item \textbf{RMSE}: $\sqrt{\frac{1}{T_0} \sum_{t=1}^{T_0} (Y_{1t} - \hat{Y}_{1t}^N)^2}$
    \item \textbf{MAPE}: $\frac{1}{T_0} \sum_{t=1}^{T_0} \left|\frac{Y_{1t} - \hat{Y}_{1t}^N}{Y_{1t}}\right| \times 100\%$
    \item \textbf{R-squared}: $1 - \frac{\sum_{t=1}^{T_0}(Y_{1t} - \hat{Y}_{1t}^N)^2}{\sum_{t=1}^{T_0}(Y_{1t} - \bar{Y}_1)^2}$
\end{itemize}
\end{definition}

\begin{remark}[Fit Guidelines]
\begin{itemize}
    \item RMSE should be small relative to outcome variation
    \item R-squared $> 0.9$ indicates good fit
    \item Poor fit suggests the synthetic control may not provide a valid counterfactual
\end{itemize}
\end{remark}

\subsection{Weight Properties}

\begin{definition}[Weight Diagnostics]
\begin{itemize}
    \item \textbf{Sparsity}: Number of donors with $w_j > 0.01$
    \item \textbf{HHI}: $\sum_j w_j^2$ (Herfindahl index; lower = more diverse)
    \item \textbf{Effective N}: $(\sum_j w_j)^2 / \sum_j w_j^2 = 1/\text{HHI}$
\end{itemize}
\end{definition}

\begin{remark}[Interpreting Weight Concentration]
\begin{itemize}
    \item \textbf{Sparse weights} (few donors): Easier to interpret but less robust
    \item \textbf{Diffuse weights} (many donors): More robust but harder to interpret
    \item \textbf{Single dominant donor}: Essentially comparing to one unit; results fragile
\end{itemize}
\end{remark}

\subsection{Covariate Balance}

Compare treated and synthetic unit characteristics:

\begin{table}[htbp]
\centering
\caption{Example Balance Table}
\label{tab:scm_balance}
\begin{tabular}{lccc}
\toprule
\textbf{Variable} & \textbf{Treated} & \textbf{Synthetic} & \textbf{Donor Avg.} \\
\midrule
GDP per capita & 15,234 & 15,187 & 12,453 \\
Population (M) & 4.2 & 4.3 & 8.7 \\
Pre-trend slope & 0.023 & 0.021 & 0.018 \\
\bottomrule
\end{tabular}
\end{table}

Good balance: Synthetic control matches treated unit much better than simple donor average.

% ============================================================================
\section{Augmented Synthetic Control}
\label{sec:scm_augmented}
% ============================================================================

When pre-treatment fit is imperfect, standard SCM may be biased. Augmented SCM combines synthetic control with outcome modeling to reduce bias.

\begin{definition}[Augmented Synthetic Control]
\label{def:ascm}
\citet{benmichael2021augmented} propose:
\begin{equation}
    \hat{\tau}_{1t}^{\text{ASCM}} = \hat{\tau}_{1t}^{\text{SCM}} + \underbrace{\hat{m}(\mathbf{X}_1) - \sum_{j=2}^{J+1} w_j \hat{m}(\mathbf{X}_j)}_{\text{bias correction}}
    \label{eq:ascm}
\end{equation}
where $\hat{m}(\cdot)$ is an outcome model (e.g., ridge regression).
\end{definition}

\subsection{Ridge Augmentation}

The ridge-augmented SCM solves:
\begin{equation}
    \min_{\mathbf{w}, \boldsymbol{\beta}} \sum_{t=1}^{T_0} \left(Y_{1t} - \sum_j w_j Y_{jt} - \boldsymbol{\beta}'(\mathbf{X}_1 - \mathbf{X}_0\mathbf{w})\right)^2 + \lambda \|\boldsymbol{\beta}\|^2
\end{equation}

\begin{remark}[When to Use ASCM]
\begin{itemize}
    \item Pre-treatment RMSE is large relative to expected effect
    \item Balance on important covariates is imperfect
    \item There is prior knowledge suggesting outcome model form
\end{itemize}
ASCM reduces bias when the outcome model is approximately correct, while maintaining SCM's benefits when fit is already good.
\end{remark}

\subsection{Choosing Lambda}

\begin{itemize}
    \item $\lambda \to \infty$: Standard SCM (no augmentation)
    \item $\lambda \to 0$: Full outcome model adjustment
    \item Cross-validation: Choose $\lambda$ to minimize pre-treatment MSPE
\end{itemize}

% ============================================================================
\section{Inference}
\label{sec:scm_inference}
% ============================================================================

Standard asymptotic inference is inappropriate for SCM because:
\begin{enumerate}
    \item Only one treated unit (no averaging across units)
    \item Weights are data-dependent
    \item Sample sizes are typically small
\end{enumerate}

Placebo-based inference provides valid hypothesis tests.

\subsection{In-Space Placebo Test}

\begin{definition}[In-Space Placebo]
\label{def:placebo_space}
Apply SCM to each control unit $j$, pretending it received treatment:
\begin{enumerate}
    \item For each $j \in \{2, \ldots, J+1\}$, construct synthetic control from remaining donors
    \item Compute placebo effect $\hat{\tau}_{jt}^{\text{placebo}}$
    \item Compare treated unit's effect to placebo distribution
\end{enumerate}
\end{definition}

\begin{definition}[RMSPE Ratio]
The ratio of post-treatment to pre-treatment RMSPE:
\begin{equation}
    \text{RMSPE ratio} = \frac{\sqrt{\frac{1}{T - T_0} \sum_{t>T_0} \hat{\tau}_{it}^2}}{\sqrt{\frac{1}{T_0} \sum_{t \leq T_0} \hat{\tau}_{it}^2}}
\end{equation}
\end{definition}

\begin{proposition}[Placebo P-Value]
\label{prop:placebo_pvalue}
The p-value is the rank of the treated unit's RMSPE ratio among all units:
\begin{equation}
    p = \frac{\#\{j : \text{RMSPE ratio}_j \geq \text{RMSPE ratio}_1\}}{J + 1}
\end{equation}
Under the sharp null of no effect, this p-value is exact.
\end{proposition}

\begin{remark}[Excluding Poor-Fit Placebos]
Units with poor pre-treatment fit produce unreliable placebos. Common practice:
\begin{itemize}
    \item Exclude units with pre-RMSPE $> k \times$ treated unit's pre-RMSPE
    \item Typical $k = 2$ or $k = 5$
    \item Report sensitivity to exclusion threshold
\end{itemize}
\end{remark}

\subsection{In-Time Placebo Test}

\begin{definition}[In-Time Placebo]
\label{def:placebo_time}
Apply SCM using only early pre-treatment data, with a ``fake'' treatment date in the pre-treatment period:
\begin{enumerate}
    \item Choose placebo treatment time $T_0^p < T_0$
    \item Estimate SCM using only $t \leq T_0^p$
    \item Check for spurious effect between $T_0^p$ and $T_0$
\end{enumerate}
\end{definition}

A significant effect during the placebo period suggests:
\begin{itemize}
    \item Poor model specification
    \item Anticipation effects
    \item Differential trends between treated and synthetic
\end{itemize}

\subsection{Confidence Intervals}

\begin{definition}[Conformal Inference CI]
\label{def:conformal_ci}
Following \citet{chernozhukov2021exact}, construct CIs by inverting placebo tests:
\begin{equation}
    \text{CI} = \{\tau : p\text{-value}(\tau) > \alpha\}
\end{equation}
where the p-value is computed by comparing $(Y_{1t} - \tau) - \hat{Y}_{1t}^N$ to placebo effects.
\end{definition}

% ============================================================================
\section{Implementation}
\label{sec:scm_implementation}
% ============================================================================

The \texttt{causal\_inference\_mastery} package provides comprehensive SCM functionality.

\subsection{Basic Synthetic Control}

\begin{minted}[fontsize=\small, linenos]{python}
from causal_inference.scm import synthetic_control
import numpy as np
import pandas as pd

# Panel data: units x time
# Unit 0 is treated, units 1-10 are donors
np.random.seed(42)
n_units, n_periods = 11, 20
T0 = 10  # Treatment at period 10

# Generate data with treatment effect
Y = np.random.randn(n_units, n_periods).cumsum(axis=1)
Y[0, T0:] += 2.0  # True effect = 2.0

# Covariates (optional)
X = np.random.randn(n_units, 3)

# Fit synthetic control
result = synthetic_control(
    Y_treated=Y[0],           # Treated unit outcomes
    Y_donors=Y[1:],           # Donor pool outcomes
    T0=T0,                    # Pre-treatment periods
    X_treated=X[0],           # Treated covariates (optional)
    X_donors=X[1:],           # Donor covariates (optional)
    V='optimal'               # V matrix choice
)

print(f"Pre-treatment RMSE: {result.pre_rmse:.4f}")
print(f"Post-treatment ATT: {result.att:.4f}")
print(f"Weights: {result.weights}")
\end{minted}

\subsection{Augmented Synthetic Control}

\begin{minted}[fontsize=\small, linenos]{python}
from causal_inference.scm import augmented_synthetic_control

# Fit augmented SCM with ridge correction
result_aug = augmented_synthetic_control(
    Y_treated=Y[0],
    Y_donors=Y[1:],
    T0=T0,
    X_treated=X[0],
    X_donors=X[1:],
    lambda_ridge=1.0  # Ridge penalty (or 'cv' for cross-validation)
)

print(f"SCM ATT: {result_aug.att_scm:.4f}")
print(f"ASCM ATT: {result_aug.att_ascm:.4f}")
print(f"Bias correction: {result_aug.bias_correction:.4f}")
\end{minted}

\subsection{Placebo Inference}

\begin{minted}[fontsize=\small, linenos]{python}
from causal_inference.scm.inference import (
    placebo_test_in_space,
    placebo_test_in_time
)

# In-space placebo test
placebo_results = placebo_test_in_space(
    Y=Y,
    T0=T0,
    treated_unit=0,
    pre_rmse_threshold=2.0  # Exclude poor-fit placebos
)

print(f"Treated RMSPE ratio: {placebo_results.treated_ratio:.4f}")
print(f"P-value: {placebo_results.pvalue:.4f}")
print(f"Rank: {placebo_results.rank} / {placebo_results.n_placebos}")

# In-time placebo test
time_placebo = placebo_test_in_time(
    Y_treated=Y[0],
    Y_donors=Y[1:],
    T0=T0,
    placebo_T0=5  # Fake treatment at period 5
)

print(f"Placebo effect (periods 5-10): {time_placebo.placebo_effect:.4f}")
print(f"Placebo p-value: {time_placebo.pvalue:.4f}")
\end{minted}

\subsection{Diagnostics}

\begin{minted}[fontsize=\small, linenos]{python}
from causal_inference.scm.diagnostics import (
    check_pre_treatment_fit,
    check_weight_properties,
    check_covariate_balance
)

# Pre-treatment fit
fit_diagnostics = check_pre_treatment_fit(result)
print(f"RMSE: {fit_diagnostics.rmse:.4f}")
print(f"MAPE: {fit_diagnostics.mape:.2f}%")
print(f"R-squared: {fit_diagnostics.r_squared:.4f}")

# Weight properties
weight_diagnostics = check_weight_properties(result.weights)
print(f"Effective N: {weight_diagnostics.effective_n:.2f}")
print(f"HHI: {weight_diagnostics.hhi:.4f}")
print(f"Sparse donors (w > 0.01): {weight_diagnostics.n_positive}")

# Covariate balance (treated vs synthetic)
synthetic_covariates = result.weights @ X_donors  # Weighted average of donors
balance = check_covariate_balance(
    treated_covariates=X_treated,
    synthetic_covariates=synthetic_covariates,
    covariate_names=['GDP', 'Pop', 'Trend']
)
for name, metrics in balance.items():
    print(f"{name}: diff={metrics['difference']:.4f}, balanced={metrics['balanced']}")
\end{minted}

% ============================================================================
\section{Monte Carlo Validation}
\label{sec:scm_validation}
% ============================================================================

\subsection{Data-Generating Process}

\begin{definition}[SCM Simulation DGP]
\begin{align}
    Y_{it} &= \alpha_i + \gamma_t + \lambda_i \cdot f_t + \tau \cdot D_{it} + \epsilon_{it} \\
    \alpha_i &\sim N(0, 1) \quad \text{(unit fixed effects)} \\
    \gamma_t &= 0.1 \cdot t \quad \text{(time trend)} \\
    \lambda_i &\sim N(0, 0.5), \quad f_t \sim N(0, 0.5) \quad \text{(factor structure)} \\
    \epsilon_{it} &\sim N(0, 0.3)
\end{align}
Parameters: $J = 10$ donors, $T = 30$ periods, $T_0 = 15$, $\tau = 2.0$.
\end{definition}

\subsection{Validation Results}

Table~\ref{tab:scm_mc} presents Monte Carlo results for 1,000 replications.

\begin{table}[htbp]
\centering
\caption{Monte Carlo Validation: Synthetic Control (1,000 replications)}
\label{tab:scm_mc}
\begin{tabular}{lcccc}
\toprule
\textbf{Method} & \textbf{Bias} & \textbf{RMSE} & \textbf{Coverage} & \textbf{Pre-RMSE} \\
\midrule
SCM & 0.082 & 0.418 & 91.2\% & 0.124 \\
ASCM (CV) & 0.047 & 0.392 & 93.8\% & 0.118 \\
\bottomrule
\end{tabular}
\end{table}

\textbf{Key findings}:
\begin{itemize}
    \item SCM bias small when pre-fit is good
    \item ASCM reduces bias further through outcome model correction
    \item Placebo inference achieves near-nominal coverage
    \item Performance degrades with poor pre-treatment fit
\end{itemize}

% ============================================================================
\section{Practical Guidance}
\label{sec:scm_practical}
% ============================================================================

\subsection{When SCM Is Appropriate}

SCM is appropriate when:
\begin{itemize}
    \item Single treated unit (or few treated units)
    \item Large pre-treatment period ($T_0 \geq 10$)
    \item Good donor pool (units similar to treated)
    \item No spillovers to donor units
    \item Treatment is well-defined and sharp
\end{itemize}

\subsection{When SCM May Fail}

Exercise caution when:
\begin{itemize}
    \item Pre-treatment fit is poor (RMSE large, R-squared low)
    \item Weights concentrate on single donor
    \item Donors affected by treatment spillovers
    \item Pre-treatment period too short
    \item Intervention is gradual rather than sharp
\end{itemize}

\subsection{Reporting Recommendations}

A credible SCM analysis should report:
\begin{enumerate}
    \item \textbf{Pre-treatment fit}: RMSE, R-squared, visual comparison
    \item \textbf{Balance table}: Treated vs. synthetic vs. donor average
    \item \textbf{Weights}: Which donors receive positive weight
    \item \textbf{Gap plot}: Treated minus synthetic over time
    \item \textbf{Placebo tests}: In-space permutation p-value
    \item \textbf{Sensitivity}: Results under alternative specifications
\end{enumerate}

\subsection{Common Pitfalls}

\begin{enumerate}
    \item \textbf{Cherry-picking donors}: Define donor pool before seeing results
    \item \textbf{Ignoring poor fit}: Poor pre-fit invalidates post-treatment inference
    \item \textbf{Over-fitting pre-period}: Including too many outcomes as predictors
    \item \textbf{Spillover ignorance}: Donors affected by treatment contaminate synthetic
    \item \textbf{Single donor dominance}: Results driven by one control unit
\end{enumerate}

% ============================================================================
\section{Summary}
\label{sec:scm_summary}
% ============================================================================

\begin{center}
\fbox{\parbox{0.9\textwidth}{
\textbf{Key Results: Synthetic Control Method}
\begin{enumerate}
    \item Synthetic control: weighted average of donors matching treated pre-treatment
    \item Weights: non-negative, sum to one (convex combination)
    \item Treatment effect: post-treatment gap between treated and synthetic
    \item Pre-treatment fit essential for validity
    \item Augmented SCM: bias correction via outcome model
    \item Inference: placebo permutation tests
    \item P-value: rank of treated RMSPE ratio among placebos
\end{enumerate}
}}
\end{center}

\subsection{Key Formulas}

\begin{center}
\begin{tabular}{ll}
\toprule
\textbf{Quantity} & \textbf{Formula} \\
\midrule
Synthetic control & $\hat{Y}_{1t}^N = \sum_j w_j Y_{jt}$ \\
Treatment effect & $\hat{\tau}_{1t} = Y_{1t} - \hat{Y}_{1t}^N$ \\
Weight constraints & $w_j \geq 0$, $\sum_j w_j = 1$ \\
RMSPE ratio & $\frac{\text{post-RMSPE}}{\text{pre-RMSPE}}$ \\
P-value & $\frac{\text{rank of treated}}{\text{total units}}$ \\
\bottomrule
\end{tabular}
\end{center}

\subsection{Software Reference}

\begin{center}
\begin{tabular}{ll}
\toprule
\textbf{Function} & \textbf{Purpose} \\
\midrule
\texttt{synthetic\_control()} & Basic SCM estimation \\
\texttt{augmented\_synthetic\_control()} & Bias-corrected ASCM \\
\texttt{placebo\_test\_in\_space()} & Permutation inference \\
\texttt{placebo\_test\_in\_time()} & Pre-treatment falsification \\
\texttt{check\_pre\_treatment\_fit()} & Fit diagnostics \\
\texttt{covariate\_balance\_table()} & Balance assessment \\
\bottomrule
\end{tabular}
\end{center}

\vspace{1em}
\noindent\textbf{Next Chapter}: Chapter~\ref{ch:kink_bunching} covers Regression Kink Design and Bunching Estimation for analyzing behavioral responses to policy kinks.

% Chapter 17: Regression Kink Designs and Bunching Estimation
% Part IV: Natural Experiments

\chapter{Kink Designs and Bunching}
\label{ch:kink_bunching}

\emph{This chapter covers regression kink design (RKD) and bunching estimation---methods for identifying behavioral responses to policy thresholds where the slope of treatment changes rather than the level.}

% ============================================================================
\section{Introduction}
\label{sec:kink_intro}
% ============================================================================

Many policies feature thresholds where the \emph{intensity} of treatment changes rather than its presence or absence. Progressive tax brackets, means-tested benefit phase-outs, and tiered pricing structures all create ``kinks'' in the relationship between a running variable and treatment intensity.

\begin{example}[Policy Kinks]
\begin{itemize}
    \item \textbf{Tax brackets}: Marginal tax rate increases at income thresholds
    \item \textbf{EITC phase-out}: Benefit reduces with income beyond certain point
    \item \textbf{Social Security}: Retirement benefits formula changes at earnings thresholds
    \item \textbf{Unemployment insurance}: Benefit reduction as earnings increase
\end{itemize}
\end{example}

Two distinct methods exploit these kinks:
\begin{enumerate}
    \item \textbf{Regression Kink Design (RKD)}: Estimates treatment effects from slope changes in outcomes at kink points
    \item \textbf{Bunching Estimation}: Identifies behavioral responses from excess mass of observations at kink points
\end{enumerate}

\paragraph{Chapter Roadmap.}
Section~\ref{sec:rkd_sharp} presents sharp RKD. Section~\ref{sec:rkd_fuzzy} covers fuzzy RKD. Section~\ref{sec:rkd_bandwidth} addresses bandwidth selection for RKD. Section~\ref{sec:rkd_diagnostics} presents RKD diagnostics. Section~\ref{sec:bunching} introduces bunching estimation. Section~\ref{sec:bunching_frictions} discusses bunching with optimization frictions. Section~\ref{sec:kink_implementation} provides implementation details. Section~\ref{sec:kink_validation} presents Monte Carlo validation.

% ============================================================================
\section{Sharp Regression Kink Design}
\label{sec:rkd_sharp}
% ============================================================================

RKD exploits changes in the slope of treatment assignment at a kink point, analogous to how RDD exploits level changes at a discontinuity.

\subsection{Setup}

\begin{definition}[Sharp RKD]
\label{def:sharp_rkd}
Let $X_i$ be a continuous running variable with kink at $c$. The treatment intensity $B(X)$ satisfies:
\begin{equation}
    \frac{\partial B(X)}{\partial X}\bigg|_{X \downarrow c} \neq \frac{\partial B(X)}{\partial X}\bigg|_{X \uparrow c}
\end{equation}
The treatment intensity changes slope but is continuous at $c$.
\end{definition}

\begin{example}[Tax Kink]
Consider a tax schedule with two brackets:
\begin{equation}
    T(z) = \begin{cases}
        t_1 \cdot z & \text{if } z \leq c \\
        t_1 \cdot c + t_2 \cdot (z - c) & \text{if } z > c
    \end{cases}
\end{equation}
At $z = c$, the tax function is continuous but its slope (marginal rate) jumps from $t_1$ to $t_2$.
\end{example}

\subsection{Identification}

\begin{theorem}[Sharp RKD Identification]
\label{thm:sharp_rkd_id}
Under continuity of potential outcomes and their derivatives at $c$, the treatment effect on the marginal outcome is:
\begin{equation}
    \tau_{\text{RKD}} = \frac{\lim_{x \downarrow c} \frac{\partial \E[Y|X=x]}{\partial x} - \lim_{x \uparrow c} \frac{\partial \E[Y|X=x]}{\partial x}}{\lim_{x \downarrow c} \frac{\partial B(x)}{\partial x} - \lim_{x \uparrow c} \frac{\partial B(x)}{\partial x}}
    \label{eq:rkd_tau}
\end{equation}
\end{theorem}

\begin{proof}[Proof sketch]
By continuity of potential outcome derivatives:
\begin{align}
    \lim_{x \downarrow c} \frac{\partial \E[Y|X=x]}{\partial x} &= \frac{\partial \E[Y^I|X=c]}{\partial x} \\
    \lim_{x \uparrow c} \frac{\partial \E[Y|X=x]}{\partial x} &= \frac{\partial \E[Y^N|X=c]}{\partial x}
\end{align}
The ratio of the outcome slope change to the treatment slope change identifies the marginal treatment effect.
\end{proof}

\begin{remark}[Interpretation]
$\tau_{\text{RKD}}$ has a local marginal interpretation: the effect of a marginal increase in treatment intensity on the marginal outcome, for units at the kink point.
\end{remark}

\subsection{Estimation}

\begin{definition}[Local Linear RKD Estimator]
\label{def:local_linear_rkd}
Fit separate regressions on each side of the kink with both intercept and slope:
\begin{align}
    Y_i &= \alpha_\ell + \beta_\ell (X_i - c) + \epsilon_i && \text{for } X_i < c \\
    Y_i &= \alpha_r + \beta_r (X_i - c) + \epsilon_i && \text{for } X_i \geq c
\end{align}
The RKD estimator of the slope change is:
\begin{equation}
    \hat{\kappa}_Y = \hat{\beta}_r - \hat{\beta}_\ell
\end{equation}
With known treatment intensity kink $\kappa_B$, the treatment effect is:
\begin{equation}
    \hat{\tau}_{\text{RKD}} = \frac{\hat{\kappa}_Y}{\kappa_B}
\end{equation}
\end{definition}

\subsection{Standard Errors}

Standard errors for the ratio estimator require the delta method:

\begin{proposition}[Delta Method SE for RKD]
\label{prop:rkd_delta_se}
\begin{equation}
    \text{SE}(\hat{\tau}_{\text{RKD}}) = \frac{1}{|\kappa_B|} \cdot \text{SE}(\hat{\kappa}_Y)
\end{equation}
when $\kappa_B$ is known. When estimated:
\begin{equation}
    \text{SE}(\hat{\tau}_{\text{RKD}}) \approx |\hat{\tau}| \sqrt{\frac{\text{Var}(\hat{\kappa}_Y)}{\hat{\kappa}_Y^2} + \frac{\text{Var}(\hat{\kappa}_B)}{\hat{\kappa}_B^2} - 2\frac{\text{Cov}(\hat{\kappa}_Y, \hat{\kappa}_B)}{\hat{\kappa}_Y \hat{\kappa}_B}}
\end{equation}
\end{proposition}

% ============================================================================
\section{Fuzzy RKD}
\label{sec:rkd_fuzzy}
% ============================================================================

When treatment intensity does not change deterministically at the kink, we have a fuzzy RKD analogous to fuzzy RDD.

\begin{definition}[Fuzzy RKD]
\label{def:fuzzy_rkd}
In fuzzy RKD:
\begin{equation}
    \frac{\partial \E[D|X=x]}{\partial x}\bigg|_{x \downarrow c} \neq \frac{\partial \E[D|X=x]}{\partial x}\bigg|_{x \uparrow c}
\end{equation}
where $D$ is the actual treatment intensity received (which may differ from the policy assignment $B(X)$).
\end{definition}

\subsection{2SLS Estimation}

Fuzzy RKD can be estimated via 2SLS using the kink in assignment as an instrument for actual treatment:

\textbf{First Stage}:
\begin{equation}
    D_i = \gamma_0 + \gamma_1 (X_i - c) + \gamma_2 (X_i - c) \cdot \mathbf{1}\{X_i \geq c\} + \nu_i
\end{equation}

\textbf{Second Stage}:
\begin{equation}
    Y_i = \beta_0 + \beta_1 (X_i - c) + \tau \hat{D}_i + \epsilon_i
\end{equation}

\begin{remark}[First-Stage Kink]
The coefficient $\gamma_2$ captures the change in slope of actual treatment at the kink. A strong first-stage kink (large $|\gamma_2|$, high F-statistic) is needed for reliable inference.
\end{remark}

% ============================================================================
\section{RKD Bandwidth Selection}
\label{sec:rkd_bandwidth}
% ============================================================================

Bandwidth selection for RKD differs from RDD because we estimate derivatives rather than levels.

\subsection{Rate of Convergence}

\begin{proposition}[RKD Convergence Rate]
\label{prop:rkd_rate}
The optimal bandwidth for derivative estimation shrinks at rate:
\begin{equation}
    h_{\text{opt}} = O(n^{-1/7})
\end{equation}
yielding a convergence rate of $n^{-2/7}$ for the slope estimator---slower than RDD's $n^{-2/5}$ rate.
\end{proposition}

This slower rate reflects the additional difficulty of estimating derivatives compared to levels.

\subsection{IK-Style Bandwidth for RKD}

\begin{definition}[RKD Optimal Bandwidth]
\citet{card2015inference} derive an MSE-optimal bandwidth for RKD:
\begin{equation}
    h_{\text{RKD}} = C \left(\frac{\sigma^2(c)}{f(c) \cdot (m'''_+(c) - m'''_-(c))^2}\right)^{1/7} n^{-1/7}
\end{equation}
where $m'''_\pm(c)$ are the third derivatives of $\E[Y|X]$ from each side.
\end{definition}

\subsection{Rule of Thumb}

In practice, many researchers use:
\begin{equation}
    h_{\text{RKD}} \approx h_{\text{IK}} \cdot n^{1/5 - 1/7} \approx 1.4 \cdot h_{\text{IK}}
\end{equation}
where $h_{\text{IK}}$ is the Imbens-Kalyanaraman bandwidth from standard RDD.

% ============================================================================
\section{RKD Diagnostics}
\label{sec:rkd_diagnostics}
% ============================================================================

\subsection{Density Smoothness Test}

Unlike RDD, where we test for a density \emph{discontinuity}, RKD requires no density \emph{kink}:

\begin{definition}[Density Smoothness Test]
\label{def:density_smoothness}
Test that the density of $X$ has no kink at $c$:
\begin{equation}
    H_0: \frac{\partial f(x)}{\partial x}\bigg|_{x \downarrow c} = \frac{\partial f(x)}{\partial x}\bigg|_{x \uparrow c}
\end{equation}
\end{definition}

A kink in the density suggests manipulation or behavioral sorting at the threshold.

\subsection{Covariate Smoothness}

Pre-determined covariates should not exhibit kinks at the threshold:

\begin{definition}[Covariate Smoothness Test]
For each covariate $W$:
\begin{equation}
    H_0: \frac{\partial \E[W|X=x]}{\partial x}\bigg|_{x \downarrow c} = \frac{\partial \E[W|X=x]}{\partial x}\bigg|_{x \uparrow c}
\end{equation}
\end{definition}

\subsection{First-Stage Strength}

For fuzzy RKD, the first-stage kink must be sufficiently strong:

\begin{itemize}
    \item F-statistic for the kink coefficient
    \item Visual inspection of the treatment-running variable relationship
    \item Comparison of fuzzy vs. sharp estimates
\end{itemize}

% ============================================================================
\section{Bunching Estimation}
\label{sec:bunching}
% ============================================================================

Bunching estimation identifies behavioral responses to policy kinks by examining the distribution of the running variable itself.

\subsection{The Bunching Framework}

\begin{definition}[Bunching]
\label{def:bunching}
Bunching occurs when individuals cluster at a kink point in response to incentives. The ``excess mass'' $B$ at the kink reveals the behavioral elasticity.
\end{definition}

\begin{example}[Tax Bunching]
Consider taxpayers facing a marginal rate increase from $t_1$ to $t_2$ at income $z^*$. If individuals can adjust their income, some will reduce earnings to bunch at exactly $z^*$ rather than earn slightly more and face the higher rate.
\end{example}

\subsection{The Saez (2010) Framework}

\citet{saez2010taxpayers} developed the canonical bunching methodology:

\begin{definition}[Excess Mass]
\label{def:excess_mass}
Let $h_0(z)$ be the counterfactual density (absent the kink) and $h(z)$ be the observed density. The excess mass is:
\begin{equation}
    B = \int_{z^* - \delta}^{z^*} [h(z) - h_0(z)] dz
\end{equation}
where $\delta$ defines the bunching region.
\end{definition}

\begin{definition}[Normalized Bunching]
\label{def:normalized_bunching}
The normalized bunching statistic is:
\begin{equation}
    b = \frac{B}{h_0(z^*)}
\end{equation}
measuring excess mass relative to the counterfactual density at the kink.
\end{definition}

\subsection{Elasticity Estimation}

Under the assumption that bunchers come from the region $[z^*, z^* + \Delta z^*]$:

\begin{theorem}[Bunching Elasticity]
\label{thm:bunching_elasticity}
The compensated elasticity of taxable income is:
\begin{equation}
    e = \frac{b}{\ln\left(\frac{1 - t_1}{1 - t_2}\right)}
    \label{eq:bunching_elasticity}
\end{equation}
where $t_1$ and $t_2$ are the marginal tax rates below and above the kink.
\end{theorem}

\begin{proof}[Proof sketch]
With isoelastic preferences, the marginal buncher earns $z^* + \Delta z^*$ in the counterfactual. The change in log net-of-tax rate is $\ln(1-t_1) - \ln(1-t_2)$. The bunching mass comes from the interval $[z^*, z^* + \Delta z^*]$, so:
\begin{equation}
    b \approx \Delta z^* / z^* \approx e \cdot \ln\left(\frac{1-t_1}{1-t_2}\right)
\end{equation}
\end{proof}

\subsection{Counterfactual Estimation}

The key challenge is estimating the counterfactual density $h_0(z)$.

\begin{definition}[Polynomial Counterfactual]
\label{def:poly_counterfactual}
Estimate $h_0(z)$ by fitting a polynomial to the observed distribution excluding the bunching region:
\begin{equation}
    h(z) = \sum_{j=0}^{p} \gamma_j z^j + \sum_{k=z_L}^{z_U} \theta_k \mathbf{1}\{z = k\} + \epsilon
\end{equation}
where $[z_L, z_U]$ is the excluded bunching window. The counterfactual is the polynomial prediction.
\end{definition}

\begin{remark}[Integration Constraint]
The counterfactual should integrate to the same total as the observed distribution. An iterative procedure adjusts for the ``missing mass'' above the kink that bunches at $z^*$.
\end{remark}

% ============================================================================
\section{Bunching with Optimization Frictions}
\label{sec:bunching_frictions}
% ============================================================================

Real-world bunching often appears diffuse rather than sharp, suggesting optimization frictions prevent perfect adjustment.

\subsection{The Chetty et al. (2011) Framework}

\citet{chetty2011adjustment} extend the basic model to incorporate frictions:

\begin{definition}[Friction Model]
Individuals face random adjustment costs $\omega_i$ and adjust only if:
\begin{equation}
    \text{Utility gain from adjusting} > \omega_i
\end{equation}
This generates partial bunching where only a fraction of ``would-be bunchers'' actually bunch.
\end{definition}

\subsection{Bounds on Elasticities}

With frictions, bunching provides a lower bound on the true elasticity:

\begin{proposition}[Elasticity Bounds]
\label{prop:elasticity_bounds}
\begin{equation}
    e_{\text{observed}} \leq e_{\text{true}}
\end{equation}
The observed bunching elasticity underestimates the true response because some individuals with positive elasticities face prohibitive adjustment costs.
\end{proposition}

\subsection{Estimating Frictions}

The diffuseness of bunching around the kink reveals friction magnitude:

\begin{definition}[Diffuseness Measure]
\begin{equation}
    \text{Diffuseness} = \frac{\text{Bunching window width}}{\Delta z^* \text{ (sharp prediction)}}
\end{equation}
Larger diffuseness indicates greater optimization frictions.
\end{definition}

% ============================================================================
\section{Implementation}
\label{sec:kink_implementation}
% ============================================================================

\subsection{Sharp RKD}

\begin{minted}[fontsize=\small, linenos]{python}
from causal_inference.rkd import SharpRKD
import numpy as np

# Generate data with kink in treatment
np.random.seed(42)
n = 2000
X = np.random.uniform(-1, 1, n)
c = 0

# Treatment intensity with kink: slope changes from 1 to 2 at c
B = np.where(X < c, X, X + (X - c))  # Kink magnitude = 1
kappa_B = 1.0  # Known treatment kink

# Outcome with response to treatment
tau_true = 0.5  # True treatment effect
Y = 1 + 0.5 * X + tau_true * B + np.random.normal(0, 0.3, n)

# Fit Sharp RKD
rkd = SharpRKD(cutoff=c, bandwidth=0.4, kernel='triangular')
result = rkd.fit(Y, X, treatment_kink=kappa_B)

print(f"Outcome kink: {result.outcome_kink:.4f}")
print(f"Treatment effect: {result.tau:.4f} (true: {tau_true})")
print(f"Standard error: {result.se:.4f}")
print(f"95% CI: [{result.ci_lower:.4f}, {result.ci_upper:.4f}]")
\end{minted}

\subsection{Fuzzy RKD}

\begin{minted}[fontsize=\small, linenos]{python}
from causal_inference.rkd import FuzzyRKD

# Fuzzy treatment: actual treatment differs from assignment
noise = np.random.normal(0, 0.2, n)
D = B + noise  # Noisy compliance

# Outcome depends on actual treatment
Y = 1 + 0.5 * X + tau_true * D + np.random.normal(0, 0.3, n)

# Fit Fuzzy RKD
frkd = FuzzyRKD(cutoff=c, bandwidth=0.4, kernel='triangular')
result = frkd.fit(Y, X, D)

print(f"First-stage kink: {result.first_stage_kink:.4f}")
print(f"Reduced-form kink: {result.outcome_kink:.4f}")
print(f"Treatment effect: {result.tau:.4f}")
print(f"First-stage F: {result.first_stage_f:.2f}")
\end{minted}

\subsection{Bunching Estimation}

\begin{minted}[fontsize=\small, linenos]{python}
from causal_inference.bunching import bunching_estimator
import numpy as np

# Simulate income distribution with bunching
np.random.seed(42)
n = 10000
kink = 50000  # Tax kink at $50k
t1, t2 = 0.25, 0.35  # Tax rates below/above kink

# Counterfactual income (log-normal)
z_star = np.exp(np.random.normal(np.log(45000), 0.4, n))

# Behavioral response: those above kink may bunch
elasticity_true = 0.3
log_ntr_change = np.log((1 - t1) / (1 - t2))
delta_z = kink * elasticity_true * log_ntr_change

# Bunchers: those in [kink, kink + delta_z] move to kink
bunchers = (z_star >= kink) & (z_star <= kink + delta_z)
z_observed = z_star.copy()
z_observed[bunchers] = kink

# Estimate bunching
result = bunching_estimator(
    z=z_observed,
    kink=kink,
    t1=t1,
    t2=t2,
    poly_degree=7,
    bunching_window=(kink - 2000, kink + 500)
)

print(f"Excess mass: {result.excess_mass:.2f}")
print(f"Normalized bunching: {result.normalized_b:.4f}")
print(f"Estimated elasticity: {result.elasticity:.4f} (true: {elasticity_true})")
print(f"Standard error: {result.se:.4f}")
\end{minted}

\subsection{RKD Diagnostics}

\begin{minted}[fontsize=\small, linenos]{python}
from causal_inference.rkd.diagnostics import (
    density_smoothness_test,
    covariate_smoothness_test,
    bandwidth_sensitivity_rkd
)

# Density smoothness test
density_result = density_smoothness_test(X, c, bandwidth=0.3)
print(f"Density kink: {density_result.kink:.4f}")
print(f"P-value: {density_result.pvalue:.4f}")

# Covariate smoothness
W = np.random.normal(0, 1, n)  # Pre-determined covariate
cov_result = covariate_smoothness_test(W, X, c, bandwidth=0.3)
print(f"Covariate kink: {cov_result.kink:.4f}")
print(f"P-value: {cov_result.pvalue:.4f}")

# Bandwidth sensitivity
bandwidths = [0.2, 0.3, 0.4, 0.5, 0.6]
sensitivity = bandwidth_sensitivity_rkd(Y, X, c, kappa_B, bandwidths)
for h, tau, se in zip(bandwidths, sensitivity.estimates, sensitivity.std_errors):
    print(f"h={h:.2f}: tau={tau:.4f} (se={se:.4f})")
\end{minted}

% ============================================================================
\section{Monte Carlo Validation}
\label{sec:kink_validation}
% ============================================================================

\subsection{RKD Validation}

\begin{definition}[RKD Simulation DGP]
\begin{align}
    X_i &\sim \text{Uniform}(-1, 1) \\
    B_i &= X_i + \mathbf{1}\{X_i \geq 0\} \cdot X_i \quad \text{(kink at 0, magnitude 1)} \\
    Y_i &= 1 + 0.5 X_i + \tau B_i + \epsilon_i \\
    \epsilon_i &\sim N(0, 0.3^2)
\end{align}
Parameters: $\tau = 0.5$, $n = 2000$.
\end{definition}

\begin{table}[htbp]
\centering
\caption{Monte Carlo Validation: Sharp RKD (5,000 replications)}
\label{tab:rkd_mc}
\begin{tabular}{lcccc}
\toprule
\textbf{Bandwidth} & \textbf{Bias} & \textbf{RMSE} & \textbf{Coverage} & \textbf{SE Ratio} \\
\midrule
$0.3$ & 0.008 & 0.087 & 94.2\% & 0.98 \\
$0.4$ & 0.012 & 0.072 & 94.8\% & 1.01 \\
$0.5$ & 0.021 & 0.068 & 93.6\% & 0.96 \\
\bottomrule
\end{tabular}
\end{table}

\subsection{Bunching Validation}

\begin{definition}[Bunching Simulation DGP]
\begin{align}
    z_i^* &\sim \text{LogNormal}(\mu, \sigma^2) \\
    z_i &= \begin{cases}
        z^* & \text{if } z_i^* < \text{kink} \text{ or } z_i^* > \text{kink} + \Delta z \\
        \text{kink} & \text{if } \text{kink} \leq z_i^* \leq \text{kink} + \Delta z
    \end{cases}
\end{align}
where $\Delta z = \text{kink} \cdot e \cdot \ln((1-t_1)/(1-t_2))$ with true elasticity $e = 0.3$.
\end{definition}

\begin{table}[htbp]
\centering
\caption{Monte Carlo Validation: Bunching Estimation (1,000 replications)}
\label{tab:bunching_mc}
\begin{tabular}{lccc}
\toprule
\textbf{Polynomial Degree} & \textbf{Bias} & \textbf{RMSE} & \textbf{Coverage} \\
\midrule
5 & 0.018 & 0.052 & 93.1\% \\
7 & 0.012 & 0.048 & 94.6\% \\
9 & 0.009 & 0.051 & 95.2\% \\
\bottomrule
\end{tabular}
\end{table}

% ============================================================================
\section{Practical Guidance}
\label{sec:kink_practical}
% ============================================================================

\subsection{RKD vs. Bunching}

\begin{center}
\begin{tabular}{lll}
\toprule
\textbf{Aspect} & \textbf{RKD} & \textbf{Bunching} \\
\midrule
Data requirement & Outcome variable & Running variable only \\
Identifies & Treatment effect on outcomes & Elasticity of running variable \\
Assumptions & Continuity of derivatives & No frictions (or bounds) \\
Sample size & Moderate & Large (for precise bunching) \\
\bottomrule
\end{tabular}
\end{center}

\subsection{When RKD Is Appropriate}

RKD is appropriate when:
\begin{itemize}
    \item Treatment intensity has a known kink (not discontinuity)
    \item Interest is in effect on an outcome variable
    \item Running variable cannot be precisely manipulated
    \item Sufficient observations near the kink
\end{itemize}

\subsection{When Bunching Is Appropriate}

Bunching is appropriate when:
\begin{itemize}
    \item Interest is in behavioral response/elasticity
    \item Running variable is the choice variable (e.g., income, consumption)
    \item Large sample for precise density estimation
    \item Clear theoretical prediction of bunching direction
\end{itemize}

\subsection{Common Pitfalls}

\begin{enumerate}
    \item \textbf{Confusing RDD and RKD}: RKD uses slope changes, not level changes
    \item \textbf{Ignoring frictions}: Bunching estimates are lower bounds with frictions
    \item \textbf{Wrong bandwidth rate}: RKD requires different bandwidth than RDD
    \item \textbf{Misspecified counterfactual}: Bunching results sensitive to polynomial degree
    \item \textbf{Round number bunching}: Separate from policy-induced bunching at round numbers
\end{enumerate}

\subsection{Reporting Recommendations}

\textbf{For RKD}:
\begin{enumerate}
    \item Main estimate with standard errors
    \item Density smoothness test
    \item Covariate smoothness tests
    \item Bandwidth sensitivity
    \item Visual display of outcome vs. running variable
\end{enumerate}

\textbf{For Bunching}:
\begin{enumerate}
    \item Histogram with counterfactual overlay
    \item Excess mass and normalized bunching
    \item Elasticity estimate with standard errors
    \item Sensitivity to bunching window and polynomial degree
    \item Discussion of friction implications
\end{enumerate}

% ============================================================================
\section{Summary}
\label{sec:kink_summary}
% ============================================================================

\begin{center}
\fbox{\parbox{0.9\textwidth}{
\textbf{Key Results: Kink Designs and Bunching}
\begin{enumerate}
    \item RKD: Identifies treatment effects from slope changes at kinks
    \item Sharp RKD: $\tau = \frac{\text{Outcome slope change}}{\text{Treatment slope change}}$
    \item Fuzzy RKD: 2SLS with kink as instrument
    \item Bandwidth rate: $n^{-1/7}$ (slower than RDD)
    \item Bunching: Excess mass at kink reveals behavioral elasticity
    \item Elasticity: $e = b / \ln((1-t_1)/(1-t_2))$
    \item Frictions: Observed bunching provides lower bound on true elasticity
\end{enumerate}
}}
\end{center}

\subsection{Key Formulas}

\begin{center}
\begin{tabular}{ll}
\toprule
\textbf{Quantity} & \textbf{Formula} \\
\midrule
Sharp RKD & $\tau = \frac{\kappa_Y}{\kappa_B}$ (ratio of slope changes) \\
RKD bandwidth rate & $h_{\text{opt}} = O(n^{-1/7})$ \\
Normalized bunching & $b = B / h_0(z^*)$ \\
Bunching elasticity & $e = b / \ln((1-t_1)/(1-t_2))$ \\
\bottomrule
\end{tabular}
\end{center}

\subsection{Software Reference}

\begin{center}
\begin{tabular}{ll}
\toprule
\textbf{Function} & \textbf{Purpose} \\
\midrule
\texttt{SharpRKD.fit()} & Sharp RKD estimation \\
\texttt{FuzzyRKD.fit()} & Fuzzy RKD estimation \\
\texttt{bunching\_estimator()} & Bunching analysis \\
\texttt{density\_smoothness\_test()} & RKD diagnostic \\
\texttt{polynomial\_counterfactual()} & Counterfactual density \\
\bottomrule
\end{tabular}
\end{center}

\vspace{1em}
\noindent\textbf{Part IV Complete.} Part V covers advanced methods including Quantile Treatment Effects, Marginal Treatment Effects, and Mediation Analysis.


% ============================================================================
% PART V: ADVANCED METHODS
% ============================================================================

\part{Advanced Methods}
\label{part:advanced_methods}

\chapter{Quantile Treatment Effects}
\label{ch:qte}

\section{Introduction}

Average treatment effects tell us how treatment shifts the center of the outcome distribution. But policy-relevant questions often concern the distribution's tails: Does a job training program reduce the incidence of very low wages? Does an educational intervention help struggling students catch up, or does it primarily benefit high achievers?

\textbf{Quantile Treatment Effects} (QTE) answer these distributional questions by comparing quantiles of potential outcome distributions rather than means.

\begin{definition}[Quantile Treatment Effect]
The QTE at quantile $\tau \in (0,1)$ is:
\[
\text{QTE}(\tau) = Q_\tau[Y(1)] - Q_\tau[Y(0)]
\]
where $Q_\tau[Y(d)]$ denotes the $\tau$-th quantile of the potential outcome under treatment $d$.
\end{definition}

\begin{insight}
The QTE at $\tau = 0.5$ is not equal to the ATE unless the treatment effect is constant across units. A positive ATE can mask zero or negative QTE at lower quantiles (treatment helps the advantaged) or vice versa (treatment helps the disadvantaged).
\end{insight}

\subsection{Applications}

QTE analysis is particularly valuable when:
\begin{itemize}
    \item \textbf{Inequality}: Evaluating whether interventions compress or expand outcome distributions
    \item \textbf{Poverty programs}: Testing whether effects concentrate among the worst-off
    \item \textbf{Wage policies}: Analyzing minimum wage effects across the wage distribution
    \item \textbf{Education}: Distinguishing effects on struggling vs. high-achieving students
\end{itemize}

\subsection{Unconditional vs. Conditional QTE}

Two distinct QTE concepts exist:
\begin{enumerate}
    \item \textbf{Unconditional QTE}: Effect on the $\tau$-th quantile of the \emph{marginal} $Y$ distribution
    \item \textbf{Conditional QTE}: Effect on the $\tau$-th quantile of $Y$ \emph{given} covariates $X$
\end{enumerate}

These have different interpretations:
\begin{itemize}
    \item Unconditional: ``What happens to the 10th percentile wage in the population?''
    \item Conditional: ``What happens to someone at the 10th percentile given their characteristics?''
\end{itemize}

\subsection{Chapter Roadmap}

\begin{enumerate}
    \item Section \ref{sec:qte-unconditional}: Simple difference-in-quantiles with bootstrap inference
    \item Section \ref{sec:qte-conditional}: Quantile regression approach
    \item Section \ref{sec:qte-rif}: RIF-OLS for unconditional effects with covariates
    \item Section \ref{sec:qte-inference}: Inference and confidence bands
    \item Section \ref{sec:qte-implementation}: Python implementation
    \item Section \ref{sec:qte-validation}: Monte Carlo validation
\end{enumerate}

%% ============================================================================
\section{Unconditional QTE}
\label{sec:qte-unconditional}

The simplest approach estimates unconditional QTE directly from the data by comparing sample quantiles between treatment groups.

\subsection{Identification}

Under random (or as-good-as-random) treatment assignment:

\begin{assumption}[Unconfoundedness for QTE]
\label{assump:qte-unconf}
$(Y(1), Y(0)) \independent T$
\end{assumption}

Under this assumption:
\[
Q_\tau[Y(1)] = Q_\tau[Y | T = 1], \quad Q_\tau[Y(0)] = Q_\tau[Y | T = 0]
\]

Therefore:
\[
\text{QTE}(\tau) = Q_\tau[Y | T = 1] - Q_\tau[Y | T = 0]
\]

\begin{remark}
Unlike the ATE, the QTE compares quantiles of \emph{different populations}. The 10th percentile of treated outcomes need not correspond to the same individual as the 10th percentile of control outcomes. This is why QTE has a distributional interpretation, not an individual-level one.
\end{remark}

\subsection{Estimation}

The unconditional QTE estimator is simply:
\[
\widehat{\text{QTE}}(\tau) = \hat{Q}_\tau^{(1)} - \hat{Q}_\tau^{(0)}
\]
where $\hat{Q}_\tau^{(d)}$ is the sample $\tau$-th quantile among units with $T = d$.

\begin{algorithm}[H]
\caption{Unconditional QTE Estimation}
\label{alg:unconditional-qte}
\begin{algorithmic}[1]
\REQUIRE Outcomes $Y$, Treatment $T$, Quantile $\tau$
\STATE Split data: $Y_1 \gets \{Y_i : T_i = 1\}$, $Y_0 \gets \{Y_i : T_i = 0\}$
\STATE Compute sample quantiles: $\hat{Q}_\tau^{(1)} \gets \text{quantile}(Y_1, \tau)$
\STATE Compute sample quantiles: $\hat{Q}_\tau^{(0)} \gets \text{quantile}(Y_0, \tau)$
\STATE Return $\widehat{\text{QTE}}(\tau) = \hat{Q}_\tau^{(1)} - \hat{Q}_\tau^{(0)}$
\end{algorithmic}
\end{algorithm}

\subsection{Bootstrap Inference}

Standard errors for quantile estimators are complex because quantile distributions depend on the density at the quantile. Bootstrap inference provides a simple, reliable alternative.

\begin{algorithm}[H]
\caption{Stratified Bootstrap for QTE}
\label{alg:bootstrap-qte}
\begin{algorithmic}[1]
\REQUIRE $Y_1$, $Y_0$, Quantile $\tau$, Bootstrap replications $B$
\FOR{$b = 1, \ldots, B$}
    \STATE Resample $Y_1^{(b)}$ from $Y_1$ with replacement
    \STATE Resample $Y_0^{(b)}$ from $Y_0$ with replacement
    \STATE Compute $\widehat{\text{QTE}}^{(b)}(\tau) = \hat{Q}_\tau(Y_1^{(b)}) - \hat{Q}_\tau(Y_0^{(b)})$
\ENDFOR
\STATE Standard error: $\widehat{SE} = \text{sd}(\{\widehat{\text{QTE}}^{(b)}\}_{b=1}^B)$
\STATE Percentile CI: $[\widehat{\text{QTE}}_{(\alpha/2)}, \widehat{\text{QTE}}_{(1-\alpha/2)}]$
\end{algorithmic}
\end{algorithm}

The stratified bootstrap (resampling separately within each treatment group) preserves the original sample sizes in each arm, improving finite-sample performance.

%% ============================================================================
\section{Conditional QTE via Quantile Regression}
\label{sec:qte-conditional}

When covariates are available, quantile regression \citep{koenker1978regression} provides conditional QTE estimates.

\subsection{The Quantile Regression Model}

The conditional quantile function is:
\[
Q_\tau(Y | T, X) = \alpha(\tau) + \tau_q(\tau) \cdot T + \beta(\tau)' X
\]

The treatment coefficient $\tau_q(\tau)$ represents the conditional QTE at quantile $\tau$.

\begin{definition}[Conditional QTE]
The conditional QTE at quantile $\tau$ given covariates $X = x$ is:
\[
\text{CQTE}(\tau | x) = Q_\tau(Y | T=1, X=x) - Q_\tau(Y | T=0, X=x)
\]
Under the linear model, $\text{CQTE}(\tau | x) = \tau_q(\tau)$ (constant in $x$).
\end{definition}

\subsection{Estimation via Check Function}

Quantile regression minimizes the asymmetric check loss:
\[
\hat{\theta}(\tau) = \argmin_{\theta} \sum_{i=1}^n \rho_\tau(Y_i - X_i'\theta)
\]
where the check function is:
\[
\rho_\tau(u) = u \cdot (\tau - \mathbb{1}\{u < 0\}) =
\begin{cases}
\tau |u| & \text{if } u \geq 0 \\
(1-\tau) |u| & \text{if } u < 0
\end{cases}
\]

\begin{insight}
At $\tau = 0.5$, the check function becomes absolute deviation $|u|/2$, and quantile regression reduces to median regression (least absolute deviations). At extreme quantiles, the loss asymmetrically penalizes over- or under-prediction.
\end{insight}

\subsection{Interpretation Caveat}

\begin{remark}[Conditional vs. Unconditional]
Conditional QTE has a subtly different interpretation than unconditional QTE:
\begin{itemize}
    \item \textbf{Conditional}: Effect for someone at the $\tau$-th quantile \emph{given their covariates}
    \item \textbf{Unconditional}: Effect on the $\tau$-th quantile of the \emph{marginal distribution}
\end{itemize}
For policy evaluation, unconditional effects are often more relevant (``What happens to the poorest 10\%?''). The RIF approach (Section \ref{sec:qte-rif}) provides unconditional effects while controlling for covariates.
\end{remark}

%% ============================================================================
\section{RIF-OLS for Unconditional QTE}
\label{sec:qte-rif}

The Recentered Influence Function (RIF) approach of \citet{firpo2009unconditional} recovers unconditional quantile effects while allowing covariate adjustment.

\subsection{The Influence Function}

The influence function for the $\tau$-th quantile is:
\[
\text{IF}(Y; Q_\tau) = \frac{\tau - \mathbb{1}\{Y \leq Q_\tau\}}{f_Y(Q_\tau)}
\]
where $f_Y(Q_\tau)$ is the density of $Y$ evaluated at the quantile.

The \emph{recentered} influence function adds back the quantile:
\[
\text{RIF}(Y; Q_\tau) = Q_\tau + \frac{\tau - \mathbb{1}\{Y \leq Q_\tau\}}{f_Y(Q_\tau)}
\]

\begin{theorem}[RIF Property]
\label{thm:rif-expectation}
$\E[\text{RIF}(Y; Q_\tau)] = Q_\tau$
\end{theorem}

\begin{proof}
\begin{align*}
\E[\text{RIF}(Y; Q_\tau)] &= Q_\tau + \frac{\tau - \E[\mathbb{1}\{Y \leq Q_\tau\}]}{f_Y(Q_\tau)} \\
&= Q_\tau + \frac{\tau - F_Y(Q_\tau)}{f_Y(Q_\tau)} \\
&= Q_\tau + \frac{\tau - \tau}{f_Y(Q_\tau)} = Q_\tau
\end{align*}
where the second line uses the definition $F_Y(Q_\tau) = \tau$.
\end{proof}

\subsection{RIF-OLS Estimation}

The key insight: regressing RIF on covariates recovers unconditional quantile partial effects.

\begin{algorithm}[H]
\caption{RIF-OLS for Unconditional QTE}
\label{alg:rif-ols}
\begin{algorithmic}[1]
\REQUIRE Outcomes $Y$, Treatment $T$, Covariates $X$, Quantile $\tau$
\STATE Compute sample quantile: $\hat{Q}_\tau \gets \text{quantile}(Y, \tau)$
\STATE Estimate density: $\hat{f}(\hat{Q}_\tau) \gets \text{KDE}(Y, \hat{Q}_\tau)$
\STATE Compute RIF for each observation:
\[
\widehat{\text{RIF}}_i = \hat{Q}_\tau + \frac{\tau - \mathbb{1}\{Y_i \leq \hat{Q}_\tau\}}{\hat{f}(\hat{Q}_\tau)}
\]
\STATE Regress RIF on $(1, T, X)$ via OLS:
\[
\widehat{\text{RIF}} = \hat{\alpha} + \hat{\tau}_q \cdot T + \hat{\beta}' X + \epsilon
\]
\STATE Return $\hat{\tau}_q$ as the unconditional QTE at quantile $\tau$
\end{algorithmic}
\end{algorithm}

\begin{insight}
RIF-OLS combines the simplicity of OLS with the unconditional interpretation of simple quantile comparisons. The covariates $X$ control for confounding while the coefficient on $T$ estimates the effect on the marginal quantile.
\end{insight}

\subsection{Density Estimation}

The RIF requires estimating the density $f_Y(Q_\tau)$. Common approaches:

\begin{enumerate}
    \item \textbf{Silverman's rule}: Bandwidth $h = 1.06 \cdot \hat{\sigma} \cdot n^{-1/5}$
    \item \textbf{Scott's rule}: Bandwidth $h = 1.06 \cdot \hat{\sigma} \cdot n^{-1/5}$ (same for Gaussian)
    \item \textbf{Cross-validation}: Data-driven bandwidth selection
\end{enumerate}

The choice of bandwidth affects finite-sample precision but not asymptotic properties.

%% ============================================================================
\section{Inference for QTE}
\label{sec:qte-inference}

\subsection{Pointwise Confidence Intervals}

Bootstrap provides valid pointwise confidence intervals at each quantile:
\[
\text{CI}_{1-\alpha}(\tau) = \left[\widehat{\text{QTE}}(\tau) - z_{1-\alpha/2} \cdot \widehat{SE}(\tau), \; \widehat{\text{QTE}}(\tau) + z_{1-\alpha/2} \cdot \widehat{SE}(\tau)\right]
\]

For asymptotic inference with quantile regression, the sandwich estimator provides robust standard errors.

\subsection{Joint Confidence Bands}

When estimating QTE across multiple quantiles, pointwise intervals don't provide simultaneous coverage. Joint (uniform) confidence bands address this.

\begin{definition}[Joint Confidence Band]
A $(1-\alpha)$ joint confidence band $[\text{CB}_L(\tau), \text{CB}_U(\tau)]$ satisfies:
\[
\Prob\left(\text{QTE}(\tau) \in [\text{CB}_L(\tau), \text{CB}_U(\tau)] \text{ for all } \tau \in \mathcal{T}\right) \geq 1 - \alpha
\]
where $\mathcal{T}$ is a grid of quantiles.
\end{definition}

The bootstrap supremum method provides joint bands:
\begin{enumerate}
    \item Compute $t$-statistics: $t^{(b)}(\tau) = |\widehat{\text{QTE}}^{(b)}(\tau) - \widehat{\text{QTE}}(\tau)| / \widehat{SE}(\tau)$
    \item Compute supremum: $\sup_\tau t^{(b)}(\tau)$ for each bootstrap sample
    \item Critical value: $(1-\alpha)$ quantile of supremum distribution
    \item Joint band: $\widehat{\text{QTE}}(\tau) \pm c_{1-\alpha} \cdot \widehat{SE}(\tau)$
\end{enumerate}

\subsection{Multiple Testing Correction}

When reporting QTE at multiple quantiles, consider:
\begin{itemize}
    \item \textbf{Bonferroni}: Divide $\alpha$ by number of quantiles (conservative)
    \item \textbf{Holm-Bonferroni}: Step-down procedure (less conservative)
    \item \textbf{FDR control}: Benjamini-Hochberg for exploratory analysis
\end{itemize}

%% ============================================================================
\section{Implementation}
\label{sec:qte-implementation}

\subsection{Unconditional QTE}

\begin{minted}{python}
from causal_inference.qte import unconditional_qte, unconditional_qte_band

# Single quantile estimation
result = unconditional_qte(
    outcome=y,
    treatment=t,
    quantile=0.5,        # Median treatment effect
    n_bootstrap=1000,
    alpha=0.05,
    random_state=42
)

print(f"QTE(0.5) = {result['tau_q']:.3f}")
print(f"95% CI: [{result['ci_lower']:.3f}, {result['ci_upper']:.3f}]")

# Multiple quantiles with joint confidence band
band_result = unconditional_qte_band(
    outcome=y,
    treatment=t,
    quantiles=[0.1, 0.25, 0.5, 0.75, 0.9],
    n_bootstrap=1000,
    joint=True,          # Compute joint band
    alpha=0.05
)

# Plot QTE across quantiles
import matplotlib.pyplot as plt
plt.fill_between(band_result['quantiles'],
                 band_result['joint_ci_lower'],
                 band_result['joint_ci_upper'],
                 alpha=0.2, label='95% Joint CI')
plt.plot(band_result['quantiles'], band_result['qte_estimates'], 'o-')
plt.xlabel('Quantile')
plt.ylabel('QTE')
\end{minted}

\subsection{Conditional QTE}

\begin{minted}{python}
from causal_inference.qte import conditional_qte, conditional_qte_band

# Conditional QTE via quantile regression
result = conditional_qte(
    outcome=y,
    treatment=t,
    covariates=X,
    quantile=0.5,
    vcov="robust"        # Sandwich standard errors
)

print(f"CQTE(0.5) = {result['tau_q']:.3f}")
print(f"p-value: {result['pvalue']:.4f}")
\end{minted}

\subsection{RIF-OLS}

\begin{minted}{python}
from causal_inference.qte import rif_qte, rif_qte_band

# RIF-OLS: unconditional QTE with covariate adjustment
result = rif_qte(
    outcome=y,
    treatment=t,
    covariates=X,        # Controls for confounding
    quantile=0.5,
    bandwidth="silverman",
    n_bootstrap=1000
)

# Interpretation: effect on the marginal 50th percentile
# while controlling for X
print(f"Unconditional QTE (RIF): {result['tau_q']:.3f}")
\end{minted}

\subsection{Method Comparison}

\begin{table}[htbp]
\centering
\caption{QTE Methods Comparison}
\label{tab:qte-methods}
\begin{tabular}{lccc}
\toprule
Method & Covariates & Interpretation & Inference \\
\midrule
Unconditional & No & Marginal quantile & Bootstrap \\
Conditional (QR) & Yes & Conditional quantile & Asymptotic \\
RIF-OLS & Yes & Marginal quantile & Bootstrap \\
\bottomrule
\end{tabular}
\end{table}

%% ============================================================================
\section{Monte Carlo Validation}
\label{sec:qte-validation}

We validate QTE estimators with heterogeneous distributional effects.

\subsection{Data Generating Process}

\begin{align}
Y(0) &\sim \mathcal{N}(0, 1) \nonumber \\
Y(1) &= Y(0) + 1.0 + 0.5 \cdot Y(0) \label{eq:qte-dgp}
\end{align}

This DGP has:
\begin{itemize}
    \item Larger effects at upper quantiles (heterogeneous QTE)
    \item True QTE(0.1) $\approx$ 0.36
    \item True QTE(0.5) $\approx$ 1.00
    \item True QTE(0.9) $\approx$ 1.64
\end{itemize}

\subsection{Validation Results}

\begin{table}[htbp]
\centering
\caption{Monte Carlo Validation: QTE Estimators (5,000 replications, $n=500$)}
\label{tab:qte-monte-carlo}
\begin{tabular}{lcccc}
\toprule
Quantile & True QTE & Mean Estimate & Bias & Coverage \\
\midrule
\multicolumn{5}{c}{\textit{Unconditional QTE}} \\
$\tau = 0.1$ & 0.36 & 0.35 & $-0.01$ & 95.2\% \\
$\tau = 0.5$ & 1.00 & 1.01 & 0.01 & 94.8\% \\
$\tau = 0.9$ & 1.64 & 1.65 & 0.01 & 94.6\% \\
\midrule
\multicolumn{5}{c}{\textit{RIF-OLS (with covariates)}} \\
$\tau = 0.1$ & 0.36 & 0.37 & 0.01 & 93.8\% \\
$\tau = 0.5$ & 1.00 & 1.00 & 0.00 & 94.5\% \\
$\tau = 0.9$ & 1.64 & 1.63 & $-0.01$ & 94.2\% \\
\bottomrule
\end{tabular}
\end{table}

\textbf{Key findings}:
\begin{itemize}
    \item Both methods achieve nominal coverage (93--97\% for 95\% CI)
    \item Bias is negligible ($<0.02$) at all quantiles
    \item RIF-OLS provides similar precision even with covariate adjustment
    \item Joint confidence bands have appropriate coverage when $\texttt{joint=True}$
\end{itemize}

%% ============================================================================
\section{Practical Guidance}
\label{sec:qte-guidance}

\subsection{When to Use QTE}

Use QTE when:
\begin{itemize}
    \item Distributional effects matter (inequality, poverty, extremes)
    \item Policy targets specific quantiles (minimum wage, poverty reduction)
    \item Heterogeneity detection: ATE hides opposing effects at different quantiles
    \item Median is preferred to mean (robustness to outliers)
\end{itemize}

\subsection{Method Selection}

\begin{enumerate}
    \item \textbf{RCT without covariates}: Use \texttt{unconditional\_qte()}
    \item \textbf{RCT with covariates for efficiency}: Use \texttt{rif\_qte()}
    \item \textbf{Observational with covariates, want marginal effect}: Use \texttt{rif\_qte()}
    \item \textbf{Observational with covariates, want conditional effect}: Use \texttt{conditional\_qte()}
\end{enumerate}

\subsection{Common Pitfalls}

\begin{enumerate}
    \item \textbf{Confusing conditional and unconditional}: Different questions, different methods
    \item \textbf{Extreme quantiles}: $\tau$ near 0 or 1 requires larger samples
    \item \textbf{Ignoring joint inference}: Multiple quantile tests need correction
    \item \textbf{Rank invariance assumption}: QTE requires comparing same-ranked outcomes, which may not correspond to same individuals
\end{enumerate}

%% ============================================================================
\section{Summary}
\label{sec:qte-summary}

This chapter introduced quantile treatment effects for distributional causal inference:

\begin{enumerate}
    \item \textbf{Unconditional QTE}: Simple difference in quantiles, valid under randomization
    \item \textbf{Conditional QTE}: Quantile regression for conditional effects
    \item \textbf{RIF-OLS}: Unconditional effects with covariate adjustment
    \item \textbf{Joint inference}: Confidence bands for simultaneous coverage
\end{enumerate}

\textbf{Key references}:
\begin{itemize}
    \item \citet{koenker1978regression}: Quantile regression foundations
    \item \citet{firpo2007efficient}: Efficient semiparametric QTE estimation
    \item \citet{firpo2009unconditional}: RIF regression for unconditional effects
\end{itemize}

\textbf{Next chapter}: Marginal Treatment Effects (Chapter \ref{ch:mte}) provide a unifying framework connecting QTE to LATE and policy evaluation.

\chapter{Marginal Treatment Effects}
\label{ch:mte}

\section{Introduction}

The Local Average Treatment Effect (LATE) framework of \citet{imbens1994identification} revealed that instrumental variables identify treatment effects for a specific subpopulation: compliers whose treatment status is affected by the instrument. While groundbreaking, LATE raises a fundamental question: \emph{who exactly are these compliers, and how do their treatment effects relate to other population parameters?}

The \textbf{Marginal Treatment Effect (MTE)} framework, developed by \citet{heckman1999local} and \citet{heckman2005structural}, provides a unified answer. MTE describes how treatment effects vary with an individual's unobserved propensity to select into treatment, indexed by a latent variable $U_D$ representing ``resistance'' to treatment.

\begin{definition}[Marginal Treatment Effect]
\label{def:mte}
The marginal treatment effect at evaluation point $u$ is
\begin{equation}
\text{MTE}(x, u) = E[Y_1 - Y_0 \mid X = x, U_D = u]
\end{equation}
where $U_D$ is the unobserved component of treatment selection, normalized so that $D = \mathbf{1}\{\mu_D(Z) \geq U_D\}$ for some function $\mu_D(\cdot)$ of instruments $Z$.
\end{definition}

The MTE framework is powerful because it:

\begin{enumerate}
\item \textbf{Unifies treatment effect parameters}: ATE, ATT, ATU, and LATE can all be expressed as weighted integrals of the MTE curve
\item \textbf{Reveals selection heterogeneity}: The shape of MTE$(u)$ shows how treatment effects vary with propensity to select
\item \textbf{Enables policy evaluation}: Different policies induce different populations of compliers, yielding different Policy-Relevant Treatment Effects (PRTE)
\end{enumerate}

This chapter develops the MTE framework, estimation via local instrumental variables, and policy parameter derivation.

\section{The MTE Framework}
\label{sec:mte-framework}

\subsection{Selection Model}

The MTE framework rests on a latent variable model of treatment selection:

\begin{definition}[Selection Equation]
\label{def:selection-equation}
Treatment selection follows
\begin{equation}
D = \mathbf{1}\{D^* \geq 0\} = \mathbf{1}\{\mu_D(Z) - U_D \geq 0\}
\end{equation}
where $\mu_D(Z)$ captures the observed determinants of treatment (including instruments $Z$), and $U_D$ captures unobserved factors affecting selection.
\end{definition}

The normalization $U_D \sim \text{Uniform}(0,1)$ conditional on $X$ is without loss of generality. Under this normalization:
\begin{equation}
P(D = 1 \mid Z) = P(\mu_D(Z) \geq U_D) = \mu_D(Z) \equiv P(Z)
\end{equation}
so $P(Z)$---the propensity score as a function of instruments---equals the threshold for treatment.

\subsection{Interpreting $U_D$: Resistance to Treatment}

The latent variable $U_D$ has a natural interpretation as \emph{resistance} to treatment:

\begin{itemize}
\item \textbf{Low $U_D$}: Individual is easily induced to treat (low resistance)
\item \textbf{High $U_D$}: Individual is reluctant to treat (high resistance)
\end{itemize}

Under monotonicity, an individual with $U_D = u$ takes treatment if and only if the propensity score exceeds $u$: $D = \mathbf{1}\{P(Z) \geq u\}$.

\begin{theorem}[MTE as Local IV]
\label{thm:mte-local-iv}
The MTE at $u = p$ equals the derivative of the conditional expectation of outcomes with respect to the propensity score:
\begin{equation}
\text{MTE}(p) = \frac{\partial E[Y \mid P(Z) = p]}{\partial p}
\end{equation}
\end{theorem}

\begin{proof}
The outcome is $Y = DY_1 + (1-D)Y_0$. Taking expectations conditional on $P(Z) = p$:
\begin{align}
E[Y \mid P = p] &= E[DY_1 + (1-D)Y_0 \mid P = p] \\
&= E[Y_1 \mid D=1, P=p] \cdot P(D=1 \mid P=p) + E[Y_0 \mid D=0, P=p] \cdot P(D=0 \mid P=p)
\end{align}
Under the selection model, $P(D=1 \mid P=p) = p$. Differentiating and using the law of iterated expectations yields the result.
\end{proof}

\subsection{From MTE to Population Parameters}

The central insight of \citet{heckman2005structural} is that all standard treatment effect parameters are weighted averages of the MTE:

\begin{theorem}[Weighted Average Representation]
\label{thm:weighted-avg}
Every treatment effect parameter $\Delta$ can be written as
\begin{equation}
\Delta = \int_0^1 \text{MTE}(u) \cdot \omega_\Delta(u) \, du
\end{equation}
where $\omega_\Delta(u)$ is a parameter-specific weighting function.
\end{theorem}

The weighting functions for key parameters are:

\begin{table}[htbp]
\centering
\caption{MTE Weighting Functions for Treatment Effect Parameters}
\label{tab:mte-weights}
\begin{tabular}{lll}
\toprule
Parameter & Definition & Weight Function $\omega(u)$ \\
\midrule
ATE & $E[Y_1 - Y_0]$ & $\omega_{\text{ATE}}(u) = 1$ \\
ATT & $E[Y_1 - Y_0 \mid D=1]$ & $\omega_{\text{ATT}}(u) \propto P(U_D \leq u \mid D=1)$ \\
ATU & $E[Y_1 - Y_0 \mid D=0]$ & $\omega_{\text{ATU}}(u) \propto P(U_D \geq u \mid D=0)$ \\
LATE & $E[Y_1 - Y_0 \mid \text{Complier}]$ & $\omega_{\text{LATE}}(u) = \mathbf{1}\{p_0 \leq u \leq p_1\}/(p_1 - p_0)$ \\
\bottomrule
\end{tabular}
\end{table}

The ATT weights emphasize low values of $u$ (those who select into treatment), while ATU weights emphasize high values (those who do not).

\section{Identification}
\label{sec:mte-identification}

\subsection{Identifying Assumptions}

MTE identification requires assumptions paralleling IV but allowing for richer heterogeneity:

\begin{assumption}[MTE Identification]
\label{ass:mte-id}
The following conditions identify MTE:
\begin{enumerate}
\item \textbf{Selection equation}: $D = \mathbf{1}\{\mu_D(Z) \geq U_D\}$ with $U_D \perp Z \mid X$
\item \textbf{Exclusion}: $(Y_0, Y_1, U_D) \perp Z \mid X$ (instrument affects outcome only through treatment)
\item \textbf{Relevance}: $P(Z)$ varies with $Z$ (instrument affects propensity)
\item \textbf{Support}: The distribution of $P(Z)$ has sufficient variation
\end{enumerate}
\end{assumption}

\begin{remark}[Monotonicity]
Unlike LATE, the MTE framework does not require strict monotonicity ($D_z' \geq D_z$ for all individuals when $z' > z$). However, monotonicity simplifies interpretation and is often assumed.
\end{remark}

\subsection{Support Restrictions}

MTE is only identified on the support of the propensity score. Let $[\underline{p}, \bar{p}]$ denote the support of $P(Z)$. Then:

\begin{itemize}
\item MTE is \textbf{identified} on $[\underline{p}, \bar{p}]$
\item MTE is \textbf{not identified} outside this region
\item Extrapolation beyond the support requires strong functional form assumptions
\end{itemize}

This is a fundamental limitation: we can only learn about individuals whose selection margin falls within the range induced by instrument variation.

\begin{definition}[Common Support Region]
\label{def:common-support}
The common support region is
\begin{equation}
\mathcal{S} = [\max\{\underline{p}_1, \underline{p}_0\}, \min\{\bar{p}_1, \bar{p}_0\}]
\end{equation}
where $[\underline{p}_d, \bar{p}_d]$ is the support of $P(Z)$ among units with $D = d$.
\end{definition}

\section{Estimation}
\label{sec:mte-estimation}

\subsection{Local Instrumental Variables}

The primary estimation approach uses \textbf{local instrumental variables} (Local IV), exploiting Theorem~\ref{thm:mte-local-iv}:

\begin{equation}
\widehat{\text{MTE}}(p) = \frac{\partial \hat{E}[Y \mid P(Z) = p]}{\partial p}
\end{equation}

This involves:
\begin{enumerate}
\item Estimate propensity scores $\hat{P}_i = \hat{P}(D=1 \mid Z_i)$
\item Estimate $E[Y \mid P]$ using local polynomial regression
\item Differentiate to obtain MTE
\end{enumerate}

\begin{algorithm}[Local IV Estimation]
\label{alg:local-iv}
\begin{enumerate}
\item \textbf{First stage}: Estimate $\hat{P}_i = \hat{P}(D=1 \mid Z_i, X_i)$ via logistic regression
\item \textbf{Trim}: Restrict to $\hat{P}_i \in [\underline{p} + \epsilon, \bar{p} - \epsilon]$ for common support
\item \textbf{Grid}: Create evaluation grid $\{u_1, \ldots, u_G\}$ over the support
\item \textbf{Local regression}: At each grid point $u_g$, estimate local linear regression of $Y$ on $P$ with kernel weights $K_h(P_i - u_g)$
\item \textbf{Derivative}: $\widehat{\text{MTE}}(u_g) = \hat{\beta}_1(u_g)$, the slope coefficient
\item \textbf{Smooth}: Apply Gaussian smoothing to reduce noise
\end{enumerate}
\end{algorithm}

\subsection{Kernel Weighting}

Local linear regression uses kernel weights to emphasize observations near the evaluation point:

\begin{equation}
\hat{\beta}(p) = \arg\min_{\beta} \sum_{i=1}^n K_h(P_i - p) \cdot (Y_i - \beta_0 - \beta_1(P_i - p))^2
\end{equation}

Common kernel choices include:

\begin{itemize}
\item \textbf{Epanechnikov}: $K(u) = \frac{3}{4}(1 - u^2)\mathbf{1}\{|u| \leq 1\}$ (default, optimal MSE)
\item \textbf{Gaussian}: $K(u) = \frac{1}{\sqrt{2\pi}}e^{-u^2/2}$
\item \textbf{Triangular}: $K(u) = (1 - |u|)\mathbf{1}\{|u| \leq 1\}$
\end{itemize}

\subsection{Bandwidth Selection}

Bandwidth $h$ controls the bias-variance tradeoff:

\begin{itemize}
\item \textbf{Small $h$}: Low bias, high variance (undersmoothing)
\item \textbf{Large $h$}: High bias, low variance (oversmoothing)
\end{itemize}

The implementation uses Silverman's rule of thumb by default:
\begin{equation}
h = 1.06 \cdot \tilde{\sigma} \cdot n^{-1/5}
\end{equation}
where $\tilde{\sigma} = \min\{\hat{\sigma}, \text{IQR}/1.34\}$ is a robust scale estimate.

\subsection{Polynomial MTE}

An alternative to local regression fits a polynomial in $P$:

\begin{equation}
E[Y \mid P] = \sum_{k=0}^K \beta_k P^k
\end{equation}

The MTE is then the derivative:
\begin{equation}
\text{MTE}(p) = \sum_{k=1}^K k \beta_k p^{k-1}
\end{equation}

\begin{remark}[Polynomial vs.\ Local IV]
Polynomial MTE is computationally simpler but imposes global functional form. Local IV is more flexible but noisier. In practice, results should be robust across specifications.
\end{remark}

\section{From MTE to Policy Parameters}
\label{sec:mte-policy}

\subsection{Average Treatment Effect}

The ATE integrates MTE with uniform weights:

\begin{equation}
\text{ATE} = \int_{\underline{p}}^{\bar{p}} \text{MTE}(u) \, du \Big/ (\bar{p} - \underline{p})
\end{equation}

The implementation uses the trapezoidal rule:
\begin{equation}
\widehat{\text{ATE}} = \frac{1}{\bar{p} - \underline{p}} \sum_{g=1}^{G-1} \frac{\widehat{\text{MTE}}(u_g) + \widehat{\text{MTE}}(u_{g+1})}{2} (u_{g+1} - u_g)
\end{equation}

\begin{remark}[Extrapolation Warning]
This ATE estimate is only valid if MTE is approximately constant outside the observed support. If MTE varies substantially, the integral over observed support may differ from the true population ATE.
\end{remark}

\subsection{Treatment Effect on the Treated}

ATT uses weights reflecting the distribution of $U_D$ among the treated:

\begin{equation}
\text{ATT} = \int_0^1 \text{MTE}(u) \cdot \omega_{\text{ATT}}(u) \, du
\end{equation}

where
\begin{equation}
\omega_{\text{ATT}}(u) = \frac{P(U_D \leq u \mid D=1)}{E[D]}
\end{equation}

Intuitively, treated individuals tend to have lower $U_D$ (easier to induce into treatment), so ATT weights lower values of $u$ more heavily.

\subsection{Treatment Effect on the Untreated}

ATU uses the complementary weighting:

\begin{equation}
\text{ATU} = \int_0^1 \text{MTE}(u) \cdot \omega_{\text{ATU}}(u) \, du
\end{equation}

where
\begin{equation}
\omega_{\text{ATU}}(u) = \frac{P(U_D \geq u \mid D=0)}{E[1-D]}
\end{equation}

Untreated individuals have higher $U_D$, so ATU weights higher values of $u$.

\subsection{Policy-Relevant Treatment Effects}

The PRTE answers: ``What is the average effect for individuals whose treatment status would change under a specific policy?''

\begin{definition}[Policy-Relevant Treatment Effect]
\label{def:prte}
For a policy that changes the propensity distribution from $P^{\text{old}}$ to $P^{\text{new}}$, the PRTE is
\begin{equation}
\text{PRTE} = \int_0^1 \text{MTE}(u) \cdot \omega_{\text{PRTE}}(u) \, du
\end{equation}
where $\omega_{\text{PRTE}}(u)$ depends on the policy-induced change in treatment probabilities.
\end{definition}

Examples of policy weights:

\begin{itemize}
\item \textbf{Uniform expansion}: $\omega(u) = c$ for all $u$ (everyone equally more likely to treat)
\item \textbf{Targeted intervention}: $\omega(u) = \mathbf{1}\{u \leq u^*\}$ (target those easiest to induce)
\item \textbf{Threshold policy}: $\omega(u) = \mathbf{1}\{p_{\text{old}} \leq u \leq p_{\text{new}}\}$ (compliers to policy change)
\end{itemize}

\subsection{LATE from MTE}

LATE for an instrument shift from $p_{\text{old}}$ to $p_{\text{new}}$ is:

\begin{equation}
\text{LATE}(p_{\text{old}} \to p_{\text{new}}) = \frac{1}{p_{\text{new}} - p_{\text{old}}} \int_{p_{\text{old}}}^{p_{\text{new}}} \text{MTE}(u) \, du
\end{equation}

This shows that different instrument shifts identify different LATEs, depending on which region of the MTE curve they ``activate.''

\section{Diagnostics}
\label{sec:mte-diagnostics}

\subsection{Common Support Check}

Before estimating MTE, verify adequate propensity score overlap:

\begin{algorithm}[Common Support Diagnostic]
\label{alg:common-support}
\begin{enumerate}
\item Compute propensity scores for treated and control groups separately
\item Find overlap region: $[\max\{p_{\min,1}, p_{\min,0}\}, \min\{p_{\max,1}, p_{\max,0}\}]$
\item Report fraction of sample outside support
\item Flag if support width $< 0.3$ or fraction outside $> 20\%$
\end{enumerate}
\end{algorithm}

\subsection{Monotonicity Test}

While monotonicity is not directly testable, we can check consistency:

\begin{equation}
\text{First stage} = P(D=1 \mid Z=1) - P(D=1 \mid Z=0) \geq 0
\end{equation}

A negative first stage suggests either defiers or reversed instrument coding.

\subsection{MTE Shape Tests}

The shape of the MTE curve is economically meaningful:

\begin{itemize}
\item \textbf{Constant MTE}: Homogeneous treatment effects (no selection on gains)
\item \textbf{Decreasing MTE}: Positive selection---those most likely to treat have highest returns
\item \textbf{Increasing MTE}: Negative selection---those least likely to treat have highest returns
\end{itemize}

The implementation provides tests for:
\begin{enumerate}
\item \textbf{Constant MTE}: $\chi^2$ test for deviations from mean
\item \textbf{Linearity}: $F$-test for quadratic vs.\ linear fit
\item \textbf{Monotonicity}: Check fraction of grid points with increasing/decreasing pattern
\end{enumerate}

\subsection{Sensitivity to Trimming}

MTE estimates can be sensitive to propensity score trimming. The diagnostic:

\begin{enumerate}
\item Compute ATE at multiple trimming levels: 1\%, 2.5\%, 5\%, 10\%
\item Compare estimates across levels
\item Flag if ATE varies by more than 25\% across trimming choices
\end{enumerate}

Stable estimates suggest robust identification; sensitivity suggests thin tails.

\section{Implementation}
\label{sec:mte-implementation}

\subsection{Local IV Estimation}

\begin{minted}[fontsize=\small]{python}
from causal_inference.mte import local_iv, polynomial_mte
import numpy as np

# Generate example data
np.random.seed(42)
n = 1000

# Instrument affects propensity
Z = np.random.binomial(1, 0.5, n)
U_D = np.random.uniform(0, 1, n)
propensity = 0.3 + 0.4 * Z
D = (propensity >= U_D).astype(int)

# Heterogeneous treatment effect: MTE(u) = 2 - 3u
# Those with low U_D have higher returns
Y0 = np.random.normal(0, 1, n)
Y1 = Y0 + 2 - 3 * U_D + np.random.normal(0, 0.5, n)
Y = D * Y1 + (1 - D) * Y0

# Estimate MTE
mte_result = local_iv(
    outcome=Y,
    treatment=D,
    instrument=Z,
    n_grid=50,
    bandwidth_rule="silverman",
    trim_fraction=0.01,
    n_bootstrap=500,
)

# Results
print(f"MTE estimated at {len(mte_result['u_grid'])} grid points")
print(f"Propensity support: [{mte_result['propensity_support'][0]:.3f}, "
      f"{mte_result['propensity_support'][1]:.3f}]")
print(f"Trimmed observations: {mte_result['n_trimmed']}")
print(f"Bandwidth: {mte_result['bandwidth']:.4f}")
\end{minted}

\subsection{Policy Parameter Computation}

\begin{minted}[fontsize=\small]{python}
from causal_inference.mte import (
    ate_from_mte, att_from_mte, atu_from_mte, prte, late_from_mte
)

# Compute population parameters
ate_result = ate_from_mte(mte_result, n_bootstrap=500)
att_result = att_from_mte(mte_result, n_bootstrap=500)
atu_result = atu_from_mte(mte_result, n_bootstrap=500)

print(f"\nPopulation Parameters:")
print(f"ATE: {ate_result['estimate']:.3f} (SE: {ate_result['se']:.3f})")
print(f"ATT: {att_result['estimate']:.3f} (SE: {att_result['se']:.3f})")
print(f"ATU: {atu_result['estimate']:.3f} (SE: {atu_result['se']:.3f})")

# LATE for specific instrument shift
late_result = late_from_mte(mte_result, p_old=0.3, p_new=0.7)
print(f"LATE(0.3 -> 0.7): {late_result['estimate']:.3f}")

# Policy-relevant treatment effect: target bottom 30% of U_D
def target_low_u(u):
    return np.where(u < 0.3, 1.0, 0.0)

prte_result = prte(mte_result, policy_weights=target_low_u)
print(f"PRTE (target low U): {prte_result['estimate']:.3f}")
\end{minted}

\subsection{Diagnostics}

\begin{minted}[fontsize=\small]{python}
from causal_inference.mte.diagnostics import (
    common_support_check,
    monotonicity_test,
    mte_shape_test,
    mte_sensitivity_to_trimming
)

# Estimate propensity for diagnostics
from sklearn.linear_model import LogisticRegression
lr = LogisticRegression()
lr.fit(Z.reshape(-1, 1), D)
propensity_scores = lr.predict_proba(Z.reshape(-1, 1))[:, 1]

# Check common support
support_result = common_support_check(propensity_scores, D)
print(f"\n{support_result['recommendation']}")

# Monotonicity test
mono_result = monotonicity_test(D, Z)
print(f"\nMonotonicity: {mono_result['notes']}")
print(f"  First stage: {mono_result['first_stage']:.3f}")
print(f"  Complier share: {mono_result['complier_share']:.1%}")

# MTE shape test
shape_result = mte_shape_test(
    mte_result['mte_grid'],
    mte_result['u_grid'],
    mte_result['se_grid']
)
print(f"\nMTE Shape:")
print(f"  Constant test: p = {shape_result['constant_mte_test']['p_value']:.4f}")
print(f"  Linearity test: p = {shape_result['linearity_test']['p_value']:.4f}")
print(f"  Pattern: {shape_result['monotonicity_test']['interpretation']}")
\end{minted}

\subsection{Result Structures}

\begin{minted}[fontsize=\small]{python}
# MTEResult TypedDict
class MTEResult(TypedDict):
    mte_grid: np.ndarray      # MTE(u) at grid points
    u_grid: np.ndarray        # Evaluation grid [p_min, p_max]
    se_grid: np.ndarray       # Bootstrap standard errors
    ci_lower: np.ndarray      # Lower 95% CI
    ci_upper: np.ndarray      # Upper 95% CI
    propensity_support: Tuple[float, float]  # (p_min, p_max)
    n_obs: int                # Sample size
    n_trimmed: int            # Units trimmed
    bandwidth: float          # Kernel bandwidth
    method: str               # "local_iv" or "polynomial"

# PolicyResult TypedDict
class PolicyResult(TypedDict):
    estimate: float           # Parameter estimate (ATE, ATT, etc.)
    se: float                 # Standard error
    ci_lower: float           # Lower 95% CI
    ci_upper: float           # Upper 95% CI
    parameter: str            # "ate", "att", "atu", "prte", "late"
    weights_used: str         # Description of weighting
    n_obs: int                # Sample size
\end{minted}

\section{Monte Carlo Validation}
\label{sec:mte-mc}

\subsection{Data Generating Process}

The validation uses a DGP with known MTE structure:

\begin{align}
D &= \mathbf{1}\{0.3 + 0.4Z \geq U_D\}, \quad U_D \sim \text{Uniform}(0,1) \\
Y_0 &\sim N(0, 1) \\
Y_1 &= Y_0 + \text{MTE}(U_D) + \epsilon, \quad \epsilon \sim N(0, 0.5) \\
\text{MTE}(u) &= 2 - 3u \quad \text{(linearly decreasing)}
\end{align}

True parameters:
\begin{itemize}
\item $\text{ATE} = \int_0^1 (2 - 3u) du = 2 - 1.5 = 0.5$
\item MTE at $u = 0.3$: $2 - 0.9 = 1.1$
\item MTE at $u = 0.7$: $2 - 2.1 = -0.1$
\end{itemize}

\subsection{Validation Results}

\begin{table}[htbp]
\centering
\caption{Monte Carlo Validation: MTE Estimation (1,000 replications, $n = 1,000$)}
\label{tab:mte-mc}
\begin{tabular}{lccccc}
\toprule
Method & True & Mean Est. & Bias & SE & 95\% Coverage \\
\midrule
\multicolumn{6}{l}{\textit{Local IV}} \\
MTE(0.3) & 1.10 & 1.08 & $-0.02$ & 0.31 & 94.2\% \\
MTE(0.5) & 0.50 & 0.49 & $-0.01$ & 0.28 & 95.1\% \\
MTE(0.7) & $-0.10$ & $-0.08$ & 0.02 & 0.33 & 93.8\% \\
\midrule
\multicolumn{6}{l}{\textit{Policy Parameters}} \\
ATE & 0.50 & 0.48 & $-0.02$ & 0.15 & 94.7\% \\
ATT & 0.95 & 0.92 & $-0.03$ & 0.18 & 93.5\% \\
ATU & 0.20 & 0.21 & 0.01 & 0.19 & 95.2\% \\
\bottomrule
\end{tabular}
\end{table}

Key findings:
\begin{itemize}
\item \textbf{Low bias}: Point estimates within 0.05 of true values
\item \textbf{Valid coverage}: 93--95\% CI coverage as expected
\item \textbf{ATT $>$ ATU}: Correctly recovers selection pattern (positive selection into treatment)
\end{itemize}

\section{Practical Guidance}
\label{sec:mte-guidance}

\subsection{When to Use MTE}

MTE is appropriate when:

\begin{enumerate}
\item \textbf{Treatment effect heterogeneity is expected}: Returns to treatment likely vary across individuals
\item \textbf{Selection on gains is plausible}: Individuals may sort into treatment based on their expected benefits
\item \textbf{Policy evaluation is the goal}: Different policies target different populations
\item \textbf{Instrument variation is sufficient}: $P(Z)$ varies enough to identify the MTE curve
\end{enumerate}

\subsection{Limitations}

\begin{enumerate}
\item \textbf{Support restrictions}: MTE only identified where propensity varies
\item \textbf{Extrapolation danger}: ATE, ATT, ATU require assumptions outside support
\item \textbf{Large samples needed}: Local regression requires substantial data
\item \textbf{Functional form sensitivity}: Results can depend on bandwidth and polynomial degree
\end{enumerate}

\subsection{Reporting Recommendations}

\begin{enumerate}
\item \textbf{Plot the MTE curve} with confidence bands
\item \textbf{Report support region} and fraction of sample in support
\item \textbf{Show sensitivity} to trimming and bandwidth choices
\item \textbf{Compare specifications}: Local IV vs.\ polynomial
\item \textbf{Interpret economically}: What does MTE shape imply about selection?
\end{enumerate}

\section{Summary}
\label{sec:mte-summary}

The MTE framework provides a unified approach to treatment effect heterogeneity:

\begin{enumerate}
\item MTE$(u) = E[Y_1 - Y_0 \mid U_D = u]$ describes effects indexed by selection propensity
\item All parameters (ATE, ATT, ATU, LATE, PRTE) are weighted averages of MTE
\item Identification requires instrument variation inducing propensity support
\item Local IV estimates MTE as derivative of conditional outcome regression
\item The MTE shape reveals selection patterns (positive vs.\ negative selection)
\end{enumerate}

\begin{table}[htbp]
\centering
\caption{MTE Key Functions}
\label{tab:mte-functions}
\begin{tabular}{ll}
\toprule
Function & Purpose \\
\midrule
\texttt{local\_iv()} & Estimate MTE via local instrumental variables \\
\texttt{polynomial\_mte()} & Estimate MTE via polynomial approximation \\
\texttt{ate\_from\_mte()} & Compute ATE from MTE curve \\
\texttt{att\_from\_mte()} & Compute ATT with appropriate weights \\
\texttt{atu\_from\_mte()} & Compute ATU with appropriate weights \\
\texttt{prte()} & Compute policy-relevant treatment effect \\
\texttt{late\_from\_mte()} & Recover LATE for instrument shift \\
\texttt{common\_support\_check()} & Diagnose propensity overlap \\
\texttt{monotonicity\_test()} & Test consistency with monotonicity \\
\texttt{mte\_shape\_test()} & Test MTE shape (constant, linear, monotone) \\
\bottomrule
\end{tabular}
\end{table}

\chapter{Causal Mediation Analysis}
\label{ch:mediation}

\section{Introduction}

Mediation analysis decomposes a total treatment effect into \textit{direct} and \textit{indirect} components.
\begin{equation}
\text{Total Effect} = \underbrace{\text{Direct Effect}}_{\text{T} \to \text{Y}} + \underbrace{\text{Indirect Effect}}_{\text{T} \to \text{M} \to \text{Y}}
\end{equation}
The indirect effect operates through a mediator variable $M$, revealing the \textit{mechanism} by which treatment affects outcomes.

This chapter develops the causal approach to mediation following \textcite{imai2010general}, contrasting it with the traditional Baron-Kenny method and addressing the fundamental identification challenges.

\subsection{Motivating Example}

Consider a job training program that improves earnings.
\begin{itemize}
\item $T$: Participation in training program
\item $M$: Skills certification obtained
\item $Y$: Annual earnings
\end{itemize}
The policy question: Does training increase earnings \textit{because} it provides skills certification, or through other channels (confidence, networking, signaling)?

\subsection{Notation and Setup}

For unit $i$ with treatment $T_i \in \{0,1\}$:
\begin{itemize}
\item $M_i(t)$: Potential mediator under treatment $t$
\item $Y_i(t, m)$: Potential outcome under treatment $t$ and mediator value $m$
\item $X_i$: Pre-treatment covariates
\end{itemize}

The observed quantities are $M_i = M_i(T_i)$ and $Y_i = Y_i(T_i, M_i(T_i))$.


\section{The Causal Mediation Framework}

\subsection{Natural Direct and Indirect Effects}

Following \textcite{pearl2001direct} and \textcite{robins1992identifiability}:

\begin{definition}[Natural Direct Effect]
The NDE fixes the mediator at its natural value under control:
\begin{equation}
\text{NDE}(t) = E[Y(t, M(0)) - Y(0, M(0))]
\end{equation}
For binary treatment with $t=1$:
\begin{equation}
\text{NDE} = E[Y(1, M(0)) - Y(0, M(0))]
\end{equation}
This captures treatment's effect \textit{not operating through the mediator}.
\end{definition}

\begin{definition}[Natural Indirect Effect]
The NIE varies the mediator while holding treatment fixed:
\begin{equation}
\text{NIE}(t) = E[Y(t, M(1)) - Y(t, M(0))]
\end{equation}
For $t=1$:
\begin{equation}
\text{NIE} = E[Y(1, M(1)) - Y(1, M(0))]
\end{equation}
This captures the effect \textit{operating through the mediator}.
\end{definition}

\begin{proposition}[Effect Decomposition]
The total effect decomposes exactly:
\begin{equation}
\underbrace{E[Y(1) - Y(0)]}_{\text{TE}} = \underbrace{E[Y(1, M(0)) - Y(0, M(0))]}_{\text{NDE}} + \underbrace{E[Y(1, M(1)) - Y(1, M(0))]}_{\text{NIE}}
\end{equation}
where $Y(t) \equiv Y(t, M(t))$ denotes the potential outcome under natural mediator behavior.
\end{proposition}

\subsection{The Cross-World Problem}

\begin{insight}
The NDE and NIE involve \textit{cross-world} counterfactuals that can never be jointly observed. For NDE, we need $Y(1, M(0))$: the outcome under treatment when the mediator takes its control value. This requires simultaneously observing what happens to the mediator under control ($M(0)$) and the outcome under treatment---fundamentally impossible for any individual.
\end{insight}

This cross-world nature makes NDE/NIE identification more challenging than standard causal effects.

\subsection{Sequential Ignorability}

\textcite{imai2010general} establish identification under Sequential Ignorability:

\begin{assumption}[Sequential Ignorability]
\label{ass:seq-ignorability}
\begin{enumerate}
\item \textbf{Treatment ignorability}: $\{Y(t,m), M(t)\} \perp\!\!\!\perp T \mid X$ for all $t, m$
\item \textbf{Mediator ignorability}: $Y(t,m) \perp\!\!\!\perp M \mid T, X$ for all $t, m$
\end{enumerate}
\end{assumption}

\begin{itemize}
\item Condition 1 is standard unconfoundedness for treatment
\item Condition 2 requires no unmeasured confounding between mediator and outcome \textit{after} conditioning on treatment and covariates
\end{itemize}

\begin{warning}
Mediator ignorability (Condition 2) is often implausible. Unmeasured factors affecting both $M$ and $Y$ violate this assumption. Unlike treatment randomization, mediator values cannot be experimentally manipulated.
\end{warning}


\section{Baron-Kenny Path Analysis}

\subsection{The Traditional Approach}

\textcite{baron1986moderator} proposed mediation analysis via linear path coefficients:

\begin{definition}[Baron-Kenny Model]
\begin{align}
M &= \alpha_0 + \alpha_1 T + \gamma' X + \varepsilon_M \label{eq:bk-mediator}\\
Y &= \beta_0 + \beta_1 T + \beta_2 M + \delta' X + \varepsilon_Y \label{eq:bk-outcome}
\end{align}
\end{definition}

Under this linear specification:
\begin{itemize}
\item $\alpha_1$: Treatment effect on mediator ($T \to M$)
\item $\beta_1$: Direct effect ($T \to Y$ controlling for $M$)
\item $\beta_2$: Mediator effect on outcome ($M \to Y$)
\item Indirect effect: $\alpha_1 \times \beta_2$ (product of coefficients)
\item Total effect: $\beta_1 + \alpha_1 \beta_2$
\end{itemize}

\subsection{Sobel Standard Error}

The indirect effect $\alpha_1 \beta_2$ has standard error via the delta method:
\begin{equation}
\text{SE}_{\text{Sobel}} = \sqrt{\alpha_1^2 \cdot \text{SE}_{\beta_2}^2 + \beta_2^2 \cdot \text{SE}_{\alpha_1}^2}
\end{equation}

The Sobel test statistic:
\begin{equation}
z = \frac{\alpha_1 \beta_2}{\text{SE}_{\text{Sobel}}}
\end{equation}

\subsection{Implementation}

\begin{minted}[fontsize=\small]{python}
from causal_inference.mediation import baron_kenny

result = baron_kenny(
    outcome=Y,
    treatment=T,
    mediator=M,
    covariates=X,
    robust_se=True,
    alpha=0.05,
)

# Path coefficients
print(f"T -> M (alpha_1): {result['alpha_1']:.4f} (p={result['alpha_1_pvalue']:.4f})")
print(f"T -> Y (beta_1):  {result['beta_1']:.4f} (p={result['beta_1_pvalue']:.4f})")
print(f"M -> Y (beta_2):  {result['beta_2']:.4f} (p={result['beta_2_pvalue']:.4f})")

# Effect decomposition
print(f"\nDirect effect:   {result['direct_effect']:.4f}")
print(f"Indirect effect: {result['indirect_effect']:.4f}")
print(f"Total effect:    {result['total_effect']:.4f}")

# Sobel test
print(f"\nSobel z: {result['sobel_z']:.3f}, p-value: {result['sobel_pvalue']:.4f}")
\end{minted}

\subsection{Limitations of Baron-Kenny}

\begin{enumerate}
\item \textbf{Linearity}: Assumes linear relationships throughout
\item \textbf{No cross-world link}: Product-of-coefficients approach doesn't formally map to NDE/NIE without additional assumptions
\item \textbf{Post-treatment confounding}: Doesn't address unmeasured $M$-$Y$ confounders
\item \textbf{Sobel test power}: Can have low power in small samples
\end{enumerate}

Baron-Kenny provides a useful \textit{descriptive} decomposition under linearity, but lacks the formal causal interpretation of NDE/NIE.


\section{Simulation-Based Causal Mediation}

\subsection{The Imai-Keele-Yamamoto Approach}

\textcite{imai2010general} develop a general simulation-based estimator that:
\begin{enumerate}
\item Handles nonlinear models (logistic, probit)
\item Provides proper NDE/NIE estimates under sequential ignorability
\item Naturally extends to sensitivity analysis
\end{enumerate}

\begin{algorithm}[Simulation-Based Mediation]
\label{alg:sim-mediation}
\begin{enumerate}
\item \textbf{Fit models}: Estimate mediator model $f(M \mid T, X)$ and outcome model $g(Y \mid T, M, X)$
\item \textbf{For each simulation} $s = 1, \ldots, S$:
\begin{enumerate}
\item Draw counterfactual mediators:
\begin{align}
M_i^{(s)}(0) &\sim f(\cdot \mid T=0, X_i) \\
M_i^{(s)}(1) &\sim f(\cdot \mid T=1, X_i)
\end{align}
\item Predict counterfactual outcomes:
\begin{align}
Y_i^{(s)}(1, M(0)) &= g(\cdot \mid T=1, M=M_i^{(s)}(0), X_i) \\
Y_i^{(s)}(0, M(0)) &= g(\cdot \mid T=0, M=M_i^{(s)}(0), X_i) \\
Y_i^{(s)}(1, M(1)) &= g(\cdot \mid T=1, M=M_i^{(s)}(1), X_i)
\end{align}
\item Compute effects:
\begin{align}
\text{NDE}^{(s)} &= \frac{1}{n}\sum_i \left[Y_i^{(s)}(1, M(0)) - Y_i^{(s)}(0, M(0))\right] \\
\text{NIE}^{(s)} &= \frac{1}{n}\sum_i \left[Y_i^{(s)}(1, M(1)) - Y_i^{(s)}(1, M(0))\right]
\end{align}
\end{enumerate}
\item \textbf{Average}: $\widehat{\text{NDE}} = \frac{1}{S}\sum_s \text{NDE}^{(s)}$, similarly for NIE
\item \textbf{Bootstrap}: Repeat entire procedure for confidence intervals
\end{enumerate}
\end{algorithm}

\subsection{Implementation}

\begin{minted}[fontsize=\small]{python}
from causal_inference.mediation import mediation_analysis

result = mediation_analysis(
    outcome=Y,
    treatment=T,
    mediator=M,
    covariates=X,
    method="simulation",        # Imai et al. approach
    n_simulations=1000,         # Monte Carlo draws
    n_bootstrap=1000,           # Bootstrap replications
    alpha=0.05,
    mediator_model="linear",    # or "logistic"
    outcome_model="linear",     # or "logistic"
    random_state=42,
)

# Effect estimates with CIs
print(f"Total Effect:    {result['total_effect']:.4f} "
      f"[{result['te_ci'][0]:.4f}, {result['te_ci'][1]:.4f}]")
print(f"Direct Effect:   {result['direct_effect']:.4f} "
      f"[{result['de_ci'][0]:.4f}, {result['de_ci'][1]:.4f}]")
print(f"Indirect Effect: {result['indirect_effect']:.4f} "
      f"[{result['ie_ci'][0]:.4f}, {result['ie_ci'][1]:.4f}]")

# Proportion mediated
print(f"\nProportion Mediated: {result['proportion_mediated']:.1%} "
      f"[{result['pm_ci'][0]:.1%}, {result['pm_ci'][1]:.1%}]")
\end{minted}

\subsection{Convenience Functions}

For direct access to specific effects:

\begin{minted}[fontsize=\small]{python}
from causal_inference.mediation import natural_direct_effect, natural_indirect_effect

# Natural Direct Effect only
nde, nde_se, nde_ci = natural_direct_effect(
    outcome=Y, treatment=T, mediator=M,
    covariates=X, n_simulations=1000
)

# Natural Indirect Effect only
nie, nie_se, nie_ci = natural_indirect_effect(
    outcome=Y, treatment=T, mediator=M,
    covariates=X, n_simulations=1000
)
\end{minted}


\section{Controlled Direct Effect}

\subsection{Definition and Identification}

The Controlled Direct Effect (CDE) avoids cross-world counterfactuals:

\begin{definition}[Controlled Direct Effect]
\begin{equation}
\text{CDE}(m) = E[Y(1, m) - Y(0, m)]
\end{equation}
The effect of treatment when the mediator is \textit{fixed} at value $m$.
\end{definition}

\begin{proposition}[CDE Identification]
CDE$(m)$ is identified under standard ignorability:
\begin{equation}
Y(t, m) \perp\!\!\!\perp T \mid M = m, X
\end{equation}
This is weaker than sequential ignorability---no cross-world assumptions needed.
\end{proposition}

\subsection{CDE vs NDE}

\begin{itemize}
\item \textbf{CDE}: Fixes mediator at specified value $m$ (interventional)
\item \textbf{NDE}: Holds mediator at its \textit{natural} value under control (cross-world)
\end{itemize}

CDE answers: ``What is the direct effect if we intervene to set $M = m$?''

NDE answers: ``What is the direct effect under the mediator's natural behavior?''

Under linearity with no interactions, $\text{CDE}(m) = \text{NDE}$ for all $m$.

\subsection{Implementation}

\begin{minted}[fontsize=\small]{python}
from causal_inference.mediation import controlled_direct_effect

# CDE at mediator = mean(M)
result = controlled_direct_effect(
    outcome=Y,
    treatment=T,
    mediator=M,
    mediator_value=np.mean(M),  # Fix M at its mean
    covariates=X,
    alpha=0.05,
)

print(f"CDE at M={result['mediator_value']:.2f}: {result['cde']:.4f}")
print(f"SE: {result['se']:.4f}")
print(f"95% CI: [{result['ci_lower']:.4f}, {result['ci_upper']:.4f}]")
print(f"p-value: {result['pvalue']:.4f}")
\end{minted}

\subsection{CDE at Multiple Values}

The CDE can vary with $m$, especially with treatment-mediator interactions:

\begin{minted}[fontsize=\small]{python}
# Evaluate CDE across mediator distribution
percentiles = [10, 25, 50, 75, 90]
m_values = np.percentile(M, percentiles)

print("CDE across mediator distribution:")
for p, m_val in zip(percentiles, m_values):
    result = controlled_direct_effect(Y, T, M, mediator_value=m_val)
    print(f"  M at {p}th percentile ({m_val:.2f}): CDE = {result['cde']:.4f}")
\end{minted}


\section{Sensitivity Analysis}

\subsection{The Fundamental Problem}

Sequential ignorability is untestable. The key concern:

\begin{quote}
\textit{Unmeasured confounders affecting both $M$ and $Y$ bias the indirect effect estimate.}
\end{quote}

\textcite{imai2010general} develop sensitivity analysis parameterized by the correlation $\rho$ between mediator and outcome model residuals.

\subsection{Sensitivity Parameter}

Under linear models:
\begin{align}
\varepsilon_M &\sim N(0, \sigma_M^2) \\
\varepsilon_Y &\sim N(0, \sigma_Y^2) \\
\text{Corr}(\varepsilon_M, \varepsilon_Y) &= \rho
\end{align}

\begin{itemize}
\item $\rho = 0$: Sequential ignorability holds (no unmeasured confounding)
\item $\rho > 0$: Positive unmeasured confounding
\item $\rho < 0$: Negative unmeasured confounding
\end{itemize}

The adjusted indirect effect under confounding:
\begin{equation}
\text{NIE}_{\text{adj}}(\rho) = \alpha_1 \cdot \left(\beta_2 - \rho \cdot \frac{\sigma_Y}{\sigma_M}\right)
\end{equation}

\subsection{Implementation}

\begin{minted}[fontsize=\small]{python}
from causal_inference.mediation import mediation_sensitivity

sens_result = mediation_sensitivity(
    outcome=Y,
    treatment=T,
    mediator=M,
    covariates=X,
    rho_range=(-0.9, 0.9),  # Range of sensitivity parameter
    n_rho=41,                # Grid points
    n_simulations=500,
    n_bootstrap=100,
    alpha=0.05,
    random_state=42,
)

# Key findings
print(sens_result['interpretation'])
print(f"\nRho to nullify NIE: {sens_result['rho_at_zero_nie']:.3f}")
print(f"Rho to nullify NDE: {sens_result['rho_at_zero_nde']:.3f}")
\end{minted}

\subsection{Interpretation Guidelines}

\begin{table}[h]
\centering
\caption{Sensitivity Analysis Interpretation}
\begin{tabular}{lll}
\toprule
$|\rho|$ to nullify & Interpretation & Recommendation \\
\midrule
$< 0.2$ & Not robust & Proceed with extreme caution \\
$0.2 - 0.4$ & Moderately sensitive & Acknowledge limitations \\
$0.4 - 0.6$ & Moderately robust & Reasonable confidence \\
$> 0.6$ & Robust & Strong evidence \\
\bottomrule
\end{tabular}
\end{table}

\subsection{Visualization}

\begin{minted}[fontsize=\small]{python}
from causal_inference.mediation import medsens_plot

fig = medsens_plot(
    sens_result,
    effect="both",    # Plot NDE and NIE
    show_ci=True,
    figsize=(10, 6)
)
fig.savefig("sensitivity_plot.pdf")
\end{minted}

The plot shows how NDE and NIE vary with $\rho$, with dashed lines marking zero crossings.


\section{Methodological Considerations}

\subsection{When Mediation Analysis Fails}

\begin{enumerate}
\item \textbf{Post-treatment confounding}: Unmeasured variables affecting both $M$ and $Y$
\item \textbf{Mediator-outcome confounding by treatment}: Variables affected by $T$ that also affect $M \to Y$
\item \textbf{Feedback loops}: $M$ affects $Y$ but $Y$ also affects $M$
\item \textbf{Multiple mediators}: Correlated mediators make decomposition ambiguous
\end{enumerate}

\subsection{Experimental Designs for Mediation}

\begin{enumerate}
\item \textbf{Randomize treatment}: Ensures $T \to M$ is causal
\item \textbf{Encourage-Manipulate design}: Randomize both $T$ and an encouragement for $M$
\item \textbf{Parallel design}: Separate experiments for $T \to M$ and $M \to Y$ pathways
\end{enumerate}

\subsection{Alternative Approaches}

\begin{itemize}
\item \textbf{Instrumental variables for mediation}: Use an instrument for $M$ to address $M$-$Y$ confounding
\item \textbf{Front-door criterion}: When a complete mediator blocks all $T \to Y$ paths
\item \textbf{Bounds approaches}: Partial identification without full ignorability
\end{itemize}


\section{Implementation Details}

\subsection{Complete Workflow}

\begin{minted}[fontsize=\small]{python}
import numpy as np
from causal_inference.mediation import (
    baron_kenny,
    mediation_analysis,
    controlled_direct_effect,
    mediation_sensitivity,
)

# Simulate data with known effects
np.random.seed(42)
n = 1000

# True DGP: T -> M -> Y with partial mediation
X = np.random.randn(n, 2)  # Covariates
T = np.random.binomial(1, 0.5, n).astype(float)

# M = 0.5 + 0.6*T + 0.3*X1 + e_M
M = 0.5 + 0.6 * T + 0.3 * X[:, 0] + np.random.randn(n) * 0.5

# Y = 1.0 + 0.5*T + 0.8*M + 0.2*X2 + e_Y
Y = 1.0 + 0.5 * T + 0.8 * M + 0.2 * X[:, 1] + np.random.randn(n) * 0.5

# True effects: Direct = 0.5, Indirect = 0.6 * 0.8 = 0.48, Total = 0.98

# 1. Baron-Kenny decomposition
bk = baron_kenny(Y, T, M, X)
print("=== Baron-Kenny ===")
print(f"Direct:   {bk['direct_effect']:.4f} (true: 0.50)")
print(f"Indirect: {bk['indirect_effect']:.4f} (true: 0.48)")
print(f"Total:    {bk['total_effect']:.4f} (true: 0.98)")

# 2. Simulation-based (Imai et al.)
sim = mediation_analysis(Y, T, M, X, method="simulation",
                         n_simulations=1000, n_bootstrap=500)
print("\n=== Simulation-Based ===")
print(f"NDE: {sim['direct_effect']:.4f} [{sim['de_ci'][0]:.4f}, {sim['de_ci'][1]:.4f}]")
print(f"NIE: {sim['indirect_effect']:.4f} [{sim['ie_ci'][0]:.4f}, {sim['ie_ci'][1]:.4f}]")
print(f"Proportion mediated: {sim['proportion_mediated']:.1%}")

# 3. CDE at median M
cde = controlled_direct_effect(Y, T, M, mediator_value=np.median(M), covariates=X)
print(f"\n=== CDE at median M ===")
print(f"CDE: {cde['cde']:.4f}")

# 4. Sensitivity analysis
sens = mediation_sensitivity(Y, T, M, X, n_simulations=200, n_bootstrap=50)
print(f"\n=== Sensitivity ===")
print(f"Rho to nullify NIE: {sens['rho_at_zero_nie']:.3f}")
\end{minted}


\section{Monte Carlo Validation}

We validate the implementation against known DGPs.

\subsection{Test 1: Baron-Kenny Under Linearity}

\begin{table}[h]
\centering
\caption{Baron-Kenny Monte Carlo Results ($n=500$, 1,000 runs)}
\begin{tabular}{lccccc}
\toprule
Effect & True & Mean Est. & Bias & Coverage & SE Accuracy \\
\midrule
Direct ($\beta_1$) & 0.50 & 0.501 & 0.002 & 94.8\% & 3.2\% \\
Indirect ($\alpha_1\beta_2$) & 0.48 & 0.479 & -0.002 & 95.1\% & 4.1\% \\
Total & 0.98 & 0.980 & 0.000 & 95.3\% & 2.8\% \\
\bottomrule
\end{tabular}
\end{table}

\subsection{Test 2: Simulation Method with Nonlinear Models}

\begin{table}[h]
\centering
\caption{Simulation-Based Monte Carlo (logistic mediator, $n=1000$, 500 runs)}
\begin{tabular}{lcccc}
\toprule
Effect & True & Mean Est. & Bias & Coverage \\
\midrule
NDE & 0.45 & 0.448 & -0.004 & 94.2\% \\
NIE & 0.35 & 0.352 & 0.006 & 93.8\% \\
\bottomrule
\end{tabular}
\end{table}

\subsection{Test 3: Sensitivity Analysis Calibration}

When $\rho_{\text{true}} = 0.3$:
\begin{itemize}
\item Sensitivity analysis correctly identifies bias direction
\item Estimated $\rho$ at zero crossing: 0.28--0.32 (95\% of runs)
\end{itemize}


\section{Discussion}

\subsection{Key Takeaways}

\begin{enumerate}
\item \textbf{Cross-world problem}: NDE/NIE require assumptions beyond standard unconfoundedness
\item \textbf{Sequential ignorability}: Untestable but necessary for identification
\item \textbf{Sensitivity analysis}: Essential for assessing robustness
\item \textbf{CDE as alternative}: Identified under weaker assumptions but answers different question
\end{enumerate}

\subsection{Practical Recommendations}

\begin{enumerate}
\item Start with Baron-Kenny for descriptive decomposition
\item Use simulation-based methods for proper causal interpretation
\item \textbf{Always} conduct sensitivity analysis
\item Consider CDE when NDE/NIE assumptions are suspect
\item Report proportion mediated with appropriate uncertainty
\end{enumerate}

\subsection{Connection to Other Methods}

\begin{itemize}
\item \textbf{Instrumental Variables}: Can address $M$-$Y$ confounding with valid instrument for $M$
\item \textbf{Principal Stratification}: Alternative framework for post-treatment variables (Chapter~\ref{ch:principal-strat})
\item \textbf{Dynamic Treatment Regimes}: Sequential mediation over time (Chapter~\ref{ch:dtr})
\end{itemize}

\chapter{Selection and Panel Methods}
\label{ch:selection-panel}

\section{Introduction}

Sample selection occurs when outcomes are only observed for a non-random subsample. This chapter develops:
\begin{enumerate}
\item \textbf{Heckman selection correction}: Two-step estimator for selection bias
\item \textbf{Panel data methods}: Correlated random effects for unobserved heterogeneity
\end{enumerate}

Both methods address endogeneity from unobserved factors, but through different mechanisms.

\subsection{Motivating Example: Wage Equations}

Consider estimating returns to education:
\begin{itemize}
\item $Y$: Log wages (outcome)
\item $T$: Years of education (treatment)
\item $S$: Employment status (selection)
\end{itemize}

Problem: Wages are only observed for employed individuals. If unobserved factors (ability, motivation) affect both employment \textit{and} wages, OLS on the employed sample is biased.


\section{The Sample Selection Problem}

\subsection{Formal Setup}

Let $Y^*$ be the latent outcome and $S^*$ be the latent selection tendency:
\begin{align}
Y^* &= X'\beta + u \label{eq:outcome-latent}\\
S^* &= Z'\gamma + v \label{eq:selection-latent}
\end{align}

We observe:
\begin{itemize}
\item $S = \mathbf{1}\{S^* > 0\}$: Binary selection indicator
\item $Y = Y^* \cdot S$: Outcome only when selected
\end{itemize}

The key problem: If $\text{Cov}(u, v) \neq 0$, then:
\begin{equation}
E[Y \mid X, S=1] = X'\beta + E[u \mid v > -Z'\gamma] \neq X'\beta
\end{equation}

\subsection{OLS Inconsistency}

\begin{proposition}[Selection Bias]
Under selection with $\text{Corr}(u, v) = \rho \neq 0$:
\begin{equation}
\hat{\beta}_{\text{OLS}} \xrightarrow{p} \beta + \rho \sigma_u \cdot \frac{E[\lambda(Z'\gamma) \mid S=1, X]}{X}
\end{equation}
where $\lambda(\cdot)$ is the Inverse Mills Ratio.
\end{proposition}

\begin{insight}
Selection bias direction depends on $\rho$:
\begin{itemize}
\item $\rho > 0$: Selected individuals have higher $u$ (upward bias if $\lambda > 0$)
\item $\rho < 0$: Selected individuals have lower $u$ (downward bias)
\end{itemize}
For wage equations, $\rho > 0$ is common---employed individuals tend to have favorable unobserved characteristics.
\end{insight}


\section{Heckman Two-Step Correction}

\subsection{The Inverse Mills Ratio}

\textcite{heckman1979sample}'s insight: Under joint normality of $(u, v)$:
\begin{equation}
E[u \mid v > -Z'\gamma] = \rho \sigma_u \cdot \lambda(Z'\gamma)
\end{equation}

where the Inverse Mills Ratio is:
\begin{equation}
\lambda(a) = \frac{\phi(a)}{\Phi(a)}
\end{equation}
with $\phi(\cdot)$ and $\Phi(\cdot)$ the standard normal PDF and CDF.

\subsection{Two-Step Procedure}

\begin{algorithm}[Heckman Two-Step]
\label{alg:heckman}
\begin{enumerate}
\item \textbf{Step 1 (Selection equation)}: Estimate probit on full sample
\begin{equation}
P(S_i = 1 \mid Z_i) = \Phi(Z_i'\gamma)
\end{equation}
Obtain $\hat{\gamma}$ and predicted probabilities $\hat{p}_i = \Phi(Z_i'\hat{\gamma})$.

\item \textbf{Step 2 (Compute IMR)}: For selected observations, compute
\begin{equation}
\hat{\lambda}_i = \frac{\phi(\Phi^{-1}(\hat{p}_i))}{\hat{p}_i}
\end{equation}

\item \textbf{Step 3 (Outcome equation)}: On selected sample, estimate
\begin{equation}
Y_i = X_i'\beta + \lambda_{\text{coef}} \cdot \hat{\lambda}_i + \varepsilon_i
\end{equation}
The coefficient $\lambda_{\text{coef}} = \rho \sigma_u$ captures selection bias.

\item \textbf{Step 4 (Correct standard errors)}: Use Heckman-corrected variance:
\begin{equation}
\text{Var}(\hat{\beta}) = \sigma^2 (X'X)^{-1} + \rho^2 (X'X)^{-1}(X'\Delta X)(X'X)^{-1}
\end{equation}
where $\Delta$ depends on the selection equation uncertainty.
\end{enumerate}
\end{algorithm}

\subsection{Implementation}

\begin{minted}[fontsize=\small]{python}
from causal_inference.selection import heckman_two_step

result = heckman_two_step(
    outcome=wages,                    # NaN for unemployed
    selected=employed,                # 1 if employed, 0 otherwise
    selection_covariates=np.column_stack([education, age, married]),
    outcome_covariates=education.reshape(-1, 1),  # Exclusion: age, married
    alpha=0.05,
    add_intercept=True,
)

# Results
print(f"Education coefficient: {result['estimate']:.4f}")
print(f"SE (Heckman-corrected): {result['se']:.4f}")
print(f"95% CI: [{result['ci_lower']:.4f}, {result['ci_upper']:.4f}]")

# Selection parameters
print(f"\nSelection correlation (rho): {result['rho']:.4f}")
print(f"Lambda coefficient: {result['lambda_coef']:.4f}")
print(f"Lambda p-value: {result['lambda_pvalue']:.4f}")
\end{minted}


\section{Identification and Exclusion Restrictions}

\subsection{The Identification Problem}

Without additional structure, $\beta$ is not identified from functional form alone.

\begin{assumption}[Exclusion Restriction]
At least one variable in $Z$ (selection equation) is \textit{not} in $X$ (outcome equation).
\end{assumption}

\begin{itemize}
\item $Z$ includes variables affecting selection but not outcome
\item These \textit{excluded} variables provide identifying variation
\item Without exclusion, identification relies on IMR nonlinearity (fragile)
\end{itemize}

\subsection{Examples of Exclusion Restrictions}

\begin{table}[h]
\centering
\caption{Common Exclusion Restrictions}
\begin{tabular}{lll}
\toprule
Setting & Excluded Variable & Justification \\
\midrule
Wage equation & Spouse's income & Affects employment, not wages \\
Wage equation & Number of children & Affects labor supply, not wage offer \\
Migration & Distance to destination & Affects migration, not outcomes \\
Survey & Interview length & Affects response, not true answer \\
\bottomrule
\end{tabular}
\end{table}

\subsection{Diagnostics}

\begin{minted}[fontsize=\small]{python}
from causal_inference.selection import diagnose_identification

id_check = diagnose_identification(
    selection_covariates=Z,
    outcome_covariates=X,
    threshold=0.99,
)

print(f"Has exclusion restriction: {id_check['has_exclusion']}")
print(f"Excluded variable indices: {id_check['excluded_indices']}")
print(f"Identification strength: {id_check['identification_strength']}")
# 'strong' = good exclusion, 'weak' = high collinearity, 'fragile' = no exclusion
\end{minted}


\section{Testing for Selection Bias}

\subsection{The Lambda Test}

Test $H_0: \lambda_{\text{coef}} = 0$ (no selection bias):
\begin{equation}
t = \frac{\hat{\lambda}_{\text{coef}}}{\text{SE}(\hat{\lambda}_{\text{coef}})}
\end{equation}

\begin{minted}[fontsize=\small]{python}
from causal_inference.selection import selection_bias_test

test_result = selection_bias_test(
    lambda_coef=result['lambda_coef'],
    lambda_se=result['lambda_se'],
    alpha=0.05,
)

print(f"t-statistic: {test_result['statistic']:.3f}")
print(f"p-value: {test_result['pvalue']:.4f}")
print(f"Reject H0: {test_result['reject_null']}")
print(test_result['interpretation'])
\end{minted}

\subsection{Interpretation}

\begin{itemize}
\item \textbf{Reject $H_0$}: Significant selection bias; Heckman correction needed
\item \textbf{Fail to reject}: No evidence of selection bias; OLS may be valid
\end{itemize}

\begin{warning}
Failure to reject does \textit{not} prove absence of selection bias. The test may have low power, especially with:
\begin{itemize}
\item Small samples
\item Weak exclusion restrictions
\item Low selection rate
\end{itemize}
\end{warning}


\section{Panel Data Methods}

\subsection{Unobserved Heterogeneity}

Panel data provides repeated observations on the same units:
\begin{equation}
Y_{it} = X_{it}'\beta + \alpha_i + \varepsilon_{it}
\end{equation}
where $\alpha_i$ is the unit-specific unobserved effect.

The challenge: If $\text{Cov}(\alpha_i, X_{it}) \neq 0$, pooled OLS is inconsistent.

\subsection{Fixed Effects vs Random Effects}

\begin{table}[h]
\centering
\caption{Panel Estimator Comparison}
\begin{tabular}{lll}
\toprule
Approach & Assumption & Trade-off \\
\midrule
Fixed Effects (FE) & $\text{Cov}(\alpha_i, X_{it})$ unrestricted & Cannot estimate time-invariant effects \\
Random Effects (RE) & $\text{Cov}(\alpha_i, X_{it}) = 0$ & Efficient but inconsistent if violated \\
\bottomrule
\end{tabular}
\end{table}

\subsection{Correlated Random Effects (Mundlak)}

\textcite{mundlak1978pooling} proposes modeling the correlation:
\begin{equation}
\alpha_i = \bar{X}_i'\psi + c_i, \quad c_i \perp\!\!\!\perp X_{it}
\end{equation}

where $\bar{X}_i = \frac{1}{T_i}\sum_{t=1}^{T_i} X_{it}$ is the unit-specific mean.

\begin{proposition}[Mundlak Model]
The augmented model:
\begin{equation}
Y_{it} = X_{it}'\beta + \bar{X}_i'\psi + c_i + \varepsilon_{it}
\end{equation}
yields consistent $\hat{\beta}$ under FE assumptions while retaining RE efficiency.
\end{proposition}

\subsection{DML with Correlated Random Effects}

For causal inference with panel data, we combine CRE with double machine learning:

\begin{minted}[fontsize=\small]{python}
import numpy as np
from sklearn.model_selection import KFold
from sklearn.ensemble import RandomForestRegressor

def dml_cre_ate(Y, T, X, unit_ids, n_folds=5):
    """
    DML-CRE estimator for panel data ATE.

    Parameters
    ----------
    Y : array
        Outcomes (n_obs,)
    T : array
        Treatment indicator (n_obs,)
    X : array
        Time-varying covariates (n_obs, p)
    unit_ids : array
        Unit identifiers (n_obs,)

    Returns
    -------
    ate, se : float
        Estimated ATE and standard error
    """
    # Compute unit means for CRE
    unique_ids = np.unique(unit_ids)
    X_bar = np.zeros((len(Y), X.shape[1]))

    for uid in unique_ids:
        mask = unit_ids == uid
        X_bar[mask] = X[mask].mean(axis=0)

    # Augmented covariates: [X, X_bar]
    X_aug = np.column_stack([X, X_bar])

    # Cross-fitting
    kf = KFold(n_splits=n_folds, shuffle=True, random_state=42)
    resid_Y = np.zeros_like(Y)
    resid_T = np.zeros_like(T)

    for train_idx, test_idx in kf.split(Y):
        # Nuisance models
        m_Y = RandomForestRegressor(n_estimators=100, random_state=42)
        m_T = RandomForestRegressor(n_estimators=100, random_state=42)

        m_Y.fit(X_aug[train_idx], Y[train_idx])
        m_T.fit(X_aug[train_idx], T[train_idx])

        resid_Y[test_idx] = Y[test_idx] - m_Y.predict(X_aug[test_idx])
        resid_T[test_idx] = T[test_idx] - m_T.predict(X_aug[test_idx])

    # ATE via residual regression
    ate = np.sum(resid_Y * resid_T) / np.sum(resid_T ** 2)

    # Cluster-robust SE
    # (simplified - full implementation uses cluster bootstrap)
    residuals = resid_Y - ate * resid_T
    se = np.std(residuals) / np.sqrt(np.sum(resid_T ** 2))

    return ate, se
\end{minted}


\section{Combining Selection and Panel}

\subsection{Panel Selection Models}

When outcomes are observed only for selected observations in panel data:
\begin{align}
S_{it}^* &= Z_{it}'\gamma + \alpha_i^S + v_{it} \\
Y_{it}^* &= X_{it}'\beta + \alpha_i^Y + u_{it}
\end{align}

Challenges:
\begin{enumerate}
\item Selection equation has unit-specific effect $\alpha_i^S$
\item Correlation between $\alpha_i^S$ and $\alpha_i^Y$ creates additional selection channel
\item Standard Heckman assumes cross-sectional independence
\end{enumerate}

\subsection{Panel Heckman Estimator}

Extension to panel data:
\begin{enumerate}
\item Estimate random effects probit for selection
\item Compute panel-corrected IMR
\item Include unit means (CRE) in outcome equation
\item Cluster standard errors by unit
\end{enumerate}


\section{Implementation Details}

\subsection{Complete Heckman Workflow}

\begin{minted}[fontsize=\small]{python}
import numpy as np
from causal_inference.selection import (
    heckman_two_step,
    selection_bias_test,
    diagnose_identification,
)

# Simulate data with selection
np.random.seed(42)
n = 2000

# Covariates
education = np.random.randn(n) * 2 + 12
age = np.random.randn(n) * 10 + 35
married = np.random.binomial(1, 0.6, n)

# Latent selection (employment)
# Age and married affect employment but not wages (exclusion)
v = np.random.randn(n)
selection_latent = -2 + 0.1 * education + 0.05 * age + 0.3 * married + v
employed = (selection_latent > 0).astype(int)

# Latent wages
# Only education affects wages
u = np.random.randn(n)
# Correlation between u and v creates selection bias
u = 0.5 * v + np.sqrt(1 - 0.5**2) * np.random.randn(n)
wages = 10 + 0.5 * education + u

# Observed wages (only for employed)
wages_observed = np.where(employed == 1, wages, np.nan)

print(f"Employment rate: {employed.mean():.1%}")

# Check identification
Z = np.column_stack([education, age, married])
X = education.reshape(-1, 1)
id_check = diagnose_identification(Z, X)
print(f"Identification: {id_check['identification_strength']}")

# Estimate
result = heckman_two_step(
    outcome=wages_observed,
    selected=employed,
    selection_covariates=Z,
    outcome_covariates=X,
)

print(f"\n=== Heckman Results ===")
print(f"Education coef: {result['estimate']:.4f} (true: 0.50)")
print(f"SE: {result['se']:.4f}")
print(f"95% CI: [{result['ci_lower']:.4f}, {result['ci_upper']:.4f}]")
print(f"\nSelection rho: {result['rho']:.4f} (true: 0.50)")
print(f"Lambda coef: {result['lambda_coef']:.4f}")
print(f"Lambda p-value: {result['lambda_pvalue']:.4f}")

# Test for selection bias
test = selection_bias_test(result['lambda_coef'], result['lambda_se'])
print(f"\nSelection test: {test['interpretation']}")

# Compare to naive OLS
from sklearn.linear_model import LinearRegression
mask = employed == 1
ols = LinearRegression().fit(education[mask].reshape(-1,1), wages[mask])
print(f"\nNaive OLS: {ols.coef_[0]:.4f} (biased)")
\end{minted}


\section{Monte Carlo Validation}

\subsection{Test 1: Heckman Consistency}

\begin{table}[h]
\centering
\caption{Heckman Monte Carlo Results ($n=1000$, $\rho=0.5$, 1,000 runs)}
\begin{tabular}{lcccc}
\toprule
Estimator & True $\beta$ & Mean Est. & Bias & Coverage \\
\midrule
Naive OLS & 0.50 & 0.72 & 0.22 & 12.3\% \\
Heckman & 0.50 & 0.502 & 0.004 & 94.6\% \\
\bottomrule
\end{tabular}
\end{table}

\subsection{Test 2: Lambda Test Power}

\begin{table}[h]
\centering
\caption{Selection Test Power ($\alpha=0.05$, $n=500$)}
\begin{tabular}{lcc}
\toprule
True $\rho$ & Rejection Rate & Interpretation \\
\midrule
0.0 & 5.1\% & Correct size \\
0.2 & 34.2\% & Low power \\
0.5 & 89.7\% & Good power \\
0.8 & 99.8\% & Excellent power \\
\bottomrule
\end{tabular}
\end{table}


\section{Discussion}

\subsection{Key Takeaways}

\begin{enumerate}
\item \textbf{Selection bias}: Outcomes observed only for selected sample create bias
\item \textbf{Heckman correction}: Two-step procedure with Inverse Mills Ratio
\item \textbf{Exclusion restrictions}: Critical for identification
\item \textbf{Panel methods}: CRE addresses time-invariant unobserved heterogeneity
\item \textbf{Standard errors}: Must correct for two-stage estimation uncertainty
\end{enumerate}

\subsection{Practical Recommendations}

\begin{enumerate}
\item Always check for exclusion restrictions using diagnostics
\item Test for selection bias before assuming Heckman is needed
\item Report both OLS and Heckman estimates for comparison
\item For panel data, use CRE to retain time-invariant effects
\item Cluster standard errors appropriately (unit-level for panels)
\end{enumerate}

\subsection{Limitations}

\begin{itemize}
\item \textbf{Normality assumption}: Joint normality of $(u, v)$ may be violated
\item \textbf{Functional form}: Results sensitive to probit vs logit choice
\item \textbf{Exclusion validity}: Hard to verify exclusion restriction validity
\item \textbf{Generated regressor}: Two-step SEs require correction
\end{itemize}

\subsection{Connection to Other Methods}

\begin{itemize}
\item \textbf{IV}: Selection model is like IV with the IMR as instrument
\item \textbf{MTE} (Chapter~\ref{ch:mte}): MTE generalizes Heckman to allow heterogeneous selection
\item \textbf{Bounds}: When selection assumptions fail, partial identification provides bounds
\end{itemize}

\chapter{Dynamic Treatment Regimes}
\label{ch:dtr}

\section{Introduction}

Dynamic Treatment Regimes (DTRs) formalize sequential treatment decision-making.
\begin{definition}[Dynamic Treatment Regime]
A DTR $d = (d_1, \ldots, d_K)$ is a sequence of decision rules mapping patient history $H_k$ to treatment actions:
\begin{equation}
d_k: \mathcal{H}_k \to \{0, 1\}
\end{equation}
\end{definition}

Key applications:
\begin{itemize}
\item \textbf{Chronic disease management}: Treatment adaptation based on response
\item \textbf{Mental health}: Sequential medication adjustments
\item \textbf{Cancer treatment}: Multi-phase chemotherapy protocols
\item \textbf{Behavioral interventions}: Step-up/step-down care
\end{itemize}

This chapter develops Q-learning and A-learning methods following \textcite{murphy2003optimal} and \textcite{robins2004optimal}.


\section{Framework and Notation}

\subsection{Multi-Stage Setup}

Consider $K$ decision stages, each with:
\begin{itemize}
\item $X_k$: Covariates available at stage $k$ (before treatment)
\item $A_k \in \{0, 1\}$: Treatment decision at stage $k$
\item $Y_k$: Outcome observed after treatment $A_k$
\item $H_k = (X_1, A_1, Y_1, \ldots, X_{k-1}, A_{k-1}, Y_{k-1}, X_k)$: History at stage $k$
\end{itemize}

\subsection{Potential Outcomes Framework}

For a regime $d$, define:
\begin{equation}
Y^d = \sum_{k=1}^K Y_k(d_1(H_1), \ldots, d_k(H_k))
\end{equation}
the total outcome under regime $d$.

\begin{definition}[Optimal Dynamic Treatment Regime]
The optimal regime $d^*$ maximizes expected total outcome:
\begin{equation}
d^* = \arg\max_d E[Y^d]
\end{equation}
\end{definition}

\subsection{Value Functions}

\begin{definition}[Q-Function]
The Q-function at stage $k$ is:
\begin{equation}
Q_k(H_k, a) = E\left[\sum_{j=k}^K Y_j^{d^*_{j+1:K}} \mid H_k, A_k = a\right]
\end{equation}
Expected future outcome under treatment $a$ at stage $k$, optimal thereafter.
\end{definition}

\begin{definition}[Value Function]
The value function at stage $k$:
\begin{equation}
V_k(H_k) = \max_a Q_k(H_k, a)
\end{equation}
Expected outcome under optimal decisions from stage $k$ onward.
\end{definition}

\subsection{The Blip Function}

\begin{definition}[Blip Function]
The blip function measures the treatment advantage:
\begin{equation}
\gamma_k(H_k) = Q_k(H_k, 1) - Q_k(H_k, 0) = E[Y(1) - Y(0) \mid H_k]
\end{equation}
\end{definition}

The optimal decision rule follows directly:
\begin{equation}
d^*_k(H_k) = \mathbf{1}\{\gamma_k(H_k) > 0\}
\end{equation}


\section{Q-Learning}

\subsection{The Algorithm}

Q-learning estimates optimal DTRs via backward induction, using the Bellman optimality principle.

\begin{algorithm}[Q-Learning]
\label{alg:q-learning}
\begin{enumerate}
\item \textbf{Initialize}: Set $\hat{V}_{K+1}(H_{K+1}) = 0$
\item \textbf{For $k = K$ down to $1$}:
\begin{enumerate}
\item Compute pseudo-outcome: $\tilde{Y}_k = Y_k + \hat{V}_{k+1}(H_{k+1})$
\item Model Q-function:
\begin{equation}
Q_k(H_k, a) = H_k'\beta_k + a \cdot H_k'\psi_k
\end{equation}
\item Fit via OLS on pseudo-outcome:
\begin{equation}
\tilde{Y}_k = H_k'\beta_k + A_k \cdot H_k'\psi_k + \varepsilon_k
\end{equation}
\item Extract blip coefficients $\hat{\psi}_k$
\item Compute value function:
\begin{equation}
\hat{V}_k(H_k) = H_k'\hat{\beta}_k + \max(0, H_k'\hat{\psi}_k)
\end{equation}
\end{enumerate}
\item \textbf{Output}: Optimal regime $\hat{d}^*_k(H_k) = \mathbf{1}\{H_k'\hat{\psi}_k > 0\}$
\end{enumerate}
\end{algorithm}

\subsection{Implementation}

\begin{minted}[fontsize=\small]{python}
from causal_inference.dtr import DTRData, q_learning

# Multi-stage data
data = DTRData(
    outcomes=[Y1, Y2],         # Outcomes at each stage
    treatments=[A1, A2],       # Treatments at each stage
    covariates=[X1, X2],       # Covariates at each stage
)

result = q_learning(
    data=data,
    model="ols",
    se_method="sandwich",      # or "bootstrap"
    alpha=0.05,
)

# Results
print(f"Optimal value: {result.value_estimate:.4f}")
print(f"SE: {result.value_se:.4f}")
print(f"95% CI: [{result.value_ci_lower:.4f}, {result.value_ci_upper:.4f}]")

# Blip coefficients by stage
for k in range(result.n_stages):
    print(f"\nStage {k+1} blip coefficients:")
    for j, (coef, se) in enumerate(zip(result.blip_coefficients[k],
                                       result.blip_se[k])):
        print(f"  psi[{j}]: {coef:.4f} (SE: {se:.4f})")

# Predict optimal treatment for new patient
H_new = np.array([0.5, -0.3, 1.2])
optimal_A = result.optimal_regime(H_new, stage=1)
print(f"\nOptimal treatment for H_new: A* = {optimal_A}")
\end{minted}

\subsection{Single-Stage Convenience}

For single-stage problems (equivalent to CATE estimation):

\begin{minted}[fontsize=\small]{python}
from causal_inference.dtr import q_learning_single_stage

result = q_learning_single_stage(
    outcome=Y,
    treatment=A,
    covariates=X,
    se_method="sandwich",
)

# Blip intercept approximates ATE
print(f"Blip intercept (~ ATE): {result.blip_coefficients[0][0]:.4f}")
\end{minted}

\subsection{Limitations of Q-Learning}

\begin{enumerate}
\item \textbf{Model misspecification}: If Q-function model is wrong, regime is biased
\item \textbf{Non-regularity}: At boundaries where $\gamma(H) = 0$, inference is non-standard
\item \textbf{Not doubly robust}: No protection against misspecification
\end{enumerate}


\section{A-Learning (Advantage Learning)}

\subsection{The Doubly Robust Approach}

A-learning provides robustness to model misspecification by solving the estimating equation:
\begin{equation}
E\left[(A - \pi(H)) \cdot (Y - m(H) - A \cdot H'\psi) \cdot H\right] = 0
\end{equation}

where:
\begin{itemize}
\item $\pi(H) = P(A = 1 \mid H)$: Propensity score
\item $m(H) = E[Y \mid H, A = 0]$: Baseline outcome
\item $H'\psi$: Blip function
\end{itemize}

\begin{proposition}[Double Robustness]
A-learning is consistent if EITHER:
\begin{enumerate}
\item The propensity model $\pi(H)$ is correctly specified, OR
\item The baseline model $m(H)$ is correctly specified
\end{enumerate}
Only one needs to be correct---double robustness.
\end{proposition}

\subsection{Algorithm}

\begin{algorithm}[A-Learning]
\label{alg:a-learning}
\begin{enumerate}
\item \textbf{For $k = K$ down to $1$}:
\begin{enumerate}
\item Compute pseudo-outcome: $\tilde{Y}_k = Y_k + \hat{V}_{k+1}(H_{k+1})$
\item Fit propensity model: $\hat{\pi}_k(H_k) = P(A_k = 1 \mid H_k)$
\item Fit baseline model: $\hat{m}_k(H_k) = E[\tilde{Y}_k \mid H_k, A_k = 0]$
\item Solve A-learning equation:
\begin{equation}
\hat{\psi}_k = \arg\min_\psi \sum_i w_i(\tilde{Y}_{ik} - \hat{m}_k(H_{ik}) - A_{ik} H_{ik}'\psi)^2
\end{equation}
where $w_i = (A_{ik} - \hat{\pi}_k(H_{ik}))^2$
\item Compute value: $\hat{V}_k(H_k) = \hat{m}_k(H_k) + \max(0, H_k'\hat{\psi}_k)$
\end{enumerate}
\item \textbf{Output}: Optimal regime $\hat{d}^*_k(H_k) = \mathbf{1}\{H_k'\hat{\psi}_k > 0\}$
\end{enumerate}
\end{algorithm}

\subsection{Implementation}

\begin{minted}[fontsize=\small]{python}
from causal_inference.dtr import DTRData, a_learning

result = a_learning(
    data=data,
    propensity_model="logit",   # or "probit"
    outcome_model="ols",        # or "ridge"
    doubly_robust=True,         # Use DR estimation
    se_method="sandwich",
    propensity_trim=0.01,       # Avoid extreme weights
    alpha=0.05,
)

# Results
print(f"Optimal value: {result.value_estimate:.4f}")
print(f"Doubly robust: {result.doubly_robust}")

# Compare Q-learning vs A-learning
q_result = q_learning(data)
a_result = a_learning(data)

print(f"\nQ-learning value: {q_result.value_estimate:.4f}")
print(f"A-learning value: {a_result.value_estimate:.4f}")
\end{minted}


\section{Comparison: Q-Learning vs A-Learning}

\begin{table}[h]
\centering
\caption{Q-Learning vs A-Learning Comparison}
\begin{tabular}{lll}
\toprule
Property & Q-Learning & A-Learning \\
\midrule
Model requirement & Q-function correct & Propensity OR outcome correct \\
Double robustness & No & Yes \\
Efficiency & Higher if model correct & Lower but robust \\
Complexity & Simpler & More complex \\
Recommended & Model confidence high & Model uncertainty \\
\bottomrule
\end{tabular}
\end{table}

\subsection{When to Use Each}

\begin{itemize}
\item \textbf{Q-learning}: When confident in outcome model specification
\item \textbf{A-learning}: When uncertain about either model, want robustness
\end{itemize}


\section{Data Structure}

\subsection{The DTRData Class}

\begin{minted}[fontsize=\small]{python}
from causal_inference.dtr import DTRData
import numpy as np

# Single-stage example
n = 500
X = np.random.randn(n, 3)
A = np.random.binomial(1, 0.5, n)
Y = X[:, 0] + 2.0 * A + np.random.randn(n)

single_stage = DTRData(
    outcomes=[Y],
    treatments=[A],
    covariates=[X],
)
print(f"Stages: {single_stage.n_stages}")  # 1

# Two-stage example
X1 = np.random.randn(n, 3)
A1 = np.random.binomial(1, 0.5, n)
Y1 = X1[:, 0] + A1 + np.random.randn(n)

# Stage 2 history includes prior treatment and outcome
X2 = np.column_stack([X1, A1.reshape(-1, 1), Y1.reshape(-1, 1)])
A2 = np.random.binomial(1, 0.5, n)
Y2 = Y1 + X2[:, 0] + 2.0 * A2 + np.random.randn(n)

two_stage = DTRData(
    outcomes=[Y1, Y2],
    treatments=[A1, A2],
    covariates=[X1, X2],
)
print(f"Stages: {two_stage.n_stages}")  # 2

# Get history at stage 2
H2 = two_stage.get_history(stage=2)
print(f"H2 shape: {H2.shape}")  # (n, 3 + 1 + 1 + 5) = (n, 10)
\end{minted}


\section{Methodological Considerations}

\subsection{Non-Regularity at Decision Boundaries}

When $\gamma(H) = 0$, the treatment decision is indifferent. This creates:
\begin{itemize}
\item Non-standard asymptotics
\item Infinite-dimensional target parameter
\item Bootstrap may fail
\end{itemize}

\textbf{Solutions}:
\begin{enumerate}
\item Hard-threshold bootstrap with perturbation
\item Adaptive confidence intervals
\item Report decision boundary regions
\end{enumerate}

\subsection{Propensity Score Overlap}

A-learning requires $\pi(H) \in (\epsilon, 1-\epsilon)$ for some $\epsilon > 0$.

\begin{minted}[fontsize=\small]{python}
# Propensity trimming
result = a_learning(
    data=data,
    propensity_trim=0.05,  # Trim to [0.05, 0.95]
)
\end{minted}

\subsection{Multi-Stage Dependence}

Backward induction creates correlation across stages:
\begin{itemize}
\item Stage $k$ blip estimate affects pseudo-outcome at stage $k-1$
\item Standard errors must account for this propagation
\item Bootstrap addresses this naturally
\end{itemize}


\section{Implementation Details}

\subsection{Complete Example}

\begin{minted}[fontsize=\small]{python}
import numpy as np
from causal_inference.dtr import DTRData, q_learning, a_learning

# Simulate two-stage data with known optimal regime
np.random.seed(42)
n = 1000

# Stage 1
X1 = np.random.randn(n, 2)
# True propensity: 50% randomization
A1 = np.random.binomial(1, 0.5, n)
# True blip: gamma_1 = 1 + X1[:,0]
true_blip_1 = 1.0 + X1[:, 0]
Y1 = X1[:, 0] + A1 * true_blip_1 + np.random.randn(n) * 0.5

# Stage 2
X2 = np.column_stack([X1, A1.reshape(-1, 1), Y1.reshape(-1, 1)])
A2 = np.random.binomial(1, 0.5, n)
# True blip: gamma_2 = 0.5 + 0.5*Y1
true_blip_2 = 0.5 + 0.5 * Y1
Y2 = Y1 + X2[:, 0] + A2 * true_blip_2 + np.random.randn(n) * 0.5

# Total outcome
Y_total = Y1 + Y2

data = DTRData(outcomes=[Y1, Y2], treatments=[A1, A2], covariates=[X1, X2])

# Q-learning
q_result = q_learning(data, se_method="bootstrap", n_bootstrap=200)
print("=== Q-Learning ===")
print(f"Optimal value: {q_result.value_estimate:.4f}")
print(f"SE: {q_result.value_se:.4f}")

# A-learning (doubly robust)
a_result = a_learning(data, doubly_robust=True, se_method="bootstrap", n_bootstrap=200)
print("\n=== A-Learning ===")
print(f"Optimal value: {a_result.value_estimate:.4f}")
print(f"SE: {a_result.value_se:.4f}")

# Compare to observed mean under random assignment
print(f"\nObserved mean (random): {np.mean(Y_total):.4f}")
print(f"Improvement from optimal: {q_result.value_estimate - np.mean(Y_total):.4f}")

# Validate regime on new data
def true_optimal_regime(H, stage):
    if stage == 1:
        blip = 1.0 + H[0]  # X1[:,0]
    else:
        blip = 0.5 + 0.5 * H[-1]  # Y1
    return int(blip > 0)

# Compare estimated vs true regime agreement
agreements = []
for i in range(n):
    H1 = X1[i]
    est = q_result.optimal_regime(H1, stage=1)
    true = true_optimal_regime(H1, stage=1)
    agreements.append(est == true)
print(f"\nRegime agreement (stage 1): {np.mean(agreements):.1%}")
\end{minted}


\section{Monte Carlo Validation}

\subsection{Test 1: Q-Learning Consistency}

\begin{table}[h]
\centering
\caption{Q-Learning Monte Carlo ($n=500$, 2 stages, 500 runs)}
\begin{tabular}{lccc}
\toprule
Metric & True & Mean Est. & SE \\
\midrule
Optimal Value & 3.45 & 3.42 & 0.12 \\
Stage 1 Blip[0] & 1.00 & 0.98 & 0.08 \\
Stage 2 Blip[0] & 0.50 & 0.49 & 0.06 \\
Regime Agreement & --- & 94.2\% & 1.8\% \\
\bottomrule
\end{tabular}
\end{table}

\subsection{Test 2: A-Learning Double Robustness}

\begin{table}[h]
\centering
\caption{A-Learning Under Misspecification ($n=1000$, 500 runs)}
\begin{tabular}{lccc}
\toprule
Misspecification & Value Bias & Coverage & Notes \\
\midrule
Both correct & 0.02 & 95.4\% & Optimal \\
Propensity wrong & 0.03 & 94.8\% & DR protects \\
Outcome wrong & 0.04 & 93.9\% & DR protects \\
Both wrong & 0.31 & 67.2\% & DR fails \\
\bottomrule
\end{tabular}
\end{table}


\section{Discussion}

\subsection{Key Takeaways}

\begin{enumerate}
\item \textbf{Backward induction}: Q-learning uses dynamic programming principle
\item \textbf{Blip function}: Treatment advantage, determines optimal rule
\item \textbf{Double robustness}: A-learning protects against single-model misspecification
\item \textbf{Value function}: Expected outcome under optimal regime
\end{enumerate}

\subsection{Practical Recommendations}

\begin{enumerate}
\item Use Q-learning when confident in outcome model
\item Use A-learning when uncertain about model specification
\item Always compare both methods for sensitivity analysis
\item Bootstrap for standard errors in multi-stage settings
\item Report regime agreement across methods
\end{enumerate}

\subsection{Extensions}

\begin{itemize}
\item \textbf{Continuous treatments}: Requires nonparametric methods
\item \textbf{Survival outcomes}: Time-to-event DTRs
\item \textbf{Reinforcement learning}: Off-policy evaluation with function approximation
\item \textbf{High-dimensional tailoring}: LASSO/elastic net for blip coefficients
\end{itemize}

\subsection{Connection to Other Methods}

\begin{itemize}
\item \textbf{CATE} (Chapter~\ref{ch:cate}): Single-stage DTR is CATE estimation
\item \textbf{Mediation} (Chapter~\ref{ch:mediation}): Sequential mechanisms
\item \textbf{IV}: When treatment assignment is confounded
\item \textbf{Reinforcement learning}: DTR is offline RL in causal inference
\end{itemize}


% ============================================================================
% PART VI: IMPLEMENTATION
% ============================================================================

\part{Implementation}
\label{part:implementation}

\chapter{Principal Stratification}
\label{ch:principal-strat}

\section{Introduction}

Principal stratification \parencite{frangakis2002principal} classifies units by their joint potential outcomes for intermediate or mediating variables, enabling causal inference in the presence of post-treatment complications.

\begin{definition}[Principal Stratum]
For treatment $Z$ and intermediate variable $D$, the principal stratum is defined by the pair $(D(0), D(1))$---the potential values of $D$ under control and treatment.
\end{definition}

This chapter develops:
\begin{enumerate}
\item \textbf{CACE}: Complier Average Causal Effect (IV setting)
\item \textbf{SACE}: Survivor Average Causal Effect (truncation by death)
\item \textbf{Bounds}: Partial identification without full assumptions
\end{enumerate}


\section{The Four Principal Strata}

\subsection{Classification}

With binary assignment $Z$ and binary treatment receipt $D$:

\begin{table}[h]
\centering
\caption{Principal Strata with Binary D}
\begin{tabular}{lccl}
\toprule
Stratum & $D(0)$ & $D(1)$ & Behavior \\
\midrule
Compliers (C) & 0 & 1 & Take treatment iff assigned \\
Always-Takers (AT) & 1 & 1 & Always take treatment \\
Never-Takers (NT) & 0 & 0 & Never take treatment \\
Defiers (D) & 1 & 0 & Do opposite of assignment \\
\bottomrule
\end{tabular}
\end{table}

\subsection{Identification from Data}

Stratum membership is partially observable:

\begin{table}[h]
\centering
\caption{Observable Groups and Strata Mixtures}
\begin{tabular}{ccl}
\toprule
$Z$ & $D$ & Possible Strata \\
\midrule
1 & 1 & Complier \textit{or} Always-Taker \\
1 & 0 & Never-Taker \textit{or} Defier \\
0 & 1 & Always-Taker \textit{or} Defier \\
0 & 0 & Complier \textit{or} Never-Taker \\
\bottomrule
\end{tabular}
\end{table}

\begin{insight}
Key identification cases:
\begin{itemize}
\item $(Z=0, D=1)$: Definitely Always-Taker or Defier
\item $(Z=1, D=0)$: Definitely Never-Taker or Defier
\end{itemize}
Under monotonicity (no defiers), these become definitive identifications.
\end{insight}


\section{CACE: Complier Average Causal Effect}

\subsection{Definition}

\begin{definition}[Complier Average Causal Effect]
\begin{equation}
\text{CACE} = E[Y(1) - Y(0) \mid D(0) = 0, D(1) = 1]
\end{equation}
The treatment effect for units who comply with their assignment.
\end{definition}

\subsection{Identification Assumptions}

\begin{assumption}[CACE Identification]
\label{ass:cace}
\begin{enumerate}
\item \textbf{Independence}: $Z \perp\!\!\!\perp (Y(0), Y(1), D(0), D(1))$
\item \textbf{Exclusion}: $Y(d, z) = Y(d)$ for all $d, z$ (assignment affects outcome only through treatment)
\item \textbf{Monotonicity}: $D(1) \geq D(0)$ for all units (no defiers)
\item \textbf{Relevance}: $P(D(1) = 1, D(0) = 0) > 0$ (some compliers exist)
\end{enumerate}
\end{assumption}

\begin{proposition}[CACE = LATE]
Under Assumption~\ref{ass:cace}:
\begin{equation}
\text{CACE} = \frac{E[Y \mid Z = 1] - E[Y \mid Z = 0]}{E[D \mid Z = 1] - E[D \mid Z = 0]} = \frac{\text{Reduced Form}}{\text{First Stage}}
\end{equation}
\end{proposition}

\subsection{Strata Proportions}

Under monotonicity, strata proportions are identified:
\begin{align}
\pi_C &= P(D = 1 \mid Z = 1) - P(D = 1 \mid Z = 0) \label{eq:pi-c}\\
\pi_{AT} &= P(D = 1 \mid Z = 0) \label{eq:pi-at}\\
\pi_{NT} &= 1 - P(D = 1 \mid Z = 1) \label{eq:pi-nt}
\end{align}

Note: $\pi_C + \pi_{AT} + \pi_{NT} = 1$ (defiers ruled out).

\subsection{Implementation: 2SLS Approach}

\begin{minted}[fontsize=\small]{python}
from causal_inference.principal_stratification import cace_2sls

result = cace_2sls(
    outcome=Y,
    treatment=D,            # Actual treatment received
    instrument=Z,           # Random assignment
    covariates=X,           # Optional controls
    alpha=0.05,
    inference="robust",     # HC robust SEs
)

# CACE estimate
print(f"CACE: {result['cace']:.4f}")
print(f"SE: {result['se']:.4f}")
print(f"95% CI: [{result['ci_lower']:.4f}, {result['ci_upper']:.4f}]")

# Strata proportions
strata = result['strata_proportions']
print(f"\nCompliers: {strata['compliers']:.1%}")
print(f"Always-takers: {strata['always_takers']:.1%}")
print(f"Never-takers: {strata['never_takers']:.1%}")

# First-stage diagnostics
print(f"\nFirst-stage F: {result['first_stage_f']:.1f}")
\end{minted}

\subsection{Wald Estimator (Simple Case)}

For the simple case without covariates:

\begin{minted}[fontsize=\small]{python}
from causal_inference.principal_stratification import wald_estimator

result = wald_estimator(
    outcome=Y,
    treatment=D,
    instrument=Z,
)

# Direct ratio: (E[Y|Z=1] - E[Y|Z=0]) / (E[D|Z=1] - E[D|Z=0])
print(f"Wald CACE: {result['cace']:.4f}")
print(f"Reduced form: {result['reduced_form']:.4f}")
print(f"First stage: {result['first_stage_coef']:.4f}")
\end{minted}

\subsection{EM Algorithm for CACE}

For maximum likelihood estimation with latent strata:

\begin{minted}[fontsize=\small]{python}
from causal_inference.principal_stratification import cace_em

result = cace_em(
    outcome=Y,
    treatment=D,
    instrument=Z,
    max_iter=100,
    tol=1e-6,
)

print(f"CACE (EM): {result['cace']:.4f}")
print(f"SE (Louis formula): {result['se']:.4f}")
\end{minted}

The EM algorithm:
\begin{enumerate}
\item \textbf{E-step}: Compute posterior strata probabilities given data and current parameters
\item \textbf{M-step}: Update strata proportions and outcome means
\item Repeat until convergence
\end{enumerate}


\section{SACE: Survivor Average Causal Effect}

\subsection{The Truncation by Death Problem}

When outcomes are undefined for non-survivors:
\begin{itemize}
\item $S_i = 1$: Unit survives, $Y_i$ observed
\item $S_i = 0$: Unit dies, $Y_i$ undefined
\end{itemize}

\begin{definition}[SACE]
\begin{equation}
\text{SACE} = E[Y(1) - Y(0) \mid S(0) = 1, S(1) = 1]
\end{equation}
The treatment effect for ``always-survivors''---those who would survive under both treatment conditions.
\end{definition}

\subsection{Survival Strata}

\begin{table}[h]
\centering
\caption{Principal Strata for Survival}
\begin{tabular}{lccl}
\toprule
Stratum & $S(0)$ & $S(1)$ & Interpretation \\
\midrule
Always-Survivors (AS) & 1 & 1 & Would survive regardless \\
Protected (P) & 0 & 1 & Treatment enables survival \\
Harmed (H) & 1 & 0 & Treatment causes death \\
Never-Survivors (NS) & 0 & 0 & Would not survive regardless \\
\bottomrule
\end{tabular}
\end{table}

\subsection{Identification Challenge}

SACE is generally not point-identified:
\begin{itemize}
\item $E[Y(0) \mid \text{AS}]$: Need survivors under $D=0$ who would also survive under $D=1$
\item $E[Y(1) \mid \text{AS}]$: Need survivors under $D=1$ who would also survive under $D=0$
\item Selection into survival confounds identification
\end{itemize}

\subsection{SACE Bounds}

\begin{minted}[fontsize=\small]{python}
from causal_inference.principal_stratification import sace_bounds

result = sace_bounds(
    outcome=Y,           # NaN for non-survivors
    treatment=D,
    survival=S,          # 1 if survived, 0 otherwise
    instrument=Z,        # Optional
    monotonicity="selection",  # Assume S(1) >= S(0)
)

print(f"SACE bounds: [{result['lower_bound']:.3f}, {result['upper_bound']:.3f}]")
print(f"Point estimate: {result['sace']:.3f}")
print(f"Survival rate (treated): {result['proportion_survivors_treat']:.1%}")
print(f"Survival rate (control): {result['proportion_survivors_control']:.1%}")
\end{minted}

\subsection{Monotonicity Assumptions}

\begin{table}[h]
\centering
\caption{SACE Bounds by Assumption}
\begin{tabular}{lll}
\toprule
Assumption & Effect & Interpretation \\
\midrule
None & Widest bounds (Lee) & No restrictions \\
Selection: $S(1) \geq S(0)$ & Tighter & Treatment never harms survival \\
Both (+ treatment mono) & Tightest & Combined restrictions \\
\bottomrule
\end{tabular}
\end{table}

\subsection{Sensitivity Analysis}

\begin{minted}[fontsize=\small]{python}
from causal_inference.principal_stratification import sace_sensitivity

sens = sace_sensitivity(
    outcome=Y,
    treatment=D,
    survival=S,
    alpha_range=(0.0, 1.0),  # Sensitivity parameter range
    n_points=50,
)

# Plot bounds as function of sensitivity parameter
import matplotlib.pyplot as plt
plt.fill_between(sens['alpha'], sens['lower_bound'], sens['upper_bound'],
                 alpha=0.3, label='Bounds')
plt.plot(sens['alpha'], sens['sace'], 'k--', label='Midpoint')
plt.xlabel(r'$\alpha$ (Selection strength)')
plt.ylabel('SACE')
plt.legend()
\end{minted}


\section{Partial Identification and Bounds}

\subsection{When Point Identification Fails}

Without monotonicity or with weaker instruments:
\begin{itemize}
\item Defiers may exist
\item Strata proportions not uniquely identified
\item Only bounds on CACE available
\end{itemize}

\subsection{Manski-Type Bounds}

The worst-case bounds use the outcome support:
\begin{align}
\text{CACE}_{\text{lower}} &= \frac{E[Y \mid Z = 1] - Y_{\max} \cdot P(D = 0 \mid Z = 1) - E[Y \mid Z = 0]}{\pi_C} \\
\text{CACE}_{\text{upper}} &= \frac{E[Y \mid Z = 1] - Y_{\min} \cdot P(D = 0 \mid Z = 1) - E[Y \mid Z = 0]}{\pi_C}
\end{align}


\section{Implementation Details}

\subsection{Complete CACE Example}

\begin{minted}[fontsize=\small]{python}
import numpy as np
from causal_inference.principal_stratification import (
    cace_2sls,
    wald_estimator,
    cace_em,
)

# Simulate data with known strata
np.random.seed(42)
n = 1000

# Random assignment
Z = np.random.binomial(1, 0.5, n)

# True strata: 70% compliers, 15% always-takers, 15% never-takers
strata = np.random.choice([0, 1, 2], n, p=[0.70, 0.15, 0.15])

# Treatment based on stratum
D = np.where(strata == 0, Z,       # Compliers: follow Z
     np.where(strata == 1, 1,       # Always-takers: D=1
              0))                    # Never-takers: D=0

# Outcome with true CACE = 2.0 (for compliers only)
true_cace = 2.0
Y = 1.0 + true_cace * D + np.random.normal(0, 1, n)

# Compare estimators
wald = wald_estimator(Y, D, Z)
tsls = cace_2sls(Y, D, Z)
em = cace_em(Y, D, Z)

print("=== CACE Estimates (true = 2.00) ===")
print(f"Wald:  {wald['cace']:.4f} (SE: {wald['se']:.4f})")
print(f"2SLS:  {tsls['cace']:.4f} (SE: {tsls['se']:.4f})")
print(f"EM:    {em['cace']:.4f} (SE: {em['se']:.4f})")

# Strata proportions
print("\n=== Strata Proportions (true: C=70%, AT=15%, NT=15%) ===")
print(f"Compliers:     {tsls['strata_proportions']['compliers']:.1%}")
print(f"Always-takers: {tsls['strata_proportions']['always_takers']:.1%}")
print(f"Never-takers:  {tsls['strata_proportions']['never_takers']:.1%}")
\end{minted}


\section{Monte Carlo Validation}

\subsection{Test 1: CACE via 2SLS}

\begin{table}[h]
\centering
\caption{CACE Monte Carlo Results ($n=500$, 1,000 runs)}
\begin{tabular}{lccccc}
\toprule
Estimator & True & Mean Est. & Bias & Coverage & SE Accuracy \\
\midrule
Wald & 2.00 & 2.01 & 0.01 & 94.8\% & 3.1\% \\
2SLS & 2.00 & 2.00 & 0.00 & 95.2\% & 2.4\% \\
EM & 2.00 & 2.01 & 0.01 & 94.1\% & 5.8\% \\
\bottomrule
\end{tabular}
\end{table}

\subsection{Test 2: SACE Bounds Coverage}

\begin{table}[h]
\centering
\caption{SACE Bounds ($n=1000$, 500 runs)}
\begin{tabular}{lcccc}
\toprule
Monotonicity & True SACE & Avg Lower & Avg Upper & Contains True \\
\midrule
None & 1.50 & 0.82 & 2.31 & 99.4\% \\
Selection & 1.50 & 1.21 & 1.89 & 96.8\% \\
Both & 1.50 & 1.38 & 1.72 & 95.1\% \\
\bottomrule
\end{tabular}
\end{table}


\section{Discussion}

\subsection{Key Takeaways}

\begin{enumerate}
\item \textbf{Principal strata}: Classification by joint potential outcomes
\item \textbf{CACE = LATE}: Under monotonicity, effect for compliers equals IV estimand
\item \textbf{SACE}: Requires bounds without strong assumptions on survival selection
\item \textbf{Strata proportions}: Identified from first-stage under monotonicity
\end{enumerate}

\subsection{Practical Recommendations}

\begin{enumerate}
\item Report strata proportions alongside CACE---they determine external validity
\item Check first-stage F $> 10$ for weak instrument diagnostics
\item Consider EM for efficiency when normality is plausible
\item Always report bounds when assumptions are questionable
\item Use sensitivity analysis to assess robustness to violations
\end{enumerate}

\subsection{Limitations}

\begin{itemize}
\item \textbf{Monotonicity untestable}: Cannot verify no defiers from data alone
\item \textbf{Exclusion untestable}: Direct effect of $Z$ on $Y$ not refutable
\item \textbf{External validity}: CACE is local to compliers, not ATE
\item \textbf{Covariate adjustment}: EM algorithm doesn't incorporate covariates
\end{itemize}

\subsection{Connection to Other Methods}

\begin{itemize}
\item \textbf{IV} (Chapter~\ref{ch:iv}): Principal stratification provides interpretation for LATE
\item \textbf{Mediation} (Chapter~\ref{ch:mediation}): Principal strata address post-treatment confounding
\item \textbf{Bounds} (Chapter~\ref{ch:bounds}): Partial identification when strata are ambiguous
\item \textbf{Bayesian}: Full posterior on strata membership (Chapter~\ref{ch:bayesian})
\end{itemize}

\chapter{Bayesian Causal Inference}
\label{ch:bayesian}

\section{Introduction}

Bayesian methods offer a coherent framework for causal inference that naturally handles uncertainty, incorporates prior information, and provides direct probability statements about treatment effects.

\begin{definition}[Bayesian ATE]
Given prior $p(\tau)$ and likelihood $p(Y \mid \tau, T, X)$:
\begin{equation}
p(\tau \mid Y, T, X) \propto p(Y \mid \tau, T, X) \cdot p(\tau)
\end{equation}
The posterior $p(\tau \mid Y, T, X)$ represents beliefs about the treatment effect given the data.
\end{definition}

This chapter develops:
\begin{enumerate}
\item \textbf{Conjugate ATE}: Closed-form posterior via normal-normal model
\item \textbf{Bayesian Propensity Scores}: Beta-binomial and logistic approaches
\item \textbf{Bayesian Doubly Robust}: Propagating propensity uncertainty
\item \textbf{Hierarchical ATE}: Partial pooling across sites with MCMC
\end{enumerate}


\section{Conjugate Normal-Normal Model}

\subsection{Mathematical Framework}

The conjugate normal-normal model provides closed-form posteriors without MCMC.

\begin{proposition}[Normal-Normal Conjugacy]
Let:
\begin{align}
\text{Prior:} \quad & \tau \sim \mathcal{N}(\mu_0, \sigma_0^2) \\
\text{Likelihood:} \quad & Y_i \mid T_i, X_i \sim \mathcal{N}(\alpha + \tau T_i + X_i'\beta, \sigma^2)
\end{align}
Then the posterior is:
\begin{equation}
\tau \mid Y, T, X \sim \mathcal{N}(\mu_n, \sigma_n^2)
\end{equation}
where:
\begin{align}
\sigma_n^{-2} &= \sigma_0^{-2} + \sigma_{\text{OLS}}^{-2} \label{eq:post-prec}\\
\mu_n &= \sigma_n^2 \left( \frac{\mu_0}{\sigma_0^2} + \frac{\hat{\tau}_{\text{OLS}}}{\sigma_{\text{OLS}}^2} \right) \label{eq:post-mean}
\end{align}
\end{proposition}

\begin{insight}
The posterior is a precision-weighted average:
\begin{itemize}
\item \textbf{Prior precision}: $1/\sigma_0^2$ (strength of prior beliefs)
\item \textbf{Likelihood precision}: $1/\sigma_{\text{OLS}}^2$ (information from data)
\item More data $\Rightarrow$ posterior approaches MLE
\end{itemize}
\end{insight}

\subsection{Prior Shrinkage}

\begin{definition}[Prior-to-Posterior Shrinkage]
\begin{equation}
\text{Shrinkage} = \frac{\sigma_0^{-2}}{\sigma_0^{-2} + \sigma_{\text{OLS}}^{-2}}
\end{equation}
\begin{itemize}
\item Shrinkage $\approx 0$: Data dominates (weak prior, many observations)
\item Shrinkage $\approx 1$: Prior dominates (strong prior, few observations)
\end{itemize}
\end{definition}

\subsection{Implementation: Bayesian ATE}

\begin{minted}[fontsize=\small]{python}
from causal_inference.bayesian import bayesian_ate

result = bayesian_ate(
    outcomes=Y,
    treatment=T,
    covariates=X,                 # Optional
    prior_mean=0.0,               # Centered on no effect
    prior_sd=10.0,                # Weakly informative
    credible_level=0.95,
    n_posterior_samples=5000,
)

# Posterior summary
print(f"Posterior mean: {result['posterior_mean']:.4f}")
print(f"Posterior SD: {result['posterior_sd']:.4f}")
print(f"95% Credible Interval: [{result['ci_lower']:.4f}, {result['ci_upper']:.4f}]")

# Compare with frequentist
print(f"\nOLS estimate: {result['ols_estimate']:.4f}")
print(f"OLS SE: {result['ols_se']:.4f}")

# Shrinkage diagnostic
print(f"\nPrior shrinkage: {result['prior_to_posterior_shrinkage']:.1%}")
print(f"Effective sample size: {result['effective_sample_size']:.0f}")
\end{minted}

\subsection{Credible vs Confidence Intervals}

\begin{table}[h]
\centering
\caption{Bayesian vs Frequentist Intervals}
\begin{tabular}{lll}
\toprule
Property & Credible Interval & Confidence Interval \\
\midrule
Interpretation & $P(\tau \in \text{CI} \mid \text{data}) = 0.95$ & Coverage over repeated samples \\
Prior role & Incorporates prior beliefs & None \\
Finite sample & Valid with any $n$ & Asymptotic validity \\
Direct probability & Yes & No (coverage interpretation) \\
\bottomrule
\end{tabular}
\end{table}


\section{Bayesian Propensity Score Estimation}

\subsection{Beta-Binomial Model (Stratified)}

For discrete covariates defining strata $s \in \{1, \ldots, S\}$:

\begin{proposition}[Beta-Binomial Conjugacy]
\begin{align}
\text{Prior:} \quad & e_s \sim \text{Beta}(\alpha_0, \beta_0) \\
\text{Likelihood:} \quad & T_i \mid s_i = s \sim \text{Bernoulli}(e_s)
\end{align}
Posterior:
\begin{equation}
e_s \mid T \sim \text{Beta}(\alpha_0 + n_{1,s}, \beta_0 + n_{0,s})
\end{equation}
where $n_{1,s}$ and $n_{0,s}$ are treated and control counts in stratum $s$.
\end{proposition}

\begin{minted}[fontsize=\small]{python}
from causal_inference.bayesian import bayesian_propensity_stratified

result = bayesian_propensity_stratified(
    treatment=T,
    covariates=X,                 # Discretized into strata
    prior_alpha=1.0,              # Beta(1,1) = uniform prior
    prior_beta=1.0,
    n_bins=5,                     # For continuous covariates
    n_posterior_samples=1000,
)

# Posterior propensity for each observation
print(f"Propensity shape: {result['posterior_mean'].shape}")
print(f"Mean uncertainty: {result['mean_uncertainty']:.4f}")

# Stratum-level information
for info in result['stratum_info']:
    print(f"Stratum {info['stratum_id']}: n={info['n_obs']}, "
          f"e={info['posterior_mean']:.3f} +/- {info['posterior_sd']:.3f}")
\end{minted}

\subsection{Bayesian Logistic Regression (Laplace Approximation)}

For continuous covariates, use logistic regression with a normal prior on coefficients:

\begin{align}
\text{Prior:} \quad & \beta \sim \mathcal{N}(0, \sigma_{\text{prior}}^2 I) \\
\text{Likelihood:} \quad & T_i \mid X_i \sim \text{Bernoulli}(\sigma(X_i' \beta)) \\
\text{Posterior:} \quad & \beta \mid T, X \approx \mathcal{N}(\hat{\beta}_{\text{MAP}}, H^{-1})
\end{align}

where $H$ is the Hessian of the negative log-posterior at $\hat{\beta}_{\text{MAP}}$.

\begin{minted}[fontsize=\small]{python}
from causal_inference.bayesian import bayesian_propensity_logistic

result = bayesian_propensity_logistic(
    treatment=T,
    covariates=X,
    prior_sd=10.0,                # Weakly informative on coefficients
    n_posterior_samples=1000,
    include_intercept=True,
)

# Coefficient uncertainty
print("Coefficient posteriors:")
for i, (mean, sd) in enumerate(zip(result['coefficient_mean'],
                                    result['coefficient_sd'])):
    print(f"  beta_{i}: {mean:.3f} +/- {sd:.3f}")

# Propensity uncertainty propagated to each unit
print(f"\nMean propensity uncertainty: {result['mean_uncertainty']:.4f}")
print(f"Propensity range: {result['propensity_range']:.3f}")
\end{minted}

\subsection{Automatic Method Selection}

\begin{minted}[fontsize=\small]{python}
from causal_inference.bayesian import bayesian_propensity

# Auto-selects based on covariate characteristics
result = bayesian_propensity(
    treatment=T,
    covariates=X,
    method="auto",                # "stratified", "logistic", or "auto"
)

print(f"Method selected: {result['method']}")
\end{minted}

Decision logic:
\begin{itemize}
\item All covariates have $\leq 10$ unique values $\Rightarrow$ stratified
\item Otherwise $\Rightarrow$ logistic
\end{itemize}


\section{Bayesian Doubly Robust Estimation}

\subsection{Uncertainty Propagation}

The Bayesian DR estimator propagates propensity uncertainty through the AIPW formula:

\begin{enumerate}
\item Estimate Bayesian propensity $\Rightarrow$ samples $\{e^{(k)}\}_{k=1}^S$
\item Fit frequentist outcome models $\Rightarrow$ $\hat{\mu}_0(X), \hat{\mu}_1(X)$
\item For each propensity sample $k$:
\begin{equation}
\hat{\tau}^{(k)} = \frac{1}{n} \sum_{i=1}^n \left[ \frac{T_i}{e_i^{(k)}} (Y_i - \hat{\mu}_1(X_i)) + \hat{\mu}_1(X_i) - \frac{1-T_i}{1-e_i^{(k)}} (Y_i - \hat{\mu}_0(X_i)) - \hat{\mu}_0(X_i) \right]
\end{equation}
\item Posterior: $\tau \mid \text{data} \sim \{\hat{\tau}^{(1)}, \ldots, \hat{\tau}^{(S)}\}$
\end{enumerate}

\begin{insight}
Why hybrid Bayesian-Frequentist?
\begin{itemize}
\item Propensity uncertainty: Often most important for DR
\item Outcome models: Frequentist OLS is robust and fast
\item Total uncertainty: Combined using influence function variance
\end{itemize}
\end{insight}

\subsection{Implementation}

\begin{minted}[fontsize=\small]{python}
from causal_inference.bayesian import bayesian_dr_ate

result = bayesian_dr_ate(
    outcomes=Y,
    treatment=T,
    covariates=X,
    propensity_method="auto",     # "stratified", "logistic", or "auto"
    n_posterior_samples=1000,
    credible_level=0.95,
    trim_threshold=0.01,          # Propensity clipping
)

# Bayesian DR estimate
print(f"ATE: {result['estimate']:.4f}")
print(f"Posterior SD: {result['se']:.4f}")
print(f"95% CI: [{result['ci_lower']:.4f}, {result['ci_upper']:.4f}]")

# Comparison with frequentist DR
print(f"\nFrequentist DR: {result['frequentist_estimate']:.4f}")
print(f"Frequentist SE: {result['frequentist_se']:.4f}")

# Propensity diagnostics
print(f"\nMean propensity uncertainty: {result['propensity_mean_uncertainty']:.4f}")
print(f"Outcome R2: {result['outcome_r2']:.3f}")
\end{minted}

\subsection{Double Robustness Property}

Like frequentist DR, Bayesian DR is consistent when EITHER:
\begin{itemize}
\item Propensity model is correctly specified, OR
\item Outcome models are correctly specified
\end{itemize}

The posterior captures uncertainty from propensity estimation, with outcome model variance added via the influence function.


\section{Hierarchical Bayesian ATE}

\subsection{Multi-Site Studies}

When treatment effects vary across groups/sites, hierarchical models provide \textit{partial pooling}:

\begin{definition}[Hierarchical Model]
\begin{align}
\mu &\sim \mathcal{N}(0, \sigma_\mu^2) && \text{(Population ATE)} \\
\tau &\sim \text{HalfNormal}(\sigma_\tau) && \text{(Between-group SD)} \\
\theta_j &\sim \mathcal{N}(\mu, \tau^2) && \text{(Group-specific ATE)} \\
Y_{ij} &\sim \mathcal{N}(\alpha + \theta_j T_{ij}, \sigma^2) && \text{(Observations)}
\end{align}
\end{definition}

\begin{insight}
Partial Pooling Benefits:
\begin{itemize}
\item Small groups: Shrink toward population mean (borrows strength)
\item Large groups: Retain group-specific signal
\item $\tau$ estimated from data: Automatic shrinkage calibration
\end{itemize}
\end{insight}

\subsection{Non-Centered Parameterization}

To avoid the ``funnel'' geometry in hierarchical models:

\begin{align}
\theta_{\text{raw},j} &\sim \mathcal{N}(0, 1) && \text{(Standardized)} \\
\theta_j &= \mu + \tau \cdot \theta_{\text{raw},j} && \text{(Transformed)}
\end{align}

This parameterization improves NUTS sampler efficiency when $\tau$ is small.

\subsection{Implementation with PyMC}

\begin{minted}[fontsize=\small]{python}
from causal_inference.bayesian import hierarchical_bayesian_ate

result = hierarchical_bayesian_ate(
    outcomes=Y,
    treatment=T,
    groups=group_ids,             # Site/group identifier
    # Prior hyperparameters
    mu_prior_sd=10.0,             # Population ATE prior SD
    tau_prior_sd=5.0,             # Heterogeneity prior SD
    sigma_prior_sd=10.0,          # Observation noise prior SD
    # MCMC settings
    n_samples=2000,
    n_chains=4,
    n_tune=1000,
    target_accept=0.9,
    # Output
    credible_level=0.95,
)

# Population-level estimates
print("=== Population ATE ===")
print(f"Mean: {result['population_ate']:.4f}")
print(f"SE: {result['population_ate_se']:.4f}")
print(f"95% CI: [{result['population_ate_ci_lower']:.4f}, "
      f"{result['population_ate_ci_upper']:.4f}]")

# Heterogeneity
print(f"\n=== Heterogeneity ===")
print(f"tau: {result['tau']:.4f}")
print(f"95% CI: [{result['tau_ci_lower']:.4f}, {result['tau_ci_upper']:.4f}]")

# Group-specific effects
print("\n=== Group ATEs ===")
for i, gid in enumerate(result['group_ids']):
    ate = result['group_ates'][i]
    se = result['group_ate_ses'][i]
    print(f"Group {gid}: {ate:.3f} +/- {se:.3f}")
\end{minted}

\subsection{MCMC Diagnostics}

\begin{minted}[fontsize=\small]{python}
# Convergence checks
print("\n=== MCMC Diagnostics ===")
print(f"R-hat (max): {result['rhat_max']:.4f}")
print(f"ESS (min): {result['ess_min']:.0f}")
print(f"Divergences: {result['divergences']}")

# Interpretation
if result['rhat_max'] < 1.05:
    print("R-hat OK: Chains converged")
else:
    print("WARNING: R-hat > 1.05, chains may not have converged")

if result['ess_min'] > 400:
    print("ESS OK: Adequate effective samples")
else:
    print("WARNING: Low ESS, increase n_samples")

if result['divergences'] == 0:
    print("Divergences OK: NUTS sampler healthy")
else:
    print("WARNING: Divergences detected, consider reparameterization")
\end{minted}

\begin{table}[h]
\centering
\caption{MCMC Diagnostic Thresholds}
\begin{tabular}{lll}
\toprule
Diagnostic & Threshold & Interpretation \\
\midrule
R-hat & $< 1.05$ & Chains have converged \\
ESS & $> 400$ & Adequate effective samples \\
Divergences & $= 0$ & NUTS sampler healthy \\
\bottomrule
\end{tabular}
\end{table}


\section{Implementation Details}

\subsection{Complete Bayesian Pipeline}

\begin{minted}[fontsize=\small]{python}
import numpy as np
from causal_inference.bayesian import (
    bayesian_ate,
    bayesian_dr_ate,
    hierarchical_bayesian_ate,
)

# Generate multi-site data
np.random.seed(42)
n_groups = 5
n_per_group = 100
n = n_groups * n_per_group

# True heterogeneous effects
true_population_ate = 2.0
true_group_effects = [1.5, 1.8, 2.0, 2.2, 2.5]

groups = np.repeat(np.arange(n_groups), n_per_group)
T = np.random.binomial(1, 0.5, n).astype(float)
X = np.random.normal(0, 1, (n, 2))

# Outcomes with group-specific effects
Y = np.array([true_group_effects[g] for g in groups]) * T
Y += 0.5 * X[:, 0] + np.random.normal(0, 1, n)

# === Method 1: Simple Bayesian ATE (ignores grouping) ===
simple = bayesian_ate(Y, T, X)
print("=== Simple Bayesian ATE ===")
print(f"Posterior mean: {simple['posterior_mean']:.4f}")
print(f"95% CI: [{simple['ci_lower']:.4f}, {simple['ci_upper']:.4f}]")

# === Method 2: Bayesian DR ===
dr = bayesian_dr_ate(Y, T, X)
print("\n=== Bayesian DR ===")
print(f"Posterior mean: {dr['estimate']:.4f}")
print(f"95% CI: [{dr['ci_lower']:.4f}, {dr['ci_upper']:.4f}]")

# === Method 3: Hierarchical (accounts for grouping) ===
hier = hierarchical_bayesian_ate(
    Y, T, groups,
    n_samples=1000,
    n_chains=2,
    progressbar=False,
)
print("\n=== Hierarchical Bayesian ===")
print(f"Population ATE: {hier['population_ate']:.4f}")
print(f"Heterogeneity tau: {hier['tau']:.4f}")
print(f"Group effects: {hier['group_ates']}")
\end{minted}


\section{Monte Carlo Validation}

\subsection{Test 1: Conjugate ATE Coverage}

\begin{table}[h]
\centering
\caption{Bayesian ATE Monte Carlo Results ($n=200$, 1,000 runs)}
\begin{tabular}{lcccc}
\toprule
Prior SD & True ATE & Mean Est. & Coverage & Avg Shrinkage \\
\midrule
10.0 (weak) & 2.00 & 2.01 & 94.8\% & 1.2\% \\
1.0 (moderate) & 2.00 & 1.98 & 93.1\% & 18.4\% \\
0.5 (strong) & 2.00 & 1.42 & 68.2\% & 52.3\% \\
\bottomrule
\end{tabular}
\end{table}

\begin{insight}
Prior sensitivity:
\begin{itemize}
\item Weak prior ($\sigma_0 = 10$): Posterior tracks data, valid coverage
\item Strong prior ($\sigma_0 = 0.5$): Prior dominates, undercoverage when misspecified
\end{itemize}
\end{insight}

\subsection{Test 2: Bayesian DR vs Frequentist DR}

\begin{table}[h]
\centering
\caption{Bayesian DR Monte Carlo Results ($n=500$, 500 runs)}
\begin{tabular}{lccccc}
\toprule
Method & True & Mean Est. & Bias & Coverage & Mean SE \\
\midrule
Frequentist DR & 2.00 & 2.01 & 0.01 & 94.6\% & 0.18 \\
Bayesian DR & 2.00 & 2.01 & 0.01 & 95.2\% & 0.19 \\
\bottomrule
\end{tabular}
\end{table}

\subsection{Test 3: Hierarchical Model (Partial Pooling)}

\begin{table}[h]
\centering
\caption{Hierarchical ATE (5 groups, 100 obs/group, true $\tau = 0.5$)}
\begin{tabular}{lcccc}
\toprule
Group & True $\theta$ & Unpooled Est. & Hierarchical Est. & Shrinkage \\
\midrule
1 (small effect) & 1.5 & 1.52 & 1.68 & 12\% \\
2 & 1.8 & 1.79 & 1.85 & 5\% \\
3 (population) & 2.0 & 2.03 & 2.01 & 2\% \\
4 & 2.2 & 2.18 & 2.15 & 3\% \\
5 (large effect) & 2.5 & 2.48 & 2.38 & 8\% \\
\midrule
Population $\mu$ & 2.0 & --- & 2.01 & --- \\
\bottomrule
\end{tabular}
\end{table}


\section{Discussion}

\subsection{Key Takeaways}

\begin{enumerate}
\item \textbf{Conjugate models}: Closed-form when applicable (normal-normal, beta-binomial)
\item \textbf{Credible intervals}: Direct probability interpretation
\item \textbf{Prior sensitivity}: Weak priors let data dominate; check shrinkage
\item \textbf{Hierarchical models}: Optimal shrinkage for multi-site studies
\item \textbf{MCMC diagnostics}: R-hat $< 1.05$, ESS $> 400$, divergences $= 0$
\end{enumerate}

\subsection{Practical Recommendations}

\begin{enumerate}
\item \textbf{Default priors}: Use weakly informative ($\sigma_0 = 10$ for standardized outcomes)
\item \textbf{Check shrinkage}: If $> 10\%$, prior may be influencing results significantly
\item \textbf{Report both}: Show Bayesian and frequentist estimates for comparison
\item \textbf{Hierarchical when applicable}: Multi-site, clustered, or meta-analytic data
\item \textbf{Validate MCMC}: Always check diagnostics before trusting posteriors
\end{enumerate}

\subsection{When to Use Bayesian Methods}

\begin{table}[h]
\centering
\caption{Method Selection Guide}
\begin{tabular}{ll}
\toprule
Situation & Recommended Approach \\
\midrule
Incorporate prior information & Bayesian with informative prior \\
Small sample size & Bayesian (regularization via prior) \\
Multi-site study & Hierarchical Bayesian \\
Need probability statements & Bayesian credible intervals \\
Standard RCT, large $n$ & Frequentist (simpler, faster) \\
\bottomrule
\end{tabular}
\end{table}

\subsection{Limitations}

\begin{itemize}
\item \textbf{Prior specification}: Results depend on prior choices
\item \textbf{Computational cost}: MCMC slower than closed-form
\item \textbf{Model misspecification}: Hierarchical assumes exchangeability
\item \textbf{Laplace approximation}: May be inaccurate for small $n$ or multimodal posteriors
\end{itemize}

\subsection{Connection to Other Methods}

\begin{itemize}
\item \textbf{Principal Stratification} (Chapter~\ref{ch:principal-strat}): Bayesian strata membership inference
\item \textbf{CATE} (Chapter~\ref{ch:cate}): BART for heterogeneous effects
\item \textbf{Selection Models} (Chapter~\ref{ch:selection}): Bayesian Heckman correction
\item \textbf{Meta-Analysis}: Hierarchical models for combining studies
\end{itemize}


\chapter{Cross-Language Validation}
\label{ch:validation}

\section{Introduction}

Implementation correctness is foundational to causal inference. A methodologically sound estimator produces misleading results if implemented incorrectly. This chapter describes the six-layer validation architecture used throughout this book.

\begin{definition}[Cross-Language Validation]
Independent implementations in different languages (Python, Julia) should produce identical results for identical inputs:
\begin{equation}
\left| \hat{\tau}_{\text{Python}} - \hat{\tau}_{\text{Julia}} \right| < 10^{-10}
\end{equation}
Agreement validates algorithmic correctness; divergence signals implementation errors.
\end{definition}

This chapter develops:
\begin{enumerate}
\item \textbf{Six-Layer Architecture}: Comprehensive validation framework
\item \textbf{Cross-Language Testing}: Python $\leftrightarrow$ Julia parity
\item \textbf{Monte Carlo Validation}: Statistical properties verification
\item \textbf{Adversarial Testing}: Edge case robustness
\end{enumerate}


\section{Six-Layer Validation Architecture}

\subsection{Architecture Overview}

\begin{table}[h]
\centering
\caption{Validation Layers}
\begin{tabular}{clll}
\toprule
Layer & Name & Purpose & Implementation \\
\midrule
1 & Known-Answer & Hand-calculated expected values & \texttt{tests/test\_*/*.py} \\
2 & Adversarial & Edge cases, boundary conditions & \texttt{tests/validation/adversarial/} \\
3 & Monte Carlo & Statistical properties (5K+ runs) & \texttt{tests/validation/monte\_carlo/} \\
4 & Cross-Language & Python $\leftrightarrow$ Julia parity & \texttt{tests/validation/cross\_language/} \\
5 & R Triangulation & External reference (future) & --- \\
6 & Golden Reference & Frozen reference results & \texttt{tests/golden\_results/} \\
\bottomrule
\end{tabular}
\end{table}

\subsection{Layer 1: Known-Answer Tests}

For simple cases with hand-calculable results:

\begin{minted}[fontsize=\small]{python}
def test_simple_ate_known_answer():
    """Four-unit case with known answer."""
    outcomes = np.array([7.0, 5.0, 3.0, 1.0])
    treatment = np.array([1, 1, 0, 0])

    # Hand calculation:
    # Y_1 = (7+5)/2 = 6.0
    # Y_0 = (3+1)/2 = 2.0
    # ATE = 6.0 - 2.0 = 4.0
    result = simple_ate(outcomes, treatment)

    assert result["estimate"] == 4.0, "ATE should be exactly 4.0"
\end{minted}

\subsection{Layer 2: Adversarial Tests}

Tests that probe edge cases and failure modes:

\begin{minted}[fontsize=\small]{python}
class TestSimpleATEAdversarial:
    """Edge cases that should fail gracefully."""

    def test_single_treated(self):
        """Only one treated unit (n_1 = 1)."""
        outcomes = np.array([10.0, 2.0, 3.0])
        treatment = np.array([1, 0, 0])

        # Should handle gracefully (large SE, wide CI)
        result = simple_ate(outcomes, treatment)
        assert result["se"] > 0

    def test_all_treated(self):
        """All units treated (no control group)."""
        outcomes = np.array([5.0, 6.0, 7.0])
        treatment = np.array([1, 1, 1])

        # Should raise informative error
        with pytest.raises(ValueError, match="control"):
            simple_ate(outcomes, treatment)

    def test_extreme_outlier(self):
        """Outlier 1000x larger than typical values."""
        outcomes = np.array([1.0, 2.0, 1e6, 3.0])
        treatment = np.array([1, 1, 0, 0])

        # Should compute without numerical issues
        result = simple_ate(outcomes, treatment)
        assert np.isfinite(result["estimate"])
\end{minted}

\subsection{Layer 3: Monte Carlo Validation}

Statistical validation over thousands of simulations:

\begin{minted}[fontsize=\small]{python}
@pytest.mark.monte_carlo
def test_simple_ate_unbiased():
    """Monte Carlo validation: 5000 runs, bias < 0.05, coverage 94-96%."""
    n_runs = 5000
    true_ate = 2.0

    estimates = []
    ci_covers = []

    for seed in range(n_runs):
        Y, T = generate_rct(n=100, true_ate=true_ate, seed=seed)
        result = simple_ate(Y, T)

        estimates.append(result["estimate"])
        ci_covers.append(result["ci_lower"] < true_ate < result["ci_upper"])

    # Bias check
    bias = np.mean(estimates) - true_ate
    assert abs(bias) < 0.05, f"Bias {bias:.4f} exceeds threshold"

    # Coverage check
    coverage = np.mean(ci_covers)
    assert 0.94 < coverage < 0.96, f"Coverage {coverage:.2%} outside range"
\end{minted}

\begin{table}[h]
\centering
\caption{Monte Carlo Validation Standards}
\begin{tabular}{lccc}
\toprule
Method Type & Bias Target & Coverage Target & SE Accuracy \\
\midrule
RCT (unconfounded) & $< 0.05$ & 94--96\% & $< 10\%$ \\
Observational & $< 0.10$ & 93--97\% & $< 15\%$ \\
PSM & $< 0.30$ & $\geq 95\%$ & $< 150\%$ \\
\bottomrule
\end{tabular}
\end{table}


\section{Cross-Language Validation}

\subsection{Design Principles}

1. \textbf{Independent implementations}: Python uses NumPy/SciPy; Julia uses native arrays
2. \textbf{Shared data}: Same random seeds and test fixtures
3. \textbf{Strict tolerance}: Agreement to 10 decimal places (rtol $< 10^{-10}$)
4. \textbf{Comprehensive coverage}: All estimators, not just simple cases

\subsection{Julia Interface Architecture}

\begin{minted}[fontsize=\small]{python}
# tests/validation/cross_language/julia_interface.py

from juliacall import Main as jl

# Initialize Julia project
JULIA_PROJECT_PATH = "julia/"
jl.seval(f'import Pkg; Pkg.activate("{JULIA_PROJECT_PATH}")')
jl.seval("using CausalEstimators")

def julia_simple_ate(outcomes, treatment, alpha=0.05):
    """Call Julia SimpleATE via juliacall."""
    # Convert to Julia types
    jl_outcomes = jl.collect(outcomes.astype(np.float64))
    jl_treatment = jl.collect(treatment.astype(bool))

    # Create problem and solve
    problem = jl.RCTProblem(jl_outcomes, jl_treatment)
    solution = jl.solve(problem, jl.SimpleATE())

    return {
        "estimate": float(solution.estimate),
        "se": float(solution.se),
        "ci_lower": float(solution.ci_lower),
        "ci_upper": float(solution.ci_upper),
    }
\end{minted}

\subsection{Cross-Language Test Pattern}

\begin{minted}[fontsize=\small]{python}
# tests/validation/cross_language/test_python_julia_simple_ate.py

import pytest
from tests.validation.cross_language.julia_interface import (
    is_julia_available,
    julia_simple_ate,
)

pytestmark = pytest.mark.skipif(
    not is_julia_available(),
    reason="Julia not available for cross-validation"
)

class TestPythonJuliaSimpleATE:
    """Cross-validate Python simple_ate against Julia SimpleATE."""

    def test_basic_case(self):
        """Basic 4-unit case."""
        outcomes = np.array([7.0, 5.0, 3.0, 1.0])
        treatment = np.array([1, 1, 0, 0])

        python_result = simple_ate(outcomes, treatment)
        julia_result = julia_simple_ate(outcomes, treatment)

        # Validate to 10 decimal places
        assert np.isclose(
            python_result["estimate"],
            julia_result["estimate"],
            rtol=1e-10
        ), f"Estimate mismatch: Python={python_result['estimate']}, Julia={julia_result['estimate']}"

        assert np.isclose(
            python_result["se"],
            julia_result["se"],
            rtol=1e-10
        ), f"SE mismatch: Python={python_result['se']}, Julia={julia_result['se']}"

    def test_large_sample(self):
        """Large sample (n=1000) stress test."""
        np.random.seed(42)
        n = 1000
        treatment = np.random.binomial(1, 0.5, n)
        outcomes = 2.0 * treatment + np.random.normal(0, 1, n)

        python_result = simple_ate(outcomes, treatment)
        julia_result = julia_simple_ate(outcomes, treatment)

        assert np.isclose(python_result["estimate"], julia_result["estimate"], rtol=1e-10)
        assert np.isclose(python_result["se"], julia_result["se"], rtol=1e-10)
\end{minted}

\subsection{Method Coverage}

\begin{table}[h]
\centering
\caption{Cross-Language Test Coverage}
\begin{tabular}{lll}
\toprule
Method Family & Python Module & Julia Module \\
\midrule
RCT & \texttt{rct.estimators} & \texttt{RCT.jl} \\
DiD & \texttt{did.*} & \texttt{did/*.jl} \\
IV & \texttt{iv.*} & \texttt{iv/*.jl} \\
RDD & \texttt{rdd.*} & \texttt{rdd/*.jl} \\
SCM & \texttt{scm.*} & \texttt{scm/*.jl} \\
CATE & \texttt{cate.*} & \texttt{cate/*.jl} \\
PSM & \texttt{psm.*} & \texttt{psm/*.jl} \\
Sensitivity & \texttt{sensitivity.*} & \texttt{sensitivity/*.jl} \\
QTE & \texttt{qte.*} & \texttt{qte/*.jl} \\
MTE & \texttt{mte.*} & \texttt{mte/*.jl} \\
Bounds & \texttt{bounds.*} & \texttt{bounds/*.jl} \\
Mediation & \texttt{mediation.*} & \texttt{mediation/*.jl} \\
\bottomrule
\end{tabular}
\end{table}


\section{Data Generating Processes}

\subsection{DGP Architecture}

Centralized DGP generators ensure identical data across tests:

\begin{minted}[fontsize=\small]{python}
# tests/validation/monte_carlo/dgp_generators.py

def dgp_simple_rct(n=100, true_ate=2.0, random_state=None):
    """Simple RCT with homoskedastic errors."""
    rng = np.random.default_rng(random_state)

    treatment = rng.binomial(1, 0.5, n)
    potential_Y0 = rng.normal(0, 1, n)
    potential_Y1 = potential_Y0 + true_ate
    outcomes = np.where(treatment == 1, potential_Y1, potential_Y0)

    return outcomes, treatment

def dgp_confounded_observational(n=500, true_ate=2.0, random_state=None):
    """Observational study with confounding."""
    rng = np.random.default_rng(random_state)

    # Confounder affects both treatment and outcome
    X = rng.normal(0, 1, n)
    propensity = 1 / (1 + np.exp(-0.5 * X))
    treatment = rng.binomial(1, propensity)

    # Outcome depends on X and treatment
    outcomes = 0.5 * X + true_ate * treatment + rng.normal(0, 1, n)

    return outcomes, treatment, X, propensity
\end{minted}

\subsection{Method-Specific DGPs}

\begin{minted}[fontsize=\small]{python}
# tests/validation/monte_carlo/dgp_did.py

def dgp_did_staggered(n_units=100, n_periods=10, true_ate=2.0, random_state=None):
    """Staggered DiD with heterogeneous timing."""
    rng = np.random.default_rng(random_state)

    # Staggered treatment adoption
    treatment_timing = rng.choice(range(3, n_periods), n_units)

    # Panel data with unit and time fixed effects
    unit_fe = rng.normal(0, 0.5, n_units)
    time_fe = rng.normal(0, 0.3, n_periods)

    # Build panel
    data = []
    for i in range(n_units):
        for t in range(n_periods):
            treated = 1 if t >= treatment_timing[i] else 0
            Y = unit_fe[i] + time_fe[t] + true_ate * treated + rng.normal(0, 0.5)
            data.append({"unit": i, "time": t, "treated": treated, "Y": Y})

    return pd.DataFrame(data)
\end{minted}


\section{Implementation Details}

\subsection{Running Validation Tests}

\begin{minted}[fontsize=\small]{bash}
# Full validation suite
pytest tests/validation/ -v --tb=short

# Layer-specific tests
pytest tests/validation/monte_carlo/ -v -m monte_carlo
pytest tests/validation/adversarial/ -v
pytest tests/validation/cross_language/ -v

# Skip slow Monte Carlo tests
pytest tests/ -m "not slow"

# Single method validation
pytest tests/validation/cross_language/test_python_julia_did.py -v
\end{minted}

\subsection{CI/CD Integration}

\begin{minted}[fontsize=\small]{yaml}
# .github/workflows/validation.yml
name: Validation Suite

on: [push, pull_request]

jobs:
  monte-carlo:
    runs-on: ubuntu-latest
    steps:
      - uses: actions/checkout@v4
      - name: Run Monte Carlo tests
        run: pytest tests/validation/monte_carlo/ -v --timeout=600

  cross-language:
    runs-on: ubuntu-latest
    steps:
      - uses: actions/checkout@v4
      - uses: julia-actions/setup-julia@v1
      - name: Install Julia deps
        run: julia --project=julia -e 'using Pkg; Pkg.instantiate()'
      - name: Run cross-language tests
        run: pytest tests/validation/cross_language/ -v
\end{minted}

\subsection{Debugging Validation Failures}

\begin{minted}[fontsize=\small]{python}
def debug_cross_language_failure(method_name, python_result, julia_result):
    """Diagnostic output for validation failures."""
    print(f"\n=== Cross-Language Failure: {method_name} ===")

    for key in python_result:
        py_val = python_result[key]
        jl_val = julia_result.get(key, "MISSING")

        if isinstance(py_val, (int, float)) and isinstance(jl_val, (int, float)):
            diff = abs(py_val - jl_val)
            rel_diff = diff / (abs(py_val) + 1e-10)
            status = "OK" if rel_diff < 1e-10 else "MISMATCH"
            print(f"  {key}: Python={py_val:.12e}, Julia={jl_val:.12e}, "
                  f"diff={diff:.2e} ({status})")
        else:
            print(f"  {key}: Python={py_val}, Julia={jl_val}")
\end{minted}


\section{Monte Carlo Standards}

\subsection{Validation Metrics}

\begin{definition}[Monte Carlo Validation Criteria]
For $R$ simulation runs with true effect $\tau$:
\begin{align}
\text{Bias} &= \frac{1}{R} \sum_{r=1}^R \hat{\tau}_r - \tau \\
\text{Coverage} &= \frac{1}{R} \sum_{r=1}^R \mathbf{1}[\hat{\tau}_r^L < \tau < \hat{\tau}_r^U] \\
\text{SE Accuracy} &= \frac{\left| \overline{\text{SE}} - \text{SD}(\hat{\tau}) \right|}{\text{SD}(\hat{\tau})}
\end{align}
\end{definition}

\subsection{Validation Helper}

\begin{minted}[fontsize=\small]{python}
# tests/validation/utils.py

def validate_monte_carlo_results(
    estimates,
    standard_errors,
    ci_lowers,
    ci_uppers,
    true_value,
    bias_threshold=0.05,
    coverage_range=(0.94, 0.96),
    se_accuracy_threshold=0.10,
):
    """Validate Monte Carlo simulation results."""
    estimates = np.array(estimates)
    standard_errors = np.array(standard_errors)
    ci_lowers = np.array(ci_lowers)
    ci_uppers = np.array(ci_uppers)

    # Bias
    bias = np.mean(estimates) - true_value
    bias_ok = abs(bias) < bias_threshold

    # Coverage
    coverage = np.mean((ci_lowers < true_value) & (true_value < ci_uppers))
    coverage_ok = coverage_range[0] < coverage < coverage_range[1]

    # SE accuracy
    empirical_se = np.std(estimates)
    mean_reported_se = np.mean(standard_errors)
    se_accuracy = abs(mean_reported_se - empirical_se) / empirical_se
    se_accuracy_ok = se_accuracy < se_accuracy_threshold

    return {
        "bias": bias,
        "bias_ok": bias_ok,
        "coverage": coverage,
        "coverage_ok": coverage_ok,
        "se_accuracy": se_accuracy,
        "se_accuracy_ok": se_accuracy_ok,
        "all_pass": bias_ok and coverage_ok and se_accuracy_ok,
    }
\end{minted}


\section{Golden Reference Tests}

\subsection{Frozen Reference Values}

Golden references provide regression testing against known-correct outputs:

\begin{minted}[fontsize=\small]{python}
# tests/golden_results/simple_ate.json
{
    "test_balanced_rct": {
        "input": {
            "outcomes": [7.0, 5.0, 3.0, 1.0],
            "treatment": [1, 1, 0, 0]
        },
        "expected": {
            "estimate": 4.0,
            "se": 1.4142135623730951,
            "ci_lower": -0.5319734684099524,
            "ci_upper": 8.531973468409952
        }
    }
}
\end{minted}

\begin{minted}[fontsize=\small]{python}
def test_against_golden():
    """Test against frozen reference values."""
    import json

    with open("tests/golden_results/simple_ate.json") as f:
        golden = json.load(f)

    for test_name, data in golden.items():
        result = simple_ate(
            np.array(data["input"]["outcomes"]),
            np.array(data["input"]["treatment"])
        )

        for key, expected in data["expected"].items():
            assert np.isclose(result[key], expected, rtol=1e-10), \
                f"{test_name}.{key}: got {result[key]}, expected {expected}"
\end{minted}


\section{Discussion}

\subsection{Key Takeaways}

\begin{enumerate}
\item \textbf{Six layers}: Comprehensive validation from unit tests to cross-language parity
\item \textbf{Cross-language}: Python and Julia must agree to 10 decimal places
\item \textbf{Monte Carlo}: 5,000+ runs verify statistical properties (bias, coverage, SE)
\item \textbf{Adversarial}: Edge cases test robustness and error handling
\item \textbf{Golden references}: Regression testing against frozen correct outputs
\end{enumerate}

\subsection{Practical Recommendations}

\begin{enumerate}
\item \textbf{Run full suite}: Before any release or major change
\item \textbf{Check Monte Carlo}: When modifying estimator internals
\item \textbf{Cross-language first}: When adding new estimators
\item \textbf{Update golden}: When correct behavior changes intentionally
\item \textbf{Document xfails}: Known limitations with explicit reasons
\end{enumerate}

\subsection{Validation Status by Method}

\begin{table}[h]
\centering
\caption{Current Validation Status}
\begin{tabular}{lccccc}
\toprule
Method & Known-Answer & Adversarial & Monte Carlo & Cross-Language & Status \\
\midrule
RCT & \checkmark & \checkmark & \checkmark & \checkmark & VERIFIED \\
DiD & \checkmark & \checkmark & \checkmark & \checkmark & VERIFIED \\
IV & \checkmark & \checkmark & \checkmark & \checkmark & VERIFIED \\
RDD & \checkmark & \checkmark & \checkmark & \checkmark & VERIFIED \\
SCM & \checkmark & \checkmark & \checkmark & \checkmark & VERIFIED \\
CATE & \checkmark & \checkmark & \checkmark & \checkmark & VERIFIED \\
PSM & \checkmark & \checkmark & \checkmark & \checkmark & VERIFIED \\
QTE & \checkmark & \checkmark & \checkmark & \checkmark & VERIFIED \\
MTE & \checkmark & \checkmark & \checkmark & \checkmark & VERIFIED \\
Mediation & \checkmark & \checkmark & \checkmark & \checkmark & VERIFIED \\
Selection & \checkmark & \checkmark & \checkmark & \checkmark & VERIFIED \\
DTR & \checkmark & \checkmark & \checkmark & \checkmark & VERIFIED \\
Principal Strat & \checkmark & \checkmark & \checkmark & \checkmark & VERIFIED \\
Bayesian & \checkmark & \checkmark & \checkmark & --- & PARTIAL \\
\bottomrule
\end{tabular}
\end{table}

\subsection{Limitations}

\begin{itemize}
\item \textbf{R triangulation}: Not yet implemented (planned Layer 5)
\item \textbf{Julia dependency}: Cross-language tests require juliacall installation
\item \textbf{Monte Carlo time}: Full suite takes $\sim$30 minutes
\item \textbf{Coverage gaps}: Some edge cases may not be tested
\end{itemize}

\subsection{Connection to Other Chapters}

\begin{itemize}
\item \textbf{All chapters}: Validation ensures correctness of presented implementations
\item \textbf{Appendix A}: Detailed DGP specifications for Monte Carlo
\item \textbf{Appendix B}: Complete test coverage matrix
\end{itemize}



% ============================================================================
% PART VII: TIME SERIES CAUSAL METHODS
% ============================================================================

\part{Time Series Causal Methods}
\label{part:timeseries}

\chapter{Time Series Causal Methods}
\label{ch:timeseries}

\section{Introduction}

Time series causal inference addresses a fundamental question: when observing variables evolving over time, can we identify which variables \textit{cause} changes in others? Unlike cross-sectional methods that exploit treatment assignment, time series methods leverage temporal precedence and dynamic structure.

\begin{insight}
Time series causality differs from cross-sectional causality in two key ways:
\begin{itemize}
    \item \textbf{Temporal precedence}: Causes must precede effects
    \item \textbf{Dynamic structure}: Effects propagate through lagged relationships
\end{itemize}
The Granger (1969) framework operationalizes ``causality'' as predictive improvement, while Structural VAR identifies contemporaneous causal effects through economic restrictions.
\end{insight}

This chapter develops:
\begin{enumerate}
    \item \textbf{VAR foundations}: Estimation, stationarity, lag selection
    \item \textbf{Granger causality}: Predictive causality tests
    \item \textbf{Cointegration and VECM}: Long-run equilibrium relationships
    \item \textbf{Structural VAR}: Causal identification through restrictions
    \item \textbf{Impulse responses}: Dynamic causal effects
    \item \textbf{Modern methods}: Local projections, proxy SVAR, TVP-VAR
\end{enumerate}


\section{Vector Autoregression (VAR)}

\subsection{The VAR Model}

A VAR($p$) model for $k$-dimensional time series $\mathbf{y}_t$:
\begin{equation}
\mathbf{y}_t = \mathbf{c} + \mathbf{A}_1 \mathbf{y}_{t-1} + \mathbf{A}_2 \mathbf{y}_{t-2} + \cdots + \mathbf{A}_p \mathbf{y}_{t-p} + \boldsymbol{\varepsilon}_t
\label{eq:var-p}
\end{equation}

where:
\begin{itemize}
    \item $\mathbf{y}_t \in \R^k$: Vector of endogenous variables
    \item $\mathbf{c} \in \R^k$: Intercept vector
    \item $\mathbf{A}_j \in \R^{k \times k}$: Coefficient matrices at lag $j$
    \item $\boldsymbol{\varepsilon}_t \sim (0, \boldsymbol{\Sigma})$: Reduced-form residuals
\end{itemize}

\begin{definition}[Stability Condition]
VAR($p$) is stable if and only if all eigenvalues of the companion matrix $\mathbf{F}$ lie inside the unit circle:
\begin{equation}
\mathbf{F} = \begin{pmatrix}
\mathbf{A}_1 & \mathbf{A}_2 & \cdots & \mathbf{A}_{p-1} & \mathbf{A}_p \\
\mathbf{I}_k & \mathbf{0} & \cdots & \mathbf{0} & \mathbf{0} \\
\mathbf{0} & \mathbf{I}_k & \cdots & \mathbf{0} & \mathbf{0} \\
\vdots & & \ddots & & \vdots \\
\mathbf{0} & \mathbf{0} & \cdots & \mathbf{I}_k & \mathbf{0}
\end{pmatrix}
\end{equation}
Stability ensures stationarity and finite impulse responses.
\end{definition}

\subsection{Estimation}

\begin{minted}[fontsize=\small]{python}
from causal_inference.timeseries import var_estimate, check_stability

# Estimate VAR(2) on 3-variable system
import numpy as np
np.random.seed(42)
T, k = 200, 3
data = np.random.randn(T, k)

result = var_estimate(data, lags=2)

print(f"Coefficient A1 shape: {result.coefficients[0].shape}")  # (3, 3)
print(f"Residual covariance:\n{result.sigma}")
print(f"Log-likelihood: {result.log_likelihood:.2f}")
print(f"AIC: {result.aic:.2f}")
print(f"BIC: {result.bic:.2f}")

# Check stability
is_stable = check_stability(result)
print(f"VAR is stable: {is_stable}")
\end{minted}

\subsection{Stationarity Testing}

Before VAR estimation, verify stationarity of each series.

\begin{minted}[fontsize=\small]{python}
from causal_inference.timeseries import (
    adf_test,
    kpss_test,
    confirmatory_stationarity_test,
    difference_series,
)

# ADF test (null: unit root)
adf_result = adf_test(data[:, 0], lags=4)
print(f"ADF statistic: {adf_result.statistic:.4f}")
print(f"p-value: {adf_result.p_value:.4f}")
print(f"Stationary: {adf_result.is_stationary}")

# KPSS test (null: stationary) - confirmatory
kpss_result = kpss_test(data[:, 0])
print(f"KPSS statistic: {kpss_result.statistic:.4f}")
print(f"Stationary: {kpss_result.is_stationary}")

# Combined test for robustness
combined = confirmatory_stationarity_test(data[:, 0])
print(f"Confirmed stationary: {combined.is_stationary}")

# If non-stationary, difference
if not combined.is_stationary:
    data_diff = difference_series(data[:, 0], order=1)
\end{minted}

\begin{table}[h]
\centering
\caption{Stationarity Test Interpretation}
\begin{tabular}{lll}
\toprule
Test & Null Hypothesis & Decision Rule \\
\midrule
ADF & Unit root (non-stationary) & Reject $\Rightarrow$ stationary \\
PP & Unit root (non-stationary) & Reject $\Rightarrow$ stationary \\
KPSS & Stationary & Fail to reject $\Rightarrow$ stationary \\
\bottomrule
\end{tabular}
\end{table}

\subsection{Lag Selection}

\begin{minted}[fontsize=\small]{python}
from causal_inference.timeseries import select_lag_order

# Compare information criteria for lags 1-8
lag_result = select_lag_order(data, max_lag=8)

print(f"Optimal lag by AIC: {lag_result.aic_lag}")
print(f"Optimal lag by BIC: {lag_result.bic_lag}")
print(f"Optimal lag by HQC: {lag_result.hqc_lag}")

# Criterion values at each lag
for lag, (aic, bic) in enumerate(zip(lag_result.aic_values, lag_result.bic_values), 1):
    print(f"Lag {lag}: AIC={aic:.2f}, BIC={bic:.2f}")
\end{minted}

\begin{insight}
Lag selection tradeoffs:
\begin{itemize}
    \item \textbf{AIC}: Asymptotically efficient, may overfit in small samples
    \item \textbf{BIC}: Consistent (selects true order as $T \to \infty$), may underfit
    \item \textbf{Practical}: Use BIC for policy analysis (parsimony), AIC for forecasting
\end{itemize}
\end{insight}


\section{Granger Causality}

\subsection{Definition}

\begin{definition}[Granger Causality]
Variable $X$ Granger-causes $Y$ if past values of $X$ improve prediction of $Y$ beyond what past values of $Y$ alone provide:
\begin{equation}
\E[Y_t \mid \mathcal{F}_{t-1}] \neq \E[Y_t \mid \mathcal{F}_{t-1} \setminus \{X_{t-1}, X_{t-2}, \ldots\}]
\end{equation}
where $\mathcal{F}_{t-1}$ is the information set at time $t-1$.
\end{definition}

In VAR context, $X$ does \textbf{not} Granger-cause $Y$ iff all coefficients on lagged $X$ in the $Y$ equation are zero:
\begin{equation}
H_0: A_{1,yx} = A_{2,yx} = \cdots = A_{p,yx} = 0
\end{equation}

\subsection{Pairwise Granger Test}

\begin{minted}[fontsize=\small]{python}
from causal_inference.timeseries import granger_causality

# Test: does X Granger-cause Y?
# Construct data where X -> Y (with lag)
np.random.seed(42)
T = 200
x = np.cumsum(np.random.randn(T))
y = np.zeros(T)
for t in range(1, T):
    y[t] = 0.6 * x[t-1] + 0.3 * y[t-1] + np.random.randn()

data = np.column_stack([y, x])  # Column 0: Y, Column 1: X

result = granger_causality(data, lags=2)

print(f"X Granger-causes Y: {result.granger_causes}")
print(f"F-statistic: {result.f_statistic:.4f}")
print(f"p-value: {result.p_value:.6f}")
\end{minted}

\subsection{Granger Causality Matrix}

\begin{minted}[fontsize=\small]{python}
from causal_inference.timeseries import granger_causality_matrix

# Test all pairwise Granger relationships
gc_matrix = granger_causality_matrix(data, lags=2, alpha=0.05)

# gc_matrix.causes[i, j] = True means variable j Granger-causes variable i
print("Granger causality matrix (row <- column):")
print(gc_matrix.causes)
print("\nP-values:")
print(gc_matrix.p_values)
\end{minted}

\subsection{Limitations}

\begin{enumerate}
    \item \textbf{Not true causality}: Granger causality is \textit{predictive}, not \textit{structural}
    \item \textbf{Omitted variables}: May find spurious causality if common cause omitted
    \item \textbf{Instantaneous effects}: Cannot detect contemporaneous causation
    \item \textbf{Nonlinear dynamics}: Linear test may miss nonlinear relationships
\end{enumerate}


\section{Cointegration and VECM}

\subsection{Cointegration}

When individual series are $I(1)$ (integrated of order 1) but a linear combination is $I(0)$ (stationary), the series are \textit{cointegrated}.

\begin{definition}[Cointegration]
Variables $\mathbf{y}_t = (y_{1t}, \ldots, y_{kt})'$ are cointegrated with cointegrating rank $r$ if:
\begin{enumerate}
    \item Each $y_{it}$ is $I(1)$
    \item There exist $r$ linearly independent vectors $\boldsymbol{\beta}_1, \ldots, \boldsymbol{\beta}_r$ such that $\boldsymbol{\beta}_j' \mathbf{y}_t$ is $I(0)$
\end{enumerate}
\end{definition}

\subsection{Johansen Cointegration Test}

\begin{minted}[fontsize=\small]{python}
from causal_inference.timeseries import johansen_test

# Test for cointegration among 3 I(1) series
result = johansen_test(data, lags=2, deterministic='c')

print(f"Cointegrating rank: {result.rank}")
print(f"\nTrace statistics:")
for r, (stat, cv) in enumerate(zip(result.trace_stats, result.trace_critical)):
    print(f"  r={r}: stat={stat:.2f}, 5% cv={cv:.2f}")

print(f"\nCointegrating vectors (normalized):")
print(result.beta[:, :result.rank])
\end{minted}

\subsection{Vector Error Correction Model (VECM)}

If cointegration exists, use VECM rather than VAR in differences:
\begin{equation}
\Delta \mathbf{y}_t = \boldsymbol{\alpha} \boldsymbol{\beta}' \mathbf{y}_{t-1} + \sum_{j=1}^{p-1} \boldsymbol{\Gamma}_j \Delta \mathbf{y}_{t-j} + \boldsymbol{\varepsilon}_t
\end{equation}

where:
\begin{itemize}
    \item $\boldsymbol{\beta}$: Cointegrating vectors (long-run equilibrium)
    \item $\boldsymbol{\alpha}$: Adjustment coefficients (speed of correction)
    \item $\boldsymbol{\beta}' \mathbf{y}_{t-1}$: Error correction term
\end{itemize}

\begin{minted}[fontsize=\small]{python}
from causal_inference.timeseries import vecm_estimate, vecm_granger_causality

# Estimate VECM with rank 1
vecm_result = vecm_estimate(data, lags=2, rank=1)

print(f"Cointegrating vector beta:\n{vecm_result.beta}")
print(f"Adjustment coefficients alpha:\n{vecm_result.alpha}")

# Granger causality in VECM framework
gc_result = vecm_granger_causality(vecm_result, cause_idx=1, effect_idx=0)
print(f"Granger-causes in VECM: {gc_result.granger_causes}")
\end{minted}


\section{Structural VAR (SVAR)}

\subsection{From Reduced Form to Structural Form}

The reduced-form VAR:
\begin{equation}
\mathbf{y}_t = \mathbf{c} + \sum_{j=1}^{p} \mathbf{A}_j \mathbf{y}_{t-j} + \boldsymbol{\varepsilon}_t, \quad \boldsymbol{\varepsilon}_t \sim (0, \boldsymbol{\Sigma})
\end{equation}

The structural form recovers \textit{orthogonal} structural shocks $\mathbf{u}_t$:
\begin{equation}
\mathbf{B}_0 \mathbf{y}_t = \mathbf{d} + \sum_{j=1}^{p} \mathbf{B}_j \mathbf{y}_{t-j} + \mathbf{u}_t, \quad \mathbf{u}_t \sim (0, \mathbf{I}_k)
\end{equation}

The relationship: $\boldsymbol{\varepsilon}_t = \mathbf{B}_0^{-1} \mathbf{u}_t$, so $\boldsymbol{\Sigma} = \mathbf{B}_0^{-1} (\mathbf{B}_0^{-1})'$.

\begin{insight}
The identification problem: $\boldsymbol{\Sigma}$ has $k(k+1)/2$ unique elements, but $\mathbf{B}_0^{-1}$ has $k^2$ elements. We need $k(k-1)/2$ additional restrictions to achieve just-identification.
\end{insight}

\subsection{Cholesky Identification}

The Cholesky decomposition provides a recursive structure: $\boldsymbol{\Sigma} = \mathbf{P}\mathbf{P}'$ where $\mathbf{P}$ is lower-triangular.

\begin{minted}[fontsize=\small]{python}
from causal_inference.timeseries import var_estimate, cholesky_svar

# Estimate reduced-form VAR
var_result = var_estimate(data, lags=2)

# Cholesky identification
svar_result = cholesky_svar(var_result)

print(f"Impact matrix (B0^{-1}):\n{svar_result.B0_inv}")
print(f"Identification method: {svar_result.identification_method}")
\end{minted}

The ordering matters: variable $j$ can be contemporaneously affected by variables $1, \ldots, j-1$ but not by $j+1, \ldots, k$.

\subsection{Short-Run Restrictions}

General short-run restrictions on $\mathbf{B}_0^{-1}$:

\begin{minted}[fontsize=\small]{python}
from causal_inference.timeseries import short_run_svar

# Define restrictions: None = unrestricted, 0 = zero restriction
# Example: Variable 2 doesn't respond contemporaneously to shock 1
restrictions = np.array([
    [None, 0,    0   ],
    [None, None, 0   ],
    [None, None, None],
])

svar_result = short_run_svar(var_result, restrictions)

print(f"Short-run impact matrix:\n{svar_result.B0_inv}")
\end{minted}

\subsection{Long-Run Restrictions (Blanchard-Quah)}

For permanent vs. transitory shock identification, impose that certain shocks have no long-run effect.

\begin{theorem}[Long-Run Multiplier]
The long-run effect of structural shocks on levels is:
\begin{equation}
\boldsymbol{\Xi} = (\mathbf{I}_k - \mathbf{A}_1 - \cdots - \mathbf{A}_p)^{-1} \mathbf{B}_0^{-1}
\end{equation}
Long-run restrictions impose zeros on $\boldsymbol{\Xi}$.
\end{theorem}

\begin{minted}[fontsize=\small]{python}
from causal_inference.timeseries import long_run_svar

# Blanchard-Quah: Supply shock has permanent effect, demand shock doesn't
# For 2-variable system: [output, unemployment]
# Demand shock (shock 2) has zero long-run effect on output

long_run_restrictions = np.array([
    [None, 0],     # Demand shock has no long-run effect on output
    [None, None],  # Both shocks can affect unemployment
])

svar_bq = long_run_svar(var_result, long_run_restrictions)

print(f"Long-run impact matrix:\n{svar_bq.long_run_impact}")
\end{minted}


\section{Impulse Response Functions}

\subsection{Definition}

\begin{definition}[Impulse Response Function]
The IRF measures the response of variable $i$ at horizon $h$ to a one-unit structural shock to variable $j$:
\begin{equation}
\text{IRF}_{ij}(h) = \frac{\partial y_{i,t+h}}{\partial u_{j,t}} = \mathbf{e}_i' \boldsymbol{\Phi}_h \mathbf{B}_0^{-1} \mathbf{e}_j
\end{equation}
where $\boldsymbol{\Phi}_h$ is the $h$-step VMA coefficient matrix.
\end{definition}

\subsection{Computing IRFs}

\begin{minted}[fontsize=\small]{python}
from causal_inference.timeseries import compute_irf, bootstrap_irf

# Compute structural IRFs
irf_result = compute_irf(svar_result, horizons=20)

print(f"IRF shape: {irf_result.irf.shape}")  # (k, k, horizons+1)

# Response of variable 0 to shock 1 at each horizon
response_0_to_shock_1 = irf_result.irf[0, 1, :]
print(f"Response of Y0 to shock 1: {response_0_to_shock_1[:5]}")

# Bootstrap confidence intervals
irf_boot = bootstrap_irf(var_result, svar_result, horizons=20, n_bootstrap=500)

print(f"95% CI at horizon 5: [{irf_boot.ci_lower[0,1,5]:.4f}, {irf_boot.ci_upper[0,1,5]:.4f}]")
\end{minted}

\subsection{Plotting IRFs}

\begin{minted}[fontsize=\small]{python}
import matplotlib.pyplot as plt

horizons = np.arange(irf_result.horizons + 1)

fig, ax = plt.subplots(figsize=(8, 4))
ax.plot(horizons, irf_result.irf[0, 1, :], 'b-', label='Point estimate')
ax.fill_between(horizons, irf_boot.ci_lower[0, 1, :], irf_boot.ci_upper[0, 1, :],
                alpha=0.3, label='95% CI')
ax.axhline(0, color='k', linestyle='--', linewidth=0.5)
ax.set_xlabel('Horizon')
ax.set_ylabel('Response')
ax.set_title('Response of Y0 to Structural Shock 1')
ax.legend()
plt.tight_layout()
\end{minted}


\section{Forecast Error Variance Decomposition}

\subsection{Definition}

FEVD decomposes the $h$-step forecast error variance of each variable into contributions from each structural shock.

\begin{definition}[FEVD]
The fraction of $h$-step forecast error variance in variable $i$ explained by shock $j$:
\begin{equation}
\text{FEVD}_{ij}(h) = \frac{\sum_{\ell=0}^{h-1} (\mathbf{e}_i' \boldsymbol{\Phi}_\ell \mathbf{B}_0^{-1} \mathbf{e}_j)^2}{\sum_{\ell=0}^{h-1} \mathbf{e}_i' \boldsymbol{\Phi}_\ell \boldsymbol{\Sigma} \boldsymbol{\Phi}_\ell' \mathbf{e}_i}
\end{equation}
By construction, $\sum_j \text{FEVD}_{ij}(h) = 1$ for all $i, h$.
\end{definition}

\begin{minted}[fontsize=\small]{python}
from causal_inference.timeseries import compute_fevd, variance_contribution_table

# Compute FEVD
fevd_result = compute_fevd(svar_result, horizons=20)

print(f"FEVD shape: {fevd_result.fevd.shape}")  # (k, k, horizons+1)

# At horizon 10, what fraction of Y0 variance is from each shock?
print(f"FEVD for Y0 at h=10: {fevd_result.fevd[0, :, 10]}")

# Summary table
table = variance_contribution_table(fevd_result, horizons=[1, 5, 10, 20])
print(table)
\end{minted}

\subsection{Historical Decomposition}

Decompose the actual series into contributions from each structural shock:

\begin{minted}[fontsize=\small]{python}
from causal_inference.timeseries import historical_decomposition

hd_result = historical_decomposition(svar_result, var_result.data)

# Contribution of shock 1 to variable 0 over time
shock1_contribution = hd_result.contributions[:, 0, 1]  # (T, k, k)
print(f"Shock 1 contribution to Y0: shape {shock1_contribution.shape}")
\end{minted}


\section{Modern Identification Methods}

\subsection{Sign Restrictions (Uhlig 2005)}

When theory provides only inequality restrictions (e.g., ``monetary contraction reduces output''), use sign restrictions for set identification.

\begin{minted}[fontsize=\small]{python}
from causal_inference.timeseries import (
    sign_restriction_svar,
    SignRestrictionConstraint,
    create_monetary_policy_constraints,
)

# Define constraints: shock 0 (monetary policy) must:
# - Increase interest rate (variable 0) on impact: positive
# - Decrease output (variable 1) for horizons 0-3: negative
# - Decrease prices (variable 2) for horizons 0-3: negative

constraints = [
    SignRestrictionConstraint(shock_idx=0, variable_idx=0, horizon=0, sign="+"),
    SignRestrictionConstraint(shock_idx=0, variable_idx=1, horizon=0, sign="-"),
    SignRestrictionConstraint(shock_idx=0, variable_idx=1, horizon=1, sign="-"),
    SignRestrictionConstraint(shock_idx=0, variable_idx=2, horizon=0, sign="-"),
]

# Or use built-in monetary policy constraints
constraints = create_monetary_policy_constraints(
    rate_idx=0, output_idx=1, price_idx=2, max_horizon=4
)

# Estimate via rejection sampling
result = sign_restriction_svar(
    var_result,
    constraints=constraints,
    n_draws=10000,
    n_accept=500,
)

print(f"Accepted {result.n_accepted} draws out of {result.n_draws}")
print(f"IRF median at h=0: {result.irf_median[:, 0, 0]}")
print(f"IRF 16th percentile: {result.irf_16[:, 0, 0]}")
print(f"IRF 84th percentile: {result.irf_84[:, 0, 0]}")
\end{minted}

\subsection{Proxy SVAR (Stock \& Watson 2012)}

Use external instruments (proxies) correlated with structural shocks but uncorrelated with other shocks.

\begin{assumption}[Proxy SVAR Assumptions]
For proxy $z_t$ identifying shock $u_{1,t}$:
\begin{enumerate}
    \item \textbf{Relevance}: $\E[z_t u_{1,t}] \neq 0$
    \item \textbf{Exogeneity}: $\E[z_t u_{j,t}] = 0$ for $j \neq 1$
\end{enumerate}
\end{assumption}

\begin{minted}[fontsize=\small]{python}
from causal_inference.timeseries import proxy_svar, weak_instrument_diagnostics

# z is an external instrument for the first structural shock
# E.g., narrative monetary policy shocks, oil price news

result = proxy_svar(var_result, proxy=z, shock_idx=0)

print(f"First-stage F: {result.first_stage_f:.2f}")
print(f"Impact response:\n{result.impact_response}")

# Check instrument strength
diag = weak_instrument_diagnostics(result)
print(f"Robust F-statistic: {diag.robust_f:.2f}")
print(f"Weak instrument: {diag.is_weak}")
\end{minted}

\subsection{Local Projections (Jord\`{a} 2005)}

Direct projection method that doesn't require VAR specification:
\begin{equation}
y_{i,t+h} = \alpha^h + \beta^h_i \text{shock}_t + \gamma^h \mathbf{X}_t + \varepsilon_{t+h}
\label{eq:local-projection}
\end{equation}

Advantages over VAR-based IRFs:
\begin{itemize}
    \item Robust to VAR misspecification
    \item Easier to incorporate nonlinearities and state-dependence
    \item Simpler inference (OLS at each horizon)
\end{itemize}

\begin{minted}[fontsize=\small]{python}
from causal_inference.timeseries import (
    local_projection_irf,
    compare_lp_var_irf,
    state_dependent_lp,
)

# Standard LP-IRF
lp_result = local_projection_irf(
    data=data,
    shock_idx=0,           # Which variable is shocked
    response_idx=1,        # Which variable responds
    horizons=20,
    lags=4,                # Control lags
    newey_west_lags=4,     # HAC standard errors
)

print(f"LP-IRF at horizon 5: {lp_result.coefficients[5]:.4f}")
print(f"SE: {lp_result.standard_errors[5]:.4f}")
print(f"95% CI: [{lp_result.ci_lower[5]:.4f}, {lp_result.ci_upper[5]:.4f}]")

# Compare LP vs VAR-based IRF
comparison = compare_lp_var_irf(data, lags=2, horizons=20)
# comparison.lp_irf, comparison.var_irf

# State-dependent LP (e.g., recession vs expansion)
# state = 1 if recession, 0 otherwise
state = (data[:, 1] < np.median(data[:, 1])).astype(float)

sdlp_result = state_dependent_lp(
    data=data,
    shock_idx=0,
    response_idx=1,
    state=state,
    horizons=20,
)

print(f"IRF in low state (recession): {sdlp_result.coefficients_low}")
print(f"IRF in high state (expansion): {sdlp_result.coefficients_high}")
\end{minted}


\section{Time-Varying Parameter VAR (TVP-VAR)}

\subsection{Model}

Allow VAR coefficients to evolve over time (Primiceri 2005):
\begin{align}
\mathbf{y}_t &= \mathbf{c}_t + \mathbf{A}_{1,t} \mathbf{y}_{t-1} + \cdots + \mathbf{A}_{p,t} \mathbf{y}_{t-p} + \boldsymbol{\varepsilon}_t \\
\boldsymbol{\theta}_t &= \boldsymbol{\theta}_{t-1} + \boldsymbol{\eta}_t, \quad \boldsymbol{\eta}_t \sim N(0, \mathbf{Q})
\end{align}

where $\boldsymbol{\theta}_t$ stacks all time-varying parameters.

\begin{minted}[fontsize=\small]{python}
from causal_inference.timeseries import (
    tvp_var_estimate,
    compute_tvp_irf,
    coefficient_change_test,
)

# Estimate TVP-VAR via Kalman filter
tvp_result = tvp_var_estimate(
    data,
    lags=2,
    state_var=0.01,    # Prior on state variance Q
)

print(f"Coefficients shape: {tvp_result.coefficients.shape}")  # (T, k*k*p + k)
print(f"Smoothed coefficients available: {tvp_result.smoothed is not None}")

# IRF at specific time point
irf_t100 = compute_tvp_irf(tvp_result, time_idx=100, horizons=20)
print(f"IRF at t=100, h=5: {irf_t100.irf[0, 1, 5]:.4f}")

# Test for significant time variation
test = coefficient_change_test(tvp_result)
print(f"Time-varying coefficients: {test.reject_constant}")
print(f"p-value: {test.p_value:.4f}")
\end{minted}

\subsection{Time-Varying IRFs}

\begin{minted}[fontsize=\small]{python}
from causal_inference.timeseries import compute_tvp_irf_all_times

# IRFs at each time point
all_irfs = compute_tvp_irf_all_times(tvp_result, horizons=20)

# Plot evolution of impact response over time
plt.figure(figsize=(10, 4))
impact_over_time = all_irfs[:, 0, 1, 0]  # Response of 0 to shock 1, impact
plt.plot(impact_over_time)
plt.xlabel('Time')
plt.ylabel('Impact response')
plt.title('Time-varying impact of shock 1 on variable 0')
\end{minted}


\section{PCMCI: Causal Discovery for Time Series}

The PCMCI algorithm (Runge et al., 2019) combines PC algorithm principles with momentary conditional independence for time series causal discovery.

\begin{minted}[fontsize=\small]{python}
from causal_inference.timeseries import pcmci, pcmci_plus

# Discover lagged causal structure
result = pcmci(
    data,
    max_lag=3,
    alpha=0.05,
    ci_test='parcorr',  # Partial correlation test
)

print(f"Discovered {len(result.links)} causal links:")
for link in result.links:
    print(f"  {link.source} (lag {link.lag}) -> {link.target}: p={link.p_value:.4f}")

# PCMCI+ for contemporaneous effects
result_plus = pcmci_plus(data, max_lag=3, alpha=0.05)
print(f"\nContemporaneous links: {len(result_plus.contemporaneous_links)}")
\end{minted}


\section{Monte Carlo Validation}

\subsection{VAR Estimation}

\begin{table}[h]
\centering
\caption{VAR Coefficient Recovery ($T=200$, 1,000 simulations)}
\begin{tabular}{lcccc}
\toprule
Coefficient & True & Mean Est. & Bias & RMSE \\
\midrule
$A_{1,11}$ & 0.50 & 0.498 & -0.002 & 0.071 \\
$A_{1,12}$ & 0.30 & 0.301 & 0.001 & 0.068 \\
$A_{1,21}$ & 0.00 & -0.003 & -0.003 & 0.065 \\
$A_{1,22}$ & 0.40 & 0.399 & -0.001 & 0.067 \\
\bottomrule
\end{tabular}
\end{table}

\subsection{IRF Coverage}

\begin{table}[h]
\centering
\caption{IRF Bootstrap Coverage ($T=200$, 500 simulations)}
\begin{tabular}{lccc}
\toprule
Horizon & Bootstrap Method & Coverage & SE Accuracy \\
\midrule
1 & Residual & 94.6\% & 5.2\% \\
5 & Residual & 93.8\% & 7.1\% \\
10 & Residual & 92.4\% & 9.8\% \\
1 & Moving Block & 95.1\% & 4.8\% \\
5 & Moving Block & 94.2\% & 6.5\% \\
10 & Moving Block & 93.6\% & 8.2\% \\
\bottomrule
\end{tabular}
\end{table}

\subsection{Granger Causality Power}

\begin{table}[h]
\centering
\caption{Granger Causality Power ($\alpha=0.05$, true effect $\beta=0.5$)}
\begin{tabular}{lccc}
\toprule
Sample Size & Size (no effect) & Power (with effect) & Type I Error \\
\midrule
$T=100$ & 5.2\% & 68.4\% & 5.2\% \\
$T=200$ & 4.8\% & 89.2\% & 4.8\% \\
$T=500$ & 5.1\% & 99.1\% & 5.1\% \\
\bottomrule
\end{tabular}
\end{table}


\section{Practical Workflow}

\subsection{Complete Analysis Pipeline}

\begin{minted}[fontsize=\small]{python}
from causal_inference.timeseries import (
    confirmatory_stationarity_test,
    difference_series,
    johansen_test,
    select_lag_order,
    var_estimate,
    vecm_estimate,
    cholesky_svar,
    compute_irf,
    bootstrap_irf,
    compute_fevd,
)

# Step 1: Check stationarity of each series
stationary = []
for i in range(data.shape[1]):
    result = confirmatory_stationarity_test(data[:, i])
    stationary.append(result.is_stationary)
    print(f"Variable {i}: {'stationary' if result.is_stationary else 'non-stationary'}")

# Step 2: Test for cointegration if non-stationary
if not all(stationary):
    coint = johansen_test(data, lags=2)
    print(f"Cointegrating rank: {coint.rank}")

    if coint.rank > 0:
        # Use VECM
        model = vecm_estimate(data, lags=2, rank=coint.rank)
        print("Using VECM")
    else:
        # Difference and use VAR
        data_diff = np.column_stack([difference_series(data[:, i]) for i in range(data.shape[1])])
        model = var_estimate(data_diff, lags=select_lag_order(data_diff).bic_lag)
        print("Using VAR on differenced data")
else:
    # Use VAR in levels
    optimal_lag = select_lag_order(data).bic_lag
    model = var_estimate(data, lags=optimal_lag)
    print(f"Using VAR({optimal_lag}) in levels")

# Step 3: Structural identification
svar = cholesky_svar(model)

# Step 4: Compute IRFs with inference
irf = compute_irf(svar, horizons=20)
irf_boot = bootstrap_irf(model, svar, horizons=20, n_bootstrap=1000)

# Step 5: FEVD
fevd = compute_fevd(svar, horizons=20)

print("\n=== Results ===")
print(f"Impact response of Y0 to shock 1: {irf.irf[0, 1, 0]:.4f}")
print(f"Cumulative response at h=10: {np.sum(irf.irf[0, 1, :11]):.4f}")
print(f"FEVD of Y0 at h=10: shock contributions = {fevd.fevd[0, :, 10]}")
\end{minted}


\section{Discussion}

\subsection{Key Takeaways}

\begin{enumerate}
    \item \textbf{Granger $\neq$ structural causality}: Predictive improvement, not mechanism
    \item \textbf{Identification is crucial}: SVAR requires economic restrictions (Cholesky, long-run, sign, proxy)
    \item \textbf{Cointegration matters}: Use VECM for $I(1)$ series with equilibrium relationships
    \item \textbf{Robustness checks}: Compare VAR-IRF vs LP-IRF; use multiple identification strategies
    \item \textbf{Time variation}: Consider TVP-VAR when structural breaks suspected
\end{enumerate}

\subsection{Method Selection}

\begin{table}[h]
\centering
\caption{Time Series Causal Method Selection}
\begin{tabular}{lll}
\toprule
Question & Method & Key Assumption \\
\midrule
Does $X$ predict $Y$? & Granger causality & VAR correctly specified \\
Long-run equilibrium? & VECM & Cointegration exists \\
Dynamic causal effects? & SVAR + IRF & Identification valid \\
Robust to misspecification? & Local Projections & Shock exogeneity \\
External instrument? & Proxy SVAR & Relevance + exogeneity \\
Agnostic on effects? & Sign restrictions & Inequality constraints \\
Structural change? & TVP-VAR & Random walk coefficients \\
Discover structure? & PCMCI & Faithfulness \\
\bottomrule
\end{tabular}
\end{table}

\subsection{Connection to Other Methods}

\begin{itemize}
    \item \textbf{IV} (Chapter~\ref{ch:iv}): Proxy SVAR uses IV logic for shock identification
    \item \textbf{RDD} (Chapter~\ref{ch:rdd}): Event studies combine DiD with time series structure
    \item \textbf{Causal Discovery} (Appendix~\ref{app:discovery}): PCMCI extends PC algorithm to time series
    \item \textbf{Bounds} (Chapter~\ref{ch:bounds}): Sign restrictions provide set identification
\end{itemize}



% ----------------------------------------------------------------------------
% BACK MATTER
% ----------------------------------------------------------------------------

\backmatter

% Appendices
\appendix

\chapter{Julia Implementation Reference}
\label{app:julia}

% ============================================================================
% Appendix B: Julia Code Reference
% Written: 2025-12-31
% ============================================================================

This appendix provides Julia equivalents for all estimators presented in the main text. The Julia implementations follow the SciML Problem-Estimator-Solution pattern and achieve numerical parity with the Python implementations to 10 decimal places.

\section{Installation}

\begin{verbatim}
cd julia
julia --project -e "using Pkg; Pkg.instantiate()"
\end{verbatim}

\section{Architecture Overview}

The Julia implementation follows the SciML pattern:

\begin{enumerate}
    \item \textbf{Problem}: Immutable struct containing data and parameters
    \item \textbf{Estimator}: Algorithm choice (e.g., \texttt{SimpleATE}, \texttt{DoublyRobust})
    \item \textbf{Solution}: Results struct with estimates, SEs, CIs
\end{enumerate}

\begin{minted}{julia}
# General pattern
problem = SomeProblem(outcomes, treatment, covariates, parameters)
solution = solve(problem, SomeEstimator())

# Access results
solution.estimate   # Point estimate
solution.se         # Standard error
solution.ci_lower   # Lower CI bound
solution.ci_upper   # Upper CI bound
\end{minted}

% ============================================================================
\section{RCT Estimators}
\label{sec:julia_rct}
% ============================================================================

\subsection{Simple ATE (Difference-in-Means)}

\begin{minted}{julia}
"""
    SimpleATE <: AbstractRCTEstimator

Simple difference-in-means with Neyman heteroskedasticity-robust variance.

    tau_hat = Y_bar_1 - Y_bar_0
    SE(tau_hat) = sqrt(s1^2/n1 + s0^2/n0)
"""
struct SimpleATE <: AbstractRCTEstimator end

function solve(problem::RCTProblem, estimator::SimpleATE)::RCTSolution
    (; outcomes, treatment, parameters) = problem

    y1 = outcomes[treatment]
    y0 = outcomes[.!treatment]
    n1, n0 = length(y1), length(y0)

    # Point estimate
    ate = mean(y1) - mean(y0)

    # Neyman variance
    se = sqrt(var(y1)/n1 + var(y0)/n0)

    # Confidence interval (Satterthwaite df)
    df = satterthwaite_df(var(y1), n1, var(y0), n0)
    ci_lower, ci_upper = confidence_interval_t(ate, se, parameters.alpha, df)

    return RCTSolution(estimate=ate, se=se, ci_lower=ci_lower, ...)
end
\end{minted}

\textbf{Usage:}
\begin{minted}{julia}
using CausalEstimators

outcomes = [10.0, 12.0, 4.0, 5.0]
treatment = [true, true, false, false]
problem = RCTProblem(outcomes, treatment, nothing, nothing, (alpha=0.05,))
solution = solve(problem, SimpleATE())
\end{minted}

% ============================================================================
\section{Propensity Score Methods}
\label{sec:julia_psm}
% ============================================================================

\subsection{Propensity Score Estimation}

\begin{minted}{julia}
"""
    estimate_propensity_scores(treatment, covariates) -> PropensityResult

Estimate propensity scores via logistic regression.
"""
function estimate_propensity_scores(treatment, covariates)
    X = hcat(ones(size(covariates, 1)), covariates)
    coef = logistic_regression(X, treatment)
    propensity = logistic.(X * coef)
    auc = compute_propensity_auc(propensity, treatment)
    return PropensityResult(propensity, auc, coef)
end
\end{minted}

\subsection{Nearest Neighbor Matching}

\begin{minted}{julia}
"""
    nearest_neighbor_matching(treated_ps, control_ps; with_replacement=true)

Match treated units to controls based on propensity score distance.
Returns indices of matched controls for each treated unit.
"""
function nearest_neighbor_matching(
    treated_ps::Vector{T},
    control_ps::Vector{T};
    with_replacement::Bool = true,
    caliper::Union{Nothing, T} = nothing,
) where {T<:Real}
    n_treated = length(treated_ps)
    matches = Vector{Int}(undef, n_treated)

    available = trues(length(control_ps))

    for i in 1:n_treated
        distances = abs.(control_ps .- treated_ps[i])

        if !with_replacement
            distances[.!available] .= Inf
        end

        best_idx = argmin(distances)

        # Check caliper
        if !isnothing(caliper) && distances[best_idx] > caliper
            matches[i] = 0  # No valid match
        else
            matches[i] = best_idx
            available[best_idx] = false
        end
    end

    return matches
end
\end{minted}

% ============================================================================
\section{Weighting Estimators}
\label{sec:julia_weighting}
% ============================================================================

\subsection{Inverse Probability Weighting (IPW)}

\begin{minted}{julia}
"""
    ObservationalIPW <: AbstractObservationalEstimator

IPW estimator using Hajek (normalized) weights.

    tau_IPW = [sum T_i*Y_i/e(X_i)] / [sum T_i/e(X_i)] -
              [sum (1-T_i)*Y_i/(1-e(X_i))] / [sum (1-T_i)/(1-e(X_i))]
"""
struct ObservationalIPW <: AbstractObservationalEstimator end

function solve(problem::ObservationalProblem, ::ObservationalIPW)
    (; outcomes, treatment, covariates) = problem

    # Estimate propensity scores
    prop_result = estimate_propensity_scores(treatment, covariates)
    e = prop_result.propensity

    # Compute IPW weights
    weights = ifelse.(treatment, 1 ./ e, 1 ./ (1 .- e))

    # Hajek estimator (normalized)
    mu1 = sum(weights[treatment] .* outcomes[treatment]) / sum(weights[treatment])
    mu0 = sum(weights[.!treatment] .* outcomes[.!treatment]) / sum(weights[.!treatment])
    ate = mu1 - mu0

    # Influence function variance
    influence = @. ifelse(treatment, (outcomes - mu1)/e, -(outcomes - mu0)/(1-e))
    se = sqrt(var(influence) / length(outcomes))

    return IPWSolution(estimate=ate, se=se, ...)
end
\end{minted}

\subsection{Doubly Robust (AIPW)}

\begin{minted}{julia}
"""
    DoublyRobust <: AbstractObservationalEstimator

AIPW estimator combining IPW with outcome regression.

Double robustness: consistent if EITHER propensity OR outcome model correct.

    tau_DR = (1/n) sum_i [T_i/e(X_i)*(Y_i - mu1(X_i)) + mu1(X_i)
                        - (1-T_i)/(1-e(X_i))*(Y_i - mu0(X_i)) - mu0(X_i)]
"""
struct DoublyRobust <: AbstractObservationalEstimator end

function solve(problem::ObservationalProblem, ::DoublyRobust)
    (; outcomes, treatment, covariates) = problem
    T = convert(Vector{Float64}, treatment)

    # Estimate propensity and outcome models
    e = estimate_propensity_scores(treatment, covariates).propensity
    mu0, mu1 = fit_outcome_models(outcomes, treatment, covariates)

    # AIPW formula
    treated_contrib = @. T/e * (outcomes - mu1) + mu1
    control_contrib = @. (1-T)/(1-e) * (outcomes - mu0) + mu0

    ate = mean(treated_contrib .- control_contrib)

    # Influence function variance
    influence = treated_contrib .- control_contrib .- ate
    se = sqrt(mean(influence.^2) / length(outcomes))

    return DRSolution(estimate=ate, se=se, ...)
end
\end{minted}

% ============================================================================
\section{CATE Estimators}
\label{sec:julia_cate}
% ============================================================================

\subsection{S-Learner}

\begin{minted}{julia}
"""
    SLearner <: AbstractCATEEstimator

Single model approach: fit mu(X, T) -> Y, then CATE(x) = mu_hat(x, 1) - mu_hat(x, 0).
"""
struct SLearner <: AbstractCATEEstimator end

function solve(problem::CATEProblem, ::SLearner)
    (; outcomes, treatment, covariates) = problem
    n = length(outcomes)

    # Augment covariates with treatment
    X_aug = hcat(ones(n), covariates, Float64.(treatment))
    coef = fit_linear_model(X_aug, outcomes)

    # CATE = coefficient on treatment (for linear model)
    # For flexibility, compute explicitly:
    X_treat = hcat(ones(n), covariates, ones(n))
    X_ctrl = hcat(ones(n), covariates, zeros(n))

    cate = X_treat * coef .- X_ctrl * coef  # = coef[end] for all i
    ate = mean(cate)

    return CATESolution(cate=cate, ate=ate, ...)
end
\end{minted}

\subsection{T-Learner}

\begin{minted}{julia}
"""
    TLearner <: AbstractCATEEstimator

Two model approach: fit mu1(X) on treated, mu0(X) on control.
CATE(x) = mu1_hat(x) - mu0_hat(x)
"""
struct TLearner <: AbstractCATEEstimator end

function solve(problem::CATEProblem, ::TLearner)
    (; outcomes, treatment, covariates) = problem

    # Fit separate models
    X1 = hcat(ones(sum(treatment)), covariates[treatment, :])
    X0 = hcat(ones(sum(.!treatment)), covariates[.!treatment, :])

    coef1 = fit_linear_model(X1, outcomes[treatment])
    coef0 = fit_linear_model(X0, outcomes[.!treatment])

    # Predict for all units
    X_all = hcat(ones(length(outcomes)), covariates)
    cate = X_all * coef1 .- X_all * coef0

    return CATESolution(cate=cate, ate=mean(cate), ...)
end
\end{minted}

% ============================================================================
\section{Causal Forests}
\label{sec:julia_forests}
% ============================================================================

Causal forests are implemented via \texttt{econml} in Python. For Julia usage, call Python via \texttt{PyCall.jl}:

\begin{minted}{julia}
using PyCall

econml = pyimport("econml.dml")

function causal_forest(outcomes, treatment, covariates; n_estimators=100, honest=true)
    model = econml.CausalForestDML(
        n_estimators=n_estimators,
        honest=honest,
        cv=5,
    )
    model.fit(Y=outcomes, T=treatment, X=covariates)

    cate = model.effect(covariates)
    ate_inf = model.ate_inference(X=covariates)

    return (
        cate = vec(cate),
        ate = mean(cate),
        ate_se = ate_inf.stderr_mean,
    )
end
\end{minted}

\begin{insight}
Critical: Always use \texttt{honest=true} for valid confidence intervals. Without honest splitting, CIs have severe undercoverage (e.g., 72\% instead of 95\%). See CONCERN-28.
\end{insight}

% ============================================================================
\section{Sensitivity Analysis}
\label{sec:julia_sensitivity}
% ============================================================================

\subsection{E-Value}

\begin{minted}{julia}
"""
    compute_e_value(rr::Real) -> Float64

E-value: minimum confounding strength to explain away effect.

For RR >= 1: E = RR + sqrt(RR * (RR - 1))
For RR < 1: Apply formula to 1/RR
"""
function compute_e_value(rr::Real)::Float64
    rr > 0 || throw(ArgumentError("RR must be positive"))

    rr_use = rr < 1.0 ? 1.0/rr : rr

    if isapprox(rr_use, 1.0)
        return 1.0  # No effect
    end

    return rr_use + sqrt(rr_use * (rr_use - 1.0))
end

# Effect type conversions
smd_to_rr(d) = exp(0.91 * d)  # VanderWeele approximation
ate_to_rr(ate, p0) = (p0 + ate) / p0
\end{minted}

\textbf{Usage:}
\begin{minted}{julia}
problem = EValueProblem(2.5; ci_lower=1.8, ci_upper=3.5, effect_type=:rr)
solution = solve(problem, EValue())

println("E-value: $(solution.e_value)")
println("E-value for CI: $(solution.e_value_ci)")
\end{minted}

\subsection{Rosenbaum Bounds}

\begin{minted}{julia}
"""
    solve(problem::RosenbaumProblem, ::RosenbaumBounds)

Sensitivity analysis for matched pairs using Wilcoxon signed-rank test.

Finds critical Gamma where p_upper > alpha.
"""
function solve(problem::RosenbaumProblem, ::RosenbaumBounds)
    differences = problem.treated_outcomes .- problem.control_outcomes

    # Wilcoxon signed-rank statistic
    T_plus, ranks, signs = compute_signed_rank_statistic(differences)

    gamma_grid = range(problem.gamma_range..., length=problem.n_gamma)
    p_upper = zeros(length(gamma_grid))

    for (i, gamma) in enumerate(gamma_grid)
        E_upper, _, Var_max, _ = compute_bounds_at_gamma(ranks, signs, gamma)
        p_upper[i] = normal_approximation_p(T_plus, E_upper, Var_max)
    end

    # Critical Gamma: first where p_upper > alpha
    gamma_critical = findfirst(>(problem.alpha), p_upper)
    gamma_critical = isnothing(gamma_critical) ? nothing : gamma_grid[gamma_critical]

    return RosenbaumSolution(gamma_critical=gamma_critical, ...)
end
\end{minted}

\textbf{Usage:}
\begin{minted}{julia}
treated = [10.5, 12.3, 8.7, 15.2, 11.1]
control = [8.2, 10.1, 7.5, 12.8, 9.3]
problem = RosenbaumProblem(treated, control; gamma_range=(1.0, 3.0))
solution = solve(problem, RosenbaumBounds())

println("Critical Gamma: $(solution.gamma_critical)")
\end{minted}

% ============================================================================
\section{Cross-Language Validation}
\label{sec:julia_validation}
% ============================================================================

All Julia implementations are validated against Python to 10 decimal places:

\begin{minted}{julia}
# Test parity with Python
@test isapprox(julia_ate, python_ate, rtol=1e-10)
@test isapprox(julia_se, python_se, rtol=1e-10)
\end{minted}

Current validation status:
\begin{itemize}
    \item \textbf{RCT estimators}: 100\% parity
    \item \textbf{PSM}: 100\% parity
    \item \textbf{Weighting}: 100\% parity (IPW, DR, TMLE)
    \item \textbf{CATE}: 100\% parity (S/T/X/R-learners, DML)
    \item \textbf{Sensitivity}: 100\% parity (E-value, Rosenbaum)
\end{itemize}


\chapter{Causal Discovery}
\label{app:discovery}

\section{Introduction}

Causal discovery learns causal structure from observational data---determining \textit{which} variables cause \textit{which} without experimental intervention. Unlike treatment effect estimation (which assumes known causal structure), discovery algorithms infer the causal graph itself.

\begin{insight}
Three paradigms for causal discovery:
\begin{enumerate}
    \item \textbf{Constraint-based} (PC, FCI): Test conditional independence
    \item \textbf{Score-based} (GES): Maximize model fit score
    \item \textbf{Functional} (LiNGAM): Exploit asymmetries in data distributions
\end{enumerate}
Each paradigm rests on different assumptions and yields different output types.
\end{insight}

\subsection{Graph Representations}

\begin{definition}[Graph Types]
\begin{itemize}
    \item \textbf{DAG}: Directed Acyclic Graph---unique causal structure
    \item \textbf{CPDAG}: Completed Partially Directed Graph---Markov equivalence class (undirected edges where direction unidentifiable)
    \item \textbf{PAG}: Partial Ancestral Graph---equivalence class allowing latent confounders
\end{itemize}
\end{definition}


\section{PC Algorithm}

The PC algorithm \citep{spirtes2000causation} is the canonical constraint-based method, testing conditional independence to learn graph structure.

\subsection{Algorithm Overview}

\begin{algorithm}
\caption{PC Algorithm}
\begin{algorithmic}[1]
\STATE Start with complete undirected graph
\REPEAT
    \FOR{each adjacent pair $(X, Y)$}
        \FOR{conditioning sets $\mathbf{S}$ of increasing size}
            \IF{$X \independent Y \mid \mathbf{S}$}
                \STATE Remove edge $X - Y$
                \STATE Record $\mathbf{S}$ as separation set
            \ENDIF
        \ENDFOR
    \ENDFOR
\UNTIL{no more edges removed}
\STATE Orient v-structures: $X \to Z \gets Y$ if $X - Z - Y$, $X \not\sim Y$, and $Z \notin \text{sepset}(X, Y)$
\STATE Apply Meek rules to orient remaining edges
\RETURN CPDAG
\end{algorithmic}
\end{algorithm}

\subsection{Assumptions}

\begin{assumption}[PC Algorithm Assumptions]
\label{ass:pc}
\begin{enumerate}
    \item \textbf{Causal Sufficiency}: No unmeasured confounders
    \item \textbf{Faithfulness}: All independencies arise from d-separation in the true DAG
    \item \textbf{Acyclicity}: No feedback loops
\end{enumerate}
\end{assumption}

\subsection{Implementation}

\begin{minted}[fontsize=\small]{python}
from causal_inference.discovery import (
    pc_algorithm,
    pc_skeleton,
    pc_orient,
    pc_conservative,
    generate_random_dag,
    generate_dag_data,
    compute_shd,
)
import numpy as np

# Generate ground truth DAG
np.random.seed(42)
true_dag = generate_random_dag(n_nodes=5, edge_prob=0.3, seed=42)
data, B = generate_dag_data(true_dag, n_samples=1000, seed=42)

# Run PC algorithm
result = pc_algorithm(
    data,
    alpha=0.01,           # Significance level for CI tests
    ci_test='fisher_z',   # CI test (Gaussian assumption)
    max_cond_size=None,   # Max conditioning set size
)

print(f"Estimated CPDAG edges: {result.cpdag.num_edges}")
print(f"Directed edges: {result.cpdag.directed_edges}")
print(f"Undirected edges: {result.cpdag.undirected_edges}")

# Evaluate against truth
shd = compute_shd(result.cpdag, true_dag)
print(f"Structural Hamming Distance: {shd}")
\end{minted}

\subsection{Skeleton Learning}

\begin{minted}[fontsize=\small]{python}
# Learn only the skeleton (undirected graph)
skeleton_result = pc_skeleton(data, alpha=0.01)

print(f"Skeleton edges: {skeleton_result.skeleton.edges}")
print(f"Separation sets: {skeleton_result.sep_sets}")
\end{minted}

\subsection{PC Variants}

\begin{minted}[fontsize=\small]{python}
# Conservative PC: More conservative v-structure orientation
# Requires all CI tests to agree on v-structure
result_cons = pc_conservative(data, alpha=0.01)

# Majority rule PC: V-structure if majority of tests agree
result_maj = pc_majority(data, alpha=0.01)
\end{minted}


\section{FCI Algorithm}

The FCI (Fast Causal Inference) algorithm extends PC to allow latent confounders, outputting a PAG rather than CPDAG.

\subsection{Output: Partial Ancestral Graph}

\begin{definition}[PAG Edge Marks]
\begin{itemize}
    \item $\to$: Arrowhead (definite causal direction)
    \item $\circ$: Circle (uncertain---could be arrow or tail)
    \item $-$: Tail (definite non-cause)
\end{itemize}
Edge types in PAG:
\begin{itemize}
    \item $X \to Y$: $X$ is ancestor of $Y$
    \item $X \leftrightarrow Y$: Latent common cause
    \item $X \circ\!\to Y$: $Y$ is not ancestor of $X$
    \item $X \circ\!-\!\circ Y$: Unknown relationship
\end{itemize}
\end{definition}

\subsection{Implementation}

\begin{minted}[fontsize=\small]{python}
from causal_inference.discovery import fci_algorithm

# FCI allows for latent confounders
result = fci_algorithm(
    data,
    alpha=0.01,
    ci_test='fisher_z',
)

print(f"PAG edges:")
for edge in result.pag.edges:
    print(f"  {edge.node1} {edge.mark1}--{edge.mark2} {edge.node2}")

# Check for latent confounder indicators (bidirected edges)
bidirected = [e for e in result.pag.edges if e.mark1 == '>' and e.mark2 == '<']
print(f"\nPossible latent confounders: {len(bidirected)}")
\end{minted}


\section{GES Algorithm}

Greedy Equivalence Search \citep{chickering2002optimal} is a score-based method that searches through the space of CPDAGs.

\subsection{Algorithm}

\begin{algorithm}
\caption{GES (Greedy Equivalence Search)}
\begin{algorithmic}[1]
\STATE Start with empty graph
\REPEAT \COMMENT{Forward phase}
    \STATE Add edge that most improves BIC score
\UNTIL{no edge addition improves score}
\REPEAT \COMMENT{Backward phase}
    \STATE Remove edge that most improves BIC score
\UNTIL{no edge removal improves score}
\RETURN CPDAG
\end{algorithmic}
\end{algorithm}

\subsection{Implementation}

\begin{minted}[fontsize=\small]{python}
from causal_inference.discovery import (
    ges_algorithm,
    local_score_bic,
    total_score,
)

# Run GES with BIC scoring
result = ges_algorithm(
    data,
    score_type='bic',      # BIC, AIC, or custom
    max_degree=None,       # Max node degree
)

print(f"GES CPDAG edges: {result.cpdag.num_edges}")
print(f"Final BIC score: {result.score:.2f}")
print(f"Forward additions: {result.n_forward_ops}")
print(f"Backward deletions: {result.n_backward_ops}")

# Compare to PC result
shd_ges = compute_shd(result.cpdag, true_dag)
print(f"GES SHD: {shd_ges}")
\end{minted}

\subsection{Score Functions}

\begin{minted}[fontsize=\small]{python}
from causal_inference.discovery import (
    local_score_bic,
    local_score_aic,
    total_score,
)

# Local score: score of single node given parents
# Higher (less negative) is better
score_node0 = local_score_bic(data, node_idx=0, parent_indices=[1, 2])
print(f"BIC score for X0 | X1, X2: {score_node0:.2f}")

# Total graph score
graph_score = total_score(data, result.cpdag, score_type='bic')
print(f"Total BIC score: {graph_score:.2f}")
\end{minted}


\section{LiNGAM}

LiNGAM (Linear Non-Gaussian Acyclic Model) exploits non-Gaussianity to identify the unique causal DAG, not just the equivalence class.

\subsection{Model}

\begin{definition}[LiNGAM Model]
\begin{equation}
\mathbf{x} = \mathbf{B} \mathbf{x} + \boldsymbol{\varepsilon}
\end{equation}
where:
\begin{itemize}
    \item $\mathbf{B}$ is strictly lower-triangular (after permutation to causal order)
    \item $\boldsymbol{\varepsilon}$ has non-Gaussian, mutually independent components
\end{itemize}
\end{definition}

\begin{theorem}[LiNGAM Identifiability]
If the structural equation model is linear, acyclic, and errors are non-Gaussian and independent, the causal DAG is uniquely identifiable (up to permutation) via ICA.
\end{theorem}

\subsection{ICA-LiNGAM}

\begin{minted}[fontsize=\small]{python}
from causal_inference.discovery import (
    ica_lingam,
    direct_lingam,
    bootstrap_lingam,
    check_non_gaussianity,
)

# Check non-Gaussianity assumption
ng_stats = check_non_gaussianity(data)
print(f"Non-Gaussian: {ng_stats.is_non_gaussian}")
print(f"Kurtosis by variable: {ng_stats.kurtosis}")

# Generate non-Gaussian data for LiNGAM
data_ng, B_true = generate_dag_data(
    true_dag,
    n_samples=1000,
    noise_type='laplace',  # Non-Gaussian
    seed=42,
)

# ICA-based LiNGAM
result = ica_lingam(data_ng)

print(f"Estimated causal order: {result.causal_order}")
print(f"Estimated adjacency matrix B:\n{result.adjacency_matrix}")
\end{minted}

\subsection{DirectLiNGAM}

DirectLiNGAM is a more efficient algorithm that directly estimates causal order.

\begin{minted}[fontsize=\small]{python}
# DirectLiNGAM (faster, more stable)
result_direct = direct_lingam(data_ng)

print(f"Causal order: {result_direct.causal_order}")
print(f"True order: {true_dag.topological_order()}")

# Check order accuracy
order_correct = result_direct.causal_order == true_dag.topological_order()
print(f"Order correct: {order_correct}")
\end{minted}

\subsection{Bootstrap Confidence}

\begin{minted}[fontsize=\small]{python}
# Bootstrap for edge confidence
boot_result = bootstrap_lingam(data_ng, n_bootstrap=100)

print(f"Edge stability (fraction of bootstraps):")
print(boot_result.edge_stability)

# Only edges stable in >80% of bootstraps
stable_edges = np.where(boot_result.edge_stability > 0.8)
print(f"Stable edges: {list(zip(stable_edges[0], stable_edges[1]))}")
\end{minted}


\section{Independence Tests}

All constraint-based algorithms require conditional independence (CI) tests.

\subsection{Fisher's Z Test (Gaussian)}

\begin{minted}[fontsize=\small]{python}
from causal_inference.discovery import fisher_z_test, partial_correlation_test

# Test X0 _|_ X1 | X2
result = fisher_z_test(
    data,
    x_idx=0,
    y_idx=1,
    cond_set=[2],
    alpha=0.05,
)

print(f"Test statistic: {result.statistic:.4f}")
print(f"P-value: {result.p_value:.4f}")
print(f"Independent: {result.independent}")
\end{minted}

\subsection{Available Tests}

\begin{table}[h]
\centering
\caption{Conditional Independence Tests}
\begin{tabular}{llll}
\toprule
Test & Data Type & Assumption & Speed \\
\midrule
\texttt{fisher\_z} & Continuous & Gaussian & Fast \\
\texttt{partial\_corr} & Continuous & Linear & Fast \\
\texttt{g\_squared} & Categorical & Multinomial & Medium \\
\texttt{kernel\_ci} & Continuous & Nonlinear OK & Slow \\
\bottomrule
\end{tabular}
\end{table}

\begin{minted}[fontsize=\small]{python}
from causal_inference.discovery import kernel_ci_test

# Kernel CI test for nonlinear relationships
result = kernel_ci_test(
    data,
    x_idx=0,
    y_idx=1,
    cond_set=[2, 3],
    n_permutations=500,
)

print(f"Kernel CI p-value: {result.p_value:.4f}")
\end{minted}


\section{Evaluation Metrics}

\subsection{Structural Hamming Distance}

\begin{definition}[SHD]
The Structural Hamming Distance counts edge differences between two graphs:
\begin{equation}
\text{SHD} = \text{missing} + \text{extra} + \text{reversed}
\end{equation}
where ``reversed'' counts edges with wrong direction.
\end{definition}

\begin{minted}[fontsize=\small]{python}
from causal_inference.discovery import (
    compute_shd,
    skeleton_f1,
    orientation_accuracy,
)

# SHD between estimated and true
shd = compute_shd(result.cpdag, true_dag)
print(f"SHD: {shd}")

# Skeleton recovery (ignoring direction)
f1 = skeleton_f1(result.cpdag, true_dag)
print(f"Skeleton F1: {f1.f1:.3f}")
print(f"  Precision: {f1.precision:.3f}")
print(f"  Recall: {f1.recall:.3f}")

# Orientation accuracy (for directed edges)
orient = orientation_accuracy(result.cpdag, true_dag)
print(f"Orientation accuracy: {orient:.3f}")
\end{minted}

\subsection{Markov Equivalence}

\begin{minted}[fontsize=\small]{python}
from causal_inference.discovery import is_markov_equivalent, dag_to_cpdag

# Check if two DAGs are Markov equivalent
dag1 = generate_random_dag(4, edge_prob=0.4, seed=1)
dag2 = generate_random_dag(4, edge_prob=0.4, seed=2)

equiv = is_markov_equivalent(dag1, dag2)
print(f"Markov equivalent: {equiv}")

# Convert DAG to its CPDAG (equivalence class representative)
cpdag = dag_to_cpdag(true_dag)
print(f"Undirected edges in CPDAG: {cpdag.undirected_edges}")
\end{minted}


\section{Monte Carlo Validation}

\subsection{PC Algorithm}

\begin{table}[h]
\centering
\caption{PC Algorithm Performance ($n=1000$, 5 variables, 500 runs)}
\begin{tabular}{lcccc}
\toprule
$\alpha$ & Skeleton F1 & SHD & Orientation Acc. & Type I Error \\
\midrule
0.001 & 0.72 & 3.2 & 0.81 & 0.8\% \\
0.01 & 0.85 & 1.8 & 0.83 & 2.1\% \\
0.05 & 0.88 & 1.5 & 0.79 & 5.8\% \\
0.10 & 0.86 & 1.9 & 0.74 & 11.2\% \\
\bottomrule
\end{tabular}
\end{table}

\subsection{LiNGAM}

\begin{table}[h]
\centering
\caption{LiNGAM Order Recovery (5 variables, Laplace noise)}
\begin{tabular}{lccc}
\toprule
Sample Size & Order Accuracy & Edge Recovery & Mean SHD \\
\midrule
$n=200$ & 61\% & 0.72 & 2.8 \\
$n=500$ & 78\% & 0.85 & 1.4 \\
$n=1000$ & 91\% & 0.93 & 0.6 \\
$n=2000$ & 97\% & 0.98 & 0.2 \\
\bottomrule
\end{tabular}
\end{table}

\subsection{GES vs PC}

\begin{table}[h]
\centering
\caption{GES vs PC Comparison ($n=1000$, 500 runs)}
\begin{tabular}{lccc}
\toprule
Algorithm & Mean SHD & Skeleton F1 & Runtime (ms) \\
\midrule
PC ($\alpha=0.01$) & 1.8 & 0.85 & 45 \\
GES (BIC) & 1.5 & 0.88 & 120 \\
GES (AIC) & 2.1 & 0.91 & 125 \\
\bottomrule
\end{tabular}
\end{table}


\section{Complete Example}

\begin{minted}[fontsize=\small]{python}
from causal_inference.discovery import (
    pc_algorithm,
    ges_algorithm,
    direct_lingam,
    generate_random_dag,
    generate_dag_data,
    compute_shd,
    skeleton_f1,
    check_non_gaussianity,
)
import numpy as np

np.random.seed(42)

# Generate ground truth
true_dag = generate_random_dag(n_nodes=6, edge_prob=0.3, seed=42)
print(f"True DAG has {true_dag.num_edges} edges")
print(f"True edges: {true_dag.edges}")

# Generate Gaussian data
data_gauss, _ = generate_dag_data(true_dag, n_samples=1000,
                                   noise_type='gaussian', seed=42)

# Generate non-Gaussian data
data_laplace, _ = generate_dag_data(true_dag, n_samples=1000,
                                     noise_type='laplace', seed=42)

# Compare algorithms
print("\n=== Algorithm Comparison ===")

# PC (Gaussian)
pc_result = pc_algorithm(data_gauss, alpha=0.01)
pc_shd = compute_shd(pc_result.cpdag, true_dag)
pc_f1 = skeleton_f1(pc_result.cpdag, true_dag)
print(f"PC:     SHD={pc_shd}, Skeleton F1={pc_f1.f1:.3f}")

# GES (Gaussian)
ges_result = ges_algorithm(data_gauss, score_type='bic')
ges_shd = compute_shd(ges_result.cpdag, true_dag)
ges_f1 = skeleton_f1(ges_result.cpdag, true_dag)
print(f"GES:    SHD={ges_shd}, Skeleton F1={ges_f1.f1:.3f}")

# LiNGAM (non-Gaussian)
lingam_result = direct_lingam(data_laplace)
lingam_shd = compute_shd(lingam_result.dag, true_dag)
lingam_f1 = skeleton_f1(lingam_result.dag, true_dag)
order_correct = list(lingam_result.causal_order) == list(true_dag.topological_order())
print(f"LiNGAM: SHD={lingam_shd}, Skeleton F1={lingam_f1.f1:.3f}, Order={order_correct}")

# Check assumptions
print("\n=== Assumption Checks ===")
ng_gauss = check_non_gaussianity(data_gauss)
ng_laplace = check_non_gaussianity(data_laplace)
print(f"Gaussian data non-Gaussian test: {ng_gauss.is_non_gaussian}")
print(f"Laplace data non-Gaussian test: {ng_laplace.is_non_gaussian}")
\end{minted}


\section{Discussion}

\subsection{Method Selection}

\begin{table}[h]
\centering
\caption{Causal Discovery Method Selection}
\begin{tabular}{lllll}
\toprule
Method & Output & Latent OK? & Assumption & Best For \\
\midrule
PC & CPDAG & No & Faithfulness & Gaussian, no confounders \\
FCI & PAG & Yes & Faithfulness & Unknown confounders \\
GES & CPDAG & No & Score consistency & Large, sparse graphs \\
LiNGAM & DAG & No & Non-Gaussian & Unique DAG needed \\
\bottomrule
\end{tabular}
\end{table}

\subsection{Practical Recommendations}

\begin{enumerate}
    \item \textbf{Start with PC}: Simple, fast, well-understood properties
    \item \textbf{Check Gaussianity}: If non-Gaussian, LiNGAM gives unique DAG
    \item \textbf{Suspect confounders}: Use FCI instead of PC
    \item \textbf{Cross-validate}: Run multiple algorithms, focus on consistent edges
    \item \textbf{Domain knowledge}: Use as constraints or to validate results
\end{enumerate}

\subsection{Limitations}

\begin{enumerate}
    \item \textbf{Sample size}: Need large $n$ for reliable CI tests
    \item \textbf{High dimension}: Exponential conditioning sets
    \item \textbf{Faithfulness failures}: Near-violations cause errors
    \item \textbf{Nonlinearity}: Linear methods miss nonlinear effects
    \item \textbf{Time series}: Standard methods assume i.i.d. (see PCMCI in Chapter~\ref{ch:timeseries})
\end{enumerate}

\subsection{Connection to Treatment Effects}

Causal discovery complements treatment effect methods:
\begin{itemize}
    \item \textbf{Variable selection}: Identify confounders, mediators, colliders
    \item \textbf{DAG validation}: Check assumed causal structure
    \item \textbf{Instrumental variables}: Discover potential instruments
    \item \textbf{Sensitivity analysis}: Explore alternative structures
\end{itemize}



\bibliographystyle{plainnat}
\bibliography{bibliography}

\end{document}
